%%============================================================================
%% thesis.tex  --  Adelaide Punt, DPhil thesis (body)
%%
%% Title: Noise and Uncertainty in Molecular Property Prediction:
%%        Robust Models and Adaptive Representations
%%============================================================================


\chapter{Introduction}\label{chap:intro}
% ---- TAKEAWAYS ----
% - A property prediction model is a molecular representation paired with a model; both limit what it can learn.
% - Labels in molecular data are noisy, and that noise is rarely homoscedastic Gaussian: its size varies with source and assay.
% - Under label noise, model architecture matters more than representation for robustness.
% - Predicted uncertainty rises with noise but does not find the corrupted labels.
% - Methods that judge a label by its nearest neighbours' labels did worst under noise. The hypothesis: the distance is computed over general-purpose representations full of irrelevant features.
% - The adaptive hybrid representation selects task-relevant features from several representations; it is tested as a model input, beyond small molecules, and as the space for similarity and noise handling.
% - The aim throughout: property prediction that stays reliable when the training labels are noisy.
% ---- END TAKEAWAYS ----
% ---- NOTES: Introduction ----
% Takeaway: property predictions filter which compounds get tested, and their training labels are noisy in ways most studies do not model.
%
% Paragraph 2: label noise in molecular data
% - Add one sentence: this noise is rarely homoscedastic Gaussian. Its size varies with source and assay, as the figures in this paragraph show.
% - Change: "Measurements from a single assay lowers" becomes "lower".
%
% Paragraph 3: the first research question
% - Change: "QSAR modelings" becomes "QSAR modelling"; \cite{Cortes2015} becomes \citet{Cortes2015}.
% - Change: no result yet supports "both model and representation choice dominate predictive performance on clean data". Chapter~\ref{chap:robustness} splits the variance only under noise.
% - Result needed: the same variance split on R^2 before any noise is added.
%
% Paragraph 4: from noise robustness to adaptive representations
% - Change: noise mitigation is a subset of learning with noisy labels (LNL). Write "learning with noisy labels, including noise mitigation".
% - Say what was tried: 12 LNL methods adapted from classification to regression.
% - Change: "all of the methods that used similarity estimates failed" is too broad. The heteroscedastic GP compares molecules through its kernel and held up.
% - Say instead: the methods that judged a label by its neighbours' labels did worst. The counts go in Section~\ref{sec:lnl_motivation}.
% - Change: "vector operations" becomes "distances". The list of distances and representations tried goes in Section~\ref{sec:lnl_motivation}, not here.
% - For the similar property principle, cite \citet{Johnson1990} as the Background does, or keep \citep{LopezPerez2024} as a review.
% - Change: typos. "simiarity" to "similarity"; "who's" to "whose"; "how it improve" to "how it improves".
% - Change: the last sentence says the hybrid improves performance. Chapters 5 and 6 test this; their results are not yet in.
%
% Paragraph 5 (to write; answers the TODO below): three to four sentences. Other theses in the field end their opener this way.
% - The aim: property prediction that stays reliable when training labels are noisy.
% - The two parts: models and representations that survive label noise, then adapting the representation itself.
% - Hand-over: "The rest of this thesis is organised as follows".
% ---- END NOTES ----

New avenues in the drug discovery pipeline have opened up thanks to advances in large dataset collection, analysis, and computation. Therapeutic targets are identified, potential drug candidates are screened and optimized, and the most promising of these go on to pre-clinical and clinical trials \citep{Lill2007}. There are various methods by which potential drug candidates can be screened. Virtual screening allows for the rapid and large-scale analysis of libraries of compounds, and is a particularly useful approach used to speed up the drug discovery process \citep{Neves2018}. One standard screening method is the prediction of bioactivities and pharmacological properties of such candidates. There is a long history of using Quantitative Structure-Activity Relationship (QSAR) models to make these predictions, to act as a filter for compounds with undesirable properties, and thus reduce the time and cost associated with experimental screening \citep{Sabando2021}. QSAR models aim to predict the relationship between a molecule's molecular features as encoded within a particular molecular representation and a desired molecular property \citep{Sabando2021} of interest. A QSAR model typically consists of a molecular representation paired with a machine learning (ML) model. An ideal molecular representation will allow a model to learn only the relevant information about each molecule with respect to the target property. Representations are either `static' hand-crafted descriptors or circular fingerprints that encode predefined structural and physicochemical features \citep{rogers2010, David2020, sns, avalon}; or learned representations, where the encoding is derived from a predictive model, often taken from the activations of the final layer of a neural network, or a latent space embedding \citep{Sabando2021, kishimoto2023}. 

QSAR models are of course, sensitive to the quality of the input data on which they are trained. Experimental noise can arise from measurement errors, variability under experimental conditions, varying provenance of the data, and intrinsic biological and chemical variability. Even the very measures of model performance may be flawed if test and validation sets are laden with noise. \citet{Kolmar2021} show that the variance associated with experimental data can contribute more to prediction error than the model’s own error. There is also a limit to how much noise the model can tolerate; \citet{jorner2021} found that once the model error approaches the magnitude of the experimental measurement error on labels, adding more training data yields diminishing returns in predictive performance. In the context of activity prediction, the inherent variability within experimental setups and assays, batch effects, measurement errors, and biological variability introduce additional layers of complexity \citep{Kolmar2021, Fourches2010}. All of this increases the likelihood of encountering noise in the data \citep{Fang2022}. Even the molecular representations in the training data themselves may be incorrect or misleading \citep{Walters2023}, and chemical datasets contain too few repeat measurements of each compound to characterize the distribution the label is drawn from. Thus, models treat a single label as the population mean \citep{Kolmar2021}. Repeat measurements of the same compound-target pair disagree by about $0.68$ log units for pIC$_{50}$, and by about $0.54$ log units for pK$_i$ \citep{Kalliokoski2013, Kramer2012}. Measurements from a single assay lowers the mean disagreement from $0.50$ to $0.27$ log units \citep{landrum2024}.  

% TODO: rephrase: However, all of the methods that used simiarity metrics failed, even when I tried using different molecular representations and vector operations
This thesis addresses the challenge of conducting molecular property prediction when the underlying experimental data contain noisy labels. My initial research question is which combinations of molecular representations and predictive model architectures maintain predictive performance under label corruption, and whether probabilistic methods that estimate uncertainty justify their additional complexity. As I will show, both model and representation choice dominate predictive performance on clean data, however, this relationship inverts under label corruption as predictive architecture becomes the primary factor in determining noise robustness.  Most existing work on noise in QSAR modelings assumes that label noise can be modeled as random error drawn from a Gaussian distribution, and they evaluate either molecular representations or predictive model architectures in isolation. \cite{Cortes2015} examined 12 algorithms across multiple noise levels but used fixed representations. Studies that compare representations are typically done on clean data. No work has systematically compared the relative importance of representation versus architecture choice for noise robustness, assessed whether probabilistic methods justify their computational overhead in noisy settings, or tested whether findings generalize across different corruption patterns and molecular property types.

Building on these findings, I began exploring techniques involving learning with noisy labels and noise mitigation. Many of these methods incorporate measures of molecular similarity, relying on the assumption that molecules which are similar in structure share similar properties \citep{LopezPerez2024}. However, all of the methods that used simiarity estimates failed, even when I tried varying the molecular representations and vector operations from the standard ECFP4 and Tanimoto coefficient. I hypothesized that to generate accurate similarity estimates, you need a specialised molecular representation who's features relate directly to similarity in that molecular domain. This brings me to my next research question, how do you select the best set of features for a given prediction task? I will outline how I developed an adaptive feature selection methodology that takes in a set of molecular representations, identifies each representation's most important features, and combines a select set of those top features into a new hybrid representation. I will then demonstrate how it improve performance across a range of predictive tasks beyond the small molecule space, and how it can be applied to the context of similarity estimation and downstream noise mitigation methods. 
% TODO: how do I end this? Check other theses


\section{The questions this thesis asks}\label{sec:intro:questions}

% ---- NOTES: The questions this thesis asks ----
% Paragraph 1: the four questions of the first half
% - Add a closing sentence: Chapter~\ref{chap:robustness} answers all four.
% - Change: the last sentence introduces NoiseInject again. Point to Chapter~\ref{chap:background}, which describes it.
%
% Paragraph 2: the second half, with the chapter that answers each question
% - Add first: do standard representations, and the similarity computed over them, fail on a given dataset? (Chapter~\ref{chap:flaws})
% - Add: is predicting each molecule's activity from its five nearest neighbours more accurate in a hybrid than in its base representations? (Chapter~\ref{chap:similarity})
% - Selecting features: Chapter~\ref{chap:method}.
% - Hybrid noise robustness: no chapter reports it yet. If it stays, "(retention percentage)" becomes "normalised area under the retention curve".
% - Distance-based noise mitigation: Chapter~\ref{chap:similarity}. Widen it to noise detection and estimation, which that chapter also reports.
% - Importance weighting: Chapter~\ref{chap:method} for hybrid construction and model training; Chapter~\ref{chap:similarity} for distances.
% - In that question, "noise mitigation tasks" becomes "distance-based noise detection and learning with noisy labels".
% - Sample weighting and label smoothing: Chapter~\ref{chap:similarity}. Reword: they are tested only on real labels with nothing injected, never under Gaussian corruption.
% ---- END NOTES ----

% TODO: review, needs manual review
The first half of this thesis addresses four questions about the behavior of QSAR models under label noise. First, how does the choice of molecular representation and model architecture explain the variance in predictive performance, and does that division change when you instead consider how much performance a QSAR model is able to retain under increasing label noise? Second, what does the choice of model provide with label noise? Within this question we compare the effects of deterministic and probabilistic modeling. Third, does the type of label noise matter, or just the quantity of it? We explore different techniques of generating and adding label noise to both clean data, as well as data with inherent experimental noise. Fourth, what is the impact of artificial label noise on a model's uncertainty estimation? Does the estimated uncertainty rise with increased label noise and can these uncertainty values point to which labels are corrupted? Together, these questions aim to identify what, if anything, makes a QSAR model robust to noise. Finally, we introduce NoiseInject, an open-source Python package for benchmarking noise robustness and uncertainty quantification in machine learning models.

% TODO: review, needs manual review
The second half asks whether the representation itself can be adapted to the task:
\begin{enumerate}
    \item How do you select the best set of features for a given prediction task?
    \item Do hybrid representations achieve higher noise robustness (retention percentage) than their base representations?
    \item Do the selected hybrid features help stabilize distance-based noise mitigation?
    \item Does feature importance weighting improve effectiveness of hybrid representation creation, model training, and noise mitigation tasks?
    \item When sample removal is suboptimal, can sample weighting or label smoothing provide superior noise handling under Gaussian corruption?
\end{enumerate}

\section{From noise robustness to adaptive representations}

% ---- NOTES: From noise robustness to adaptive representations ----
% Takeaway: the evidence behind each step from the first half to the hybrid. Three short paragraphs; method detail stays in the chapters cited.
%
% Paragraph 1: what the first half found (Chapter~\ref{chap:robustness})
% - Robustness is the normalised area under the R^2 retention curve as noise rises.
% - Under Gaussian noise on QM9, model architecture explains far more of the variance in robustness than representation. The percentages stay in Chapter~\ref{chap:robustness}.
% - Noise shared within a scaffold family, or clipped at an assay limit, did the most damage. The shape of the noise distribution barely mattered.
% - Mean predicted uncertainty rose with noise. Per-molecule uncertainty did not pick out the corrupted labels (Section~\ref{sec:uq}).
% - So uncertainty alone cannot say which labels to fix. That leads to methods that judge each label another way.
%
% Paragraph 2: the exploratory study of learning with noisy labels (Section~\ref{sec:lnl_motivation})
% - Twelve LNL methods were adapted to regression and tested on the QM9 HOMO-LUMO gap with Gaussian noise.
% - The worst configurations all judged a label by its neighbours' labels, using a distance between molecules. This held for each distance tried. Counts stay in Section~\ref{sec:lnl_motivation}.
% - The methods that held up estimated noise from the model's own loss or uncertainty.
% - One dataset and one noise type: state it as motivation, not a general result.
%
% Paragraph 3: why the distance failed, and the step to the hybrid
% - Hypothesis: on a given dataset, most features of a general-purpose representation do not vary or do not relate to the property. Chapter~\ref{chap:flaws} measures this.
% - Each such feature adds a term to the distance, so irrelevant features dominate it.
% - Proposed fix: select a fixed number of task-relevant features from several representations, and compute distances in that space.
% - A first version showed these selected features also work as ordinary model inputs. How to choose them became Chapter~\ref{chap:method}.
% - Chapter~\ref{chap:similarity} returns to distances, noise detection and LNL with the hybrid as the feature space.
% ---- END NOTES ----
\todo[inline, color=yellow!40]{Write Sec 1.4 from the outline bullets above. The raw material is Section~\ref{sec:lnl_motivation}.}


\section{Outline of the thesis and contributions}

% ---- NOTES: Outline of the thesis and contributions ----
% Takeaway: one short paragraph per chapter: what it does, what it found (Chapters 3 and 4) or will test (Chapters 5 and 6), and where it is published or planned.
% - Close each chapter's paragraph with one sentence on what is new in it compared with earlier work. Findings in words, no numbers.
%
% Chapter 2 (Background)
% - What the later chapters rely on, including the seven noise conditions of Chapter 3 and NoiseInject, the open-source benchmarking framework released with this work.
%
% Chapter 3 (Noise robustness across molecular representations and model architectures)
% - Adapted from "NoiseInject: A Benchmarking Framework to Study Noise Robustness in QSAR Models" (Punt, Hanser, Werner and Morris), in preparation for the Journal of Cheminformatics.
% - Main task: the QM9 HOMO-LUMO gap, whose computed labels carry negligible measurement noise. Repeated on LogD, Caco-2 and hERG Ki.
% - Findings in words only: those of Paragraph 1 in the section above, plus one. Random forests were the most consistent model across the four datasets.
%
% Chapter 4 (The flaws of the similarity assumption)
% - A motivation chapter: small measurements, no new method.
% - On a given dataset, standard representations carry many features that do not vary, and some map different molecules to the same vector.
% - Similarity computed over them inherits both problems. Methods built on it fail, including the label cleaning from the LNL study in the confirmation report.
% - Ends with what a representation would need to fix this. Chapter 5 builds it.
%
% Chapter 5 (Building adaptive hybrid representations). First planned method paper; venue to be confirmed.
% - The method: select each base representation's most important features, split a fixed budget of features between representations, and combine them.
% - Will test when it helps: against the best single representation and an ensemble of one model per representation.
% - Will test why: against other feature-selection methods and random features at the same number of features, and by reading the selected features against the chemistry.
% - Will test whether it extends beyond small molecules (antibodies in most depth) and to data held at separate sites.
% - Result needed: one sentence of findings.
%
% Chapter 6 (Hybrid representations for similarity and label noise). Second planned method paper; venue to be confirmed.
% - Asks whether a hybrid repairs the failures of Chapter 4: similarity, noise detection and estimation, and LNL including noise mitigation, one noise type at a time.
% - Result needed: one sentence of findings.
%
% Chapter 7 (Discussion and future directions)
% - Summarises each chapter, then argues: model the noise before choosing a method. No single method fits every dataset, and simple methods are often as good.
% - Then limitations and future directions.
%
% Contributions (short)
% - Publications: point to the Research Papers list in the Summary of Research Contributions.
% - Software: NoiseInject, and the KIRBy repository that implements the hybrid search. Repositories, licences and data: Appendix~\ref{app:availability}.
% ---- END NOTES ----
\todo[inline, color=yellow!40]{Write Sec 1.5 from the outline bullets above.}

\section{Collaborations and funding}

% ---- NOTES: Collaborations and funding ----
% Takeaway: who funded the work, what Lhasa Limited and OPIG contributed, and what is the author's own.
% - Funding: EPSRC Centre for Doctoral Training in Sustainable Approaches to Biomedical Science: Responsible and Reproducible Research (EP/S024093/1), with support from Lhasa Limited, a not-for-profit organisation.
% - Lhasa: Thierry Hanser and Stephane Werner co-supervised, with Garrett Morris as supervisor. Chapter~\ref{chap:robustness} is adapted from the jointly authored NoiseInject manuscript; contributions are in the Statement of Originality.
% - Lhasa supplied the hERG data used in Chapter~\ref{chap:method} as data held at separate sites. Name its federated-learning platform, Effiris, in one clause \citep{Bassani2023, Hanser2023}.
% - OPIG pointed to antibody and protein datasets used in Chapters~\ref{chap:method} and~\ref{chap:similarity} (list in Appendix~\ref{app:datasets}).
% - The OPIG collaboration also produced five posts on its blog, Blopig. Point to the Blogposts section.
% - One sentence: the hybrid, its software and the experiments in Chapters~\ref{chap:flaws} to~\ref{chap:similarity} are the author's own work.
% ---- END NOTES ----
\todo[inline, color=yellow!40]{Write Sec 1.6.}

\pagebreak
\chapter{Background}\label{chap:background}\label{sec2}
% ---- TAKEAWAYS ----
% - A representation turns a molecule into a fixed-length vector of features. Each family captures different information and has its own weakness, so no single one captures everything.
% - Model families differ in how far one wrong label spreads: a tree keeps it in one leaf, an ensemble dilutes it, a neural network can memorise it.
% - Similarity only exists relative to a representation and a coefficient. Activity cliffs and uninformative features break the similar property principle.
% - Real label noise is not independent, identically distributed or Gaussian. It has heavy tails, shared offsets between laboratories and censoring at assay limits.
% - Learning with noisy labels (LNL) covers every method for training on wrong labels. Noise mitigation is the subset that changes the labels. No single method wins.
% - Every split in the thesis is group-aware, never random. NoiseInject supplies the noise and the robustness and uncertainty metrics used in later chapters.
% ---- END TAKEAWAYS ----

\section{Introduction}

% ---- NOTES: Introduction ----
% Takeaway: this chapter introduces, once, the ideas every later chapter refers back to.
% Paragraph 1: one-sentence roadmap, in section order: representations, models, similarity and activity cliffs, noise in molecular data, LNL, NoiseInject, evaluation.
% - Change: the sentence below names only two parts, and ends in a colon with no list after it.
% ---- END NOTES ----

I will now break down the two components of QSAR models: molecular representations and predictive models:


\section{Molecular Representations in QSAR}

% ---- NOTES: Molecular Representations in QSAR ----
% Takeaway: a representation maps a molecule to a fixed-length vector of features that a model takes as input.
% Takeaway: each family encodes different information and fails in a different way. This is what the hybrid in Chapter 5 builds on.
% Paragraph 1: what a representation is.
% - Definition 2.1, numbered: a molecular representation maps a molecule to a fixed-length numeric vector. Computing it is featurization; each position is one feature.
% - Features can be bits, counts, real-valued descriptors or embedding dimensions. The same definition covers an antibody sequence mapped to a vector (Chapter 5).
% - Change: the opening paragraph says "the type of input labels". The subject is the input features, not the labels.
% - Change: a numbered definition needs \usepackage{amsthm} and \newtheorem{definition}{Definition}[chapter] in the preamble.
% Paragraph 2: canonical and randomized one-hot SMILES.
% - The canonical string is the one a toolkit picks by fixed rules. Randomized SMILES start the atom walk at a random atom \citep{ArusPous2019}.
% - Chapter 4's exploratory LNL study used both.
% - Do not cite Bertoni2020 for one-hot SMILES. It is about bioactivity signatures inferred by neural networks, not SMILES encoding.
% Paragraph 3: MACCS and Avalon.
% - Change: cite Durant2002 for MACCS and rogers2010 for ECFP. The paragraph cites only Xie2020 for both.
% - Change: ECFP4 in this thesis ignores chirality (radius 2, 2,048 bits). The paragraph says "active chirality flags".
% - MACCS: 166 predefined substructure questions, 167 bits as computed here. Each bit's meaning is known. Substructures not on the list are not seen.
% - Avalon: a second hashed fingerprint, 2,048 bits here, built through RDKit \citep{Gedeck2006}. Hashing folds different substructures onto one bit, as in ECFP.
% - What Avalon hashes is not yet sourced. Check the Gedeck2006 full text before writing that sentence.
% Paragraph 4: mordred.
% - More than 1,800 2D and 3D descriptors \citep{Moriwaki2018}. This thesis uses the 2D ones only, about 1,613 features.
% - Weakness: multicollinearity, as for descriptor vectors. Chapter 4 shows that most features do not vary on a given dataset.
% Paragraph 5: MolFormer.
% - A transformer encoder trained on the SMILES of 1.1 billion unlabelled molecules from PubChem and ZINC \citep{ross2022}. Used as a fixed encoder, without fine-tuning.
% Paragraph 6: MHG-GNN, one or two sentences.
% - Combines a molecular hypergraph grammar with a graph neural network \citep{kishimoto2023}. Gives a fixed-length vector read from the molecular graph.
% Paragraph 7 (new, short): protein and antibody language models, so Chapter 5 can point back.
% - Same idea as the chemical language models, on amino-acid sequences. One vector per residue; a chain's vector is usually their average.
% - The individual protein and antibody models, with their citations, belong in the Chapter 5 table "Representations by domain", not here.
% - Strength: dense vectors learned from far more sequences than any labelled dataset. Weakness: single features cannot be read, and computing them is costly.
% Closing: a comparison table, one row per family: what it encodes, fixed or learned, length in features, whether a feature can be read, main weakness.
% ---- END NOTES ----
The success of machine learning hinges on the quality of input data. Although the quality can refer to noise, which the majority of my work is centered around, in this case I'm talking about the type of input labels. Molecular representations provide different features with which QSAR models and noise detection techniques can be learned. However, a potential bottleneck on larger data sets should be noted; the more complex the data structure, the greater the computational cost. For smaller data sets, the number of parameters is limited by the size of the training set, to prevent overfitting. 

\begin{figure}[h]
	\centering 
	\includegraphics[width=1.05\textwidth]{molecular_representations.png}	
    \caption{\label{fig:molecular_representations}Visualisation of the molecule \lstinline{O=C1C=CN(C1)C=O} using three different methods: a) ECFP4 showing the molecule broken into substructures with circular neighbourhoods for each atom. Each substructure is hashed to produce unique identifiers, which are then compiled into a binary bit vector representing the presence or absence of specific substructures. b) SMILES tokenization and one-hot encoding of the molecule. The SMILES string is tokenized, and using a vocabulary of the most common tokens, these tokens are converted into a one-hot encoded matrix. c) Molecular graph representation where nodes represent atoms and edges represent bonds.}
\end{figure}

One of the most common and intuitive molecular representations are Simplified Molecular-Input Line-Entry System (SMILES) strings which consist of a sequence of ASCII characters generated by a fixed set of rules to specify a given chemical identity \citep{Weininger1988}. A common way to use SMILES strings as inputs to QSAR models is to tokenize and one-hot-encode them rather than to use the strings themselves \citep{TeachOpenCADD2023}. The strings are tokenized into atomic branches and symbols and mapped to indices of a vocabulary, and finally converted into binary vectors using one-hot encodings \citep{TeachOpenCADD2023}. This representation, referred to as one-hot-encoded SMILES, allows models to process molecular structures as bit-vector-based fingerprints or as sequences of vectors. Sequence-based models such as recurrent neural networks and transformers have proven to be particularly successful in capturing text semantics within a sequence. When applied to raw or one hot encoded (OHE) SMILES strings, these models can effectively capture the structure of the molecule in its textual representation \citep{TeachOpenCADD2023}. 

Another common representation is the molecular fingerprint, a vector containing information about the properties of the molecule in question. In a typical bit vector-based fingerprint, each bit accounts for the presence (1) or absence (0) of a given chemical feature or substructure \citep{Fassio2022}. One common fingerprint is the Molecular ACCess System (MACCS) key, a substructure key-based fingerprint that encodes predefined atom symbols and properties, as well as bond types \citep{Xie2020}. Another is the Extended Connectivity Fingerprint (ECFP), which encodes local neighborhoods around each atom and bond connectivity in molecules \citep{Xie2020}. ECFPs are typically constructed with a radius of 2, a length of 2048 bits, and active chirality flags. These existing molecular representations are well established and have been quite successful when coupled with QSAR models \citep{CortsCiriano2020, Sabando2021}. They are relatively simple, and as the bits can be mapped back to one or more substructures, they are human-interpretable, and furthermore easy-to-compute. However, they are subject to information loss due to bit collisions, which happens when two different substructures are encoded at the same bit position \citep{Fassio2022}. They are also based on a fixed procedure and are not trainable or adaptable to a particular dataset. Alternative representations that capture different structural, chemical, or physical features are able to overcome these limitations. 

Recent advances have led to the development of new molecular fingerprints. One such method is the Sort \& Slice (SNS) technique, an alternative to hash-based folding for pooling ECFP substructures \citep{sns}. Sort \& Slice begins by sorting all circular substructures in the training set according to their frequency of occurrence. It then "slices" away all but the $N$ most frequent substructures, which are subsequently used to generate an $N$-dimensional, bit-collision-free binary fingerprint. This approach solves the ambiguity problem of hashed fingerprints, where multiple substructures might be encoded at the same bit position, making it very interpretable \citep{sns}. Sort \& Slice fingerprints outperform ECFP4 with standard hash-based folding methods across various molecular property prediction tasks. These experiments consisted of different data splitting techniques, machine learning models, fingerprint diameters, fingerprint dimensions, and atomic invariants \citep{sns}. Furthermore, Sort \& Slice outperformed advanced substructure selection schemes such as filtered fingerprints and mutual information maximization \citep{sns}. 


Neural fingerprints are a category of learned  representations that overcome the fixed, hard-coded nature of traditional fingerprints. Unlike ECFP or MACCS keys, which rely on predefined hashing schemes or substructure catalogues, neural fingerprints train on molecular data to capture chemical information \citep{Jaeger2018, Chithrananda2020}. These representations are the result of deep learning architectures that have learned dense, continuous embeddings which can be used for downstream tasks. 

% TODO: review, needs manual review
Mol2vec, which adapts the Word2vec algorithm from NLP to molecular substructures \citep{Jaeger2018}, is common neural fingerprint. Mol2vec treats molecules like "sentences" and circular substructures derived from Morgan fingerprints like "words". The model learns vector embeddings for each substructure through unsupervised training, such that chemically similar substructures are mapped to nearby points in the embedding space \citep{Jaeger2018}. The embeddings of all present substructures are summed to obtain a full representation. This approach produces 300-dimensional dense vectors \citep{Jaeger2018}. Mol2vec has performed strongly on various QSAR benchmarks, matching or exceeding Morgan fingerprints on solubility and bioactivity prediction tasks \citep{Jaeger2018}.

Transformer-based architectures have also been applied to molecular representation learning, similarly stemming from NLPs ideas. MolBERT applies a bidirectional transformer to learn molecular representations from SMILES strings through self-supervised pre-training \citep{Fabian2020}. The model has domain-relevant auxiliary adaptations, such as the prediction of calculated physicochemical properties alongside the masked token prediction, to improve the chemical fidelity of the learned embeddings \citep{Fabian2020}. MolBERT demonstrated improvements over traditional fingerprints and graph-based methods on Virtual Screening and QSAR benchmarks \citep{Fabian2020}. ChemBERTa is another transformer-based approach which pre-trains masked language models on SMILES data from PubChem \citep{Chithrananda2020}.  ChemBERTa uses a tokenization strategy adapted for the SMILES notation, which allows the model to learn representations of atoms and functional groups through self-supervised learning \citep{Chithrananda2020}. Attention visualization reveals that ChemBERTa learns chemically meaningful patterns, with different attention heads specializing in functional groups, ring structures, or SMILES syntax \citep{Chithrananda2020}. ChemBERTa-2, an optimized variant pre-trained in 77 million compounds, achieved even stronger performance than the original ChemBERTa in MoleculeNet benchmarks \citep{Ahmad2022}.

Neural fingerprints provide several advantages over traditional fingerprints. They produce dense continuous vectors that avoid sparseness and information loss from bit collisions \citep{Jaeger2018}. They can also be pre-trained once on large unlabeled datasets and then fine-tuned for specific downstream tasks, opening the door to transfer learning \citep{Chithrananda2020, Fabian2020}. However, these methods are often computationally expensive \citep{Jaeger2018} and due to their black-box nature, they are not as interpretable.  \citep{Chithrananda2020}. However, attention mechanisms can have some degree of interpretability through visualization \citep{Chithrananda2020}.

Beyond fingerprints, physicochemical-descriptor vectors (PDVs) are a class of continuous-value molecular descriptors that encode structural and chemical properties of molecules \citep{Todeschini2009}. Unlike binary fingerprints or graph-based representations, PDVs aggregate atom and bond-level quantities into fixed-length feature vectors that capture a wide variety of information about a molecule and its structure \citep{Mauri2006}. PDVs typically consist of hundreds of properties including constitutional descriptors such as molecular weight, atom and bond counts, and heteroatom ratios; topological indices encoding molecular connectivity patterns; physicochemical properties such as logP (lipophilicity), molar refractivity, and topological polar surface area; and electronic descriptors derived from partial charges and polarizabilities \citep{Todeschini2009, Cherkasov2014}. The \texttt{RDKit} cheminformatics toolkit provides a comprehensive suite of more than 200 molecular descriptors that cover these categories \citep{rdkit}. PDVs are very interpretable; each descriptor has a clear physicochemical meaning \citep{Cherkasov2014}. They are also less computationally expensive than some 3D methods or neural fingerprints, with many descriptors simply requiring a 2D molecular structure without conformer generation or extensive pre-training. This efficiency makes PDVs particularly appropriate for high-throughput applications. However, it is also possible to use descriptors based on 3D components \citep{Mauri2006}.

However, PDVs have limitations. The dimensionality of PDV feature spaces can lead to problems with multicollinearity \citep{Todeschini2009}. This can be mitigated by feature selection or dimensionality reduction techniques \citep{Cherkasov2014}. Additionally, PDVs encode molecular properties as aggregated global features, which may not capture fine-grained structural details such as specific substructure arrangements or local chemical environments, information important for particular endpoints \citep{Mauri2006}.

One of the more intuitive representations are molecular graphs. These graphs can be categorized into 2D and 3D graphs. 2D graphs rely on adjacency matrices that encode topological atom- and bond-level features but leave out geometric properties such as bond angles or torsion, meaning they may fail to capture stereochemistry. 3D graphs incorporate these spatial features, but are often limited by remaining in a single conformation \citep{Zhichun2023}. A standard 3D molecular graph is constructed such that each molecule is represented as $G = (V,E)$ with node features $\mathbf{x}_i$ and edge features $\mathbf{e}_{ij}$. Nodes encode atom type, degree, hybridization, aromaticity, and formal charge; edges encode bond order and conjugation \citep{Zhichun2023}. 

% TODO: review, needs manual review
Another common molecular graph is the standard distance compatibility molecular graph, in which atoms and bonds are mapped to sets of nodes and edges \citep{David2020, Hnig2022}. These molecular graphs are typically 2D objects that represent 3D information such as atomic coordinate, bond angles, chirality, etc. \citep{David2020}. However, there are alternative molecular graph representations such as graphs with quantum 3D embeddings, weight alignment graphs, association graphs containing pairs of molecules, and geometric interaction graph neural network that incorporate 3D structures and physical interactions \citep{Hnig2022, yan2023, Yang2023GIGN}. \cite{Du2022} developed hybrid 2D and 3D graphs, combining 3D data, guaranteeing precision and completeness by providing unabridged information, and 2D data, providing chemical intuitions for guidance, in a zipper-like arrangement. These graphs pose various challenges, such as the atoms being localized at voxels creating issues with the points being too sparse, leading to poor performance. \cite{Kuzminykh2018} developed a novel approach in which atoms are extended to fill nearby voxels with a wave transformation.
      
These graph representations do not highlight specific substructures or specific arrangements of atoms and bonds within each molecule. The substructures contain functional groups that are related to a range of chemical properties and are potentially involved in downstream tasks related to the function of molecules. \cite{Bouritsas2023} developed graph substructure networks, a topologically aware message passing scheme based on substructure encoding. The method involves encoding structural roles based on the appearance of specific substructures and using them in message passing updates. Since message-passing GNNs are at most as powerful as the Weissfeiler Leman test, they cannot isolate sufficient features from the graph structure \citep{Bouritsas2023}. Graph substructure networks, on the other hand, capture these substructure-related features. And recent research has highlighted the abilities of pretrained GNN encoders produce embeddings on unlabeled datasets that capture structural features \citep{kishimoto2023, ross2022}. 

% TODO: review, needs manual review
\cite{Kolmar2021} examined the predictive limits of QSAR in the presence of noise, showing that by learning the underlying trends rather than fitting to noise, the models can achieve predictions more accurate than those of their training data. They assumed experimental noise could be modeled as Gaussian-distributed, using the Central Limit Theorem to argue that aggregated experimental measurements from multiple independent error sources exhibit a Gaussian distribution. However, \cite{Heid2023} used mean-variance estimation and bias-variance decomposition to show that noise can be tied to specific modalities, for example structure dependence. As an example, nitrogen-containing functional groups introduce measurement deviations, implying the presence of systematic noise. 

\section{Models for molecular property prediction}
% ---- NOTES: Models for molecular property prediction ----
% Takeaway: model families differ in how far one wrong label spreads. A tree keeps it in one leaf, an ensemble dilutes it, a network can memorise it.
% Paragraph 1: models treat labels as true values.
% - The evidence (too few repeat measurements per molecule, Kolmar2021) is in the Chapter 1 introduction. Point there.
% Paragraphs 2-5: trees and tree ensembles.
% - A noisy label only affects predictions for molecules in its own leaf.
% - Bootstrapping and averaging over many trees dilute any one noisy label.
% Paragraphs 6-7: feed-forward neural networks.
% - Change: backpropagation is cited to Hornik1989, which proves universal approximation. Cite Rumelhart1986.
% - Change: "it's layers" should be "its layers".
% - Change: wu2018 does not state the "2-3 hidden layers with 128-512 units" recommendation.
% Paragraph 8: why networks struggle on QSAR data and under label noise.
% - Change: "don't rely on a static 3D molecular structure like NNs do" is wrong for the networks in this thesis. They take fixed-length vectors, never 3D coordinates.
% - Networks can fit randomly assigned labels perfectly \citep{Zhang2017}. So they can also fit noisy labels.
% - Change: the last sentence cites Zhang2017 for robust losses, mixup and label smoothing. Cite robust losses \citep{Ghosh2017}, mixup \citep{Zhang2018mixup}, label smoothing \citep{Szegedy2015}. Northcutt2021 for confident learning is correct.
% - All four are LNL methods. Label smoothing and removing flagged examples change the labels, so they are noise mitigation.
% Paragraph 9: graph neural networks.
% - Change: add Kipf2017 for GCNs. "Convolutional Networks (GCNs)" is missing "Graph". Write "activity cliffs" in lower case.
% - No graph network is trained as a predictor in this thesis. MHG-GNN, a GIN-based autoencoder, is used only as a fixed-length embedding \citep{kishimoto2023}.
% Paragraph 10: which chapters use which families, one sentence.
% - Chapter 3: tree ensembles, SVMs, Gaussian processes and feed-forward networks, with probabilistic versions. Chapters 5 and 6: tree ensembles, an SVM, a feed-forward network, and a Gaussian process on datasets under 1,000 molecules.
% ---- END NOTES ----

One key issue with QSAR models is the standard assumption that labels represent true values. Most ML models rely on the assumption that the training data is, in fact, an accurate reflection of the population. Therefore, they treat these potentially noisy labels as discrete quantities rather than distributions, which opens the door to overfitting to noise and giving unduly optimistic assessments of the models.

There are a wide variety of model architectures that can be used for QSAR modeling. Starting off with simpler and more explainable models, we have tree-based methods. Decision trees are non-parametric supervised learning algorithms that perform recursive binary partitioning of the feature space \citep{Breiman1984}. A decision tree builds a predictive model by learning a series of if-then rules that recursively split the data based on given feature values, creating a tree  where internal nodes represent decision rules, branches represent the results of those decisions, and leaf nodes contain predicted values or class labels \citep{Breiman1984}. Each recursive split is chosen to maximize  homogeneity within the resulting partitions. For classification, Gini impurity is commonly used and for regression, mean squared error reduction \citep{Breiman1984}.

There are several properties that make decision trees a good foundation for QSAR modeling. First, they handle non-linear relationships between molecular descriptors and target properties without requiring explicit feature engineering or transformation \citep{Breiman1984}. They also do not make assumptions about the underlying distribution of the data, making them robust to violations of parametric assumptions that can limit linear regression or other classical methods \citep{Svetnik2003}. 

Decision trees handle data in a "contributory" rather than "consensus" manner. In consensus methods such as linear regression or neural networks, every data point influences the prediction at every other point in the feature space. Tree-based methods make predictions based on the subset of training samples that fall into the same leaf node as the test point \citep{Breiman1984}. This mechanism essentially quarantines outliers in their respective leaf nodes, limiting their impact on predictions for samples in other regions or chemical domains. 

Individual decision trees tend to have high variance and often overfit, particularly at large depths \citep{Breiman1984}. Small modifications made to the training data can result in large differences in the structure of the tree, resulting in unstable predictions \citep{Breiman2001}. Ensemble methods are used to counteract this instability, aggregating predictions from many trees. Random forests (RFs) aggregate predictions from trees trained on bootstrap samples of data with random subsets of features at each split \citep{Breiman2001}. Gradient boosting methods such as XGBoost sequentially fit trees to the residuals of previous iterations, progressively reducing the prediction error \citep{Chen2016}. Bootstrapping ensures that each tree is trained on a slightly different subset of data, which effectively reduce the influence of any individual sample \citep{Breiman2001}. This is especially beneficial with noisy samples. The selection of random features at every node decorrelates the trees, preventing systematic biases from propagating through the ensemble \citep{Breiman2001}. The aggregation of trees further prevents overfitting to noise by smoothing out prediction errors from individual noisy labels. 

Neural networks (NNs), a very popular family of models, can also be used for QSAR. A NN is a system of interconnected nodes which learn by passing information through it's layers. There are weights and biases attached to the nodes and various other parameters which the model optimizes to minimize the difference between the predicted output and the actual target value. During training, the network parameters are updated using backpropagation, in which gradients of the loss function with respect to each weight are propagated backward through the network \citep{Hornik1989}. The structure of these networks and the hyperparameters used depend on the type of NN. In QSAR, one of the most commonly used architectures is the multilayer perceptron (MLP), a NN containing fully-connected layers with non-linear activation functions which transforms molecular feature vectors (i.e., fingerprints, descriptors) into a property or activity prediction \citep{Ma2015}.

For molecular property prediction, it is recommended that architectures use 2-3 hidden layers with 128-512 units each, ReLU activations, and dropout regularization to prevent overfitting. Training is performed using backpropagation with gradient descent optimizers such as Adam, to minimize a loss function (MSE for regression, cross-entropy for classification) on batches of training data \citep{wu2018}. 

However, NN performance on QSAR-related tasks has been very mixed. They are often outperformed by more traditional methods: decision trees and related ensemble models tend to perform stronger on smaller datasets, overfit less, and don't rely on a static 3D molecular structure like NNs do \citep{Baskin2008, Koutsoukas2017}. Their complexity and use of gradients make them prone to label memorization \citep{Zhang2017}, which in the context of noise is problematic, as it means that models may memorize the noise themselves. Although models that are memorizing labels may initially appear to have strong predictive performance, they cannot generalize out of distribution samples properly and often provide inaccurate uncertainty estimates. Various approaches have been proposed to improve neural network robustness to label noise, including noise-robust loss functions \citep{Zhang2017}, confident learning to identify mislabelled examples \citep{Northcutt2021}, and regularization techniques such as mixup or label smoothing \citep{Zhang2017}. 

All the models up to this point have assumed a vector-based input. Molecular graphs, on the other hand, require different types of predict modeling. Two standard ML methods on molecular graphs include Convolutional Networks (GCNs) and Graph Isomorphism Networks (GINs). Graph Convolutional Networks (GCNs) apply the concept of convolution to graph-structured data. Instead of operating on grid-based data, such as images, GCNs aggregate information from neighbors of a node to update its representation \citep{TeachOpenCADD2023, Xu2019}. In contrast, Graph Isomorphism Networks (GINs) address the limitations in GCNs’ ability to differentiate between non-isomorphic graphs. Unlike GCNs, which normalize and weight the sum of neighboring node features, GINs use a simple sum followed by a trainable MLP, giving the model flexibility to learn complex transformations of node features. The GIN model’s design is informed by the Weisfeiler-Lehman test for graph isomorphism, enabling it to distinguish between molecular graphs that are structurally similar but differ in subtle ways \citep{TeachOpenCADD2023, Xu2019}. This makes GINs particularly suited to Activity Cliffs (ACs), where small variations in molecular structure can lead to significant changes in chemical or biological activity \citep{TeachOpenCADD2023}. 


\subsection{Tree ensembles}
% ---- NOTES: Tree ensembles ----
% Takeaway: quantile regression forests and NGBoost turn tree ensembles into models that report their own uncertainty.
% Takeaway: boosting may keep fitting a wrong label, because its residual stays large.
% Paragraph 1: the quantile regression forest.
% - Same trees as a random forest, but each leaf keeps all its training labels, not only their mean. Quantiles come from those labels, weighted over all trees \citep{Meinshausen2006}.
% - A wrong label stays in its leaf and widens that leaf's spread. Chapter 3 tests whether the interval tracks injected noise.
% Paragraph 2: gradient boosting, LightGBM and NGBoost.
% - Argument only: a wrong label keeps a large residual, so each boosting round keeps fitting it. LightGBM's one-side sampling keeps large-gradient molecules, so it never drops them \citep{ke2017lightgbm}.
% - Evidence: 12 algorithms on 12 bioactivity datasets, with Gaussian noise added to training pIC50 values. Gradient boosting had low noise tolerance; the other 11 degraded smoothly \citep{Cortes2015}.
% - NGBoost predicts a distribution per molecule, for example a mean and a variance \citep{Duan2020}. A molecule with a large residual can get a larger variance instead of a shifted mean.
% ---- END NOTES ----


% TODO: review, needs manual review
RFs are one of the most common choices for QSAR modeling thanks to their robustness and interpretability \citep{Svetnik2003}. Quantile Regression Forests (QRFs) hold the distribution of training labels in each leaf and compute quantiles over all trees \citep{Meinshausen2006}.

% TODO: review, needs manual review
Gradient-boosted tree ensembles include eXtreme Gradient Boosting (XGBoost) \citep{Mustapha2016, Tian2022}, Light Gradient-Boosting Machine (LightGBM) \citep{ke2017lightgbm}, and Natural Gradient Boosting (NGBoost) \citep{Duan2020}. NGBoost extends deterministic gradient boosting by treating the parameters of a chosen parametric distribution as regression targets and learns them through boosting with a natural gradient update rule \citep{Duan2020}.


\subsection{Support vector machines}
% ---- NOTES: Support vector machines ----
% Takeaway: support vector regression ignores errors inside a tube and caps each molecule's pull, so it may resist label noise.
% Paragraph 1: what a support vector machine is.
% - Two sentences at most. Support vector regression ignores training errors smaller than epsilon and caps each molecule's weight at a constant C \citep{Vapnik1995, Smola2004}.
% Paragraph 2: how it uses the features.
% - Features enter only through the kernel. With the RBF kernel every feature enters every Euclidean distance, so uninformative features blur it.
% Paragraph 3: why it might resist label noise (answers your TODO).
% - Inside the tube a label error changes nothing. Outside it the loss grows linearly, not quadratically. Each molecule's weight is capped at C \citep{Smola2004}.
% - In the 12-algorithm study the RBF-kernel SVM degraded smoothly and linearly \citep{Cortes2015}.
% - Chapter 3 reports how the SVM holds up under each noise condition. Point there; do not restate the numbers.
% Paragraph 4: the limitation on large datasets (answers your "Limitations" note).
% - The kernel matrix has one row and one column per training molecule. It grows with the square of the number of molecules \citep{Chang2011}.
% ---- END NOTES ----

Support vector machines (SVMs) are a well-established baseline in QSAR modeling \citep{Vapnik1995, Svetnik2003}.
% TODO: anything interesting to say on SVMs? Background? promising results with noise (it's actually one of the top performing in my results)
% Limitations - compute too high on large N

\subsection{Gaussian processes}
% ---- NOTES: Gaussian processes ----
% Takeaway: a GP sees molecules only through a kernel. A standard GP has one noise variance for the whole dataset.
% Paragraph 1: what a GP is and what limits it.
% - Change: the first sentence is missing its verb: "Gaussian processes (GPs) are a non-parametric class of models ...".
% Paragraph 2: the two kernels.
% - Define $\sigma_f^2$ once: the signal variance, a learned scale on the kernel. Both equations use it without naming it.
% - The Tanimoto kernel takes binary vectors. Continuous features, such as the physicochemical descriptor vector, use the RBF kernel.
% Paragraph 3 (new, short): why a GP might absorb label noise.
% - The likelihood adds one observation-noise variance to the diagonal of the kernel matrix \citep{Rasmussen2005}. It cannot absorb one molecule's large error.
% - A heteroscedastic GP learns a noise variance per molecule \citep{le2005heteroscedastic, kersting2007most}.
% - The two halves of the thesis use different GPs. Point to the Chapter 3 methods and to "Models and metrics" in Chapter 5.
% ---- END NOTES ----

Gaussian processes (GPs) a non-parametric class of models that use a kernel to produce a Gaussian predictive distribution over every data point \citep{Obrezanova2007, Rasmussen2005}. One of the major limitations of using GPs in practice is computational scaling; inference has $O(N^3)$ time complexity and $O(N^2)$ memory, where $N$ is the number of training points \citep{Rasmussen2005}. This can be mitigated with sparse GP approximations, but these introduce approximation error and can be more difficult to tune \citep{quinonero2005}.


% TODO: review, needs manual review
The \texttt{Gauche} framework \citep{gauche} defines kernels for strings, fingerprints, and graphs that operate on a range of widely used molecular representations. The Tanimoto kernel \citep{Ralaivola2005, moss2020} is defined for binary vectors $\mathbf{x}, \mathbf{x}' \in \{0, 1\}^d$ for $d \geq 1$ as:

$$
k_{\text{Tanimoto}}(\mathbf{x}, \mathbf{x}') := \sigma_f^2 \cdot \frac{\langle \mathbf{x}, \mathbf{x}' \rangle}{\|\mathbf{x}\|^2 + \|\mathbf{x}'\|^2 - \langle \mathbf{x}, \mathbf{x}' \rangle},
$$

where $\|\cdot\|$ represents the Euclidean norm \citep{gauche}. 

% TODO: review, needs manual review
The radial basis function (RBF) kernel is defined for
$\mathbf{x}, \mathbf{x}' \in \mathbb{R}^d$ as:

$$
k_{\text{RBF}}(\mathbf{x}, \mathbf{x}') := \sigma_f^2 \cdot \exp\left(-\frac{\|\mathbf{x} - \mathbf{x}'\|^2}{2\ell^2}\right),
$$

where $\ell$ is the lengthscale.


\subsection{Bayesian neural networks}
% ---- NOTES: Bayesian neural networks ----
% Takeaway: a network can be made to report its own uncertainty, at different costs.
% Takeaway: aleatoric and epistemic uncertainty are defined here, once, for the whole thesis.
% Paragraph 1: the full Bayesian network (Figure \ref{fig:bayesian_transformation}).
% - Each weight gets a probability distribution instead of one value, fitted by a variational approximation \citep{Neal1995, Blundell2015}. The spread over many forward passes is the uncertainty.
% Paragraph 2: cheaper versions.
% - Variance head: a second output, a variance per molecule \citep{kendall2017}. Variational Bayesian last layer (VBLL): only the final layer gets a weight distribution \citep{Harrison2024}.
% - Which uncertainty each version reports, and its training loss, belongs in the Chapter 3 methods.
% Paragraph 3: other ways of getting uncertainty.
% - MC dropout: dropout stays on at prediction time \citep{gal2016}. Deep ensembles: networks from different random starts, their disagreement used as uncertainty \citep{lakshminarayanan2017}.
% Paragraph 4: aleatoric against epistemic uncertainty.
% - Answer to your TODO: keep the paragraph. It is the only definition of the two parts in the thesis, and Chapter 3 uses both.
% - Epistemic: uncertainty about the model's fit. It shrinks with more training data, most where few similar molecules were seen \citep{kendall2017}.
% - Aleatoric: noise in the labels themselves, such as assay error. More data from the same assay does not reduce it \citep{kendall2017}.
% - Change: "higher quality training data" is ambiguous. More precise labels lower the aleatoric part. Wider coverage of chemical space lowers the epistemic part.
% - The total predictive variance is the sum of the two. The split also applies to Gaussian processes, quantile regression forests and NGBoost.
% ---- END NOTES ----

NNs are often deterministic; however, they can be transformed into probabilistic models.

\begin{figure}[htbp]
    \centering
    \includegraphics[width=0.85\textwidth]{bayesian_transformation.png}
    \caption{Deterministic versus Bayesian neural network architectures. Deterministic networks use fixed weights $W$ to produce point predictions, while Bayesian networks treat weights as distributions $\mathcal{N}(\mu, \sigma^2)$, generating prediction histograms through multiple stochastic forward passes to quantify epistemic uncertainty.}
    \label{fig:bayesian_transformation}
\end{figure}

Unlike standard deterministic ML methods, Bayesian methods model distributions over predictions using Bayes' rule, allowing for the quantification of uncertainty within parameters and labels. Bayesian methods treat the model parameters as random variables with distributions. Uncertainty can also be estimated from other non-deterministic methods, including deep ensembles \citep{lakshminarayanan2017}, bootstrap modeling \citep{palmer2022}, and federated learning \citep{Hanser2023}.

% TODO: this paragraph may need to go
Epistemic uncertainty is model-driven and can be reduced by collecting either more or higher quality training data. Aleatoric is data-driven and reflects the intrinsic noise in the labels.

\section{Molecular similarity and activity cliffs}
% ---- NOTES: Molecular similarity and activity cliffs ----
% Takeaway: similarity only exists relative to a representation and a coefficient; the representation is the larger choice.
% Takeaway: activity cliffs and uninformative features break the similar property principle, and neighbour-based noise detectors inherit the break.
% Paragraph 1: the similar property principle.
% - Numbered definition: structurally similar molecules tend to have similar properties \citep{Johnson1990, vanTilborg2022}.
% - This thesis relies on it in nearest-neighbour prediction, the Tanimoto-kernel GP \citep{Ralaivola2005, moss2020}, and noise detectors that compare a label with its neighbours' (Chapter 6). State no result here.
% Paragraph 2: similarity depends on the measure.
% - A similarity value is fixed by two choices: the representation and the coefficient \citep{Maggiora2014}.
% - \citet{Bajusz2015}: Tanimoto, Dice, Cosine and Soergel were the best coefficients on fingerprints, and roughly equivalent.
% - Change: "changing the similarity formula yields minimal improvement if the underlying representation is weak" claims more than Bajusz2015 found. They did not test weak representations.
% - That the representation matters more than the coefficient is this thesis's own argument, not a cited result.
% - Change: "Distance-based reasoning rely" should be "relies". "\citet{Bajusz2015} et al." prints "et al." twice, and the sentence cites Bajusz2015 a second time.
% Paragraph 3: activity cliffs.
% - Numbered definition: a pair of structurally similar molecules with a large potency difference, from a small chemical change \citep{Maggiora2006, Stumpfe2012, Dablander2023}.
% - Change: the prose defines activity cliffs twice (opening paragraph and "The sharpest failure" paragraph). Keep one.
% - Change: drop "the benchmark used in this thesis". MoleculeACE's datasets are not thesis datasets. Cite it only for its cliff rule and cliff shares.
% - This thesis's own cliff rule (its similarity threshold and activity gap) belongs in the Chapter 6 methods, where it is used.
% Paragraph 4: why cliffs and label noise are hard to tell apart.
% - To a detector that compares a molecule with its neighbours, a cliff looks like label noise. Removing cliffs as noise removes the molecules with the most structure-activity information.
% Paragraph 5: high dimensionality.
% - Change: the prose credits the converging nearest and farthest distances to Aggarwal2001 alone. Cite \citet{Beyer1999} as well; Aggarwal2001 extend it to each $L_k$ norm.
% - That uninformative features dominate the distance is this thesis's argument. Chapter 4 measures it.
% - Answer to your TODO on the bib: the entries are in refs.bib. Still to fix: Dablander2023 lists "Michael Dablander"; it should be Markus Dablander.
% ---- END NOTES ----

Distance-based reasoning rely on the assumption that the data embeddings  contain quality information. When the labels are corrupted or the important features are scarce, similarity scores over-rely on compounding false or useless information, and the assumption that similar molecules have similar properties breaks down. \citet{Bajusz2015} et al. found that while the Tanimoto coefficient can be effective with binary fingerprints, changing the similarity formula yields minimal improvement if the underlying representation is weak \citep{Bajusz2015}.
This problem is further complicated by activity cliffs (ACs), which are defined as structurally similar molecules with large property differences, often due to a subtle chemical change \citep{Stumpfe2012, Dablander2023}. These molecules, which potentially contain the most important information for model prediction, can be confused with noise. Another barrier to distance-based reasoning is persistent noise across labels. The task of trying to determine if a molecule's value matches that of its neighbors becomes a lot harder when all of the neighbors have some form of noise as well. 


% TODO: hand edit
% TODO: review, needs manual review
Much of cheminformatics rests on the similar property principle: molecules that are structurally similar tend to have similar properties \citep{Johnson1990, vanTilborg2022}. 
One of the clearest applications of this principle is the use of nearest-neighbour prediction which shows up across cheminformatics workflows and in the case of this research, in the kernels used by GPs and at the basis of many noise identification algorithms.

% TODO: review, needs manual review
The difficulty with this assumption is that similarity is defined only relative to a similarity measure, and that measure is fixed by the choice of representation and
coefficient \citep{Maggiora2014}. There is no standard definition of similarity between molecules, especially when handling different molecular embeddings in a computational setting. Two molecules that are near-identical under a substructure fingerprint may be far apart under a physicochemical descriptor vector. \citet{Maggiora2014} caution that quantified similarity can be over-interpreted in this manner. 

There are many downstream tasks which rely on the concept of similarity. The second half of this thesis relies on many such tasks, and as such focuses on identifying representations that agree with this principle. 


% TODO: review, needs manual review
The sharpest failure of the similar property principle is the activity cliff: a pair of structurally similar molecules with a large difference in potency, produced by a small
chemical change \citep{Maggiora2006, Stumpfe2012}. Activity cliffs inherit the measure-dependence of the principle they violate, and there is no widely accepted definition of one
\citep{vanTilborg2022}. The MoleculeACE benchmark used in this thesis makes the choice explicit: a pair is ``highly similar'' if it exceeds 90\% similarity under at least one of
three measures --- Tanimoto on extended connectivity fingerprints, Tanimoto on scaffold fingerprints, or a scaled Levenshtein distance on SMILES strings --- and forms a cliff if
it also differs by tenfold or more in $K_i$ or EC$_{50}$ \citep{vanTilborg2022}. The soft consensus over three measures is used precisely because each captures only part of the
underlying chemistry. Cliff compounds account for between 7\% and 52\% of a dataset across the thirty targets.

% TODO: review, needs manual review
A further difficulty is dimensionality itself. \citet{Aggarwal2001} show that as the number of dimensions grows, the distance to the nearest point and the distance to the furthest
point become increasingly similar, leaving pairwise distances progressively less discriminative. Molecular representations are routinely high-dimensional and sparse, and when a
large proportion of their dimensions carry no information relevant to the property being modelled, distances computed over them are dominated by those uninformative dimensions.


% TODO: make changes to bib

\section{Noise in molecular data}
% ---- NOTES: Noise in molecular data ----
% Takeaway: real label noise comes from many sources and is not independent, identically distributed or Gaussian.
% Takeaway: label noise is not always harmful; it can act as regularisation.
% Paragraph 3: independence and zero mean fail too.
% - Change: the last sentence ("In this study, we tested noise conditions...") should point to Section~\ref{sec:noise_injection} and Chapter~\ref{chap:robustness}.
% Paragraph 5: instruments and detection limits.
% - Change: the sentence on quenching giving false negatives for potent inhibitors does not match Simeonov2008. Rewrite to its finding, below.
% - \citet{Simeonov2008}: more than 5\% of over 70,000 library samples fluoresced themselves, causing both false positives and false negatives.
% - \citet{Zhang1999} defines the Z-factor, which scores an assay by its signal range and variation. It does not cover detection-limit or radiometric error.
% - Change: "systematic bia" should be "systematic bias".
% Paragraph 7: aggregation across databases.
% - Add the current ChEMBL release paper \citep{Zdrazil2023} next to Gaulton2017.
% - Change: "between 10 and 20 reported values... two orders of magnitude" is not what \citet{Kramer2012} reports. Its repeat ChEMBL Ki pairs differ by 0.44 pKi units on average.
% - \citet{landrum2024}, IC$_{50}$ pairs: 27\% differ by more than one log unit with minimal curation, 13\% with assay metadata matched. Ki assays were similar.
% Paragraph 8: federated learning.
% - \citet{Hanser2025} describe federated learning by knowledge distillation, the teacher-student set-up of Effiris. Noisy, differing clients are a known problem \citep{Fang2022}.
% - Change: "presents" should be "present"; "confidental" should be "confidential".
% Paragraph 9: annotation and structure errors.
% - Change: the prose says 1-5\% structure errors. \citet{Young2008} found 0.1\% to 3.4\% across six databases, with no breakdown by cause. Drop the list of causes or cite another source.
% - Cite the structure-curation workflow \citep{Fourches2010}.
% - Protonation states and counterions: point to Chapter 4, where molecules the representation cannot tell apart are counted.
% Paragraph 10: real noise is not i.i.d. or symmetric.
% - Survey of label-noise sources and consequences: \citet{Frenay2014}.
% - One sentence before the Heid2023 sentence: \citet{Kolmar2021} modelled experimental noise as Gaussian, arguing from the central limit theorem.
% - Chapter 6 also measures noise on real labels: repeat measurements, disagreement between laboratories, curator flags.
% Paragraph 11: label noise can regularise.
% - Change: "noise analysis and mitigation literature" should be "the label-noise literature, including learning with noisy labels". Noise mitigation is one subset of LNL.
% ---- END NOTES ----

% TODO: review, needs manual review
Real-world cheminformatics datasets contain label noise arising from multiple sources across the drug discovery pipeline. To properly model experimental noise, we must first understand where the noise is coming from.

Studies involving noise, and specifically the quantification of noise robustness, typically approach label noise by adding independent, zero-mean, Gaussian error to the labels. However, there's evidence to suggest that noise from experimental assays takes on different forms. \cite{Kruger2012} rejected normality of bioactivity differences by Anderson-Darling test and fitted a Laplace distribution instead. \citet{Kalliokoski2013} examined the widest disagreements among 16,844 pairs of repeat pIC$_{50}$ measurements and found that in almost all cases, the wide gap between repeat values was due to annotation error rather than to an imprecise assay.

Beyond just the shape, the assumptions of independence and zero mean may fail as well. \citet{Bentz2013} examined P-glycoprotein transport measurements across 23 laboratories and found that $62\%$ of the variance in log efflux ratio was due to differences between laboratories. When you pool data from different sources, each source has the potential to carry a shared offset. These sources usually contain families of related compounds, breaking the independence assumption. The zero-mean assumption is not assured either, especially when true values fall above or below an assay's detection limit, called censoring. \citet{Svensson2025} count censored labels in fifteen industrial assays, of which thirteen carry censored labels and eight sit between a quarter and two thirds censored. In this study, we tested noise conditions that vary the shape, the correlation, and the bias of injected error separately at a matched quantity.  

Assay-related variability represents a primary source of label noise in bioactivity data. High-throughput screening assays are often laden with errors from physical phenomena such as edge effects, where compounds in wells exhibit different activities at plate boundaries due to evaporation or temperature gradients \citep{Macarron2011}. Other examples include compound carryover between wells, incomplete liquid transfers, and air bubbles  \citep{Inglese2006}.  Detection reagents or unstable compounds can degrade over time, which leads to drift within a single plate run. This turns into heteroscedastic noise correlated with plate position or read time \citep{Brideau2003}.

Instrument calibration and detection limits can impact measurement accuracy. Spectrophotometric assays near detection limits exhibit higher relative error, such that signal-to-noise ratios degrade for weak binders \citep{Zhang1999}. Fluorescence quenching by test compounds introduces a systematic bia in fluorescence-based assays, resulting in false negatives for potent inhibitors that accidentally get absorbed at emission wavelengths \citep{Simeonov2008}. Radiometric assays require background subtraction and decay corrections that accumulate error, particularly for low-activity compounds where the signal barely exceeds background \citep{Zhang1999}.

Another source of noise is biological variability. Even minor variations in cell culture conditions such as seeding density, plating efficiency, growth rate, and cell cycle distribution have the potential to affect downstream tasks such as cell metabolism and drug responsiveness. Even serum properties may introduce variability \citep{Hatzis2014}. The speed of cell culture passages is even found to be a significant factor in behavior \citep{Postnikova2018}. 

Data aggregation across multiple sources compounds these individual error sources. Common databases like ChEMBL, BindingDB, and PubChem aggregate measurements from hundreds of laboratories using a diverse set of protocols, compound purities, and quality control standards \citep{Gaulton2017, Gilson2016}. The $K_i$ or IC$_{50}$ values from the same compound-target pair with between 10 and 20 reported values, often end up spanning two orders of magnitude apart, making it challenging to identify the underlying ground truth \citep{Kramer2012}. All the sources of error mentioned above regarding inconsistent experimental protocols cannot be distinguished easily from the genuine underlying value without extensive, and often unrealistic amounts of metadata \citep{Fourches2016}.


Federated datasets, a method of data aggregation centered around privacy protection, presents a particularly challenging situation for noise. In pharmaceutical research and beyond, companies often want to merge collective information without exposing proprietary data to competitors. One example of such platforms is Lhasa Limited's federated learning platform Effiris. Teacher models are first trained separately on confidental data, and then used to label a shared public chemical data set. Those results are consolidated into a federated student model \citep{Bassani2023}. While this is an effective approach of protecting intellectual property, it introduces additional label noise through the prediction errors of all the teacher models. There also exists aleatoric uncertainty from the underlying experimental measurements, many of which likely come from differing protocols, resulting in compounding label noise. 

Beyond the experimental setups, annotation errors in the data itself can arise from mis-registered structures, incorrect stereochemistry assignments, and salt form ambiguity \citep{Williams2017}. Public databases are estimated to contain 1-5\% structure errors from incorrect SMILES parsing, tautomer misassignment, or manual transcription mistakes during data entry \citep{Young2008}. Ambiguous protonation states and counterions create representation-level noise where chemically distinct samples are classified as the same \citep{Bento2020}.

These noise sources can interact with each other and present in more ways than a simple Gaussian. Unlike artificial noise injection, real-world noise is neither independent, identically distributed, nor symmetric  \citep{Heid2023}. I wanted to address this complexity in my research by examining a variety of ways of both modeling and adding noise to data sets. 

Although the general assumption in many of the noise analysis and mitigation literature is that noise, specifically label noise, will have a negative effect on the model, this is not necessarily the case. There are multiple instances of using noise as a method of regularization. Although not the optimal approach to counteracting overfitting in linear regression, adding artificial noise to labels is considered a form of regularization \citep{Ullah2024}. Furthermore, adding artificial noise drawn from a Gaussian distribution to the activation layers of a neural network is considered another form of regularization \citep{Camuto2021}. These studies indicate that the presence of noise may act as a form of regularization in some contexts. It is important to address this potential of regularization from label noise in QSAR and QSPR models.

\subsection{Artificial noise injection}\label{sec:noise_injection}
% ---- NOTES: Artificial noise injection ----
% Takeaway: the seven conditions deliver the same expected amount of noise, so a difference between conditions is a difference in pattern, not amount.
% Takeaway: each parameter comes from a published measurement error where one exists.
% Paragraph 1: what gets noise.
% - Change: the first sentence says validation data stay free of noise; the next says they are noised. Drop "validation and" from the first. Test labels stay clean.
% Paragraph 2: the seven conditions.
% - Change: use "grouped-shifted" throughout. Paragraph 2 writes "Group shifted".
% - Change: "on specific subset" should be "on a specific subset"; "every label extends past" should be "every label that extends past".
% Paragraph 5: the parameters and where each comes from.
% - Change: "62\% of the variance in Caco-2 across 23 laboratories" does not match Bentz2013. Only 11 of the 23 laboratories used Caco-2 cells.
% - The 62\% is the laboratory share of variance in log efflux ratio, cell-line data only, Bentz2013 Table 7. Word it as in "Noise in molecular data".
% - Answer to your TODO on the 20\%: no published estimate exists, so 20\% was chosen. Scaffold families are drawn in random order until about 20\% of the split is covered.
% - Censoring has no parameter here: its level is the fraction of labels clipped.
% ---- END NOTES ----

Our objective is to evaluate the robustness of QSAR models against label noise. To do so, we ``inject'' artificial noise into the labels during training while validation and test data remain free of noise. Validation labels are given noise of the same type and quantity as the training labels, drawn independently. This artificial noise simulates real-world experimental noise. The clean target label $y_i$ is replaced by a noisy label $\tilde{y}_i = y_i + \epsilon_i$, where $\epsilon_i$ is the artificial noise injected at the index $i$, the value of which is determined by a specified noise strategy. As we increased the quantity of noise, we measured how much predictive performance fell. 

We explored seven methodologies of adding noise, referred to as noise conditions (Figure~\ref{fig:noise_conditions}). Three vary the shape of the distribution from which the error is drawn. These shapes include Gaussian, Student-$t$ with $\nu = 5$, and Laplace distributions. Two conditions concentrate error on specific subset of molecules; outlier isolates a random tenth of them and grouped-wider concentrates on Bemis--Murcko scaffold families. Group shifted also applies noise to scaffolds, but instead gives each scaffold its own offset rather than a wider error. The seventh, censoring, sets a limit such that every label extends past that limit is reassigned to the limit itself. 

For a given label, every condition except censoring delivers the same expected quantity of noise at a set noise level. These levels are defined as a fraction of the standard deviation of clean training labels. The noise added to the label of molecule $i$ is
\begin{equation}
\epsilon_i = \frac{\tau}{G}\,s_i z_i,
\qquad
G = \sqrt{\frac{1}{n}\sum_{j=1}^{n} s_j^{2}}\;\;\sigma_z,
\label{eq:dose_matching}
\end{equation}
% TODO: review, needs manual review
Here $\tau$ is the target standard deviation of the noise, in the units of the label, and $z_i$ is
drawn from the condition's distribution at unit scale. The multiplier $s_i$ is $1$ for most
molecules and $\lambda = 3$ for those the condition affects, which are a random $10\%$ of molecules
under outlier and whole scaffold families covering $20\%$ of the split under grouped-wider. The
constant $\sigma_z$ is the standard deviation of $z$, equal to $1$ for the Gaussian,
$\sqrt{\nu/(\nu-2)}$ for the Student-$t$ and $\sqrt{2}$ for the Laplace. The factor $G$ divides both
of these out, so that $\epsilon$ has standard deviation $\tau$ in expectation whatever the
distribution and whichever molecules are affected.

% TODO: review, needs manual review
Grouped-shifted widens no molecule's error, and instead gives every scaffold family $g$ an offset of
its own,
\begin{equation}
\epsilon_i = \frac{\tau}{G}\left(\sqrt{\rho}\,b_{g(i)} + \sqrt{1-\rho}\,w_i\right),
\label{eq:grouped_shifted}
\end{equation}
with one draw $b$ per family and one draw $w$ per molecule. The two variances sum to $\tau^{2}$, and
the family offset carries a share $\rho$ of that total.

% TODO: No published figure exists for the fraction of scaffold families a difference between laboratories would cover, so we chose the $20\%$ above rather than sourcing it. - need to rephrase this 
We selected parameters from published measurement errors, wherever possible. As inter-laboratory error is estimated to be approximately three times intra-laboratory error, we set $\lambda = 3$ \citep{Avdeef2019}. \cite{Hampel2001} reports that $1-10\%$ of routine scientific data contains gross errors; we selected the upper bound of $10\%$. Laboratory identity accounts for $62\%$ of the variance in Caco-2 across 23 laboratories, prompting us to set $\rho = 0.62$ \citep{Bentz2013}. We took the Student-$t$ and Laplace shapes from the distributions fitted to repeated public bioactivity measurements \citep{Kruger2012}, and contaminated a fraction of the samples with a wider error according to Huber's contamination model \citep{huber1964robust}. As we found no published estimate of how large a share of a given dataset inter-laboratory differences affect, we chose to randomly select scaffolds until $20\%$ of the molecules had been covered. 

% TODO: review, needs manual review (caption)
\begin{figure}[htbp]
\centering
\includegraphics[width=\textwidth]{F1_noise_conditions.png}
\caption{The seven noise conditions applied to one set of labels. The labels are 2,000 synthetic
values rather than a study dataset, split at random into 40 groups that stand in for scaffold
families. a)--g) Clean labels in gray and noised labels in color, at a noise level of 0.5 of the
clean label spread. Censoring is drawn with 25\% of labels clipped, since its level is the fraction
clipped. The bottom axis is the label value and the side axis is the share of labels in each bin.
h) For each condition except censoring, the standard deviation of the mean error within a group,
divided by the standard deviation of all errors. A value near 1 means a group's labels move
together. With about 50 labels per group, noise that ignores the groups gives about 0.14.}
\label{fig:noise_conditions}
\end{figure}

% ---- NOTES: Learning with noisy labels (planned section, no heading yet) ----
% Takeaway: LNL covers every method for training a model when some labels are wrong. Noise mitigation is the subset that changes the labels before training.
% Takeaway: no single method wins, and methods that compare a molecule with its neighbours fail where similarity fails.
% Paragraph 1: what LNL is.
% - Noise mitigation changes, drops or reweights the training labels once, before training; any model can then be fitted. Examples: removal, down-weighting, smoothing, relabelling, correction.
% - Noise-robust training leaves the labels and changes how the model trains: robust losses, noise-modelling layers, regularisation, reweighting inside the training loop.
% - A sample-selection step used only to delete molecules before a refit counts as noise mitigation.
% - Noise detection and noise estimation find or size noisy labels. They feed noise mitigation but are not LNL methods themselves.
% - Taxonomy of 62 deep-learning methods in five groups \citep{Song2022}.
% Paragraph 2: why no single method wins.
% - No free lunch: averaged over all problems, no learning algorithm beats another \citep{Wolpert1996, Wolpert1997}. Chemistry analogue under simulated noise: \citep{Cortes2015}.
% - Neighbour-based cleaning goes back to editing a training set by its nearest neighbours \citep{Wilson1972}. It relies on similar molecules having similar labels, which activity cliffs break.
% - QSAR models frequently fail to predict activity cliffs \citep{Dablander2023}. Chapter 6 tests this.
% ---- END NOTES ----

\section{NoiseInject: a benchmarking framework}
% ---- NOTES: NoiseInject: a benchmarking framework ----
% Takeaway: NoiseInject is introduced here, before Chapters 3 and 6, because both use it.
% - NoiseInject 1.0.0 is an open-source Python package, MIT licence \citep{noiseinject}.
% - Change: add that the label-noise experiments on the hybrid in Chapter 6 use it too.
% - QM9 is the exception: its noise comes from a separate Rust implementation of the same seven conditions. Point to Chapter 3 for that pipeline.
% ---- END NOTES ----

% TODO: review, needs manual review
The noise conditions described above are implemented twice. The three experimental datasets and
every uncertainty result use NoiseInject, an open-source Python package. QM9 uses a separate
Python--Rust precompute pipeline, which carries its own Rust implementation of the same conditions.
An executable cross-check compares the two implementations before any cluster run. Rust's and
NumPy's random number generators produce different streams, so the check compares statistics rather
than individual draws: on the same labels, groups and seeds, it compares the delivered noise, the
unit dose $G$, the median error, the number of badly corrupted labels and the number of molecules
affected.

\subsection{NoiseInject}
% ---- NOTES: NoiseInject ----
% Takeaway: the package injects each noise condition at a matched amount and reports how much performance and uncertainty quality survive.
% - Change: "In the KIRBy experiments" should be "In the experiments with the hybrid". KIRBy names only the repository.
% - The main text needs the conditions, the robustness and uncertainty metrics, and the retention summary. The rest belongs in an Additional file on NoiseInject: the check between the Python and Rust implementations, the registry's extra settings, the dose-tolerance check (Equation~\ref{eq:dose_tolerance}), the scale-without-noise method and the residual-sized injector.
% - Classification: Chapter 6 states its levels as the fraction of training labels flipped to the wrong class (0 to 0.40). That fraction is the target the binary search hits; point forward to "Noise conditions" in Chapter 6.
% ---- END NOTES ----

% TODO: review, needs manual review
NoiseInject (version 1.0.0) evaluates model robustness and uncertainty calibration under artificial
label noise, for both regression and classification. It operates on label arrays and on the
predictions of any externally trained model, so it is model-agnostic. The core package depends only
on NumPy, pandas, SciPy and scikit-learn. Results are returned as \texttt{pandas} DataFrames, one
row per noise level and one aggregate summary row. The software is available under an MIT license
and archived on Zenodo \citep{noiseinject}.

% TODO: review, needs manual review
\paragraph{Regression noise.} The regression injector takes a noise shape (Gaussian, Student-$t$ or
Laplace) and a targeting rule (uniform, grouped-wider, grouped-shifted, outlier or censoring)
separately. Each named condition in the package's registry is one combination of the two with its
parameters, and the same name is written on every results row. The registry holds more settings
than this study runs: Student-$t$ at $\nu \in \{3, 5, 10\}$, and outlier contamination at $1\%$,
$5\%$ and $10\%$. It also holds an outlier variant in which the unselected molecules receive no
noise at all, rather than the base-scale noise of Equation~\ref{eq:dose_matching}.

% TODO: review, needs manual review
The requested level is the root-mean-square noise to deliver, $\tau$, in the label's own units. The
injector computes the per-molecule multipliers $s_i$ and solves for the scale $\tau / G$ in closed
form (Equation~\ref{eq:dose_matching}); there is no calibration search. A level of zero returns the
labels exactly unchanged. Which molecules are affected is drawn from a second random generator,
seeded separately from the noise draws and re-seeded on every call. The affected molecules are
therefore the same at every noise level, including zero, while the draws themselves change. Under
grouped-wider, whole scaffold families are added in random order until the fraction of affected
molecules is as close as possible to the target of $20\%$, because a fixed fraction of families
gives a very different fraction of molecules when family sizes are uneven.

% TODO: review, needs manual review
Censoring is parameterised by the fraction $f$ of labels clipped. The limit is the $(1-f)$ quantile
of the reference labels, which are the clean training labels. On the training split, the censored
set is exactly the $k = \mathrm{round}(fn)$ largest labels. On any other split it is whichever
molecules lie past the training limit.

% TODO: review, needs manual review
Each injection returns its provenance alongside the noisy labels: the target dose, the unit dose
$G$, the solved scale, the delivered root-mean-square noise (in label units and as a fraction of the
clean label standard deviation), the mean of $\epsilon$, the fraction of molecules affected, the
number and largest share of groups, and the seeds. The injector also checks the delivered amount.
The delivered dose is an estimate from a finite sample, so it is compared against a band of three
standard errors,
\begin{equation}
\left|\frac{\hat{\tau}}{\tau} - 1\right| \le \max\!\left(3\sqrt{\frac{\kappa - 1}{4\,n_{\mathrm{eff}}}},\; 0.005\right),
\qquad
n_{\mathrm{eff}} = \frac{\left(\sum_i s_i^2\right)^2}{\sum_i s_i^4},
\label{eq:dose_tolerance}
\end{equation}
where $\hat{\tau}$ is the delivered root-mean-square noise and $\kappa$ is the sample kurtosis of the
unscaled draws $z$, bounded between 3 and 60. The effective count $n_{\mathrm{eff}}$ is the number of
independent contributions the delivered dose averages over. It falls below $n$ when the noise is
concentrated on a few molecules. Under grouped-shifted it is
$n_{\mathrm{eff}} = \left(\rho^2/g_{\mathrm{eff}} + (1-\rho)^2/n\right)^{-1}$, with
$g_{\mathrm{eff}} = \left(\sum_g n_g\right)^2 / \sum_g n_g^2$ for family sizes $n_g$. At
$\nu \le 4$ the Student-$t$ has no fourth moment, so the standard error is instead
$1.3\,n_{\mathrm{eff}}^{\,2/\nu - 1}$, with the constant calibrated on the measured spread. A draw
outside the band raises a warning rather than an error. At three standard errors a correct injector
falls outside the band on about $1\%$ of draws, and the delivered amount is recorded on the row
either way.

% TODO: review, needs manual review
A second method returns the scale each molecule's region would receive without drawing any noise.
We use it to score held-out molecules against the pattern the training labels were exposed to. For
uniform and grouped-shifted noise this scale is the same for every molecule, so the package flags it
as degenerate. A rank correlation against a constant is undefined, and must not be reported as zero.

% TODO: review, needs manual review
A separate injector sizes noise against a unit the caller supplies rather than the label standard
deviation. In the KIRBy experiments the unit is the standard deviation of a model's out-of-fold residual on
the clean training labels, so the noise is sized against the error the model already makes. Its
conditions add Gaussian noise with variance equal to the level times the unit squared; move $2\%$,
$5\%$ or $10\%$ of labels by the level times the unit, with a random sign or always upwards; or add
Gaussian noise whose standard deviation follows either a per-molecule covariate or a hidden
log-normal draw. Both of the last two are scaled to the same mean variance as the first.

% TODO: review, needs manual review
\paragraph{Classification noise.} Six label-flipping strategies are provided: uniform flips; flip
rates set by class frequency; asymmetric flip rates for binary tasks; a random per-sample flip rate;
a separate flip rate for each class; and flips directed by a confusion matrix. Because the realised
flip rate differs from the nominal flip probability under most of these strategies, a binary search
finds the flip probability that produces a target fraction of changed labels.

% TODO: review, needs manual review
\paragraph{Robustness metrics.} For regression the package computes $R^2$, RMSE and mean absolute
error at every noise level. For classification it computes accuracy, and precision, recall and F1
score, both macro- and weighted-averaged, with a per-class breakdown. Robustness is read off the
retention curve $\mathrm{ret}(x) = m(x)/m(0)$, where $m$ is a bounded, higher-is-better metric
($R^2$ for regression; accuracy, F1, precision or recall for classification) and $x$ is the noise
level. The primary summary is the normalised area under that curve,
\begin{equation}
\mathrm{AUC}_{\mathrm{norm}} = \frac{1}{x_{\max} - x_{\min}} \int_{x_{\min}}^{x_{\max}} \mathrm{ret}(x)\,\mathrm{d}x,
\label{eq:auc_norm}
\end{equation}
evaluated by the trapezoid rule. It is the mean fraction of clean performance retained across the
sweep, lies near $[0, 1]$, and does not depend on how good the model was on clean labels. A
secondary summary fits $\mathrm{ret}(x) \approx \exp\!\left(-(x/\tau_W)^{\beta}\right)$. The shape
$\beta$ separates curves that hold and then collapse ($\beta > 1$) from curves that fall early and
then plateau ($\beta < 1$), and $\tau_W$ sets the scale of the decay. The fit needs at least four
noise levels. These two summaries replace the linear slope of performance against noise used in
earlier versions, which misdescribed a nonlinear curve and was coupled to clean performance. The
package also reports the clean baseline and the retention at the highest level. A flag marks any
curve that falls below zero retention, and an optional threshold withholds the curve summaries for
models whose clean baseline is too low to degrade meaningfully.

% TODO: review, needs manual review
\paragraph{Uncertainty metrics.} For a model that returns a per-sample predicted standard deviation
$u_i$, the package computes, at each noise level: the Spearman correlation between $u_i$ and the
absolute error $|y_i - \hat{y}_i|$; the empirical coverage $\frac{1}{n}\sum_i \mathbb{1}\left[|y_i -
\hat{y}_i| \le k u_i\right]$ at $k = 1$ and $k = 2$, against Gaussian targets of $68.3\%$ and
$95.4\%$; the mean interval width $\frac{1}{n}\sum_i 2u_i$; and the expected calibration error,
\begin{equation}
\mathrm{ECE} = \sum_{b=1}^{10} \frac{|B_b|}{n}\,\left|\bar{u}_b - \overline{|y - \hat{y}|}_b\right|,
\label{eq:ece}
\end{equation}
over ten equal-count bins $B_b$ of $u$. These are computed on the clean test labels. The Spearman
correlation between $u_i$ and the injected noise $|\epsilon_i|$ can only be computed where the noise
is known, which is the corrupted molecules, so it takes a separate pair of inputs. It is reported
within each noise level. A pooled correlation across all levels is also returned, under a separate
name. The pooled value mostly reflects that uncertainty and noise both rise with the level, so it
does not measure whether a model can pick out which molecules are corrupted.

% TODO: review, needs manual review
Three optional wrappers produce $u_i$ from a point predictor: split-conformal prediction for any
scikit-learn-compatible model, optionally scaled by the label spread of each molecule's nearest
training neighbours; an exact Gaussian process with a Tanimoto or RBF kernel, via \texttt{Gauche}
\citep{gauche}; and Monte Carlo dropout for any \texttt{PyTorch} network \citep{gal2016}. Example
notebooks and scripts cover regression and classification end to end.


\section{Evaluation methodology}

% ---- NOTES: Evaluation methodology ----
% Takeaway: every split and every fold in the thesis is group-aware; no random split is used anywhere.
% Takeaway: the metric families are named once here; each chapter says which one it uses.
% Paragraph 1 (new, before the metrics): why never at random.
% - All members of a group land on the same side of every split and fold.
% - A random split puts close structural analogues on both sides, so the score rewards recall of near-copies \citep{Wallach2018, Yang2019}. Sequence data has the same problem \citep{Notin2023, Vieira2026}.
% Paragraph 2: the small-molecule group.
% - Group = the Bemis-Murcko scaffold: ring systems plus the linkers between them, side chains removed \citep{Bemis1996}. Computed with RDKit \citep{rdkit}.
% - How stereoisomers and acyclic molecules are grouped differs between Chapter 3 and Chapters 5-6. Each chapter's methods states its own rule.
% Paragraph 3: how folds are formed.
% - Regression: group k-fold. Classification: stratified group k-fold. How groups are assigned to folds belongs in an Additional file.
% - Non-molecule datasets use their own groups. Point to Chapter 5, "Datasets and splitting".
% Paragraph 4: predictive metrics.
% - Regression: RMSE and R$^2$, plus Spearman's $\rho$ where ranking matters. Classification: accuracy, F1, AUROC and MCC. Chapter 5, "Models and metrics", sets which each dataset uses.
% ---- END NOTES ----

Root mean squared error (RMSE) and the coefficient of determination ($R^2$) were primarily used to compare QSAR predictive performance.


When evaluating uncertainty, we assessed the relationship between predicted uncertainty and prediction quality. Spearman's rank correlation coefficient ($\rho$) was used to quantify the relationship between predicted uncertainty ($u_i$) and absolute error ($|y_i - \hat{y}_i|$), as well as between predicted uncertainty and injected noise magnitude. Higher correlations indicate that the model's uncertainty estimates reliably identify predictions that are likely to be incorrect.

% TODO: review, needs manual review
For probabilistic models, empirical coverage at $1\sigma$ and $2\sigma$ and the expected calibration error (Equation~\ref{eq:ece}) are defined with the NoiseInject uncertainty metrics above.

\section{Conclusions}

% ---- NOTES: Conclusions ----
% Takeaway: a predictor is a representation paired with a model; Chapter 3 asks which of the two governs how badly a predictor degrades under label noise.
% Paragraph 1: 3-4 sentences.
% - Sentence 1: both halves of a predictor, the representation and the model, have now been introduced.
% - Sentence 2: real assay noise is not independent, identically distributed or symmetric. The artificial noise model has seven conditions. LNL has been set out; no single method wins.
% - Sentence 3: NoiseInject injects the conditions and reports robustness and uncertainty metrics. The evaluation section fixes the splits and metrics.
% - Sentence 4: Chapter 3 asks whether the representation or the model architecture governs how badly a predictor degrades under label noise. Point to Chapter 1 for its questions.
% ---- END NOTES ----
\todo[inline, color=yellow!40]{Write the Background chapter's closing paragraph.}
\chapter{Noise Robustness Across Molecular Representations and Model Architectures}\label{chap:robustness}
% ---- TAKEAWAYS ----
% - The model, not the representation, sets noise robustness: 49.5% against 9.2% of the variance in AUC_norm (QM9, Gaussian noise).
% - Clean accuracy does not predict robustness. Random forests were the most consistent across the four datasets; the neural networks depended most on the representation.
% - The pattern of the noise matters more than the shape of its distribution. Grouped-shifted noise and censoring cost the most; Laplace, Student-t and outlier noise behave like Gaussian noise.
% - The representation still matters in two places: the best representation for a model changes with the noise condition, and it decides whether uncertainty splits cleanly into aleatoric and epistemic parts.
% - Mean predicted uncertainty rises with label noise, but no model's per-molecule uncertainty identifies the corrupted labels. Under censoring, models are most confident on the clipped labels.
% ---- END TAKEAWAYS ----

% ---- NOTES: Preamble ----
% Takeaway: the chapter is adapted from a jointly authored paper.
% - Change: the front matter has no "Statement of Authorship" page. Contributions and signatures are on the "Statement of Originality" page.
% ---- END NOTES ----
% TODO: review, needs manual review
\textit{Preamble.} The work in this chapter is adapted from Punt, A., Hanser, T., Werner, S., and Morris, G.M., in preparation for the \textit{Journal of Cheminformatics}. Most of the figures, tables, and text in this chapter are therefore either identical or similar to the content of that manuscript. A signed Statement of Authorship for this jointly authored work is included in the front matter.

\section{Introduction}
% ---- NOTES: Introduction ----
% Takeaway: few studies compare how representations respond to label noise. This chapter does, with noise patterns beyond a matched mean and spread.
% - Change: "direction comparisons" to "direct comparisons"; "responds" to "respond"; drop the second "fingerprints" after "(ECFP)".
% - Open with one clause: a predictor is a representation paired with a model (Chapter~\ref{chap:background}).
% - After the Deng2023 sentence, one sentence: \citet{Cortes2015} varied algorithm, noise level and dataset on bioactivity data, but held the representation fixed. This chapter varies it.
% - Add: one sentence saying what is new in this chapter, beside the list of questions it answers.
% ---- END NOTES ----


While there have been many investigations into the effect of noise on molecular property prediction, direction comparisons between how different molecular representations responds to label noise remains limited. The choice of molecular representation certainly affects predictive performance. \cite{Deng2023} evaluated 62,820 models and found that extended-connectivity fingerprints (ECFP) fingerprints frequently outperform learned graph embeddings across models. Whether these conclusions translate to noise robustness is our question. Between different types of label noise, \cite{Heid2023} found no difference in model performance between error distributions with the same mean and standard deviation. We expand on this by considering alternative patterns of noise, namely noise concentrated in structural groups or on a sparse subset of samples.

% TODO: review, needs manual review
This chapter answers the first four questions set out in Section~\ref{sec:intro:questions}: how representation and model architecture divide the variance in predictive performance and in noise robustness, what probabilistic models add, whether the type of label noise matters, and whether uncertainty estimates track label noise.

\section{Methods}\label{sec:robustness:methods}
% ---- NOTES: Methods ----
% Takeaway: four datasets, six representations, 20 model configurations, seven noise conditions, seven noise levels.
% - The 20 configurations are 13 base models and 7 variant models, defined under Models.
% ---- END NOTES ----


\subsection{Datasets}

Although we are investigating the effects of label noise stemming from experimental settings, for the purposes of testing we inject controlled amounts of artificial noise into clean labels. Data were extracted from the QM9 molecular property dataset distributed by PyTorch Geometric \citep{Ramakrishnan2014}, which contains 130,831 small organic molecules with pre-computed quantum-mechanical properties obtained from density functional theory (DFT) calculations at the B3LYP/6-31G(2df,p) level of theory. We selected the gap between the highest occupied and the lowest unoccupied molecular orbital (HOMO--LUMO gap) as the prediction target in QM9, as it captures a molecule's electronic excitability, charge transfer capability, and chemical stability \citep{Fediai2023}. We assumed that QM9 labels were free of noise, notwithstanding approximations in the level of theory. Each experiment was replicated ten times with a distinct random seed and a randomly selected subset of $N = 10{,}000$ molecules, split 80/10/10 into 8,000 training, 1,000 validation and 1,000 test molecules. All models were fitted on those 8,000 training molecules, except the Gaussian processes, which were fitted on a random subset of 5,000 of them; the validation and test molecules were unchanged.

We repeated these experiments on three molecular property datasets with experimentally determined endpoints. From the OpenADMET initiative \citep{openadmet}, we used LogD and Caco-2 efflux permeability, with $N = 5{,}039$ and $N = 2{,}161$ respectively. We also selected hERG $K_i$ data from ChEMBL \citep{Zdrazil2023}, following a protocol inspired by \citet{landrum2024}: filter for binding assays, deduplicate by median pChEMBL value, and remove compounds with an inter-assay standard deviation greater than 1.0 log unit, resulting in $N = 1{,}415$ compounds. Each of the three was partitioned into five folds by grouped $k$-fold cross-validation, and each fold's training block gave up a further grouped fifth as a validation set for early stopping. A model therefore fits between 3,194 and 3,364 of the 5,039 LogD molecules, depending on the fold. Unlike QM9's ten replicates, which are ten independent draws from one pool, each assay dataset is a single fixed set in which every molecule is tested exactly once and nothing is redrawn.

All splits and cross-validation folds were grouped by the chirality-blind Bemis--Murcko scaffold, implemented such that each distinct acyclic molecule forms its own group, and groups are filled in random order. Prior to modeling, all molecular structures were sanitized and canonicalized using \texttt{RDKit} to remove invalid valence states and standardize atom and bond typing \citep{rdkit}. Prediction targets were standardized to zero mean and unit variance. Experiments that track uncertainty were run once rather than in replicates.

\subsection{Representations}

In this study, we evaluated six molecular representations: extended-connectivity fingerprints of diameter four (ECFP4), Avalon, Sort \& Slice (SNS), physicochemical descriptors (PDV), ChemBERTa and Molecular Hypergraph Grammar GNNs (MHG-GNNs). The same six were used on QM9 and on the three assay datasets. ECFP4 is a Morgan fingerprint with radius $r = 2$ with $d = 2048$ bits \citep{rogers2010, rdkit}. Sort \& Slice takes the same substructures as ECFP4, but it ranks them by prevalence and selects the top $L$ ($L=1024$) \citep{sns}. We fitted the Sort \& Slice vocabulary on the training molecules alone, refitting it for each QM9 replicate and for each cross-validation fold.  Avalon is a 2,048-bit fingerprint, generated with the Avalon toolkit through RDKit \citep{avalon, rdkit}. The PDV representation is a vector of 200 physicochemical descriptors computed with RDKit's \texttt{MolecularDescriptorCalculator} \citep{rdkit, Cherkasov2014}; the 200 are listed in Additional file~\ref{app:noise:sec:pdv_descriptors}, and the list is fixed in our code rather than taken from RDKit's default descriptor set, which grows between releases. ChemBERTa embeddings are the mean over the non-padding token embeddings of the \texttt{DeepChem/ChemBERTa-77M-MTR} checkpoint, producing 384 features \citep{Ahmad2022}. MHG-GNNs are a GIN-based autoencoder pre-trained on 1.34 million PubChem molecules using $\beta$-VAE loss, producing 1024-dimensional embeddings through iterative message passing and found in the published \texttt{mhggnn\_pretrained\_model\_0724\_2023} checkpoint \citep{kishimoto2023}. We z-score normalized PDV, ChemBERTa and MHG-GNN per feature, using the training molecules' mean and standard deviation; ECFP4, Avalon and Sort \& Slice reach the model as built.

\subsection{Models}
% ---- NOTES: Models ----
% Takeaway: 13 base models carry every comparison across models; 7 variant models appear only where they are the subject.
% - Base models: RF, QRF, XGBoost, LightGBM, NGBoost, SVM, GP (RBF kernel), NN-$\alpha$, BNN-$\alpha$, VBLL-$\alpha$, NN-$\beta$, BNN-$\beta$, VBLL-$\beta$.
% - Variants: RF with 300 trees, GP with a Tanimoto kernel, heteroscedastic GP, BNN-$\alpha$ and BNN-$\beta$ with a variance head, VBLL-$\alpha$ and VBLL-$\beta$ with input-dependent noise.
% Table: one row per configuration (20 rows): family, name, base or variant, what it adds, representations it runs on.
% - Change: "idential" to "identical"; "accomodates" to "accommodates".
% - Change: "As such, we used a standard GP" follows nothing. Give the reason: an exact GP costs time cubic in the number of training molecules.
% - Change: move the sentence on preliminary kernel tests next to the kernel sentence of the GP or the SVM.
% - Change: the PyTorch citation is \citep{Paszke2019}, not \citep{pytorchGeometric}.
% - Change: "We used Bayesian transformations ..." repeats the next paragraph's opening. Keep one.
% - Add: one tuned setting per network per dataset, shared across the six representations (Additional file~\ref{app:noise:sec:hyperparameters_tuned}). Every non-network model uses defaults.
% - Change: the variance head is added to the BNN, not to the plain network. So "a cheaper alternative to the BNN" does not hold.
% - Change: VBLL here makes every linear layer variational, not only the last. Keep "replaces every linear layer"; say the name comes from \citet{Harrison2024}.
% - Change: "VBLL networks report only epistemic uncertainty" is wrong. VBLL gives both parts, the aleatoric part one number for the whole fit. The BNN without a head is the epistemic-only one.
% - Add after the last sentence: the VBLL networks were tuned separately, so NN against VBLL differs by more than the variational layers. VBLL-$\beta$ kept the defaults on Caco-2 and hERG K$_i$.
% ---- END NOTES ----

% TODO: review, needs manual review
To understand model behavior in the presence of label noise, we explored a range of model families. We fitted tree ensembles (random forests, quantile regression forests, XGBoost, LightGBM and NGBoost), a support vector machine, Gaussian processes and neural networks. Several are variants of one another, differing in the probabilistic machinery they add. All hyperparameters for the models below are included in Additional file~\ref{app:noise:sec:hyperparameters_default}.

While we used 100 trees for RF, we observed a noisy quantile estimate with 100 trees for QRF, prompting us to use 300 instead. To observe the effect of the quantile mechanism, we then fit a second RF at 300 trees so the two models were idential in all other settings.

NGBoost relies on scoring off validation data, which differs by dataset. On QM9 the validation split is already held out of every fit. However, on the three assay datasets the validation set is a scaffold-grouped fifth of the rows NGBoost is handed, which already exclude the validation molecules.

We used a radial basis function (RBF) with $C = 1.0$ and $\gamma = \texttt{scale}$ as the kernel, as it accomodates continuous vectors.

We fitted three variants of GPs, each with a constant mean and a scaled kernel. The first uses an RBF kernel and runs on all six representations. The second uses the Tanimoto kernel \citep{Ralaivola2005, moss2020} and was run on ECFP4 alone. The third is a heteroscedastic variant in which the RBF process is given a second network that predicts the observation noise of each molecule. As such, we used a standard GP and only fit on at most $5{,}000$ training molecules, drawn by a seeded sample. Due to the size of the datasets, this cap only affected QM9. We implemented these GP models in \texttt{GPyTorch} \citep{gpytorch}, with the Tanimoto kernel taken from the \texttt{Gauche} framework \citep{gauche}.

While preliminary tests indicated that the Tanimoto kernel provides GP with substantial benefits over other kernels, for SVM, the RBF kernel was sufficient for all types of data.

% TODO: review, needs manual review
We tested two deterministic feed-forward architectures, both fully connected and both implemented in PyTorch \citep{pytorchGeometric}: (i) NN-$\alpha$ with two hidden layers of sizes 128 and 64, and dropout after each hidden layer and (ii) NN-$\beta$ with two hidden layers, each of size 128, with dropout applied before the output layer. Each model was trained using the Adam optimizer with early stopping on validation loss. We used Bayesian transformations to convert NNs into Bayesian neural networks (BNNs), with ReLU activations and dropout regularization ($p=0.2$). Additional file~\ref{app:noise:sec:hyperparameters_default} gives their full widths, dropout fractions, learning rates and batch size.

To get these networks producing uncertainty estimates, we applied four probabilistic transformations to each architecture, producing eight probabilistic networks. The first replaces every linear layer with a Bayesian layer carrying a Gaussian prior $\mathcal{N}(0, 0.1^2)$ on its weights, producing BNN-$\alpha$ and BNN-$\beta$. They are trained on an evidence lower bound whose  Kullback–Leibler (KL) divergence term is divided by the number of training molecules. While this is a straightforward approach, it is computationally expensive. We explored alternative techniques like our second variant, where we add an output head to the standard NN and fit the negative log likelihood of a Gaussian with that variance \citep{kendall2017}. The third, Variational Bayesian Last Layer (VBLL), replaces every linear layer with a layer that produces mean-field variational posteriors $q(\mathbf{W}) = \mathcal{N}(\boldsymbol{\mu}_W, \text{diag}(\boldsymbol{\sigma}_W^2))$ over the last-layer weights and trains by maximizing the evidence lower bound. This can be thought of as minimizing the reconstruction loss plus KL divergence $D_{\text{KL}}(q(\mathbf{W}) \| p(\mathbf{W}))$ from a standard normal prior, scaled by $1/N$ \citep{Harrison2024}. Because only its output layer keeps a learned observation-noise parameter, VBLL networks report only epistemic uncertainty. The fourth makes observation noise a function of the input. All eight probabilistic networks estimate their predictive distributions from 100 stochastic forward passes. Each network and its Bayesian counterpart were tuned to the same set of parameters so that NN-$\alpha$ against BNN-$\alpha$ and NN-$\beta$ against BNN-$\beta$ differ by the Bayesian layers alone.

\subsection{Noise conditions and levels}
% ---- NOTES: Noise conditions and levels ----
% Takeaway: how much noise each level adds, and which conditions ran on which models and representations.
% - Change: give the unit. A level is a multiple of the standard deviation of the clean training labels; level 0.5 adds noise with half that spread.
% - Change: say how censoring clips: every training label above the $k$th largest clean training label is set to that label's value.
% - Change: "selected pairs" is never spelled out. Point to the table below. A pair is one model on one representation.
% Table: which conditions ran on which pairs, all four datasets.
% - Gaussian, grouped-wider, grouped-shifted: every configuration on every representation, the Tanimoto GP on ECFP4 only (115 pairs).
% - Student-$t$, Laplace, outlier: 8 models on ECFP4, PDV and ChemBERTa (24 pairs). Chosen to span the most and least noise-tolerant models under Gaussian noise, plus the models with a per-molecule uncertainty.
% - Censoring: NGBoost, GP (RBF) and BNN-$\alpha$ (variance head) on PDV; RF and the heteroscedastic GP on ECFP4 (5 pairs). Four report a per-molecule uncertainty; RF is the contrast.
% Result needed: confirm 115 pairs. VBLL-$\alpha$ with input-dependent noise has no Student-$t$, Laplace or outlier runs on the assay datasets (21 of 24 pairs there).
% ---- END NOTES ----
% TODO: review, needs manual review
The seven noise conditions are defined in Section~\ref{sec:noise_injection}.

We ran experiments across seven noise levels, $0$, $0.2$, $0.3$, $0.5$, $0.75$, $1.0$ and $1.5$. We added the noise to the raw label and standardized the target afterwards, using the clean training mean and spread. For censoring, the level is the fraction of labels clipped, at $0$, $0.1$, $0.2$, $0.25$, $0.3$, $0.4$ and $0.5$ \citep{Svensson2025}.

We ran Gaussian, grouped-wider and grouped-shifted noise on every model and representation pairing. The
remaining four conditions ran on selected pairs. 


\subsection{The QM9 precompute pipeline}
% ---- NOTES: The QM9 precompute pipeline ----
% Takeaway: how QM9 goes from featurisation through noise injection to training.
% - Add to "Models": the plain networks train on squared error; the Bayesian and VBLL networks on the evidence lower bound; with a variance head, squared error becomes the Gaussian negative log-likelihood \citep{kendall2017}.
% - Change: the figure and caption show one-hot SMILES and domain labels, which no run writes. Redraw with the six representations; drop both from the caption.
% - Suggest: the byte layout of each entry and the provenance-row paragraph belong in an Additional file. The main text keeps what the pipeline does, not its file format.
% ---- END NOTES ----

% TODO: review, needs manual review
QM9 \citep{Ramakrishnan2014} runs through a separate pipeline, loaded via PyTorch Geometric. Python
loads the data, splits it, computes every representation, and trains and scores the models. A Rust
program sits between featurisation and training. It injects the noise, standardises the labels and
writes the files the models are trained from.

% TODO: review, needs manual review
Each replicate draws a fresh scaffold split, as described under Datasets. Each acyclic molecule's
group is keyed on its canonical SMILES, and a group is never divided between two splits. The same
grouping defines the scaffold families used by the grouped noise conditions.

% TODO: review, needs manual review
Python computes the six representations described above. Binary fingerprints are bit-packed, Sort \&
Slice counts are stored as 16-bit integers, and the descriptors and learned embeddings are stored
as 32-bit floats. A molecule whose ECFP4 or Avalon fingerprint is entirely zero stops the run, so it
cannot enter training as if it had been featurised. One-hot encoded SMILES can still be built but
are not used in this thesis.

% TODO: review, needs manual review
Traditional approaches load entire datasets into memory from CSV or Parquet files, creating RAM bottlenecks at scale and redundant deserialization overhead across experimental replicates. Our pipeline instead serializes preprocessed data into custom memory-mapped binary files (\texttt{.mmap}) that enable on-demand access with constant low memory footprint regardless of dataset size (Figure~\ref{fig:fig_mmap}). Python writes one file per split, the Rust program reads each one and writes a new file with the processed labels in its place, and Python then reads the new files back for training.

% TODO: review, needs manual review
Each mmap file stores a sequence of fixed-format entries. Entry structure includes: length-prefixed UTF-8 SMILES strings (isomeric and canonical), the raw target value (float32, 4 bytes), the requested molecular representations as fixed-width blocks (Sort \& Slice: uint16[1024] = 2,048 bytes, PDV: float32[200] = 800 bytes, ChemBERTa: float32[384] = 1,536 bytes, MHG-GNN: float32[1024] = 4,096 bytes, Avalon: 256 bytes, ECFP4: 256 bytes), and the processed target value incorporating noise and normalization (float32, 4 bytes). Only the requested representations are written, so the byte offset of every field depends on the list of representations. This format enables zero-copy memory mapping—Python and Rust access the same disk-backed memory pages without deserialization, maintaining constant RAM usage regardless of dataset size.

% TODO: review, needs manual review
The Rust program builds a noise plan for each split separately, in that split's own record order.
The training plan defines the condition. The dose, the censoring limit and the affected scaffold
families are all computed from the clean training labels. Validation labels receive an independent
draw, from a seed derived from the training seed, at the same dose in label units, the same limit
and the same affected families. The program checks that validation received the same amount of
noise as training. Test labels are never noised. Every split still records the per-molecule scale
its region would have received, so held-out molecules can be scored against the training pattern.

% TODO: review, needs manual review
For every molecule in every split, the Rust program writes a provenance row: its canonical SMILES,
clean label, injected noise, noise scale, the pattern its region was exposed to, noisy label and the
standardised value written. It also writes a run-level manifest of the condition's parameters and
the delivered dose. The record at each row index must hold the same molecule on the way out as when
the noise was drawn, or the run stops. A self-test mode runs the same dose and targeting checks on a
label file without the rest of the pipeline.

% TODO: review, needs manual review
\paragraph{Models.} The pipeline trains the models described under Models above.
Hyperparameters are either package defaults, values tuned by Optuna \citep{akiba2019optuna}, or
supplied from a JSON file. For models that return uncertainty, the pipeline can also score the
training molecules out of fold, so that each training molecule's uncertainty comes from a model that
never saw its corrupted label.

\begin{figure}[htbp]
    \centering
    \includegraphics[width=0.95\textwidth]{figure_mmap_structure.png}
    \caption{Memory-mapped file architecture contrasted with traditional in-memory loading. Top: CSV/Parquet approaches load entire datasets into RAM, creating bottlenecks at scale. Bottom: Memory-mapped binary files enable on-demand line-by-line access with constant low RAM usage. Right: Entry-level structure showing byte-aligned fields with fixed-size representations (target values, fingerprints, domain labels) and length-prefixed variable-length fields (SMILES strings, one-hot encodings).}
    \label{fig:fig_mmap}
\end{figure}


\subsection{Uncertainty}
% ---- NOTES: Uncertainty ----
% Takeaway: how each model's uncertainty splits into aleatoric and epistemic parts, and why comparisons with the injected noise use out-of-fold training molecules.
% - Change: "the variance the labels the molecule's given leaf" to "the variance of the training labels in the molecule's leaf, in each tree"; "are those variances" to "is"; "the the" to "the".
% ---- END NOTES ----

Up to this point we have been concerned with how noise affects the predictive accuracy of QSAR models. However, label noise should also affect the uncertainty these models report. Some models return only a point prediction, while others also return an estimate of uncertainty, either as a predicted distribution or as the spread across an ensemble. Furthermore, some report a single value of uncertainty while others report one that can be decomposed into an aleatoric and an epistemic component \citep{kendall2017}. The aleatoric component describes underlying noise the model associates with the data, while the epistemic component represents the uncertainty within the model's fit. We sought to determine if, as expected, the aleatoric component rises with injected label noise while the epistemic component does not.

For the BNN and VBLL transformations, we took the two components over $T = 100$ stochastic forward
passes \citep{kendall2017}:
\begin{equation}
\hat{u}^2_{\text{ale}}(x) = \frac{1}{T}\sum_{t=1}^{T}\hat{v}_t(x),
\qquad
\hat{u}^2_{\text{epi}}(x) = \frac{1}{T}\sum_{t=1}^{T}\big(\hat{\mu}_t(x) - \bar{\mu}(x)\big)^2 ,
\label{eq:sampling_split}
\end{equation}
where $\hat{\mu}_t(x)$ and $\hat{v}_t(x)$ are the mean and the variance returned for molecule $x$ on
pass $t$, and $\bar{\mu}(x)$ is the mean over passes. Here $\hat{v}_t(x)$ is the predicted observation
noise variance. It varies per molecule in the variance-head and heteroscedastic variants, and is one
number for the whole fit in the VBLL transformation. A full-BNN without a variance head returns no
variance, so its uncertainty is epistemic alone.

The other models follow the same split. For GPs, the epistemic component is the posterior variance of the latent function and the aleatoric component is the learned observation noise \citep{Rasmussen2005}. All variants of the GP produce one number for the whole fit except heteroscedastic. For QRF models, each tree holds multiple training molecules. So for a given molecule, the aleatoric component is the variance the labels the molecule's given leaf. The total aleatoric uncertainty are those variances averaged over the trees. The epistemic component is the variance of the leaf means across the trees. NGBoost predicts one variance for each molecule from a single fit, which is aleatoric alone \citep{Duan2020}. In every model, the total uncertainty is the sum of the the aleatoric and epistemic components.

No held-out label carries injected noise, so every comparison against the injected noise was made on
training molecules. Each training set was divided into five folds grouped on Murcko scaffolds, and each
molecule was scored by a refit that excluded its own fold. For NGBoost on QM9, three of the five folds
were scored to limit its cost. No uncertainty was calibrated after fitting.


\subsection{Evaluation}
% ---- NOTES: Evaluation ----
% Takeaway: robustness is AUC_norm, read beside clean R^2. An ANOVA per noise condition splits the variance between model and representation.
% Result needed: the excluded-configurations Additional file still uses the old gate (clean R^2 <= 0.6). Regenerate at 0.3.
% - Change: "probabilistic models with its" to "each probabilistic model with its".
% - Change: the ANOVA and the two exclusion paragraphs are out of date. Her current paper has them rewritten; copy those in. What changed:
% - The ANOVA now has a replicate block and models nested in eight model families. The interval on each share is a 95% bootstrap over whole replicates.
% - No model or representation is dropped on a correlation. Near-identical models are grouped into one family instead (RF with QRF, XGBoost with LightGBM, each NN with its BNN, the two VBLLs).
% - ECFP4 and Sort & Slice stay separate: their clean R^2 differs by more than the replicate spread for all 13 base models.
% - Drop the ICC sentences. Cite the paper's family, representation-group and Tukey tables, not the older redundancy and ICC Additional files.
% - Change: "whether uncertainty rise" to "rises"; "models who" to "models that".
% - Change: the second uncertainty question divides the epistemic rise by the rise in mean aleatoric uncertainty, not by its mean.
% ---- END NOTES ----
We compute statistics on the held-out molecules, for QM9 the test split and for three assay datasets the held-out fold of the scaffold cross-validation. Different combinations of dataset, model, representation, noise condition, and noise level are examined in each scenario. We used the coefficient of determination ($R^2$) throughout. For molecule $i$, let $y_i$ be the clean label, $\hat{y}_i$ the prediction, $\epsilon_i$ the noise injected into its label and $u_i$ the predicted uncertainty.

To evaluate the effect of label noise, we tracked predictive performance retention under increasing artificial noise. For each QSAR pairing we recorded $R^2(\tau)$ at each noise level $\tau$, and normalized it to obtain the retention ratio $R^2(\tau)/R^2(0)$. To quantify robustness, we used the normalized area under this retention curve, \aucnorm (Equation~\ref{eq:auc_norm}), evaluated using the trapezoidal rule over the seven noise levels. While we tested other metrics of robustness, we found \aucnorm to be less directly correlated to clean-label $R^2$ predictions. For censoring, we sweep over a set of fractions to clip the labels prior to normalization. 

An \aucnorm value near 1 indicates that the QSAR model in question has retained predictive performance with artificial noise, meaning it is more robust to noise. Since we normalize the curve by its own clean-label predictive performance, a model with a low clean $R^2$ has less to lose, so \aucnorm is reported beside clean $R^2$. We excluded replicates with a clean $R^2 < 0.3$ from the robustness analysis (Additional file~\ref{app:noise:sec:excluded_configurations}), as performance retention ratios become unstable when the clean-label denominator approaches zero. On QM9 one \aucnorm is computed per replicate, and we report the median over the ten replicates. On the three assay datasets one value is computed per scaffold fold, and the median over the five folds is reported. 

To compare the relative effect of the thirteen base models across datasets and noise conditions, we compared their rankings using Spearman's rank correlation coefficient ($\rho$) between clean $R^2$ and \aucnorm. We compared probabilistic models  with its deterministic counterpart using a two-sided Wilcoxon signed-rank test paired on the replicate, at $\alpha = 0.05$, separately for each representation and noise condition. This test was run only on QM9, as with five scaffold folds the
smallest attainable $p$ is 0.0625.

We conducted a separate two-way analysis of variance (ANOVA) for each noise condition to identify the
relative contributions of molecular representation and model architecture to both predictive
performance and noise robustness. For a given metric $y$, either $R^2$ at a noise level of 1.0 or
\aucnorm:
\begin{equation}
y_{ijr} = \mu + \alpha_i + \beta_j + (\alpha\beta)_{ij} + e_{ijr},
\end{equation}
where $\alpha_i$ represents the effect of model architecture $i$, $\beta_j$ represents the effect of
molecular representation $j$, $(\alpha\beta)_{ij}$ is the interaction term, and $e_{ijr}$ is the
residual for replicate $r$. The proportion of variance explained by each term was calculated as the
effect size
\begin{equation}
\eta^2_{\text{term}} = \frac{SS_{\text{term}}}{SS_{\text{total}}},
\end{equation}
using Type I (sequential) sums of squares, entered in the order model architecture, molecular
representation, interaction. The ten QM9 replicates provide the residual. The uncertainty on each share is the range of $\eta^2$ when the decomposition is repeated with each replicate left out in turn. 

ANOVA relies on the assumption that observations both within and across groups are independence. Across the whole set of models and representations in this study, some pairs produced similar predictions that are not independent. This causes the residual degrees of freedom to increase, while the residual sum of squares barely moves. To handle this issue, we computed pairwise Spearman rank correlations between all models and between all representations across noise strategies. Models whose profiles correlated at $\rho > 0.99$ with another model were excluded, including quantile regression forests (QRF), which were very similar to RF, as well as different BNN variants. 

Across representations, SNS correlated at $\rho > 0.90$ with the ECFP4 fingerprint and was excluded to avoid inflating representation degrees of freedom, as both encode overlapping circular substructure information. We also computed intraclass correlation coefficients ICC(1,1) for all model pairs. ICC(1,1) measures the proportion of total variance attributable to between-subject differences, with values near 1.0 indicating that two models rank configurations almost identically. The full redundancy and ICC tables are reported in Additional files~\ref{app:noise:sec:model_redundancy}--\ref{app:noise:sec:model_icc}.

To assess the interaction between uncertainty estimation and label noise, we posed three questions. First, whether uncertainty rise with label noise, as measured by the mean $u_i$ over held-out molecules at the highest noise level
divided by its value with no added noise, in each replicate or fold. The second pertains to models who report separate aleatoric and epistemic components. We divided the increase of mean epistemic uncertainty between the highest level of noise and clean labels by the mean of the aleatoric component to determine how each responds to label noise. A ratio near 0 means the added noise was attributed to the labels while a ratio of 1 or above means the epistemic component grew as much or more. We also compared the mean aleatoric component with the standard deviation of the injected noise. The third question we had was whether per-sample uncertainty values track injected label noise, as measured by the Spearman correlation between $u_i$ and $|\epsilon_i|$ over out-of-fold training molecules, per fold and per noise level. We also report the Spearman correlation between $u_i$ and the absolute error of the prediction. Finally, we asked whether dividing a molecule's out-of-fold absolute error by $u_i$ ranks the corrupted labels better than the error alone, and report the difference between the two Spearman correlations with $|\epsilon_i|$.


\section{Representation governs clean performance, architecture governs robustness}
% ---- NOTES: Representation governs clean performance, architecture governs robustness ----
% Takeaway: model architecture explains most of the variance in robustness. For R^2 under noise, model and representation are comparable.
% - Change: the prose, table and figures are from an earlier study (six noise strategies, NDS, gate 0.6). Rewrite from these bullets.
% Result needed: regenerate the table and figure with the replicate block and the eight families; every share will move.
% - Data: QM9, 13 base models x 6 representations x 10 replicates; Gaussian, grouped-wider, grouped-shifted.
% - Gaussian, AUC_norm: model 49.5%, representation 9.2%. R^2 at level 1.0: model 29.1%, representation 26.6%. Grouped-wider gives the same picture.
% - The interaction is larger for robustness than for R^2 (numbers in the table).
% - Grouped-shifted is dominated by the residual: its offset per scaffold group differs between replicates.
% Result needed: the clean-R^2 row, and the decomposition on LogD, Caco-2 and hERG K$_i$ under AUC_norm.
% Table (replaces the one below): rows condition x outcome (AUC_norm, R^2 at level 1.0), plus the clean row; columns model, representation, interaction, residual.
% - Change: the interaction figure (NDS) has no counterpart; drop it or rebuild on AUC_norm. The last paragraph and the overview figure move to the clean-against-robustness comparison.
% ---- END NOTES ----


Considering the two main components of a QSAR model, the molecular representation and the model architecture, we sought to understand how each contributes to predictive performance and noise robustness. ANOVA decomposition across six noise strategies on models to predict HOMO--LUMO gap from QM9 demonstrates that for predictive performance (R$^2$ at $\sigma = 0.3$), the model--representation interaction term is the  largest source of variance across all six strategies, followed by model architecture and representation (Table~\ref{tab:anova_decomposition}; Figure~\ref{fig:anova_decomposition}a). However, for noise robustness as measured by the noise degradation slope (NDS), this ordering shifts. The choice of model architecture is instead the largest source of variance, while the interaction term drops to second and representation third. This observed inversion makes sense because the molecular representation dictates the chemical information the model receives, while the architecture determines how efficiently that information is used. The interaction term reflects the idea that these two are not independent. In the context of label noise robustness, representation has less of an impact since the features are not corrupted by noise. Meanwhile, the model's regularization, ensembling, and uncertainty handling have a stronger impact on how it handles noise. 

The impact varied by type of noise, represented by noise strategy (Figure~\ref{fig:anova_decomposition}b). The strategies which degraded in performance the most, threshold and value-proportional, both corrupted label regions, forcing models to rely on their inductive biases. This translates to the model architecture explaining a significant portion of the variance. On the other hand, outlier noise only affects statistical outliers. For outlier noise, model and interaction effects were both smaller.

\begin{table}[htbp]
\centering
\caption{ANOVA variance decomposition by noise strategy on the QM9 HOMO--LUMO gap. Each cell reports $\eta^2$ (\%), the share of variance explained by that factor. Columns: \textbf{Model} = model architecture effect; \textbf{Rep} = molecular representation effect; \textbf{Inter.}\ = model$\times$representation interaction; \textbf{Resid.}\ = residual unexplained variance. ``Performance'' uses R$^2$ at $\sigma = 0.3$ as the outcome; ``Robustness'' uses the noise degradation slope (NDS). Bold entries highlight the dominant factor in each row.}
\label{tab:anova_decomposition}
\small
\begin{tabular}{lcccccccc}
\toprule
& \multicolumn{4}{c}{\textbf{Performance ($\eta^2$, \%)}} & \multicolumn{4}{c}{\textbf{Robustness ($\eta^2$, \%)}} \\
\cmidrule(lr){2-5} \cmidrule(lr){6-9}
\textbf{Strategy} & \textbf{Model} & \textbf{Rep} & \textbf{Inter.} & \textbf{Resid.} & \textbf{Model} & \textbf{Rep} & \textbf{Inter.} & \textbf{Resid.} \\
\midrule
Gaussian        & 24.6 & 22.8 & \textbf{48.1} &  4.5 & \textbf{48.7} &  8.2 & 11.9 & 31.2 \\
Quantile        & 25.1 & 22.4 & \textbf{46.9} &  5.6 & \textbf{47.9} &  6.8 & 14.9 & 30.4 \\
Threshold       & 25.0 & 21.2 & \textbf{48.0} &  5.8 & \textbf{48.0} & 10.1 & 28.0 & 13.9 \\
Heteroscedastic & 27.9 & 23.5 & \textbf{46.2} &  2.5 & \textbf{37.0} &  4.5 & 17.5 & 41.0 \\
Value-prop.     & 27.8 & 22.0 & \textbf{46.5} &  3.8 & \textbf{52.8} &  6.9 & 22.2 & 18.2 \\
Outlier         & 25.3 & 22.0 & \textbf{49.2} &  3.5 & 12.7 &  3.2 &  4.7 & \textbf{79.3} \\
\bottomrule
\end{tabular}
\end{table}

\begin{figure}[htbp]
    \centering
    \includegraphics[width=\textwidth]{fig2_anova_decomposition.png}
    \caption{ANOVA variance decomposition ($\eta^2$, \%) for a) predictive performance (R$^2$ at $\sigma=0.3$) and b) noise robustness, quantified by the noise degradation slope (NDS), by noise strategy. Lower (more negative) NDS indicates faster performance loss under noise.}
    \label{fig:anova_decomposition}
\end{figure}

The interaction between model architecture and representation is seen in Figure~\ref{fig:interaction}. While some model architectures, like SVM and full BNNs maintain consistent noise robustness scores across all representations, others like NN-$\beta$ display greater variation. Model robustness on one representation strongly predicts robustness on another (Spearman $\rho = 0.73$, $p = 0.007$ for ECFP4 vs PDV; 11 models). However, the scatter around the diagonal reflects the interaction term identified by the ANOVA. 

\begin{figure}[htbp]
\centering
\includegraphics[width=\textwidth]{fig_interaction.png}
\caption{Representation--model interaction effects on the QM9 HOMO--LUMO gap. a) Noise degradation slope (NDS) heatmap showing model robustness across representations under Gaussian noise; lower (more negative) NDS indicates faster degradation. ``N/A'' indicates configurations excluded due to baseline $R^2 < 0.6$ or kernel incompatibility (GP uses a Tanimoto kernel defined on binary vectors, which is incompatible with PDV). b) NDS on PDV versus ECFP4 fingerprint for each model under Gaussian noise, with Spearman correlation.}
\label{fig:interaction}
\end{figure}

We filtered out model--representation configurations with baseline predictive performance R$^2 < 0.6$, excluding 66 configurations where poor clean-data performance would produce misleadingly shallow degradation slopes (Additional file~\ref{app:noise:sec:excluded_configurations}). We evaluated all model architectures under increasing label noise on the PDV representation across six strategies, considering both predictive performance as measured by $R^2$ as well as the noise degradation slope (NDS). We observe that NDS varies substantially across architectures and noise strategies, as seen in Figure~\ref{fig:global_overview}. 

\begin{figure}[htbp]
\centering
\includegraphics[width=\textwidth]{fig1_global_overview.png}
\caption{Global overview of noise robustness on the QM9 HOMO--LUMO gap. a) Performance degradation curves (PDV, Gaussian noise) showing R$^2$ versus noise level $\sigma$ for representative models. b) Noise degradation slope (NDS) heatmap across all model architectures and noise strategies on PDV; lower (more negative) NDS indicates faster degradation under noise. The Tanimoto GP is absent because its kernel is incompatible with PDV (see Section~\ref{sec2}); a separate GP with an RBF kernel is included for PDV-specific analyses.}
\label{fig:global_overview}
\end{figure}

\section{Probabilistic models are no more noise-robust than deterministic ones}
% ---- NOTES: Probabilistic models are no more noise-robust than deterministic ones ----
% Takeaway: the Bayesian networks kept more AUC_norm than the plain networks on all four datasets, at a cost in clean R^2.
% - Change: the prose is from the earlier study (NN against VBLL, NDS, QM9 only, mol2vec). Rewrite from these bullets.
% - Change: the heading says probabilistic models are no more robust; the takeaway above says the Bayesian networks kept more. One of them has to change.
% - Pairs: NN-$\alpha$ with BNN-$\alpha$, NN-$\beta$ with BNN-$\beta$, RF (300 trees) with QRF. Each shares one setting. VBLL is a base model, not a pair.
% - BNN-$\beta$ beat NN-$\beta$ in most representation and condition combinations on every dataset; BNN-$\alpha$ on Caco-2 and logD, and about half on hERG K$_i$.
% - The variance-head BNN runs at a different setting from its BNN, so the head alone cannot be credited with a difference.
% - Change: drop "QRF was significantly less robust than RF". It compared 100-tree RF with 300-tree QRF.
% Result needed: QRF minus RF (300 trees), and variance-head BNN minus BNN, in AUC_norm and clean R^2 per dataset.
% Result needed: whether VBLL still fails on ChemBERTa and MHG-GNN under the 0.3 gate. If so, the point moves to the base-model results.
% Table (replaces the Wilcoxon table): pair x condition rows, representation columns; median change in AUC_norm (clean R^2 in brackets), counterpart minus base, QM9. No p-values.
% Result needed: Additional file~\ref{app:noise:sec:bayesian_transformation_tests} still has slope-based Wilcoxon tests; replace it.
% ---- END NOTES ----

To further explore the impact of model architecture on noise robustness, we examined variants of the same core architecture. Specifically, Bayesian transformations were used to explore the impact of the Bayesian framework on noise. By comparing models with the same core architecture, we observe that with full Bayesian transformations, representation explains less than $11\%$ of NDS variance ($p > 0.25$) while variational Bayesian inference on the same NNs is representation-dependent. On MHG-GNN and mol2vec representations, VBLL models were unstable, such that all 24 VBLL $\times$ \{MHG-GNN, mol2vec\} configurations were excluded due to baseline R$^2 < 0.6$. Further work could involve more thoroughly-tuning these models and exploring different core architectures to adapt to these embedding-style inputs. 

In general, the full Bayesian transformations reduced the degradation in performance with Gaussian noise the most, as seen in Figure~\ref{fig:nn_family_comparison}. Both full BNN and VBLL transformations significantly improved robustness for both NN-$\alpha$ and NN-$\beta$ (Table~\ref{tab:wilcoxon_bnn}). Full transformations are more computationally expensive, making VBLL a potentially stronger approach to handling noise with NNs when working with larger data sets or modeling architectures. The choice of transformation varies by representation, as seen in Additional file~\ref{app:noise:sec:bayesian_transformation_tests}, with SMILES receiving the largest benefits from Bayesian transformations. QRF was significantly less robust than RF (Table~\ref{tab:wilcoxon_bnn}), suggesting that quantile regression adds complexity and flexibility which may result in overfitting.

\begin{table}[htbp]
\centering
\caption{Wilcoxon signed-rank tests for Bayesian and probabilistic transformations on the QM9 HOMO--LUMO gap. \textbf{Base} is the deterministic reference model; \textbf{Variant} is the probabilistic/Bayesian transformation tested against it. \textbf{$\Delta$ NDS} = mean change in noise degradation slope (NDS) from base to variant; positive values mean the variant is more robust (slope closer to zero). \textbf{Sig.}\ marks pairs that are statistically significant at $p < 0.05$ (asterisk).}
\label{tab:wilcoxon_bnn}
\small
\begin{tabular}{llrrl}
\toprule
\textbf{Base} & \textbf{Variant} & \textbf{$\Delta$ NDS} & \textbf{$p$-value} & \textbf{Sig.} \\
\midrule
NN-$\alpha$ & BNN-$\alpha$ & $+0.056$ & $2.9 \times 10^{-10}$ & \textbf{*} \\
NN-$\alpha$ & VBLL-$\alpha$ & $+0.061$ & $1.2 \times 10^{-6}$ & \textbf{*} \\
NN-$\beta$ & BNN-$\beta$ & $+0.096$ & $2.9 \times 10^{-11}$ & \textbf{*} \\
NN-$\beta$ & VBLL-$\beta$ & $\mathbf{+0.124}$ & $1.2 \times 10^{-7}$ & \textbf{*} \\
RF  & QRF  & $\mathbf{-0.022}$ & $1.6 \times 10^{-10}$ & \textbf{*} \\
\bottomrule
\end{tabular}
\end{table}

\begin{figure}[htbp]
    \centering
    \includegraphics[width=\textwidth]{fig_nn_family_comparison.png}
    \caption{Neural network family comparison on the QM9 HOMO--LUMO gap with N=10k: R$^2$ versus $\sigma$ under Gaussian noise on the physicochemical descriptor vector (PDV) representation for a) NN-$\alpha$ family, b) NN-$\beta$ family, and c) RF versus QRF. Heteroscedastic results (not shown) are qualitatively consistent.}
    \label{fig:nn_family_comparison}
\end{figure}

\section{Relative model robustness is largely strategy-independent}
% ---- NOTES: Relative model robustness is largely strategy-independent ----
% Takeaway: clean accuracy does not predict robustness.
% Takeaway: systematic noise costs more than random noise; the shape of the distribution barely matters.
% Takeaway: the best representation for a model, and the order of the models, both change with the noise condition and level.
% - Change: the prose and table are from the earlier study (six strategies, NDS, Kendall's W, SMILES). Rewrite from these bullets.
% Paragraph 1: clean accuracy against robustness (ECFP4, Gaussian).
% - Clean R^2 against AUC_norm: Spearman -0.18 over the 13 base models. Result needed: the latest run gives -0.15; take the reported run.
% - The four plain and Bayesian networks are among the most accurate on clean labels and the four least robust (AUC_norm 0.881 to 0.897).
% - NGBoost is the reverse: the highest AUC_norm on all six representations, the least accurate on four.
% - The networks' AUC_norm spans more across representations than any other model's: the model-level face of the interaction term.
% Paragraph 2: the shape of the distribution.
% - Laplace, Student-$t$ and outlier noise against Gaussian: none of the 63 changes exceeds that pair's replicate spread. Agrees with \citet{Heid2023}.
% Paragraph 3: systematic noise.
% - Grouped-shifted lowered AUC_norm on 77 of 78 base-model and representation pairs, beyond the replicate spread on 50.
% - Censoring gave the largest drop: median AUC_norm 0.813, against 0.951 for the same five pairs under Gaussian noise.
% Paragraph 4: the best representation moves.
% - From Gaussian to grouped-shifted, five of 13 base models change their best representation. Give the AUC_norm gap beside each; some are 0.002.
% Result needed: the fifth model is NN-$\beta$ in the paper and BNN-$\beta$ in the latest run.
% Paragraph 5: ranks move with the level.
% - ECFP4, Gaussian: RF goes from seventh on clean labels to first by level 1.0.
% - AUC_norm weights all levels equally, so it can understate what the pattern of noise costs.
% Table: change in AUC_norm from Gaussian to grouped-shifted, one row per base model, one column per representation, plus the replicate spread.
% Result needed: regenerate the two Additional files cited in the prose (ECFP4 overview; clean R^2 against robustness) on AUC_norm.
% ---- END NOTES ----

NGBoost and SVM showed the smallest degradation in predictive performance across all strategies, while NN-$\beta$ showed the steepest (Table~\ref{tab:nds_ranking}). When sorting models by NDS, their ranks across noise strategies were consistent: Kendall's $W = 0.92$ ($p = 2.7 \times 10^{-8}$; 11 models, 6 strategies), indicating that a model's relative robustness is largely strategy-independent. The same ranking pattern holds on the ECFP4 fingerprint (Additional file~\ref{app:noise:sec:ecfp4_overview}). Robustness is also largely decoupled from baseline predictive performance: across models on PDV under Gaussian noise, NDS clusters near $-0.38$ regardless of baseline R$^2$ (Additional file~\ref{app:noise:sec:clean_r2_vs_robustness}), so a high-performing model is not automatically a more robust one.

\begin{table}[htbp]
\centering
\caption{Noise degradation slope (NDS) by model on PDV (QM9 HOMO--LUMO gap), ranked by mean across six strategies. Lower $|$NDS$|$ (slope closer to zero) indicates greater robustness. Strategy abbreviations: \textbf{Gauss.}\ = Gaussian, \textbf{Quant.}\ = Quantile, \textbf{Thresh.}\ = Threshold, \textbf{Hetero.}\ = Heteroscedastic, \textbf{Val.-P.}\ = Value-Proportional. Bold rows mark the two most robust models; bold cell entries mark each strategy's most robust configuration.}
\label{tab:nds_ranking}
\small
\begin{tabular}{lccccccc}
\toprule
\textbf{Model} & \textbf{Gauss.} & \textbf{Outlier} & \textbf{Quant.} & \textbf{Thresh.} & \textbf{Hetero.} & \textbf{Val.-P.} & \textbf{Mean} \\
\midrule
\textbf{NGBoost}    & $\mathbf{-0.33}$ & $-0.12$ & $\mathbf{-0.29}$ & $\mathbf{-0.65}$ & $\mathbf{-0.18}$ & $\mathbf{-0.58}$ & $\mathbf{-0.36}$ \\
\textbf{SVM}        & $-0.36$ & $\mathbf{-0.11}$ & $-0.30$ & $-0.71$ & $-0.20$ & $-0.63$ & $-0.38$ \\
RF                  & $-0.36$ & $-0.12$ & $-0.30$ & $-0.71$ & $-0.20$ & $-0.63$ & $-0.39$ \\
LightGBM            & $-0.36$ & $-0.13$ & $-0.32$ & $-0.72$ & $-0.20$ & $-0.64$ & $-0.40$ \\
VBLL-$\alpha$       & $-0.38$ & $-0.16$ & $-0.35$ & $-0.68$ & $-0.22$ & $-0.60$ & $-0.40$ \\
XGBoost             & $-0.37$ & $-0.13$ & $-0.33$ & $-0.72$ & $-0.20$ & $-0.65$ & $-0.40$ \\
VBLL-$\beta$        & $-0.39$ & $-0.14$ & $-0.34$ & $-0.70$ & $-0.23$ & $-0.64$ & $-0.41$ \\
BNN-$\alpha$        & $-0.38$ & $-0.15$ & $-0.35$ & $-0.72$ & $-0.21$ & $-0.64$ & $-0.41$ \\
BNN-$\beta$         & $-0.38$ & $-0.15$ & $-0.36$ & $-0.72$ & $-0.21$ & $-0.65$ & $-0.41$ \\
NN-$\alpha$         & $-0.41$ & $-0.16$ & $-0.37$ & $-0.77$ & $-0.23$ & $-0.68$ & $-0.44$ \\
NN-$\beta$          & $-0.49$ & $-0.17$ & $-0.42$ & $-0.88$ & $-0.26$ & $-0.79$ & $-0.50$ \\
\bottomrule
\end{tabular}
\end{table}

The strategies differed not only in predictive performance severity, but also in their ability to discriminate between architectures. The spread of NDS across models differed by noise strategy, with outlier noise having a spread of 0.06 and threshold noise having 0.23, presumably exposing architectural differences that are masked by other types of noise. As seen in Table~\ref{tab:nds_ranking}, NGBoost was the most resistant to noise across strategies. With respect to the Variational Bayesian Last Layer (VBLL) models, VBLL-$\alpha$ showed a similar pattern, outperforming all tree-based models except NGBoost under threshold and value-proportional noise despite ranking below them under Gaussian noise. This suggests that its learned noise variance is most beneficial when label corruption is systematic rather than uniform. This observation with VBLL is not alone; predictive performance on clean data is not well-correlated with noise robustness. In fact, some of the models with the lowest NDS across strategies such as NGBoost and SVM did not perform particularly well on clean data. These models may be more prone to picking up on noise rather than underlying signals in the data. 

% TODO: review, needs manual review
As established by the ANOVA, model architecture and the model--representation interaction term dominate NDS variance. A handful of models stood out as being noise-robust regardless of representation, namely SVM and full BNNs. Both methods rely on inductive biases, SVM on margin maximization and BNN on weight priors. However, other models like RF and NN-$\beta$ show the opposite pattern. These models can be robust to noise, but only when paired with particular representations. The pairings differ by model. 

% TODO: review, needs manual review
While representation may not be a dominant determinant in NDS variance, the model--representation interaction term does have an impact. For SMILES, model choice explains over 91\% of robustness variance, while on PDV, model choice explains 72\%. Across all six noise strategies, PDV configurations produced the shallowest mean NDSs across representations, with its best performance coming with models with strong inductive biases like SVM and full BNN. This suggests that PDVs are preferable when dealing with noisy data. It appears that PDV's compact feature space complements margin-based and Bayesian regularization. Embeddings struggled, particularly when paired with NNs. These patterns are aligned with the importance of the model--representation interaction term, indicating that representation choice modulates noise robustness within a given architecture. 

Taking a closer look at model choice for noise robustness, we ranked all model--representation configurations on QM9 by their Gaussian NDS. Tree-based ensembles dominate the top 10: NGBoost appears five times, XGBoost and RF twice each, and LightGBM once. No neural-network configuration appears the top 10. Rankings are consistent across strategies and the top QM9 configurations maintain their advantage on ADME datasets, though SVM configurations rank higher on these ADME datasets than NGBoost.

\section{The pattern holds across prediction targets}
% ---- NOTES: The pattern holds across prediction targets ----
% Takeaway: on measured labels RF keeps its QM9 standing; NGBoost, the boosted trees and VBLL-$\alpha$ do not.
% Takeaway: grouped-shifted noise is the hardest condition on every dataset.
% - Change: the prose and both figures are from the earlier study (NDS, NGBoost first everywhere). Rewrite from these bullets.
% - Change: the hERG K$_i$ caption says 1,482 molecules; the dataset has 1,415.
% - These labels already carry measurement error, so injected noise adds to an unknown noise floor. Repeat pK$_i$ measurements differ by about 0.54 log units \citep{Kalliokoski2013, Kramer2012}.
% - Clean R^2: the model term leads on all four datasets. Representation takes 6.7 to 16.8% on the assay datasets against 29.2% on QM9.
% Result needed: the AUC_norm decomposition on the assay datasets, before claiming representation has little effect there.
% - Grouped-shifted gives the lowest median AUC_norm on every dataset (Caco-2 0.733, hERG K$_i$ 0.748, LogD 0.829, QM9 0.896).
% - Non-Gaussian shapes: RF and the GP never lose more than 0.031. SVM never loses; NGBoost loses up to 0.15 on Caco-2. So shape matters for SVM and NGBoost only.
% - Ranks: RF never below 7th on any assay dataset. NGBoost falls to 10th on Caco-2. LightGBM and XGBoost sit near the bottom on Caco-2 and hERG K$_i$. VBLL-$\alpha$ is 1st or 2nd on hERG K$_i$ and last on LogD.
% - Change: drop "BNN variants suffer the most"; the BNNs do not fall as a group.
% - Change: "QRF was consistently less robust than RF on every external data set" compared 100-tree RF with 300-tree QRF. Against RF at 300 trees, QRF is lower under grouped-shifted noise on every representation and dataset, and mixed otherwise.
% - The rank range of each model over the six representations, per dataset, goes in an Additional file table; figure (b) carries the ranks in the main text.
% Figures (replace both): (a) AUC_norm per model x condition on ECFP4, one panel per assay dataset; (b) AUC_norm rank per model on each dataset, ECFP4, Gaussian.
% ---- END NOTES ----

\begin{figure*}[htbp]
\centering
\includegraphics[width=\textwidth]{fig_validation_overview.png}
\caption{Noise degradation slope (NDS) heatmaps for three external validation datasets: a) LogD, b) Caco-2 Efflux, and c) hERG K$_i$. Lower (more negative) NDS indicates faster degradation under noise. Black cells indicate baseline R$^2 < 0.3$ for that (model, strategy) pair; ``N/A'' indicates filtered extreme values ($|$NDS$| > 2$). Models whose baseline R$^2$ fell below 0.3 across all strategies are omitted entirely, accounting for the smaller set of rows on hERG K$_i$.}
\label{fig:validation_overview}
\end{figure*}

In addition to working with ``clean'' quantum molecular properties, we tested our findings on experimentally-obtained data, replicating our process of adding artificial noise. Unlike QM9, where $\sigma = 0$ corresponds to a noise-free baseline, the experimental endpoints carry their own measurement noise in both training and test labels; the injected $\sigma$ therefore stacks on top of an unknown noise floor. Per-dataset NDS heatmaps across all model architectures and noise strategies are shown in Figure~\ref{fig:validation_overview}. The ANOVA on these additional datasets confirms that model architecture dominates robustness variance on all three datasets, as seen in Additional file~\ref{app:noise:sec:assay_variance_decomposition}, and that the trends with respect to noise robustness generalize. NGBoost ranks first on both QM9 and external datasets, SVM also transfer quite well. Some models' predictive capabilities are limited on these external datasets, likely due to a combination of smaller training sets, narrower chemical coverage, and the unknown experimental noise in the labels. XGBoost and BNN variants suffer the most here, as seen in Figure~\ref{fig:validation_combined}a.

\begin{figure*}[htbp]
\centering
\includegraphics[width=\textwidth]{fig_validation_combined.png}
\caption{Comparison of noise robustness on QM9 (HOMO--LUMO gap, $N = 10{,}000$) and three external validation datasets: LogD ($N = 5{,}039$), Caco-2 Efflux ($N = 2{,}161$), and hERG K$_i$ ($N = 1{,}482$). a) Model robustness measured by the noise degradation slope (NDS) across all four datasets; lower (more negative) NDS indicates faster degradation under noise. b) QM9 vs external robustness: models with inherent robustness mechanisms (NGBoost, SVM) maintain their relative standing, while ensemble-dependent models (XGBoost) degrade on the external datasets.}
\label{fig:validation_combined}
\end{figure*}

QRF was consistently less robust than RF on every external data set (Additional file~\ref{app:noise:sec:rf_vs_qrf_assay}). As seen with QM9, choice of molecular representation has a minimal effect on noise robustness across all three external datasets, reinforcing the finding that model architecture, not representation, is the main driver of noise robustness. 

\section{Uncertainty under label noise}\label{sec:uq}
% ---- NOTES: Uncertainty under label noise ----
% Takeaway: for models that report uncertainty, does it rise with label noise, and does it point to the corrupted labels?
% - End the existing paragraph with those two questions.
% - One or two sentences of context: earlier comparisons of uncertainty methods on molecules found no method best everywhere \citep{Hirschfeld2020, Scalia2020}. None injected label noise.
% ---- END NOTES ----
\todo[inline, color=yellow!40]{Write the opening of Section~\ref{sec:uq}. Hirschfeld, Scalia and Norinder are cited in no source -- add to refs.bib.}

The sections above showcased the importance of model selection when working with noisy labels. We take this one step further and ask, for the models which produce uncertainty estimates, how those estimates are affected by noise and if they can be used to track noise. 

\subsection{Models with an explicit noise term track prediction error under noise}
% ---- NOTES: Models with an explicit noise term track prediction error under noise ----
% Takeaway: mean predicted uncertainty rises with label noise, but per molecule it does not track the injected noise.
% Takeaway: censoring is the exception, and the sign is negative: the clipped labels get the lowest uncertainty.
% - Change: the prose and table say per-molecule uncertainty tracks injected noise. Their 0.56 (GP, SNS) pools noise levels; within one level the same pair gives -0.04.
% - Mean uncertainty rises with Gaussian noise for every model and representation on QM9. On the assay datasets the exceptions are networks on learned embeddings.
% - Within one noise level, under six of seven conditions, the Spearman correlation between uncertainty and injected noise stays between -0.16 and 0.16 for every model, representation and dataset.
% - Under censoring the noise depends on the label's value. QM9, half the labels clipped, PDV: -0.84 to -0.80 for NGBoost, -0.47 to -0.41 for the GP.
% Table (replaces the one below): per model and condition, QM9, PDV: slope of mean uncertainty against added noise, uncertainty against error, change in ranking corrupted labels.
% Result needed: mean uncertainty at no added noise and at the highest level (PDV, Gaussian, held-out molecules).
% ---- END NOTES ----

\citet{Kolmar2021} found that the mean predicted uncertainty derived from the GPs increases with the amount of label noise in the training data. We go beyond the population level and instead ask if a model's per-sample uncertainty tracks label noise. One key sign that a model is able to handle noise is its ability to track it, represented by the per-sample Spearman correlation between predicted uncertainty and noise magnitude (Table~\ref{tab:top_unc_noise}). While GPs and NGBoost gave the strongest correlations, BNNs were moderate. Learned embeddings (MHG-GNN and mol2vec) struggled, so much so that their correlations dropped to near zero or below on every model tested. 

\begin{table}[htbp]
\centering
\caption{Strongest and weakest model--representation combinations for uncertainty--noise correlation on the QM9 HOMO--LUMO gap (Gaussian noise strategy). Columns: \textbf{Unc-Noise $\rho$} = Spearman correlation between predicted uncertainty and injected noise magnitude $|\epsilon|$ (higher = better noise detection); \textbf{Unc-Error $\rho$} = Spearman correlation between predicted uncertainty and absolute prediction error; \textbf{ECE} = expected calibration error (lower = better-calibrated); \textbf{Cov.\ $1\sigma$} / \textbf{$2\sigma$} = empirical coverage at the $1\sigma$ (68\% target) and $2\sigma$ (95\% target) intervals.}
\label{tab:top_unc_noise}
\small
\begin{tabular}{llrrrrr}
\toprule
\textbf{Model} & \textbf{Rep.} & \textbf{Unc-Noise $\rho$} & \textbf{Unc-Error $\rho$} & \textbf{ECE} & \textbf{Cov. $1\sigma$} & \textbf{Cov. $2\sigma$} \\
\midrule
\textbf{GP}    & \textbf{SNS}    & \textbf{0.56} & 0.23 & 0.14 & 71.2\% & 94.6\% \\
GP             & Morgan          & 0.53 & 0.22 & 0.14 & 71.4\% & 95.4\% \\
\textbf{NGBoost} & \textbf{SNS}  & \textbf{0.47} & 0.16 & 0.13 & 69.3\% & 95.4\% \\
\textbf{NGBoost} & \textbf{PDV}  & \textbf{0.47} & \textbf{0.29} & \textbf{0.11} & 69.0\% & 95.2\% \\
NGBoost        & PDV (binary)    & 0.44 & 0.25 & 0.13 & 69.8\% & 95.4\% \\
NGBoost        & Morgan          & 0.44 & 0.20 & 0.14 & 69.8\% & 95.3\% \\
BNN-$\alpha$   & SNS             & 0.42 & 0.21 & 0.17 & 73.4\% & 95.1\% \\
BNN-$\beta$    & SNS             & 0.40 & 0.21 & 0.18 & 73.5\% & 95.1\% \\
BNN-$\beta$    & Morgan          & 0.38 & 0.20 & 0.18 & 72.1\% & 95.0\% \\
BNN-$\beta$    & PDV             & 0.38 & 0.24 & 0.19 & 72.1\% & 94.0\% \\
\midrule
\multicolumn{7}{l}{\textit{Contrast: lowest correlations (graph/embedding representations)}} \\
\midrule
BNN-$\beta$ & mol2vec & $-0.05$ & 0.02 & 1.33 & 50.6\% & 76.4\% \\
VBLL-$\alpha$ & MHG-GNN & $-0.05$ & 0.09 & 13.41 & 26.6\% & 38.6\% \\
VBLL-$\beta$ & MHG-GNN & $-0.03$ & 0.10 & 15.34 & 31.1\% & 47.1\% \\
VBLL-$\beta$ & mol2vec & $-0.03$ & 0.15 & 23.86 & 44.8\% & 58.0\% \\
\bottomrule
\end{tabular}
\end{table}

\subsection{NGBoost and the GP stay near target coverage while Bayesian variants miscalibrate}
% ---- NOTES: NGBoost and the GP stay near target coverage while Bayesian variants miscalibrate ----
% Takeaway: each model absorbs the large residuals of noisy labels in its own way, and that decides where its uncertainty comes from.
% - Her paragraph covers only the global-noise GP. The heteroscedastic GP's noise term depends on the input.
% - Change: drop "both the aleatoric and epistemic components increased" for VBLL; its aleatoric part is one number for the whole fit.
% - Change: "QRF ... produced poor label predictions" does not hold; QRF is among the more robust models.
% - Change: drop the per-molecule correlations and the citations of Table~\ref{tab:top_unc_noise}, Additional file~\ref{app:noise:sec:uncertainty_metrics_ecfp4} and Figure~\ref{fig:uncertainty_combined}. All three are from the earlier study.
% Result needed (only if the coverage heading stays): coverage at 1 and 2 standard deviations per probabilistic model and noise level.
% ---- END NOTES ----

Samples containing label noise tend to produce large residuals in training, and every model has its own way of handling them. NGBoost contains predicted scales that increase to absorb residual \citep{Duan2020}. The GP contains has a single global observation noise term. It does not absorb individual residuals, so its per-sample uncertainty derives primarily from the posterior variance rather than label noise  \citep{Rasmussen2005, Obrezanova2007}. BNNs treat these large residuals differently; they broaden their weight posteriors, increasing predicted uncertainty without explicitly modeling observation noise \citep{gal2016, kendall2017}. VBLL adds a learned noise variance to the BNN loss \citep{Harrison2024}, but this value is a global scalar, independent of the input. As seen in Table~\ref{tab:top_unc_noise} and Additional file~\ref{app:noise:sec:uncertainty_metrics_ecfp4}, both BNNs and VBLLs achieve moderate correlations which improve when paired with fingerprint or descriptor-based representations. Among models that learn separate distributional parameters during training, GP and NGBoost produced the strongest correlations between uncertainty and label noise, though QRF, which separates parameters using quantile distributions, produced both poor label predictions and uncertainty-noise correlations. At the population level, mean predicted uncertainty increases with artificial noise. Although we expect to see aleatoric uncertainty increase with injected noise, for VBLL both the aleatoric and epistemic components increased (Figure~\ref{fig:uncertainty_combined}b).

\begin{figure}[htbp]
\centering
\includegraphics[width=\textwidth]{fig_uncertainty_combined.png}
\caption{Uncertainty estimates under label noise on the QM9 HOMO--LUMO gap. a) Mean predicted uncertainty versus $\sigma$ for probabilistic models (PDV, Gaussian strategy). b) Aleatoric and epistemic uncertainty components versus $\sigma$ for VBLL models.}
\label{fig:uncertainty_combined}
\end{figure}

\subsection{Uncertainty attribution is itself representation-dependent}
% ---- NOTES: Uncertainty attribution is itself representation-dependent ----
% Takeaway: the variance-head BNNs put added noise into the aleatoric part. QRF puts much of it into the epistemic part, and the representation decides how much.
% Takeaway: the aleatoric part shows that noise rose, not how much. Under censoring, uncertainty shrinks.
% - Models with a per-molecule split: the BNNs with a variance head, the heteroscedastic GP, QRF. The other GPs have one noise value per fit.
% - Measure: rise in mean epistemic divided by rise in mean aleatoric, no added noise to level 1.5, held-out molecules. Near 0: noise went to the aleatoric part.
% - Variance-head BNNs: at most 0.43 everywhere.
% - QRF: 0.63 or above on every representation, lowest on the fingerprints and highest on ChemBERTa and MHG-GNN (ranges in the table).
% - Heteroscedastic GP: at most 0.36 under Gaussian noise. Grouped-shifted raises it, most on ECFP4. A possible reading: a shared offset per scaffold family looks like structure, not noise.
% - Slope of mean aleatoric uncertainty against the added noise, out-of-fold training molecules, ECFP4 and PDV: 0.36 to 1.09. So it compares noise between datasets; it does not estimate a dataset's own error.
% - Change: the paper gives 0.36 to 0.99; the fitted slopes give 1.09. Re-check the paper's source.
% - Censoring narrows the label spread. NGBoost and QRF shrink their total uncertainty with it on every dataset: more confident, less informed.
% Table: epistemic-to-aleatoric ratio, rows model x condition (Gaussian, grouped-shifted), columns the six representations, entries the range over datasets.
% ---- END NOTES ----
\todo[inline, color=yellow!40]{MISSING RESULT: the GP uncertainty decomposition by representation.}

\subsection{Fingerprint representations carry the usable uncertainty signal}
% ---- NOTES: Fingerprint representations carry the usable uncertainty signal ----
% Takeaway: the representation changes how cleanly a model splits its uncertainty, not whether per-molecule uncertainty tracks the noise.
% - Change: the paragraph below says fingerprints and PDV let the GP and NGBoost follow each molecule's noise. No representation gives that signal outside censoring.
% - The direction differs by model: fingerprints give QRF its cleanest split, but give the heteroscedastic GP its least clean split under grouped-shifted noise.
% - Change: mol2vec is not tested; the two learned embeddings are ChemBERTa and MHG-GNN.
% ---- END NOTES ----

The choice of representation has a much stronger impact on per-sample uncertainty tracking than on noise robustness. Fingerprints and physicochemical descriptors provided a clear signal with which models such as GP and NGBoost could separate noise-induced residuals from structural variation. However, learned embeddings such as MHG-GNN and mol2vec did not produce viable uncertainty-noise correlations (Table~\ref{tab:top_unc_noise}). Although these learned embeddings can be effective representations for many tasks, this research suggests that they provide challenges for per-sample noise tracking. However, there is room for further investigation, as only two embeddings were tested.

\section{Conclusions}
% ---- NOTES: Conclusions ----
% Takeaway: architecture governs robustness, systematic noise costs most, RF is the most consistent choice, and no model's per-molecule uncertainty identifies the corrupted labels.
% Takeaway: the representation still mattered, and that opens the next chapter.
% - Change: several paragraphs predate the current results. Her paper's current Conclusion rewrites Paragraphs 2, 4 and 5 and "When working with noisy experimental data"; copy those in.
% - Change: "These results lead to two main conclusions" repeats "We investigated" and announces a second conclusion that no longer holds. Merge or cut.
% - Change: drop NDS, the strategy names, Kendall's W, mol2vec, "QRF significantly less robust than RF", and every claim that NGBoost, GPs or fingerprints track per-molecule label noise.
% - Change: "QM9 served" names only XGBoost as degrading on the assay datasets; LightGBM did too.
% - Uncertainty in one line: mean uncertainty rose with noise; no model's per-molecule uncertainty found the corrupted labels; under censoring the models were most confident on clipped labels.
% - Practical advice: RF was the most consistent across the four datasets; Bayesian networks traded some clean R^2 for robustness.
% Last paragraph, hand-off to Chapter~\ref{chap:flaws}: the model decided robustness, but the representation still mattered. The best representation changed with the noise condition for five of 13 models, and it decided QRF's uncertainty split.
% ---- END NOTES ----

These results lead to two main conclusions. The first is that noise robustness, defined as the rate at which predictive performance degrades with increasing label noise, is primarily determined by the model's training mechanism. The margin maximization of SVM, the weight priors of BNNs, and the ensembling in tree methods all appear to keep those models from fitting to noise. The choice of representation has a stronger influence on predictive performance than noise robustness. The interaction term between representation and model has a stronger impact; some models like SVM and full BNNs remain robust across all representations, while RF and NNs are robust only with certain ones. While the pattern of noise has a strong impact on predictive performance, it does not influence noise robustness rankings between models or representations much.

We investigated how a wide variety of molecular representations and model architectures perform under increasing label noise in QSAR modeling. We used ANOVA variance decomposition across three noise conditions on the QM9 HOMO--LUMO gap task to show that with regards to predictive performance, model architecture and molecular representation account for comparable shares of the variance. However, for robustness to label noise, measured as \aucnorm, the normalized area under the $R^2$ retention curve, model architecture is the dominant factor, while representation explains less than 10\% of the variance.

Threshold and value-proportional noise produce the widest spread of NDS across models; outlier noise barely separates them. Model rankings are highly concordant across strategies (Kendall's $W = 0.92$), suggesting that a model's robustness to noise is independent of type of noise. Predictive performance on clean data is not correlated with noise robustness; NGBoost and SVM, the most noise-robust models, did not perform particularly well on clean data. Embeddings (MHG-GNN, mol2vec) struggled with regards to noise robustness, especially with neural network architectures. On the flip side, QRF was significantly less noise-robust than RF across datasets, suggesting that quantile regression may be overfitting to noisy labels rather than absorbing them.

% TODO: review, needs manual review
The second conclusion concerns whether a model's uncertainty estimates track per-sample label noise. Some models struggle, particularly those like BNNs whose uncertainty comes from a posterior over weights or QRF which does produce a full quantile, but is unable to successfully track noise. 
Models that fit a separate, per-sample scale or observation-noise parameter during training are able to track that label noise more effectively. NGBoost identifies noisy labels through a per-sample predicted scale that absorbs large residuals. The GP does not: it has a single global observation noise term, so its per-sample uncertainty comes from the posterior variance, which depends on where a molecule sits relative to the training molecules and not on its label. Among representations, uncertainty tracking collapses with the learned embeddings MHG-GNN and mol2vec, while fingerprints and descriptors produce a viable signal. Neither strong predictive performance nor noise robustness implies noise tracking; again it is up to the model's architecture and how it handles loss. 

Although the choice of representation carries less weight than that of models, we observed an interesting phenomenon. PDVs produce the strongest predictive performance on clean data, but are the least noise-resistant. Embeddings, particularly mol2vec, are the most robust to noise yet perform the worst noise tracking. At the end of the day, the choice of both model and representation comes down to your problem, including the amount of noise in your data, your objectives, and your compute limits. 

QM9 served as a relatively clean baseline; it is computationally-derived with negligible measurement noise. The validation datasets, LogD, Caco-2 \citep{openadmet}, and hERG Ki \citep{Zdrazil2023}, tested are smaller and cover a narrow range of the chemical and biological space, and have an unknown inherent amount of noise from experimentally-determined values, on top of the added artificial noise. These validation datasets supported the QM9 findings, with XGBoost the notable exception, degrading on the external datasets. Future work should extend this type of noise robustness benchmarking to classification tasks and larger experimental datasets with better-known uncertainty estimates.

When working with noisy experimental data in QSAR settings, it is important to keep in mind that noise robustness and noise-aware uncertainty estimation are primarily driven by the model's mechanism of training, and to a certain extent the choice of molecular representation, usually with respect to the model. While Bayesian transformations on neural networks improve noise robustness, they do not produce uncertainty estimates that correlate particularly strongly with noise. Models like NGBoost and GPs, which include a learned noise or scale parameter in the training objective, are the most robust to noise across both various types of noise and datasets and often produce uncertainty estimates which track per-sample label noise. Pairing these models with fingerprint-based representations provides additional benefits in restricting, and potentially detecting and mitigating label noise. 

\chapter{The Flaws of the Similarity Assumption}\label{chap:flaws}
% ---- TAKEAWAYS ----
% - General-purpose representations are built for all of chemistry. On one dataset, most of their features barely vary.
% - Some representations give different molecules exactly the same vector, mostly stereoisomers. No model can tell those molecules apart.
% - A similarity or distance compares every feature, so features irrelevant to the task dilute it. Different representations also disagree about which molecules are close.
% - Methods that judge a molecule by its neighbours did worst under label noise: LNL methods that use a distance, and noise detection by neighbours.
% - What is needed is a compact set of features chosen for the task from several representations. Chapter 5 builds it: the hybrid.
% ---- END TAKEAWAYS ----
% ---- NOTES: The Flaws of the Similarity Assumption ----
% Paragraph 1: a short preamble, opening with \textit{Preamble.} as in Chapter~\ref{chap:robustness}.
% - Section \ref{sec:lnl_motivation} reuses the exploratory learning-with-noisy-labels study from the DPhil confirmation report, "Investigating the Effects of Noise on Molecular Property Prediction" (unpublished, January 2026).
% - Cite it once here as \citet{Punt2026confirmation}.
% ---- END NOTES ----

\section{Introduction}
% ---- NOTES: Introduction ----
% Takeaway: Chapter 3 used the representation to predict. This chapter asks what it does to methods that compare molecules with each other.
% Takeaway: the chapter introduces no method. It shows flaws with small measurements, and closes on three research questions.
%
% Paragraph 1: the bridge from Chapter 3
% - Chapter~\ref{chap:robustness} in one sentence: model architecture decided how much performance survived label noise; the representation mattered for performance but little for robustness.
% - Change: quote the same pair of variance shares Chapter 3 prints: 48.7% and 8.2%.
% - Turn: methods that compare molecules have no model in between to down-weight features that do not matter. For them the representation is the whole method.
%
% Paragraph 2: the similarity assumption, restated as what is tested
% - Restate it in one sentence; Chapter 2, "Molecular similarity and activity cliffs", defines it.
% - A similarity value depends on the representation and the coefficient. The coefficient matters little \citep{Bajusz2015}, so this chapter questions the representation.
%
% Paragraph 3: related work, short, not a review
% - MODI scores whether a dataset can be modelled from its neighbours \citep{Golbraikh2014}. It scores dataset and representation together, never the representation alone.
% - Fingerprints agree with each other more than expected \citep{Riniker2013}, and learned embeddings mostly do no better than ECFP \citep{Praski2025} (preprint; say so).
% - Representations disagree about which pairs are activity cliffs, so choosing one defines what a cliff is \citep{DabrowskiTumanski2026} (preprint; say so).
% - The gap: all of these describe one fixed feature set for every target. None asks which features matter for the property at hand.
%
% Paragraph 4: what the chapter does, in order
% - A motivation chapter with no new method. Each flaw is shown with one small measurement.
% - Order: flaws in the representation, then in similarity computed on it, then in methods that use that similarity (the exploratory LNL study, Section~\ref{sec:lnl_motivation}).
% - Framing: the argument is about general-purpose representations used on one narrow dataset. ECFP4 is not called a bad fingerprint; it was built to cover all of chemistry.
%
% Paragraph 5: research questions, as an itemised list closing the section
% - Q1: on a given dataset, how many of a standard representation's features vary, and how many molecules does it map to the same vector?
% - Q2: how closely does distance in each representation follow differences in the property?
% - Q3: do methods that rely on that distance do better or worse than methods that do not?
% Result needed: one sentence per research question previewing its answer, from the full runs.
%
% Figure: the chapter's argument as a chain of boxes. "Representation built for all of chemistry" -> "distance dominated by features irrelevant to this task" -> "methods that use neighbours usually do worse" -> "select the task-relevant features" (Chapter 5).
% ---- END NOTES ----
\todo[inline, color=yellow!40]{Write 4.1: bridge from Chapter 3, research questions, modelability literature.}


\section{Flaws in standard base representations}
% ---- NOTES: Flaws in standard base representations ----
% Takeaway: standard representations are built before any dataset or task is known, so on one narrow dataset much of each describes chemistry that is absent.
% Takeaway: different representations disagree about which molecules are close, so each one sees only part of the chemistry.
%
% Paragraph 1: why the chapter starts with the representation
% - Similarity is computed on a representation, so any flaw in it is carried into every similarity value.
% - One dataset usually covers a few chemical series. Much of a general-purpose representation then describes chemistry that is absent.
% - Which features matter also changes from task to task.
% - Each flaw gets one simple statistic on the full feature matrix. No model is fitted. Small molecules only; families as defined in Chapter 2.
%
% Paragraph 2: closing point, after the learned-embeddings point
% - Different representations disagree about which molecules are close. Each sees only part of the chemistry, which is the reason to combine them (Chapter 5).
% - Measure: overlap of each molecule's five nearest neighbours between two representations, or the Spearman correlation between their pairwise distances.
% Result needed: the least and most alike pair of representations on hERG Ki. The agreement for every pair and dataset belongs in an Additional file.
%
% Table: one row per small-molecule representation family. Columns: what it encodes, share of features that barely vary, molecules it cannot tell apart, whether its features have a chemical name. Short words only; the counts stay in the two tables below.
% Figure: one schematic. A molecule mapped to a long feature vector, with three small panels: most features constant; two stereoisomers with the same vector; an embedding feature labelled only by an index.
% ---- END NOTES ----


\subsection{Most features do not vary on a given dataset}
% ---- NOTES: Most features do not vary on a given dataset ----
% Takeaway: on one dataset, most bits of a 2048-bit hashed fingerprint are near-constant. The learned embeddings counted have none.
%
% Paragraph 1: the argument
% - ECFP4 folded to 2048 bits is built to describe any molecule \citep{rogers2010}. On a dataset of a few chemical series, most bits are off in almost every molecule.
% - Scope: this counts only how many features vary. Whether the varying ones matter for the task is Chapter~\ref{chap:method}'s question.
%
% Paragraph 2: what is counted
% - Constant feature: one value in every molecule. Near-constant: one value in at least 95% of molecules (includes the constant ones).
% - Datasets: hERG Ki, Caco-2, LogD and QM9. The other small-molecule datasets belong in an Additional file.
%
% Paragraph 3: the numbers, as a pattern
% - Hashed substructure fingerprints have the most near-constant bits; the two learned embeddings counted have none (on these datasets only).
% - Example: ECFP4 on hERG Ki (1,412 molecules), 1,815 of 2,048 features near-constant. MolFormer (768 features) and mol2vec (300 features): 0 (preliminary).
% - The near-constant count is the one that matters for distances; the next section says why.
% - One sentence: after folding, one bit can be shared by unrelated substructures (Chapter 2), so a chosen bit is harder to read.
% - Close: if too few bits vary, different molecules can get identical vectors.
% Result needed: constant and near-constant counts for the four datasets above, on the full-run feature matrices.
% Result needed: Uni-Mol v2, once its features are rebuilt (every batch failed on hERG Ki and LogD, so every row is zero).
%
% Table (label tab:constant_features): one row per dataset, one column per representation. Each entry reads "constant / near-constant / total features". A line above says what one entry counts.
% - MACCS: 167 features; one line below the table says RDKit adds an unused bit 0 that is always off.
% - Preliminary hERG Ki row: ECFP4 105 / 1,815 / 2,048. Topological torsion 1,012 / 1,747 / 2,048. Atom pair 201 / 977 / 2,048. MACCS 24 / 60 / 167. mordred 257 / 351 / 1,613. PDV 21 / 55 / 200. RDKit fingerprint 0 / 18 / 2,048. MolFormer 0 / 0 / 768. mol2vec 0 / 0 / 300.
% ---- END NOTES ----
\todo[inline, color=yellow!40]{4.2.1: constant and near-constant feature counts (Table 4.1).}




\subsection{Molecules the representation cannot tell apart}
% ---- NOTES: Molecules the representation cannot tell apart ----
% Takeaway: some representations give two different molecules exactly the same vector, mostly stereoisomers.
% Takeaway: every model must then give both one prediction. Only a change of representation can remove that error.
%
% Paragraph 1: the flaw
% - Two molecules with one vector get one prediction from every model. If their labels differ, part of the error is fixed before training.
% - More data, tuning or a better model cannot remove it.
%
% Paragraph 2: why stereoisomers tie
% - Most common fingerprints do not encode stereochemistry \citep{orsi2024map4c}. The Morgan fingerprints here use the library default, chirality off. PDV and Mordred also ignore it.
% - Stereoisomers can differ in measured activity \citep{Nguyen2006}.
% - Learned embeddings are not immune: SMI-TED ties 52 hERG Ki molecules, all stereo-only (preliminary).
%
% Paragraph 3: the measurement
% - Group molecules with exactly identical vectors. Count separately: the same molecule twice (dataset fault), a row of zeros (featurisation fault), two different molecules (the representation's fault), and of those, stereo-only.
% - Best score any model could reach: each tied group gets its mean label. $\mathrm{RMSE}_{\min} = \sqrt{\frac{1}{n}\sum_{G}\sum_{i \in G}(y_i - \bar{y}_G)^2}$, in label units.
% - Also the largest label spread inside one group. Scale: repeat pKi measurements differ by about 0.54 log units \citep{Kramer2012}.
% - Datasets: hERG Ki (1,412 molecules), LogD (5,028) and Caco-2 (2,158).
%
% Paragraph 4: what was found, as a pattern (preliminary)
% - Every tie is a merge by the representation; none is a repeated molecule or a row of zeros.
% - ECFP4, atom pair, PDV and Mordred tie almost only stereoisomers. MACCS ties the most, and on hERG Ki most of its ties are not stereo-only.
% - 3D representations and most learned embeddings tie none. Joining every representation side by side leaves no molecule tied.
% Result needed: ChemBERT's tie count on hERG Ki and LogD, on features rebuilt from the current code.
%
% Paragraph 5: consequence
% - A tied pair with different labels has infinite SALI \citep{Guha2008}: a cliff made by the representation, not the chemistry.
% - A neighbour method sees that molecule at distance 0 with a different label, and may flag a correct label as noise.
%
% Table: one row per dataset x representation that ties any molecule. Columns: molecules tied (count and %), stereo-only, best reachable RMSE and $R^2$, largest spread (label units).
% - Preliminary hERG Ki rows (label standard deviation 0.915 pKi): ECFP4 56 tied (4.0%), all stereo-only, best RMSE 0.041 pKi, largest spread 0.57 pKi. MACCS 377 tied (26.7%), 30 stereo-only, best RMSE 0.134 pKi, largest spread 0.78 pKi.
% Result needed: every row on the final representation matrices.
% Figure: one pair of hERG Ki stereoisomers, their identical ECFP4 bits and their two pKi values. Caption: the prediction every model must give is their mean.
% ---- END NOTES ----
\todo[inline, color=yellow!40]{4.2.3: shared-vector census (herg\_ki, logd, caco2).}


\paragraph{Learned embeddings cannot be read chemically.}
% ---- NOTES: Learned embeddings cannot be read chemically. ----
% Takeaway: a descriptor feature has a chemical name; an embedding feature does not. A short point, 2-3 sentences.
% - A blueprint (Chapter 5) can name chosen descriptors or substructure bits. Chosen embedding features stay unnamed.
% - Change: typos: "PDV's" should be "PDV", "it's" should be "its", "useful informative" should be "useful information".
% ---- END NOTES ----

It is worth noting that some of these features are more straightforward than others. For example, PDV's carry very clear physicochemical properties such that if you examined one of it's features, you'd know exactly the type of chemical information it contains. Other representations are frozen embeddings, or the last layer of a network. While these features may be found to contain useful informative for the prediction task at hand, they do not tell us what type of chemical information, if any, they relate to.


\section{Flaws in similarity and other pairwise computations}
% ---- NOTES: Flaws in similarity and other pairwise computations ----
% Takeaway: a similarity or distance compares every feature and cannot tell which ones matter for the property. Features irrelevant to the task dilute it, and their share grows linearly with their number.
% Takeaway: the most common choice, ECFP4 with Tanimoto, cannot tell many molecules apart.
%
% Paragraph 1: the usual way to compare two small molecules (her paragraph below)
% - Sources for her sentence: Tanimoto is the most popular coefficient for binary fingerprints \citep{LopezPerez2024}; ECFP4 is the most popular fingerprint \citep{Capecchi2020}.
% - Change: say "the most common choice", not "the most standard way". No source names the pair as the default.
% - Her "ECFPs do not individually distinguish each molecule": in 96,456 human metabolites, 72.9% had an ECFP4 nearest neighbour at Tanimoto distance 0 \citep{Capecchi2020}. Those are distinct structures even ignoring stereochemistry.
%
% Paragraph 2: a fixed measure compares every feature (the main point)
% - A similarity or distance compares every feature, and cannot tell which matter for the property.
% - Split features into task-relevant $R$ and the rest $I$: $d^2(\mathbf{x},\mathbf{y}) = \sum_{j\in R}(x_j-y_j)^2 + \sum_{j\in I}(x_j-y_j)^2$.
% - Each irrelevant feature that varies adds one term. The irrelevant part grows linearly with $|I|$, not exponentially. The share carrying signal is roughly $|R|/(|R|+|I|)$.
% - Constant features: bits off in both molecules are ignored by Tanimoto (defined in Chapter 2), and add 0 to Euclidean distance.
% - Bits on in nearly every molecule push every Tanimoto similarity toward 1 and narrow its range.
% - After standardising, a rare bit takes large values on the few molecules that have it, so it can outweigh many common bits. This is why the near-constant count matters.
% - Distance concentration (Chapter 2) makes this worse in high dimension; cross-reference only.
%
% Paragraph 3: errors compound
% - An error in the representation gives a wrong similarity, then wrong neighbours, then a wrong prediction, noise flag or applicability-domain call.
% - Part of what looks like an activity cliff may be the measure calling two molecules similar on features that do not matter.
% Result needed: molecules in activity cliffs per dataset and representation, under the rule Chapter 2 names.
%
% Paragraph 4: does distance follow the property
% - All pairs: rank correlation between pair distance and absolute property difference (the neighbourhood-behaviour plot of \citet{Patterson1996}, as one number).
% - Nearest neighbours only: accuracy of predicting each molecule's activity from its five nearest neighbours.
% - These are two questions. Neighbours can carry signal when the all-pairs correlation is near zero; say which one each claim is about.
% Result needed: both numbers, per dataset and representation.
%
% Paragraph 5: prior work and what follows
% - Choosing descriptors for a nearest-neighbour model of one endpoint has been done \citep{Zheng2000, Itskowitz2005}. The hybrid selects across several representations.
% - "Similarity in a hybrid" (Chapter 6) tests whether a feature set chosen for the task gives a better similarity.
%
% Figure: absolute activity difference against ECFP4 Tanimoto similarity for sampled hERG Ki pairs, binned, beside the same pairs under a second representation.
% ---- END NOTES ----

% TODO: need sources for this
Within small molecule research, the most standard way to measure how similar two molecules are is the Tanimoto coefficient between two ECFPs representations of those molecules \citep{Capecchi2020, LopezPerez2024}. There's a handful of problems with this approach. Firstly, ECFPs do not indvidually distinguish each molecule.

\subsubsection{Distance metrics for continuous molecular descriptors}
% ---- NOTES: Distance metrics for continuous molecular descriptors ----
% Takeaway: defines the distances used on continuous features in the thesis. Standardising does not give every feature an equal say under every metric.
%
% Corrections to her list and closing sentence:
% - Mahalanobis needs a formula: $d(x,y)=\sqrt{(x-y)^\top S^{-1}(x-y)}$, $S$ the feature covariance on the training molecules. Cite \cite{Mahalanobis1936}.
% - Bray-Curtis is defined only for non-negative features. The comparison shifted the matrix so its smallest value is zero; say so. Cite \cite{BrayCurtis1957}.
% - Standardising equalises scales for Euclidean, Manhattan and Chebyshev only. How it acts on each other metric belongs in an Additional file.
% - Equal say is also the cost argued above: an irrelevant feature counts as much as a relevant one.
% ---- END NOTES ----

The choice of distance, or similarity metric has a strong impact on performance of distance-based noise estimation methods. We evaluate eight distance metrics appropriate for continuous molecular descriptors:

\begin{itemize}
    \item Euclidean ($L_2$): Standard geometric distance, $d(x, y) = \sqrt{\sum_i (x_i - y_i)^2}$
    \item Manhattan ($L_1$): Sum of absolute differences, $d(x, y) = \sum_i |x_i - y_i|$, more robust to outlier features
    \item Cosine: Angle-based similarity, $d(x, y) = 1 - \frac{x \cdot y}{\|x\| \|y\|}$, invariant to magnitude
    \item Chebyshev ($L_\infty$): Maximum absolute difference, $d(x, y) = \max_i |x_i - y_i|$
    \item Canberra: Weighted $L_1$ normalized by magnitude, $d(x, y) = \sum_i \frac{|x_i - y_i|}{|x_i| + |y_i|}$
    \item Correlation: $d(x, y) = 1 - \rho(x, y)$, where $\rho$ is the Pearson correlation
    \item Bray-Curtis: Ecological dissimilarity measure, $d(x, y) = \frac{\sum_i |x_i - y_i|}{\sum_i |x_i + y_i|}$
    \item Mahalanobis: Covariance-adjusted distance accounting for feature correlations
\end{itemize}

All features are standardised (zero mean, unit variance) before distance computation to ensure equal contribution across features with different relative scales.

\todo[inline, color=yellow!40]{4.3: distance against activity difference, local against global, cliffs. [DIAGRAM D4.1]}


\section{Consequences for methods that rely on similarity}
% ---- NOTES: Consequences for methods that rely on similarity ----
% Takeaway: any method that uses nearby molecules as evidence inherits both flaws.
% Takeaway: changing the representation can fix wrong neighbours; it cannot fix real activity cliffs.
%
% - One neighbour set, from one representation and one distance, feeds several methods (listed under "Other standard uses of similarity").
% - Two claims, kept apart. (a) The similarity principle is false for some pairs: activity cliffs \citep{Maggiora2006, Stumpfe2012}. (b) The standard representation and distance pick the wrong neighbours even where it holds.
% - Only (b) can be fixed by changing the representation. That is why the thesis changes the representation, not the downstream methods.
% - The evidence: the exploratory LNL study, then noise detection by neighbours.
% ---- END NOTES ----

\subsection{The exploratory learning-with-noisy-labels study}\label{sec:lnl_motivation}
% ---- NOTES: The exploratory learning-with-noisy-labels study ----
% Takeaway: in the confirmation report's LNL study, all 20 configurations with the lowest noise retention used a distance between molecules.
% Takeaway: the methods that held up estimated noise from the model's own loss or uncertainty. None compared a label with its neighbours' labels.
%
% Paragraph 1: her first paragraph
% - Open by saying this reuses the confirmation report's study; cite \citet{Punt2026confirmation} once.
% - Change: LNL and the five groups of \citet{Song2022} are defined in Section~\ref{sec:lnl_background}. Replace the five-group list with one clause pointing there, and delete the second \citet{Song2022} sentence ("There are a wide variety...").
% - Change: \citet{Cortes2015} found the other algorithms degraded smoothly and linearly with noise, with gradient boosting worst. Replace "no single algorithm was overall best", drop "10 noise levels", and say "including QSAR \citep{Cortes2015}".
% - One sentence: most LNL methods were built for image classification with flipped labels \citep{Song2022}; the labels here are continuous.
%
% Paragraph 2: what the confirmation report tested
% - Change: "In this chapter, we explore" becomes "The confirmation report tested"; "classical noise mitigation techniques" becomes "the 12 LNL methods".
% - QM9 HOMO-LUMO gap, 10,000 molecules, scaffold splits. Gaussian noise on training labels only, standard deviation 0 to 1.0.
% - Five representations (ECFP4, Sort \& Slice, PDV, canonical and randomised SMILES). 12 LNL methods adapted to regression, 221 configurations.
% - Distances: Tanimoto on ECFP4, Euclidean on PDV, Mahalanobis. Change: list Manhattan and cosine only if the run records show they were used.
% - Four of the 12 remove suspected molecules and retrain; these are noise mitigation. The other eight are noise-robust training.
% - Change: move the list of methods that held up after the distance result.
%
% Paragraph 3: the distance result (her paragraph, the core)
% - Define retention first: $R^2$ under noise divided by clean $R^2$, as a percentage.
% - Change: "significantly worse" becomes "did worse"; drop the ANOVA sentence (no effect size, and the thesis reports no significance tests); fix the doubled full stop.
% - The failure is in judging a label by its neighbours' labels, not in all similarity: the heteroscedastic GP compares molecules through its kernel and held up.
%
% Paragraph 4: what held up
% - DivideMix kept 86.7% \citep{li2020dividemix}, small-loss selection 84.2% \citep{Han2018}, heteroscedastic GP 79.0% \citep{liu2020}. All estimate noise from the model's own loss or uncertainty.
% - Caveats: one dataset, one property, Gaussian noise on every label; and the distance step was added to methods designed without one. Chapter~\ref{chap:similarity} repeats the question on more datasets and noise types.
%
% Paragraph 5: why the distance failed (her two last paragraphs)
% - Two possible reasons: one representation carries too little information; a sparse one carries many features irrelevant to the dataset. The hybrid addresses the second. The evidence does not separate them; say so once.
% - Change: MHG-GNN does not "capture 3D structural features"; it is built from the 2D graph \citep{kishimoto2023}. Say "encode the molecular graph".
% - Change: "exponential growth in the number of irrelevant pairwise interactions" becomes: each uninformative feature that varies adds one term to the distance, so their share grows linearly with their number.
% ---- END NOTES ----

Numerous studies have explored Learning with Noisy Labels (LNL) in the broader context of ML, as well as specifically for QSAR models. \citet{Song2022} provide a comprehensive review of Deep LNL, identifying five main approaches: (i) modifying architectures to better handle noise in training; (ii) employing regularization to reduce overfitting to noise; (iii) adjusting loss functions to be less sensitive to outliers or label flips; (iv) post-processing to adjust the weights for noisy data; and (v) sample selection to select presumably clean data to pre-train. \citet{Cortes2015} examined 12 different algorithms across 12 datasets and 10 noise levels, finding that no single algorithm was overall best at maintaining predictive performance in the presence of noise; certain algorithms perform better or worse at different levels of noise. This is a recurring theme in LNL research: there is no universal solution to handling noise with experimental endpoints.
% TODO make this flow
There are a wide variety of approaches to LNL, spanning sample selection, meta-learning, curriculum learning, and iterative refinement \citep{Song2022}. Cheminformatics data presents additional challenges in the form of high-dimensional molecular representations and dealing with a wide variety of chemical domains with various similarities and relationships. In this chapter, we explore a set of 12 LNL methods in the context of QSAR modeling. We examine the impact of uncertainty quantification, molecular distance metrics, and alternative loss functions in the hopes of improving performance under label noise. We tested 221 total configurations across LNL methods, loss functions, distance metrics, and molecular representations. Notably, we found that molecular representation choice explains even less variance under LNL methods. We also identified the following LNL techniques which showed promise on the tasks at hand: heteroscedastic GP, curriculum learning, and sample selection methods. The work was meant to connect the representation and uncertainty quantification insights found in Chapter~\ref{chap:robustness} with more classical noise mitigation techniques as a means of determining broader research directions. 

One of the most prominent outcomes was that distance-augmented methods performed significantly worse than ones that did not incorporate any sense of distance or molecular similarity. An ANOVA test was run on the different adapted factors and shows that distance metric usage is the dominant factor in variance in $R^2$.. When ranking the performance of methods, all 20 worst-performing configurations with respect to noise robustness used distance metrics, not a single one achieved both high baseline ($R^2 > 0.6$) and high noise retention ($> 80\%$). This pattern was observed across standard distance types (Euclidean, Manhattan, Cosine, Mahalanobis) and a chemistry-aware metric (Tanimoto coefficient). 

Building upon the exploratory study above, we decided to investigate alternative featurization techniques. There are multiple possible explanations behind the poor performance of distance metrics, first being that a singular molecular representation lacks sufficient information for distance calculations, with and possibly even without noise. These representations all encode different types of molecular information, PDVs encode global molecular properties, ECFPs contain substructure information, and MPNNs such as MHG-GNNs capture 3D structural features. These single sources of data may not be sufficient for robust similarity estimation. Additionally, some of these representations are sparse and contain a large number of features that are not relevant to molecules of a given domain. 

The main issue we identified here is that when a large proportion of dimensions are uninformative, similarity measures become dominated by those uninformative features, causing exponential growth in the number of irrelevant pairwise interactions as dimensionality increases. This is what we sought to address by developing a representation with informative features and low dimensionality. We do so by creating a budget-constrained hybrid representation, which takes high-importance features from a variety of sources. These hybrid representations concatenate the top-ranked features by importance from a selection of base representations. 


\subsection{Noise detection by neighbours}
% ---- NOTES: Noise detection by neighbours ----
% Takeaway: judging a label by its neighbours' labels usually made the model worse. Of 860 cleaning configurations, 8 helped at both noise levels.
% Takeaway: it needs the right neighbours and clean neighbour labels at once. Under Gaussian noise on every label, the second fails.
%
% Paragraph 1: the idea (new paragraph, before her prose)
% - A label far from its neighbours' labels is treated as probably wrong. Origin: editing a nearest-neighbour training set \citep{Wilson1972}; adapted to regression by \citet{Martin2021}.
% - Source needed: a QSAR source for removing or down-weighting a molecule whose label disagrees with its neighbours'.
% - It needs the distance to pick the right neighbours and those neighbours' labels to be clean. A large gap can mean a noisy label, a noisy neighbour, or wrong neighbours; the score cannot tell which. Activity cliffs also look like noise.
%
% Paragraph 2: what the confirmation report measured \citep{Punt2026confirmation}
% - 5,000 QM9 molecules, HOMO-LUMO gap, scaffold split, Gaussian noise at sigma 0.3 and 0.6. Each method judged by the recovery rate in $R^2$ of a random forest on ECFP4 (defined in Chapter~\ref{chap:similarity}).
% - Scores compared each label with its $k$ nearest training molecules' labels. Cite only the single-representation rows; the hybrid rows used an allocation method no longer in the thesis.
% - This is not the 10,000-molecule removal study in "Sample removal methods"; keep their numbers apart.
% - Used as sample weights, mean recovery was negative for every representation at both noise levels (from -0.1% to -7.3%).
% - Of 860 cleaning configurations, 8 had positive recovery at both noise levels.
% - It did not always fail: the best single configuration at sigma 0.6 was neighbour-based (PDV, Euclidean, +17.2%). Say "usually made things worse".
% - Caveat: every training label was noisy, so the neighbour mean is itself a noisy target.
%
% Paragraph 3: her two paragraphs below
% - Change: "Sampling weighting methods generally proved effective" is stronger than the report's table: positive recovery in 14.2% of configurations at sigma 0.3 and 4.7% at 0.6.
% - Change: the hybrid paragraph is out of date. That hybrid (PDV plus MHG-GNN, 100 features) used greedy allocation, now cut; the current search sweeps 10 to 2,000 features. "Consistently positive recovery" is also too strong: neighbour-based mean recovery was -0.9% and -3.0%.
% - Close: Chapter~\ref{chap:similarity} repeats this with neighbours found in a hybrid and under other noise types.
% Result needed: where removal by distance to neighbours ranks among the twelve LNL methods, per noise level.
% ---- END NOTES ----

Overall, we see in agreement with conclusions of Section~\ref{sec:lnl_motivation} on the use of distance metrics in LNL methods, the concept of using a sample's neighbors to determine it's ground truth is flawed. It is likely all labels are too uniformly noisy for a methodology like this. Uncertainty-based methods show potential at lower noise, but their estimation performance degrades with higher noise. Sampling weighting methods generally proved effective, and one estimator using exponential sample weighting produced very strong recovery across both noise levels.

Hybrid representations had mixed results. For sample removal, hybrids showed advantages at high noise, but struggled at low noise. They showed consistently positive recovery rates for neighbor-based estimation, which roughly falls in line with the idea that a low-dimension, high-quality information representation benefits neighbor-based reasoning. The feature budget in this experiment was restricted to 50-100 features, which may limit its advantages. This range was successful at standard molecular property prediction, but potentially not for certain noise mitigation or estimation tasks. Distance-based methods remain challenging under Gaussian noise regardless of representation, suggesting the core issue is the noise distribution rather than representation quality alone. This merits an investigation into mitigation of noise from other distributions or noise injection strategies.


\paragraph{Other standard uses of similarity.}
% ---- NOTES: Other standard uses of similarity ----
% Takeaway: the same neighbours are used well beyond noise cleaning, and each use inherits the flaws above. One sentence per use.
% - Nearest-neighbour prediction: its accuracy is a direct test of the distance.
% - Applicability domain \citep{Tropsha2010}: if irrelevant features dominate the distance, it measures novelty, not how likely the model is to be wrong.
% - Activity-cliff detection \citep{vanTilborg2022}: which pairs count depends on the distance.
% - Clustering for data splits by Tanimoto similarity \citep{Butina1999}.
% - Federated learning: each site's prediction is weighted by the query's distance to that site's data \citep{Hanser2025}.
% ---- END NOTES ----

\section{Conclusions: what a representation would need}
% ---- NOTES: Conclusions: what a representation would need ----
% Takeaway: answer each research question, then list what a representation would need. No new evidence.
% Takeaway: the answer is the hybrid: task-chosen, compact features from several representations, usable as model input and for similarity.
% Takeaway: open with what the chapter adds: per dataset, how many of each representation's features vary and how many molecules it merges. The modelability work cited in the Introduction does not report these.
%
% Paragraph 1: answers, one sentence each
% - Q1: many features of a general-purpose fingerprint barely vary on one dataset, and several representations merge different molecules.
% - Q2: a distance sums over all features, so it is dominated by the ones that do not matter.
% - Q3: methods that judge a label by its neighbours did worst under label noise, and usually made the model worse.
% Result needed: the share of near-constant features, the count of tied molecules and the distance-property correlation from the full runs.
%
% Paragraph 2: what a representation would need (one sentence each)
% - Features chosen for the task, from the training labels of each fold.
% - Compact: features that do not vary or do not matter are removed; the number kept is searched, not fixed.
% - Drawn from several representations. Plain concatenation meets this but not the previous one (Chapter 5).
% - Usable as a model input and as a space for similarity.
% - Traceable to its representation, and read chemically where the base feature can be.
%
% Paragraph 3: what feature selection cannot fix
% - Molecules identical in every base representation stay identical. A chosen hashed bit still merges every substructure folded into it. Selection needs enough labelled molecules.
%
% Paragraph 4: bridge to Chapter 5
% - The hypothesis: ECFP4 with Tanimoto and its variants carry many features irrelevant to the task, and methods built on it, such as noise detection and mitigation, suffer.
% - The first version selected task-relevant features from several representations only to compute similarity. They also worked as ordinary model inputs, so the selection became the method: the adaptive hybrid representation.
% - Questions left for Chapter 5, one clause each: how to score a feature (importance), and how to split the budget across representations (allocation methods).
% ---- END NOTES ----

% Similarity background
% TODO: find compass file on this

% Just need to formalise this

% TODO: Go back to original LNL experiments (explained in confirmation report, find this) - copy over the evidence that similarity-based approaches just automatically fail
 % Context for similarity and my motivating hypothesis: I think that molecular similarity as we often do it (ECFP4/Tanimoto or other variations of that) is flawed and applications that this concept of similarity (in my case noise detection and mitigation) suffer as a result
 %    - Take a look at ECFP4 as an example. Let's say you're working with molecules in a particular domain, how many of those features are going to useless? First, you can get numbers on in a given dataset how many features have no difference between the molecules. Some will differ, but will have nothing to do with the task at hand. There's definitely useful features in that set of 2048, but they're probably quite sparse, and when you take a similarity metrix that does a vector by vector comparison, and some of your features are irrelevant, you have exponentially more irrelevant pieces of information (is this true? phrase this better). I thought a much more efficient approach would be to select useful features for that task (I define task as some kind of supervised model fit on a training set) and you specify how many features you want, and you use that to calculate similarity. Now I implemented a proof-of-concept version of this and realised that these customised feature sets (hybrids) work as standard inputs to the model, irregardless of similarity. Then the main question was how on earth do I select the best features? How do I score how meaningful a feature is (feature importance), if I have lots of base reps, how do I know how many features from each to select (allocation strategies), once I have a candidate, how do I score it to see if its useful (internal evaluator). Then the other axes I test for the purposes of seeing if this works: does it work on different downstream evaluators, and what happens when I modify/control the budget (number of features).
 %        - Compare to other feature selection algorithms but explain why this is different - you could just pool all the features (and I've tried this, pull up my results on it), but it's not actually as effective as looking at them on a representation-level. It turns out if you have 10 reps you're considering, and you limit it to the top 2,3, 5, SOME datasets prefer that. Others prefer all 10. It's more complex than a simple feature selection algorithm if you really want to optimise it, and my results show why
 %        - This use case is specifically taking existing reps that say vastly different things about the molecule (PDV vs ECFP4 vs graph based embeddings vs sequence based embeddings) and tailor selecting a set of features from all of these to the very specific task at hand
 %    - Two concepts of using similarity: make a normal hybrid and use it for similarity OR modify the objective that the hybrid search optimises for to make it similarity-aligned (called "similarity hybrids")
\chapter{Building Adaptive Hybrid Representations}\label{chap:method}
% ---- TAKEAWAYS ----
% - Joining whole representations adds their irrelevant features along with the useful ones. The hybrid combines and filters at once.
% - A hybrid is an input, not a model: up to a budget of features, selected from several base representations by a search for one dataset and one downstream model.
% - No setting of the search wins on every dataset or model, so importance method, allocation method, representation count, internal evaluator, budget and downstream model are all searched.
% - Every selected feature traces back to its representation, so a blueprint can show a chemist which substructures and descriptors the model relies on.
% - The method needs only a supervised label and several vector representations of the same samples, so it carries over to antibodies, reactions, nucleotides and tumour data.
% - Where one representation already dominates, the hybrid should show little gain. The search is costly and must be rerun when the data, representations or model change.
% ---- END TAKEAWAYS ----

\section{Introduction}
% ---- NOTES: Introduction ----
% Takeaway: Standard representations carry many features that do not matter for a given task, and the best representation changes with the dataset and the model.
% Takeaway: This chapter builds the hybrid and reports when and why it helps. Distances and label noise are left to Chapter 6.
% Paragraph 1: the bridge from the earlier chapters
% - Chapter 4: standard representations carry many features that do not vary, or do not matter, on a given dataset. A distance over all of them inherits the problem.
% - Chapter 3: the best representation changed with the dataset and the model. That is the practical reason to search rather than fix one choice.
% - One clause on the origin: the selection was first built for similarity, then found to work as model input.
% Paragraph 2: what this chapter does
% - It builds the adaptive hybrid representation ("hybrid"), described in one sentence, and reports first results on when and why it helps.
% - It measures prediction only. Distances in a hybrid, noise detection and learning with noisy labels are in Chapter 6.
% - Result needed: on how many datasets the hybrid beats the best single representation, and whether the preferred search settings vary between datasets.
% Paragraph 3: what is new in this chapter, as a short numbered list: selecting features across several representations with a per-representation budget; searching the selection settings instead of fixing them; blueprints that trace each feature to its representation; the same search on antibodies, reactions, nucleotides and tumour data.
% - Figure: chapter-opening overview. Left: several base representations as matrices of different widths (descriptors, fingerprint, graph embedding, sequence embedding). Middle: a few highlighted features from each. Right: the narrow hybrid going into any downstream model.
% ---- END NOTES ----


\subsection{Combining representations and the alternatives}
% ---- NOTES: Combining representations and the alternatives ----
% Takeaway: Most prior work joins whole representations or pools the predictions of one model per representation. The hybrid selects individual features across representations.
% Takeaway: No single representation wins across tasks, and new representations keep arriving, so the choice has to be made per task.
% Changes to the prose below, by its paragraph:
% - Paragraph 1: Turkina et al. is the closest prior work. It combines and filters at once, but only across six non-hashed fingerprints.
% - TODO(lit): read the Turkina2025 full text for whether a fixed number of bits is kept per fingerprint, or all bits are ranked together.
% - Paragraph 2: "noise mitigation tasks" should read "noise detection and learning with noisy labels". "concatenate features" should read "select and join features". Augmentation gets one clause.
% - Paragraph 2: the GNN sentence holds. MHG-GNN and GROVER are base representations on every small-molecule property dataset except QM9; cite Table~\ref{tab:reps_smallmol}.
% - Paragraph 3: the authors are Deng, Yang, Wang, Ojima, Samaras and Wang \citep{Deng2023}, not "Li and Deng". Cite them for concatenation hiding which representation earned the score, not for concatenation degrading.
% - Paragraph 4: FernandezDelgado2014 never adds irrelevant features, so it does not support "distance-based models are affected most severely". Why distances suffer is argued in Chapter 4; one clause.
% - Paragraph 4: say that a concatenation of two representations scores between its two parts only if the concatenation test in "Comparators, controls and ablations" shows it. Proposed reason: it adds the irrelevant features of both.
% - Paragraph 5: typo "farmework". The point in two sentences: a new embedding enters as one more base representation; a new downstream model gets its own search.
% - Paragraph 7: replace the domain-by-domain list with one sentence citing the representation tables. Keep \citep{Song2026} (over 40 molecular foundation models, 2020 to 2024) and \citep{Bommasani2021}. Remove "presumably" and "somehow". Move the El Nahhas sentence to the ensembling paragraph.
% - Paragraph 8: cut the first two sentences; they repeat Paragraph 3. Write "the hybrid search", not "KIRBy", and "all base representations side by side", not "the full pool".
% - Paragraph 8: move the consensus sentence \citep{Valsecchi2020} to the ensembling paragraph. Move the information bottleneck sentences to the Discussion; Tishby2015 does not show that compression helps.
% - Paragraph 8 keeps its framing: automated machine learning chooses an algorithm and its hyperparameters together \citep{Thornton2013}. The hybrid search adds a third choice, which features go in.
% New paragraph: ensembling
% - One model per representation; their predictions are pooled. Consensus QSAR does this \citep{Valsecchi2020}.
% - An ensemble of four pathology foundation models, 31 clinical tasks: mean AUROC 1.2\% above the best single model \citep{ElNahhas2025}.
% - No single model sees features from two representations together, and every representation and model must be run for each new molecule.
% New paragraph: feature selection (her TODO)
% - Feature selection keeps a subset of the original features. Three families: filter, wrapper and embedded \citep{Guyon2003}.
% - Where the hybrid sits: it ranks features by a model's importance (embedded), then keeps a fixed number (like a filter).
% - Classical methods treat all features as one flat list. The nearest relative, the group lasso, keeps or drops predefined groups whole \citep{Yuan2006}. The hybrid instead gives each representation a share of the budget.
% - Only original features are kept, so every feature traces back to one representation. This is what makes blueprints possible.
% New paragraph: dimensionality reduction (her TODO)
% - PCA replaces the features with fewer linear combinations of all of them, ordered by variance \citep{Jolliffe2016}. It never sees the labels. Partial least squares is the supervised counterpart \citep{Wold2001}.
% - Both give mixed components. A component cannot be traced to one substructure or representation, so no blueprint can be read from it.
% New paragraph: shrinking the feature space
% - Every alternative except ensembling also reduces the number of features.
% - Models may want different numbers of features, and different features. Distance-based models use every feature in every distance; tree ensembles choose one feature per split.
% - Selection that uses labels may overfit. Taken up in the Discussion.
% - Table: the alternatives side by side. Rows: concatenation, late fusion and ensembling, intermediate fusion, multiple kernel learning, PCA and PLS, flat feature selection, group lasso, the hybrid. Columns: uses the labels; keeps original features; one model sees several representations; adding a representation needs retraining. Each row carries its citation.
% ---- END NOTES ----

% TODO: review
Prior research in this field shows promise, with conjoint fingerprints achieving stronger performance on binding affinity prediction than alone across a range of models \citep{Xie2020}. Augmenting structural-based fingerprints like ECFP4 with bio-activity based fingerprints has also been successful at improving both predictive performance and scaffold-hopping \citep{Laufktter2019}. A very similar fusion strategy, in which features are selected from six non-hashed fingerprints and are combined by feature importance score within individual models, consistently improves predictive accuracy over the original fingerprints \citep{Turkina2025}. 

% TODO: review
The majority of the work in this space deals with combining relatively standard fingerprints and linear representations like SMILES \citep{Xie2020, Laufktter2019}. However, graph-based representations capture important structural information that is often lost in fingerprints. We explore using GNN embeddings within these hybrid, fusion methodologies. We further expand on the mechanisms used to compute feature importance, concatenate features, and use these hybrids within similarity quantification and noise mitigation tasks. We also explore adding on smaller sets of structural features to a single base representation, referred to as augmentation. 

% TODO: review and the li and deng defini
% TODO: review, needs manual review
So far, the most common strategy for combining feature sets has been concatenation, which requires only a single downstream model \citep{Stahlschmidt2022}. It works best when feature sets are homogeneous. When they are not, concatenation degrades \citep{Stahlschmidt2022}, and \cite{Deng2023} called the practice ``misleading'' when the goal is to measure the contribution of individual representation types. Late fusion trains a separate model per source and aggregates predictions \citep{Steyaert2023}, avoiding the dimensionality problem. But cross-source feature interactions are lost. To recover cross-source interactions requires intermediate fusion (MOGONET \citep{Wang2021}, cross-attention \citep{Zhang2024}), but at the cost of end-to-end training on all sources; adding a new source means retraining from scratch. No interpretable feature subsets emerge from multiple kernel learning \citep{Gonen2011}, which operates entirely at the kernel level. There remains, as Steyaert et al.\ noted, no decisive evidence that any of these strategies is broadly superior \citep{Steyaert2023}.

% TODO: review, needs manual review
In practice, most research groups default to one representation and ignore the rest. ECFP is common in small molecules, while deep learning groups may opt for embeddings like MoLFormer or ChemBERTa. ESM-2 embeddings are common for antibodies, and can be broken up by chain. While there is some research that compares these representations, there is not much that focuses on combining them beyond simple concatenation. Concatenating just three standard sources, however (ECFP4 at 2,048 dimensions, ESM-2 at 1,280, Mordred descriptors at ${\sim}1{,}600$), gives a ${\sim}5{,}000$-dimensional input for a dataset that may contain only a few hundred labeled samples. At these sample-to-feature ratios, more features mean worse predictions \citep{Hughes1968, Bellman1961}. Models sensitive to distance, such as kNN and SVM, are affected most severely \citep{FernandezDelgado2014}.

% TODO: review, needs manual review
The pace of new model development shows no sign of slowing, and extends well beyond small molecules. First, these new models may prefer different feature sets, and that's what this methodology can provide. But beyond that, new models can be fit and converted into frozen embeddings and used as further input into this farmework. 

% TODO: review, needs manual review
A comparison of 62,820 models in molecular property prediction benchmarks showed that classical RDKit2D descriptors matched or outperformed learned features from graph neural networks and pretrained language models on the majority of datasets \citep{Deng2023}. Descriptor-based models occupied 73\% of top-three positions on 11 benchmarks \citep{Jiang2021}, and in 1,310 ChEMBL assays, feedforward networks on ECFP fingerprints outperformed graph convolutions and SMILES-based recurrent models \citep{Mayr2018}. The input often matters more than the model. Yet we still lack a systematic way to exploit this. No single representation wins. The same pattern holds outside small-molecule chemistry: no single protein language model consistently achieves the highest performance on all tasks \citep{Kenlay2024, Schmirler2024}; for antibodies, general-purpose models sometimes outperform those built for the domain \citep{Kenlay2024}. In computational pathology, no individual foundation model dominates on clinical endpoints \citep{ElNahhas2025}, and in cancer genomics, integrating additional omics layers can negatively impact performance when the data types are not properly combined \citep{Steyaert2023}.

% TODO: review, needs manual review
For small molecules alone, the landscape now includes extended-connectivity fingerprints, physicochemical descriptor sets exceeding 1,600 dimensions, graph neural network embeddings \citep{Yang2019, Fang2022}, language models trained on billions of SMILES strings \citep{ross2022}, image-based encodings (molecules rendered as 2D images) \citep{Zeng2022}, and stereoelectronics-infused molecular graphs \citep{Boiko2025}. The same is true for proteins (ESM-2 \citep{Lin2023}, ProtTrans \citep{Elnaggar2022}), antibodies (AntiBERTy \citep{Ruffolo2021}, AbLang \citep{Olsen2022}), histopathology \citep{ElNahhas2025}, and multi-omics cancer profiles \citep{Steyaert2023}. A recent review counted over 40 molecular foundation models published between 2020 and 2024 \citep{Song2026}. Increasingly, new models appear faster than we can compare them \citep{Bommasani2021}. The choice is not getting simpler. Different sources presumably encode distinct aspects of the same biology. Noisy or incomplete on their own, single data modalities somehow give more sensitive and robust predictions when combined with independent signals from other modalities \citep{Steyaert2023}. An ensemble of frozen embeddings from the top four pathology foundation models, El Nahhas et al.\ found, achieved a mean AUROC 1.2\% higher than the best individual model on 31 clinical tasks \citep{ElNahhas2025}.

% TODO: review, needs manual review
Concatenation, the most common method for combining feature sets, gives the full pool that KIRBy starts from, but on biological sample sizes it can degrade rather than help \citep{Bellman1961, Stahlschmidt2022}. At the kernel level, multiple kernel learning \citep{Gonen2011} (e.g., pairwiseMKL for drug bioactivity \citep{Cichonska2018}) can learn source weights, but no interpretable feature-level selections result. Consensus methods \citep{Valsecchi2020} combine predictions, not features, and so cannot exploit cases where one source's noise is another's signal. Instead, KIRBy works at the feature level: it selects features that span source boundaries. Perhaps the closest analogy in the optimization literature is the CASH problem \citep{Thornton2013} (the joint selection of algorithm and hyperparameters), extended here to include the input encoding as an additional search dimension. Tishby and Zaslavsky \citep{Tishby2015} showed why compression helps: the information bottleneck principle states that optimal representations maximize mutual information with the target while minimizing mutual information with the input. In effect, the budget forces the model to keep only what matters. The feature count controls the trade-off between informativeness and dimensionality.

% TODO: background on feature selection algorithms
% But more of a direct focus on feature selection

% TODO: background on PCA and other dimensionality reduction techniques


\subsection{The adaptive hybrid representation}
% ---- NOTES: The adaptive hybrid representation ----
% Takeaway: A hybrid is up to a budget of features selected from several base representations of the same molecules, chosen by a search for one dataset and one downstream model.
% Takeaway: Filtering out useless features matters as much as picking useful ones, because pooled irrelevant features act as noise.
% Paragraph 1: no single representation is best for every dataset, even for one model (answers her TODO on NGBoost)
% - Change: name the method "adaptive hybrid representation", "hybrid" for short. Drop "KIRBy (Knowledge Integration by Representation)" and "framework". "the noise robustness study above" should read "in Chapter~\ref{chap:robustness}".
% - Her NGBoost example holds on QM9 and fails on Caco-2 as written.
% - QM9, clean $R^2$ (median of 10 replicates): PDV 0.865 is highest, but MHG-GNN 0.852 is close. "Only strong with PDV" overstates it.
% - Caco-2: Avalon leads only after 3 of its 5 folds are dropped. Over all 5 folds PDV is highest, 0.370 clean $R^2$ against 0.287 for Avalon.
% - hERG Ki gives a clean second example: MHG-GNN is highest there (0.511 clean $R^2$). So NGBoost's best representation changes with the dataset: PDV on QM9 and logD, MHG-GNN on hERG Ki.
% - Cut the sentence guessing that Avalon captures features relevant to efflux. Nothing tested it.
% - Typos: "overidden"; "its own type" should be "their own type"; "data set" is "dataset".
% Paragraph 2: what a hybrid is, in words
% - Up to a fixed number of features (the budget), selected from several base representations of the same molecules.
% - Chosen by a search, for one dataset and one downstream model. It is an input, not a model: any supervised model can be trained on it.
% - Intended use: representations that say different things about the molecule (PDV, ECFP4, graph embeddings, sequence embeddings), with the features for this task picked from all of them.
% Paragraph 3: why select (her second paragraph)
% - Add one sentence: distance-based noise handling did worst in an earlier exploratory study (Section~\ref{sec:lnl_motivation}). That study covered one dataset and one noise type; Chapter~\ref{chap:similarity} measures the question properly.
% - Typo: "the feature which carry" should be "the features which carry".
% Paragraph 4: why the budget is split per representation
% - Representations differ in width by orders of magnitude, and their importance scores are not on a common scale. In a single ranking, wide ones can crowd out narrow ones.
% - One pooled allocation method is also searched, so per-representation and pooled selection are compared, not assumed.
% - One sentence listing what the idea leaves open. The search settings answer each: how useful each feature is, how many each representation gets, how many representations take part, how a candidate is judged, the budget, and the downstream model.
% - One sentence: the two ways to use a hybrid for similarity are defined in Chapter~\ref{chap:similarity}.
% - Figure: base representation matrices of different widths; features ranked by importance within each; the budget as a bar split between representations; the top features joined into a narrow hybrid. Side panel: the full concatenation, for contrast.
% ---- END NOTES ----

% TODO: is this right??? I'm making shit up with NGBoost.
Here, we present KIRBy (Knowledge Integration by Representation), a framework that takes a collection of feature matrices for a given data set and outputs a customised feature set optimised for a specific prediction task. The inputs can be any type of frozen feature vector, such as fingerprints, descriptor sets, foundation model embeddings, gene expression signatures, histopathology features. Within a given domain, these individual vectors typically contain its own type of information. It's a well-established phenomenon when it comes to machine learning, there's rarely ever a one-size-fits-all approach. Take the noise robustness study above as evidence, NGBoost's predictive performance on QM9 was only ever strong with given PDV input features. However, on Caco-2, NGBoost preferred Avalon, with PDV as a close second. The model NGBoost itself produces better fits when given physicochemical properties, yet this preference can be overidden by a different prediction task. Avalon likely captures features that are pertinent to the log efflux ratio.

But as important as the selection of useful features is, it is also important to filter out the useless features. We could throw the kitchen sink at the models and pool together features across a wide array of representations. Not only would that pose computational challenges, but it also introduces a huge set of irrelevant features that function as noise in downstream applications (Chapter~\ref{chap:flaws}). With all of that in mind, I decided to develop a method that selects the best features for a given data set and model. In theory, it should effectively filter out noisy features and isolate the feature which carry the most relevant information.

% TODO: review, needs manual review
The questions this chapter and Chapter~\ref{chap:similarity} address are set out in Section~\ref{sec:intro:questions}.


\subsection{Blueprints and what they tell a chemist}
% ---- NOTES: Blueprints and what they tell a chemist ----
% Takeaway: A finished search is saved as a blueprint, and every selected feature traces back to its representation, so the search says something about the problem, not only its score.
% Takeaway: Named features let a chemist look for new chemistry or confirm the selection makes sense. Learned embedding features cannot be read this way.
% Paragraph 1: the output is more than a score
% - A finished search is saved as a blueprint: the settings needed to rebuild the hybrid. Every selected feature traces back to its representation.
% - There is one blueprint per downstream model and budget. Comparing them shows whether tree ensembles and neural networks draw on different features.
% Paragraph 2: discovery and confirmation
% - Readable features: descriptor names, the atom environment behind an ECFP bit, the substructure behind a MACCS key. Outside small molecules: a sequence position, or a gene pathway.
% - Learned embedding features say which representation was useful, not which chemistry. One clause; Chapter~\ref{chap:flaws} covers why.
% - Discovery: substructures or descriptors that recur among the kept features point to parts of the molecule that may drive activity. A starting point for a chemist, not a finding.
% - Confirmation: kept features that make chemical sense support the claim that the selection finds real signal.
% - Caveat, one sentence: selected means predictive, not causal; correlated features can stand in for each other.
% - Context: explainability in drug discovery \citep{JimenezLuna2020}; interpretable inputs instead of explaining a black box afterwards \citep{Rudin2019}.
% Paragraph 3: other uses of a similarity chosen for the task (one or two sentences)
% - A hybrid also gives a distance chosen for one task. Name the uses from Chapter~\ref{chap:flaws}: noise detection, nearest-neighbour prediction, applicability domain, activity-cliff detection, clustering for group-aware splits, weighting predictions from separate sites.
% - Future work: the predictor that scores molecules in generative design \citep{Olivecrona2017}.
% ---- END NOTES ----
\todo[inline, color=yellow!40]{5.1.4: explainability motivation; other applications of a task-specific similarity.}


\subsection{Beyond small molecules}
% ---- NOTES: Beyond small molecules ----
% Takeaway: The search needs only a supervised task and several vector representations of the same samples, so it applies to tabular data in general.
% Takeaway: A dataset needs enough samples, and base representations whose features carry importance.
% Paragraph 1: what the method needs
% - A supervised task and several vector representations of the same samples, one feature matrix each.
% - The default objective scores prediction only and reads no chemical structure, so it applies to tabular data in general.
% - What does not carry over: reading substructures from a blueprint, the optional shared-substructure objective, and the distance and label-noise work of Chapter~\ref{chap:similarity}.
% Paragraph 2: one-hot encodings against learned sequence embeddings
% - One-hot encodings are sparse. At most positions the residue never changes across a dataset, so those features carry no importance. Recommend language-model embeddings.
% - One-hot can still work, as on some antibody datasets. Result needed: those datasets, with the score in the dataset's metric.
% Paragraph 3: the other domains, one or two sentences each
% - Antibodies get the most weight after small molecules. Reactions and nucleotide sequences are proofs of concept.
% - Cancer genomics: each representation is a different measurement of the patient, not a different encoder of one input. The search treats both the same way.
% Paragraph 4: what a dataset needs
% - Enough samples. The threshold (about 200) waits on the full runs.
% - Base representations whose features carry importance. Where every representation scores poorly, there is too little to improve on.
% - Figure: one row per domain (small molecule, antibody, reaction, DNA sequence, patient sample). Left: the raw input. Middle: several representation matrices. Right: the same search box, outputting a hybrid. Only the left column changes between rows.
% ---- END NOTES ----
\todo[inline, color=yellow!40]{5.1.5: other domains; requirements for a dataset.}


\section{Methods: the hybrid search}
% ---- NOTES: Methods: the hybrid search ----
% Takeaway: The search builds and scores every combination of setting values inside cross-validation, keeps a winner for each downstream model and budget, and scores it once on an untouched outer fold.
% Lead-in, 3-5 sentences, no definitions:
% - The search is a grid: every combination of setting values is built and scored inside the cross-validation folds.
% - A winner is kept for each pair of downstream model and budget. The budget is a maximum number of features.
% - The reported comparison runs the whole search inside each outer training fold and scores the winner once on the untouched outer test fold.
% - Two options make it cheaper: the fast tier and adaptive pruning.
% ---- END NOTES ----


\subsection{From base representations to a hybrid}
% ---- NOTES: From base representations to a hybrid ----
% Takeaway: Each representation's features are ranked by importance, the budget is split between representations by how well each predicts alone, and the top features from each are joined.
% Takeaway: Every step uses the training molecules of the fold only. The outer test molecules score only the finished hybrid.
% Paragraph 1: inputs and output (her overview paragraph; fix it, do not add to it)
% - Inputs: $K$ base representations of the same $n$ molecules, $\mathbf{X}^{(k)} \in \mathbb{R}^{n\times d_k}$, with widths $d_k$ from tens to thousands of features. Also the labels $\mathbf{y}$, one group per molecule, the metric, a downstream model and a list of budgets.
% - Output: a hybrid of at most $B$ features. The budget is a maximum, not a target.
% - Changes: "at most $b$ columns" should read "at most $B$ features". "sources" should read "representations". "Blueprint" is lower case.
% - Change: features outside the budget come only from uncapped augmentation.
% - Change: "according to baseline performance" holds for four of the five allocation methods. The global pool ranks the offered features together; say so in one clause.
% Paragraph 2: preprocessing inside each fold (her paragraph)
% - Change: "drop columns" should read "drop features".
% - Change: binary features stay 0/1 only when a single representation is scored to pick the comparator. Every hybrid feature is standardised with training-fold statistics.
% - One sentence: a representation wider than 5,000 features is cut to its 500 most-split features (LightGBM) for that single-representation scoring only. The hybrid search always sees the full width.
% Paragraph 3: which molecules each step sees
% - Five outer folds, group-aware. Classification folds are also stratified by class.
% - Inside each training fold, one group-aware split holds out about 20% of the molecules. Importance and each representation's standalone score come from this fold only.
% - The outer test molecules score only the finished hybrid.
% Paragraph 4: importance and standalone score per representation
% - The importance method gives each feature of representation $k$ a score $s^{(k)}_j \ge 0$. Scores rank features within one representation only.
% - The internal evaluator is fitted on representation $k$ alone and scored on the held-out part, giving $r_k$.
% - $r_k$ is rescaled to $\tilde r_k$ so that 0 means no better than a trivial predictor, whatever the metric.
% - Spearman, Pearson, MCC and $R^2$: $\max(0, r_k)$. AUROC and concordance index: $\max(0, 2(r_k - 0.5))$.
% - RMSE: $\max(0, 1 - \mathrm{RMSE}_k/\mathrm{RMSE}_0)$, where $\mathrm{RMSE}_0$ comes from predicting the training mean. MAE, accuracy and weighted F1 follow the same pattern.
% Paragraph 5: splitting the budget
% - Optional representation count $N$: only the $N$ representations with the highest $\tilde r_k$ go on.
% - Each representation gets a weight $w_k = g(\tilde r_k)$, set by the allocation method.
% - Its share is $b_k \approx B\, w_k / \sum_l w_l$, rounded so that $\sum_k b_k = \min(B, \sum_k d_k)$ and $b_k \le d_k$.
% - The global pool does not use this share. One joint ranking keeps at most $B$.
% Paragraph 6: selecting features
% - $S_k$ = the top $b_k$ features of representation $k$, among those scoring above 0.
% - Why the cut at 0: tree importance gives 0 to every feature no tree split on. Filling the budget from those would pick features by their position in the matrix.
% - Hybrid: $\mathbf{H} = [\mathbf{X}^{(1)}_{:,S_1}, \dots, \mathbf{X}^{(K)}_{:,S_K}]$. So a hybrid can be narrower than its budget.
% Paragraph 7: scoring, and the reported comparison
% - A configuration is one value on every searched setting. Each is scored by fitting the downstream model on the hybrid and averaging the held-out score over folds.
% - The winner is picked per downstream model and budget. Its score chose among many configurations, so it is optimistic \citep{Varma2006, Cawley2010}.
% - The reported comparison is nested: the whole search runs inside each outer training fold, and the winner is scored once on the untouched outer test fold.
% - Result needed: how much higher the inner search score is than the outer test score, in each dataset's metric.
% Paragraph 8: two variants used in Chapter 6, one sentence each
% - Under label noise, every training-side quantity uses the noisy labels. Held-out labels are clean.
% - The nearest-neighbour objective replaces both scores with predicting each molecule's activity from its nearest neighbours in the space.
% - Figure: the search pipeline. Left: matrices of different widths. Middle, inside one outer training fold: the internal 80/20 group split; importance bars per representation; standalone scores feeding the budget split; top features pulled from each. Right: the hybrid, the downstream model, the held-out score. Mark where the blueprint is saved, and that the search never sees the outer test fold.
% ---- END NOTES ----

% TODO: review
% TODO: review, needs manual review
Given $K$ frozen representation matrices $\mathbf{R}_1, \dots, \mathbf{R}_K$ for the same $n$ samples, a label vector $\mathbf{y}$, and optional group assignments for cross-validation, the search finds which features from which sources best predict for a specified metric and downstream model. The output is a hybrid feature matrix of at most $b$ columns, where $b$ is the feature budget (plus any augmentation features added outside it), together with a Blueprint that records all parameters needed to reproduce the identical feature selection on new data. The search first scores each source independently, then scores feature importance on training data within each cross-validation fold, allocates the budget among sources according to baseline performance, and scores candidate hybrid configurations via cross-validation. We retain the best configuration for each combination of downstream model and budget. Every selected feature traces back to its source and importance rank.

% TODO: review, needs manual review
Inputs from different sources differ in scale, distribution, and missingness patterns, and naively pooling them before splitting would leak validation-set statistics into the training procedure, a well-documented source of optimistic bias in feature selection studies \citep{Ambroise2002, Vabalas2019}. All preprocessing therefore runs within each cross-validation fold. We impute missing values with per-feature medians computed on the training fold and drop columns that are entirely missing. Every feature of the hybrid, binary fingerprint bits included, is standardized; only when a single representation is scored on its own are binary features left as 0/1. We standardize to zero mean and unit variance using training-fold statistics only.

\todo[inline, color=yellow!40]{Figure 5.1: the search pipeline schematic. Does not exist yet.}


\paragraph{Blueprints.}
% ---- NOTES: Blueprints ----
% Takeaway: A blueprint stores the recipe, not the features. Rebuilt on other labelled molecules it can pick other features; for unlabelled molecules, the features chosen by a labelled fit are reused.
% Paragraph 1: what a blueprint stores (her first paragraph; fix it)
% - Add to her list, in one clause: the objective, the metric's direction and the augmentation settings.
% - Drop the step size from her list. No searched allocation method reads it.
% - Typo: "it's results".
% - Change: the recipe does not recreate "this exact hybrid representation on any data". On other molecules a rebuild re-scores the representations and re-picks features.
% Paragraph 2: two ways to reuse a blueprint (her second paragraph; fix it)
% - Re-run the recipe on new labelled data. The features adapt to the new dataset, and may differ.
% - Reuse the exact features from a labelled fit. This needs no labels, so it is the only option for unlabelled molecules. It assumes the same features matter.
% - Result needed: which of the two transfers better from a large dataset to a smaller related one.
% Paragraph 3: naming the selected features (her third paragraph; fix it)
% - Change: the names come from a separate explanation step run on a finished hybrid. The blueprint does not store them.
% - Change: ECFP bits are named by the atom environment that sets them, written as a SMILES fragment, not SMARTS. MACCS keys are named by their SMARTS pattern. Descriptor features carry their descriptor names.
% - Change: her prose names only learned embeddings as unreadable. Also given only as representation name and index: MAP4, DRFP, Avalon, Mordred, atom-pair and graph-kernel features.
% - Figure: two paths. Left, "re-run the recipe": new labelled data, score representations, split the budget, pick features. Right, "reuse the exact features": stored features taken from unlabelled molecules. Mark which path needs labels.
% ---- END NOTES ----

% TODO: review, needs manual review
To be able to better examine the selected features and to regenerate them on new data sets, we produce blueprints. After a hybrid feature selection search is completed, it's results are stored as a blueprint containing the ``recipe'' required to recreate this exact hybrid representation on any data. These blueprints store the selected picks for the pre-specified budget,  the full list of input representations searched and how many of them the recipe used, feature importance scoring method, allocation strategy, and internal evaluator. It contains other runtime information like the step size, number of folds, the seed, the metric being used to score, the task type (regression or classification), whether augmented features were being used, and the per-representation scores.

The blueprint does not store selected feature indices, those can only be recovered when the blueprint is refit on the data, which re-runs the feature selection process deterministically from that recipe and seed.

It also provides a registry of selected features, stating where that feature came from. On some representations, we added an extension where features are mapped back to their function. For instance, ECFP and MACCS bits are mapped to SMARTS substructures. This blueprint will tell you which of those substructures were selected. This is not possible for learned embeddings for which the blueprint will return the bare representation name and dimension index. This mapping provides a layer of explainability. It can be used both to discover what features are useful for your task and to confirm that the features being selected are in fact lining up with the underlying known chemistry.

\subsection{The search space and why each axis is swept}
% ---- NOTES: The search space and why each axis is swept ----
% Takeaway: Six settings are searched as a full grid: downstream model, budget, representation count, importance method, allocation method and internal evaluator.
% Takeaway: Each is searched because different datasets and models prefer different values, and the settings interact.
% Paragraph 1: motivation
% - Move her paragraph "When I began to implement this idea..." here, as the opening. It now sits at the end of "Search configuration and internal evaluation".
% - Her three questions map onto settings. How to measure importance: the importance method. How many features from each representation: the allocation method, budget and representation count. How to score a feature set: the internal evaluator.
% Paragraph 2: why a grid
% - Every combination of values is built and scored. The claim that no single setting suits every dataset is then measured, not assumed.
% - A full product, not one setting at a time, because the settings interact: the best budget depends on the downstream model (example below).
% - Motivation only: no learning method is best on every problem \citep{Wolpert1996}.
% Paragraph 3: the settings
% - Six settings: downstream model, budget, representation count, importance method, allocation method, internal evaluator. Values in Table~\ref{tab:search_tiers}; the full breakdown in Additional files, Section~\ref{app:searchspace}.
% - Not searched inside one search: the objective and augmentation. Each is a separate search.
% - Downstream model: it decides what counts as a useful feature and how many it can use.
% - Budget: eight values, 10 to 2,000 features. Too few leave informative features out; too many add uninformative ones.
% - Example (preliminary), nanobody thermostability, 5 folds: SVM did best at 300 features (0.6384 Spearman). DNN, LightGBM and random forest did best at 50 features, XGBoost at 100.
% - Representation count: 2, 3, 5 or all of a dataset's base representations (usually 10). It is the only step that drops whole representations.
% - Pooled feature selection has no such setting. Some datasets do best with 2 or 3 representations, others with all 10.
% ---- END NOTES ----

\subsubsection{Feature importance scoring}
% ---- NOTES: Feature importance scoring ----
% Takeaway: Three importance methods from different model families rank features differently, and which one builds the best hybrid depends on the dataset and the downstream model.
% Paragraph 1: what the search does with an importance score
% - Each representation gets its own freshly fitted model, on its own features, on the internal training part of the fold. The downstream model is not used.
% - Only the order within a representation is used, so scores are never normalised or compared across representations.
% - Scores are computed once per fold and method, and reused by every other setting.
% - Change: "Tree-based feature importance measures favor features with high split gain" does not match \citep{Strobl2007}. That paper shows split-based importance favours features with many possible split points. This matters when binary bits sit beside continuous embeddings.
% Paragraph 2: the three methods (answers her TODO on naming)
% - Names in text: "LightGBM gain", "TreeSHAP", "DeepLIFT" \citep{Shrikumar2017}. Change: the prose spells it "DeepLift".
% - Change: "explanation-based" does not separate the three. Use family labels: split gain, Shapley values, reference-based attribution.
% - LightGBM gain: the total loss reduction of every split on the feature, summed over 50 boosted trees \citep{ke2017lightgbm}.
% - TreeSHAP: the mean absolute Shapley value over the training molecules, from a 50-tree random forest \citep{Lundberg2020}.
% - DeepLIFT: the mean absolute attribution against an all-zero reference, which is each feature at its training mean. It comes from a small two-layer network; call it that, to keep it apart from the DNN downstream model.
% - Model settings are in Table~\ref{tab:app:model_settings} (Additional files).
% Paragraph 3: why these three (her paragraph "We tested if...")
% - Change: the prose says the methods' "predictions did not correlate". It is their feature rankings.
% - Selection rule: a method must score every feature for regression and classification. Of two methods whose rankings agree above Spearman 0.95, the faster is kept. Methods that rank features differently are kept.
% - Under that rule DeepLIFT replaced integrated gradients and gradient SHAP.
% - The three kept methods rank features differently (preliminary, 7 datasets, fitted once on all molecules). Spearman between LightGBM gain and either other method: 0.20 to 0.67. Between TreeSHAP and DeepLIFT: 0.56 to 0.99.
% Paragraph 4: no single method suits every dataset and model
% - Preliminary: 5 datasets, 5 folds, 5 budgets (25 to 300 features), 5 downstream models. A win is one (downstream model, budget) pair where that method gave the best configuration, out of 25 per dataset.
% - LightGBM gain won most pairs on four datasets. DeepLIFT won 17 of 25 on LogD, and 2 of 100 elsewhere.
% - Change: "networks preferred tree-based measures" holds on four of five datasets. With the DNN downstream, LightGBM gain won 17 of 20 budgets outside LogD.
% - Change: "tree ensembles preferred gradient-based approaches" holds on LogD only, for XGBoost and LightGBM (5 of 5 budgets each). Random forest took LightGBM gain there.
% - Table: one row per dataset. Columns: dataset | metric | wins out of 25: LightGBM gain, TreeSHAP, DeepLIFT | best score reached with each, same order (preliminary).
%   Nanobody thermostability | Spearman | 15, 10, 0 | 0.6696, 0.6471, 0.6440
%   hERG Ki | RMSE (pKi), lower is better | 14, 11, 0 | 0.5589, 0.5605, 0.5728
%   Humanness | AUROC | 12, 12, 1 | 0.9723, 0.9713, 0.9689
%   LogD | RMSE (log units), lower is better | 8, 0, 17 | 0.4250, 0.4345, 0.4170
%   BRCA PAM50 subtype | weighted F1 | 21, 3, 1 | 0.9097, 0.9018, 0.9038
% - Result needed: the same win counts from the full runs.
% - The full search sweeps all three. The fast tier uses LightGBM gain only.
% ---- END NOTES ----

% TODO: this is out of date

% What were the qualifications for correlation/filtering allocation strategies? Can I formalize this a bit?
% TODO: review, needs manual review
Not all features from a given base representation are equally informative, and there are a variety of techniques that can assess each feature's relative importance. Tree-based feature importance measures favor features with high split gain, while neural attributions such as DeepLIFT measure how much each feature moves a neural network's output away from its output on a reference input \citep{Strobl2007, Ancona2018}.
We tested if, for the purposes of this task, there was a single feature importance method that worked across data sets and model architectures. Not only was there no clear winner, but there was evidence that certain networks preferred tree-based measures of importance while certain tree ensembles preferred gradient-based approaches. We landed on a set of three who's predictions did not correlate strongly with each other and who's selected features performed well as model input.

% TODO: background on each of these, also how do I name these in text
% TODO: review, needs manual review
We score features using LightGBM gain, DeepLift, and TreeSHAP. LightGBM gain is a tree based method which.... TreeSHAP is explanation-based, meaning it... DeepLift relies on neural attribution which...

When searching for the features to use in the hybrid, thse feature importance scoring methods are swept.


\subsubsection{Feature allocation strategies}
% ---- NOTES: Feature allocation strategies ----
% Takeaway: An allocation method decides how many of the budget's features each representation contributes. Four give fixed shares from each representation's standalone score; the fifth lets candidates compete in one pooled ranking.
% Takeaway: No allocation method was best in most comparisons, so all five are searched.
% Paragraph 1: what an allocation method decides
% - How many of the budget's $B$ features each representation contributes. The importance method then decides which ones.
% - Five methods. Linear, squared, softmax and rank decay give fixed shares from each representation's standalone score. The global pool lets the features each representation offers compete in one ranking.
% - Change: "select the top features from each source and concatenate" holds for the four fixed-share methods only.
% Paragraph 2: the four fixed-share methods (replaces her numbered list, which repeats "From base representations to a hybrid" and writes $n_r$ for $b_k$)
% - Change: the standalone score comes from the internal validation part of the training fold, not "the training set".
% - They differ only in the weight $w_k$ computed from the rescaled score $\tilde r_k$.
% - Linear: $\tilde r_k$. Squared: $\tilde r_k^{\,2}$. Softmax: $\exp(\tilde r_k - \max_l \tilde r_l)$. Rank decay: $e^{-0.3\rho_k}$, where $\rho_k$ is the rank by score (0 for the best).
% - Linear and squared give 0 features to a representation scoring 0. Softmax and rank decay never do.
% - Worked example: budget 100 features, three representations scoring 0.6, 0.5 and 0.2. Linear 46/39/15; squared 55/39/6; softmax 39/35/26; rank decay 44/32/24 features.
% - Change: "variants of performance-weighted allocation ... proportional to performance" fits linear only. Say "fixed-share methods".
% - Table: the worked example, one row per method, one column per representation.
% Paragraph 3: the global pool
% - Each representation offers candidates, sized by squared weights as shares of $2B$. The same importance method is fitted once on all offered features together, and the top $B$ are kept.
% - Each representation's count is then an outcome of the joint ranking. A feature's rank depends on what the other representations offered.
% Paragraph 4: what fixed shares assume (her two sentences; merge them)
% - A strong representation may owe its score to a few features, so a large share can fill with weak ones.
% - Fixed shares cannot see complementarity: two representations carrying the same information both get a share. Only the global pool ranks features across representations. State it as a difference in mechanism, without a verdict.
% Paragraph 5: greedy forward selection, tried and cut (answers her TODO; three sentences)
% - It added 10 features at a time from whichever representation most improved the internal score, so it could pick features that work together \citep{Kohavi1997}.
% - At 500 features and 10 representations it cost up to 500 fits per configuration and fold, against 10 for the score-based methods. It placed 16th of 17 methods on both datasets it was compared on (preliminary).
% - Drop the commented sentence that it "works better on larger datasets". The PAM50 set has 821 tumours and greedy still placed 16th.
% Paragraph 6: how the five were chosen (answers her TODO; two sentences)
% - Preliminary: 8 candidate methods on 9 regression datasets, 3 downstream models and 2 budgets, 54 combinations in all.
% - A fixed-share method was best in 38 of the 54, the global pool in 16, so all five stay.
% ---- END NOTES ----

% TODO: allocation strategies completely out of date
% Check framework doc?
% TODO: review, needs manual review
How to divide the total feature budget among sources is the role of the allocation strategy. We search five. Four divide the budget in fixed shares from transformed standalone scores (linear, squared, softmax, rank decay). The fifth, a global pool, lets features from all sources compete in one importance ranking. Once we have the allocation, we select the top features from each source by importance score and concatenate them to form the hybrid.

% TODO: completely replace this is SO out of date, I need a version of this for all the allocation strategies I'm using, as well as the knobs

% \begin{enumerate}
%     \item \textbf{Split:} Divide training data into allocation-train (80\%) and validation (20\%) subsets using scaffold splitting.

%     \item \textbf{Initialize:} Set current feature set to empty, remaining budget to $B$.

%     \item \textbf{Rank features:} For each representation $r \in \{base_1, base_2, \dots base_m\}$, train RFs using the base representations individually, then rank features by mean decrease in impurity.

%     \item \textbf{Iterate:} While remaining budget $\geq$ step\_size and patience not exhausted:
%     \begin{enumerate}
%         \item For each representation, try adding the top step\_size candidates from its ordered list to the current feature set.
%         \item Train a RF using the tentative feature set.
%         \item Evaluate $R^2$ on the validation set.
%         \item Select the representation whose addition yields the highest validation $R^2$.
%         \item Permanently add step\_size features from the winning representation; remove from its candidate list.
%         \item If validation $R^2$ does not improve for patience consecutive rounds, terminate early.
%     \end{enumerate}

%     \item Re-compute importance scores for the features of the new hybrid on the full training data. For each representation, select the top $n_r$ features according to the allocation determined above, where $n_r$ is the number of features allocated to representation $r$ during iteration.
% \end{enumerate}

% This greedy approach works better on larger datasets, as the size of the validation set may not be enough to gain meaningful insights from.

% TODO: this is not what's happening
% TODO: review, needs manual review
The four fixed-share methods are variants of performance-weighted allocation, where the distribution of features from each base representation is proportional to that representation's standalone predictive performance on the training set. This works under the assumption that representations with stronger predictive performance contribute more valuable features. It is worth noting that this is not always the case, as high-quality, large representations may rely on a small subset of valuable features.

% TODO: review, needs manual review
\begin{enumerate}
    \item Fit the internal evaluator on each base representation independently, on the internal training part of the fold.
    \item Score each representation on the internal validation part in the dataset's own metric, rescaled so 0 means no better than a trivial predictor.
    \item Allocate features proportionally:
    \[
    n_r = B \times \frac{w_r}{\sum_{r'} w_{r'}}, \quad w_r = g(\tilde r_r)
    \]
    \item For each representation, rank features by the configuration's importance method and select the top $n_r$.
    \item Concatenate selected features to form the hybrid representation.
\end{enumerate}

% TODO: review, needs manual review
This method cannot adapt allocation based on feature complementarity, and it relies on the assumption that many features in a well-performing representation are valuable.

% TODO: explain why the greedy approach was tried, but it's too computationally intensive and actually quite ineffective
% TODO: explain framework experiments and how they helped here

\subsubsection{Search configuration and internal evaluation}
% ---- NOTES: Search configuration and internal evaluation ----
% Takeaway: The internal evaluator scores each representation on its own during the search. That score splits the budget and decides which representations take part.
% Takeaway: The evaluator that best serves a downstream model is not always the same model family, so four evaluators are searched, matching included.
% Paragraph 1: what an internal evaluator is (her paragraphs; fix them)
% - A model fitted inside the search on one representation at a time, separate from the downstream model.
% - Its score sets each representation's share of the budget and ranks representations for the representation count. It does not rank features inside a representation.
% - Judging a feature set by a model's score is the wrapper idea \citep{Kohavi1997}.
% - Change: "a model whose inductive bias best fits the downstream model" is out of date. All four evaluators are searched and the best configuration wins.
% - Change: kNN is not a downstream model; use SVM or DNN in the opening example. "as we found in our feature importance scoring analysis" names no result; cite one or drop it.
% - Typos: "evalutors"; "another axes" should be "another setting".
% Paragraph 2: the four values (answers her TODO on architectures)
% - Matched (the downstream model's own family), random forest, LightGBM and DNN.
% - So random forest, LightGBM and DNN downstream models are searched with 3 evaluators; XGBoost, SVM and the Gaussian process with 4.
% - Same family is not the same settings. The tree evaluators use 50 trees against 100 downstream. The SVM and Gaussian process evaluators see at most 500 training molecules.
% - All settings are in Table~\ref{tab:app:model_settings} (Additional files).
% Paragraph 3: a downstream model the search does not know (answers her TODO; one or two sentences)
% - A user's own model gets no matched evaluator and no warning; the other evaluators score it. State it as a limit: a user with their own model should name an evaluator.
% Paragraph 4: why the evaluator is searched
% - Variable importance depends on the model \citep{Dong2020}. Matching is the obvious default, but another model may rank representations better for the downstream model.
% - A mismatch winning is informative, not a negative result.
% - Preliminary, 5 datasets, 5 folds, the search's own scores. The matched evaluator was best in 43 of 125 combinations of dataset, downstream model and budget. With no effect, about 38 would be expected.
% - Every gap between evaluators was smaller than the range of matched's own fold scores. The per-dataset gaps belong in Additional files, Section~\ref{app:searchspace}.
% - Result needed: per dataset and downstream model, how often the winner used the matched evaluator, and its gap to the best other evaluator.
% Paragraph 5: how the setting values were chosen (her paragraph "We ran a series of experiments...")
% - Smaller design runs came before the full runs, each dataset scored in its own metric.
% - A value stayed if it won for some dataset, downstream model and budget. A value that won but cost many hours stayed in the full search and left the fast tier.
% - A setting was fixed at one value if the spread between values was smaller than the spread across folds.
% - Example: the DNN was added as a downstream model because it scored best on 4 of 5 datasets (preliminary).
% - The full list of design runs is in Additional files, Section~\ref{app:searchspace}.
% Paragraph 6: size and caching (her paragraphs "The full search tests..." and "The complete search result is cached...")
% - 8,160 configurations per dataset holds where all five downstream models run, fewer where a dataset drops one. Breakdown in Additional files.
% - Importance and standalone scores are computed once per fold and reused, so cost grows more slowly than the configuration count.
% - A fast search can be extended to a full one.
% ---- END NOTES ----

% TODO: review, needs manual review
Features that improve predictions under RF may not necessarily help a kNN model, and vice versa. This is because variable importance depends fundamentally on the model, as Dong and Rudin showed \citep{Dong2020}, and as we found in our feature importance scoring analysis. To address this, we use an internal evaluator which scores each base representation on its own with a model whose inductive bias best fits the downstream model. We use the following evaluators: RF, DNN, LightGBM, and the special case of matched, which matches the internal evaluator to the final downstream evalutors. We fit each on the internal training part of the training fold and score it on the internal validation part.

% TODO: review, needs manual review
This internal evaluator is used to score each representation, and those scores set how the budget is split and which representations are kept within the larger hybrid search. Thus, it is another axes we sweep while searching for the best set of features.

% TODO: need additional files to specify architectures

% TODO: does match always work? custom downstream model?

% TODO: these numbers are outdated
% TODO: review, needs manual review
The full search tests all combinations of budgets, allocations, importance methods, representation counts, downstream models, and internal evaluators under 5-fold cross-validation: 8,160 configurations per dataset when using 5 downstream models. We approach this problem combinatorially as it allows us to measure which factors contribute most to performance variance (model choice, budget, allocation, or their interactions) rather than assuming any single axis dominates.

% TODO: mention representation knobs

% TODO: review, needs manual review
The complete search result is cached to disk. The search runs once; everything after that is a lookup.

When I began to implement this idea, I figured that there would be a fairly straightforward way to select the most informative features. However, I began to ask some questions. How do you measure feature importance? If you're starting with a handful of input representations, how many features do you take from each? How do you score these potential feature sets?  Unfortunately, my one-size-fits-all hypothesis rang true on multiple fronts. As I found out, different prediction tasks prefer different answers to all of those questions. So to figure out how to even go about selecting the best features, I realised I had to sweep a search space.

% TODO: explain framework experiments
We ran a series of experiments to determine the best variables to include in the search space.
When making a selection, we took into account how many potential hybrid representations would be gained by including that variable, as well as the compute time


\subsection{Augmentation}
% ---- NOTES: Augmentation ----
% Takeaway: Small or specialised feature sets cannot compete fairly in the budget split, so they are offered afterwards as extras, ranked by how well each explains what the hybrid gets wrong.
% Takeaway: Augmentation also lets a chemist keep one trusted representation and learn which few outside features help, and how many.
% Paragraph 1: why some feature sets sit outside the budget split (her first paragraph)
% - Some are too small to compete for budget: MACCS keys are 167 features, ESM-2 log-likelihood ratios are 3.
% - Others mean something only next to a full representation, such as mutation features.
% - Putting them in the budget split distorts it; leaving them out loses signal. So they are offered after the hybrid is built.
% - Change: the prose describes only the uncapped mode. Add the capped mode.
% Paragraph 2: which feature sets are offered (answers her TODO for a better list)
% - Set per dataset: the Augmentation column of Table~\ref{tab:dataset_reps}.
% - MACCS is usually an augmentation set; on MDR1-MDCK it is a base representation. QM9 offers topological torsions instead.
% - Common thread: each is small or describes one aspect only (shape, charge, the mutated positions, one tumour property). A plausible supplement, and weak on its own.
% - Change: the Coulomb matrix set is not a fingerprint. It is the 50 largest eigenvalues of a matrix built from nuclear charges and 3D interatomic distances.
% - Spelling: "Coloumb"; "sames" should be "samples" or "variants".
% - Why 3D sets are worth offering: joining 2D and 3D descriptors gave many more significant models on six protein--ligand series \citep{Bahia2023}. Their 3D descriptors came from bioactive conformations; here one generated conformer is used.
% Paragraph 3: ranking the extra features by what the hybrid gets wrong
% - Inside each outer training fold, a small random forest is fitted on the hybrid's features, and its out-of-fold residuals $\mathbf{r}$ are kept.
% - Every offered feature is scored by $|\rho_{\mathrm{P}}(\mathbf{x}_j, \mathbf{r})|$, its absolute Pearson correlation with the residuals. The top 50 go on.
% - Why residuals and not labels: ranking against $\mathbf{y}$ favours features that repeat what the hybrid already knows.
% - Same idea as each stage of gradient boosting fitting the current residual \citep{friedman2001greedy}.
% Paragraph 4: how many extras to keep
% - The top 0, 5, 10, 20 or 40 are each scored with the downstream model on the internal validation molecules.
% - 0 is scored first and wins ties, so "no extras" is the default outcome, not a failure.
% - Five whole-block choices leave less room to overfit a small validation set than 50 one-by-one decisions \citep{Cawley2010}.
% Paragraph 5: capped and uncapped
% - Uncapped: extras are appended, so the width is at most $B + 40$ features.
% - Capped: each extra replaces the weakest remaining hybrid feature by random forest importance, so the width stays at $B$.
% - Why both: uncapped against no augmentation mixes "better features" with "more features". Capped removes the width difference.
% - Each is a separate search, not a searched setting. 14 of the 16 datasets run all three conditions; humanness and trastuzumab--HER2 have no augmentation sets.
% - Table: the three conditions. Columns: final width (features); what an extra feature must beat; datasets that run it.
% Paragraph 6: the second use, one strong representation
% - Keep one trusted representation at full width and add a few extras, uncapped. A chemist learns which few outside features help, and how many.
% - It is tested with decoy features planted among the candidates (Comparators, controls and ablations).
% - Limit, one sentence: the correlation sees only linear relations, one feature at a time, and not redundancy among the extras.
% - Figure: flow diagram. Left: the hybrid and the augmentation sets. Middle: forest on the hybrid, residuals, augmentation features ranked by correlation with them, top 50. Right: 0 to 40 extras chosen on internal validation, then two branches: uncapped appends; capped swaps out the weakest hybrid features.
% - Result needed: how often 0 extras is chosen, per dataset and downstream model; which augmentation sets are kept most often, per domain.
% ---- END NOTES ----

% TODO: review, also I think this is out of date
% TODO: review, needs manual review
Some features are too low-dimensional or too specialized to participate in the budget allocation stage: MACCS fingerprints (167 features), for instance, or log-likelihood ratios (3 features). Excluding them from the allocation would lose potentially useful signal; including them would distort the budget distribution because they cannot contribute many features. We handle these as augmentation: after the base hybrid is constructed, we pool the augmentation features, rank them by how strongly each correlates with the out-of-fold errors of a small random forest fitted on the base features, and keep the top 0, 5, 10, 20 or 40 of them, whichever scores best with the downstream model on the internal validation set.

% TODO: I need a better list to describe augmentations
Augmentation features range from smaller fingerprints like MACCS or Coloumb matrices, sets of structural information like 3D pharmacophores, and features describing only the differences between the sames like mutation features for antibodies.
% TODO: add augmentation subsection


\subsection{Fast tier and adaptive pruning}
% ---- NOTES: Fast tier and adaptive pruning ----
% Takeaway: The fast tier searches shorter value lists, fixed in advance, and accepts a possibly worse answer. Adaptive pruning drops values that clearly lose on the first two folds and aims to find the same winner.
% Paragraph 1: two ways to make the search cheaper
% - The fast tier uses shorter value lists, fixed in advance and the same on every fast-tier dataset. It accepts a possibly worse answer.
% - Adaptive pruning decides during the search, from this dataset's own fold scores. It aims to find the same winner.
% Paragraph 2: which datasets run the fast tier
% - Six datasets, marked in Table~\ref{tab:dataset_reps}. Their full search would take more than about 40 hours, or they have more than 10,000 samples. Example: QM9's full search was estimated at about 103 hours.
% - The dropped values are still searched on every other dataset. The fast tier is a compute compromise, not a claim that those values are useless.
% - A fast run can be extended to a full one later, scoring only the skipped configurations.
% Paragraph 3: how the fast lists were chosen (preliminary)
% - A value stays if it is cheap and adds to the best score. The run times and scores behind each dropped value belong in Additional files, Section~\ref{app:searchspace}.
% - Result needed: the share of run time the fast tier and pruning saved on the full runs.
% ---- END NOTES ----

\subsubsection{Search tier settings}
% ---- NOTES: Search tier settings ----
% Takeaway: Every dataset is searched at one of two tiers, fixed before the search. The fast tier shortens three settings and keeps the rest.
% - Every dataset is searched at one of two tiers, full or fast, fixed per dataset before the search starts.
% - The fast tier shortens three settings: allocation methods, importance methods and representation counts. Downstream models, internal evaluators, budgets and the 5 folds are the same at both tiers.
% ---- END NOTES ----

% TODO: review, needs manual review
Table~\ref{tab:search_tiers} lists, for each search setting, the values tried at the full tier and at the fast tier.

% TODO: review, needs manual review
\begin{table}[htbp]
\centering
\caption{Values tried for each search setting at the two search tiers. Both tiers use five folds; configurations removed by adaptive pruning are scored on the first two folds only.}
\label{tab:search_tiers}
\begin{tabular}{@{}p{0.22\linewidth}p{0.40\linewidth}p{0.30\linewidth}@{}}
\toprule
Setting & Full tier & Fast tier \\
\midrule
Downstream models & Random forest, XGBoost, LightGBM, SVM$^{a}$, DNN; Gaussian process$^{b}$ & Random forest, XGBoost, LightGBM, SVM$^{a}$, DNN \\
Internal evaluators & The downstream model itself, random forest, LightGBM, DNN & Same as full \\
Feature budgets (features) & 10, 25, 50, 100, 250, 500, 1{,}000, 2{,}000 & Same as full \\
Allocation methods & Linear, squared, softmax, rank decay, global pool & Linear, squared, rank decay \\
Importance methods & LightGBM gain, DeepLIFT, TreeSHAP & LightGBM gain \\
Representation counts & 2, 3, 5, all representations & 3, all representations \\
\bottomrule
\end{tabular}

\smallskip
{\footnotesize $^{a}$SVM is not fitted above 10{,}000 molecules (5{,}000 when augmentation features are added outside the budget). $^{b}$Added only on Humanness. Of the 16 datasets, it is the only one that is under 1{,}000 molecules, has its size recorded and is not excluded for run time.}
\end{table}

\subsubsection{Adaptive pruning}
% ---- NOTES: Adaptive pruning ----
% Takeaway: Pruning scores every configuration on two folds, drops values that are rarely among the best within groups that fix every other setting, and never removes every value of a setting.
% Takeaway: Its cut-offs were set on five datasets so that no best score changed.
% Paragraph 1: what pruning is for
% - Some configurations are so far behind after two folds that the other three would only confirm a loss. Pruning stops scoring those.
% - Its one criterion: never remove a configuration that would have been the best. Time saved is secondary.
% - Related: successive halving \citep{Jamieson2016} and Hyperband \citep{Li2018Hyperband} drop configurations on partial evaluations. Here the partial evaluation is two of the five folds.
% - As submitted, pruning is on for every dataset.
% Paragraph 2: the steps (her paragraph; one change)
% - Change: configurations are ranked by their mean score over folds 1 and 2, in the dataset's own metric.
% - Why grouped: a value that is best for a random forest at 25 features is not removed for losing to a DNN at 500.
% - Why two folds: after one fold on PAM50 subtype, 726 of 1,020 groups had every value tied. A one-fold vote would remove values by list order.
% - A vote can never remove every value of a setting. The winner is still picked by the five-fold mean.
% Paragraph 3: how the cut-offs were set (preliminary)
% - Five datasets were scored unpruned on all five folds, and each rule was replayed on the stored scores.
% - A cut-off counts as safe if the best score is unchanged for every (downstream model, budget) pair on every dataset.
% - Limits: five datasets is few. The calibration applied the votes one after another; the search applies them at once.
% Paragraph 4: effect on the reported score
% - Folds 1 and 2 both choose the survivors and score them, so the winner's five-fold mean is slightly optimistic \citep{Cawley2010}.
% - Single representations are never pruned. A fair comparison leaves both voting folds out of both sides, or re-scores without them.
% - The calibration (run times and the share of configurations left at each vote) belongs in Additional files, Section~\ref{app:searchspace}.
% ---- END NOTES ----

% TODO: review, needs manual review
Adaptive pruning: every configuration is first scored on the first two folds. At budget 100, within each (model, evaluator, representation count) group, (allocation, importance) pairs with no configuration in the better half are not searched at larger budgets. Each axis is then voted on within groups that fix every other axis: an allocation method is pruned if it is in the best 2 in fewer than 20\% of groups, a representation count if it is best in fewer than 10\%, an importance method if it is in the best 2 in fewer than 20\%, and an internal evaluator if it is in the best 3 in fewer than 10\%. Survivors are scored on the remaining folds.

\section{Methods: experimental setup}
% ---- NOTES: Methods: experimental setup ----
% Takeaway: Every experiment reuses one protocol: 5 outer group folds, the search on training molecules only, and a paired comparison with the best single representation on the same held-out molecules.
% Paragraph 1: the shared protocol
% - Each of the 16 datasets is split into 5 outer folds, and no group is split. The antibody expression scan in the transfer results used 3; say so there.
% - The search sees only the training molecules of the outer fold. The selected features are rebuilt on them and applied unchanged to the held-out molecules.
% Paragraph 2: the comparator
% - The best single base representation at its full width: what a user would otherwise run.
% - It is chosen from the search's own single-representation scores on the training molecules, then scored on the same held-out molecules. It can differ between folds.
% - It is often many times wider than the hybrid. Record both widths per fold.
% - Per fold: hybrid score minus best-single score, signed so that positive means the hybrid is better.
% - One sentence: Chapter 6 reuses this protocol with noise added to the training labels only.
% ---- END NOTES ----


\subsection{Datasets and splitting}
% ---- NOTES: Datasets and splitting ----
% Takeaway: 16 datasets from six domains, each grouped so that a held-out fold means what "unseen" means in its field.
% - Change: the metrics paragraph disagrees with Table~\ref{tab:datasets} in four ways. It gives $R^2$ for reaction yield; both reaction datasets use RMSE. It names concordance index for drug-target affinity; none of the 16 datasets is one. It does not name MCC, used on the two nucleotide sets. It lists tumour subtype under AUROC; PAM50 subtype is four classes, scored with weighted F1.
% - Change: the metrics paragraph could shrink to its first sentence plus one sentence saying each dataset's metric is in Table~\ref{tab:datasets}, with the reasons in "Models and metrics". That also removes the four disagreements.
% - Change: the no-random-splits paragraph says "5-fold GroupKFold throughout". Classification datasets use stratified group $k$-fold.
% - Change: the no-random-splits paragraph repeats Chapter 2, "Evaluation methodology", with the same four citations. It could shrink to one sentence citing that paragraph.
% - Change: the internal-split paragraph repeats "From base representations to a hybrid". It could shrink to one sentence.
% - Result needed: a last column for Table~\ref{tab:datasets}, the best single base representation per dataset and its score.
% ---- END NOTES ----

% TODO: review, needs manual review
The hybrid experiments use 16 datasets from six domains: five small-molecule ADME and toxicity endpoints, one quantum-chemical property, two reaction-yield sets, four antibody sets, two nucleotide-sequence tasks and two cancer classification tasks from TCGA (Table~\ref{tab:datasets}). Every score in this chapter is computed on held-out folds that keep each group whole. Each dataset is grouped to match what ``unseen'' means in its field. The curation of each dataset, with counts before and after each step, is given in Appendix~\ref{app:datasets}; the representations each dataset uses are listed in Table~\ref{tab:dataset_reps}.

% TODO: review, needs manual review
Different biological domains have established different conventions for what constitutes a good prediction, and applying a single metric everywhere would misrepresent performance in at least some of them. To this end, we assign each task a primary metric matched to its domain: RMSE for molecular property prediction, Spearman rank correlation for antibody regression and developability, AUROC for binary classification (including antibody classification and tumor subtype), weighted F1 for multiclass classification, concordance index for drug-target affinity, and $R^2$ for reaction yield.

% TODO: review, needs manual review
No random splits are used. Random splitting inflates performance estimates for molecular data by placing structurally similar compounds on both sides of the train/test boundary \citep{Yang2019, Wallach2018}, and position-level leakage in deep mutational scanning data allows models to memorise site-specific effects rather than learning generalisable sequence-fitness relationships \citep{Notin2023, Vieira2026}. Splits therefore match domain conventions and use 5-fold GroupKFold throughout.

% TODO: review, needs manual review
\begin{table}[htbp]
\centering
\small
\setlength{\tabcolsep}{3pt}
\caption{The 16 datasets used in the hybrid experiments. One row is one dataset after curation. $n$ is the number of samples; the last column gives what the folds are grouped by and the number of groups. $\downarrow$: lower is better; $\uparrow$: higher is better.}
\label{tab:datasets}
\begin{tabular}{p{3.0cm}p{3.3cm}p{2.3cm}p{1.1cm}p{3.5cm}}
\toprule
Dataset & Label (unit) & Task, metric & $n$ & Folds grouped by (groups) \\
\midrule
\multicolumn{5}{l}{\textit{Small molecules: ADME and toxicity}} \\
hERG $K_i$ \citep{Zdrazil2023} & p$K_i$ & regression, RMSE $\downarrow$ & 1,412 & scaffold (769) \\
LogD \citep{openadmet} & LogD (log units) & regression, RMSE $\downarrow$ & 5,028 & scaffold (2,413) \\
Caco-2 efflux \citep{openadmet} & log$_{10}$ efflux ratio & regression, RMSE $\downarrow$ & 2,158 & scaffold (1,086) \\
Rat liver microsome clearance \citep{fang2023, polaris} & log$_{10}$ CL$_\text{int}$ (mL/min/kg) & regression, RMSE $\downarrow$ & 3,053 & scaffold (2,625) \\
MDR1--MDCK efflux \citep{fang2023, polaris} & log$_{10}$ efflux ratio & regression, RMSE $\downarrow$ & 2,641 & scaffold (2,318) \\
\midrule
\multicolumn{5}{l}{\textit{Quantum chemistry}} \\
QM9 \citep{Ramakrishnan2014, qm9_dataset} & HOMO--LUMO gap (eV) & regression, RMSE $\downarrow$ & 9,978 & scaffold (3,452) \\
\midrule
\multicolumn{5}{l}{\textit{Chemical reactions}} \\
Buchwald--Hartwig yield \citep{Ahneman2018} & yield (\%) & regression, RMSE $\downarrow$ & 3,955 & additive (22) \\
Buchwald--Hartwig plates \citep{Ahneman2018} & yield (\%) & regression, RMSE $\downarrow$ & 4,312 & additive (24) \\
\midrule
\multicolumn{5}{l}{\textit{Antibodies}} \\
Koenig expression \citep{Koenig2017, Chungyoun2024} & expression enrichment ratio (no unit) & regression, Spearman $\uparrow$ & 4,275 & mutated position modulo 5 (5) \\
Nanobody thermostability \citep{ValdesTresanco2023, Chungyoun2024} & melting temperature ($^\circ$C) & regression, Spearman $\uparrow$ & 672 & 95\% identity cluster (309) \\
Humanness \citep{Prihoda2022} & human or not (198 : 366) & binary, AUROC $\uparrow$ & 564 & 95\% identity cluster (461) \\
Trastuzumab--HER2 binding \citep{Chinery2024} & binder or not (1,639 : 3,361) & binary, AUROC $\uparrow$ & 5,000 & mutations from wild type (10) \\
\midrule
\multicolumn{5}{l}{\textit{Nucleotide sequences}} \\
H3K4me3 \citep{DallaTorre2025} & mark present or not (9,122 : 9,122) & binary, MCC $\uparrow$ & 18,244 & chromosome (23) \\
TATA promoter \citep{DallaTorre2025} & promoter or not (2,637 : 2,637) & binary, MCC $\uparrow$ & 5,274 & chromosome (23) \\
\midrule
\multicolumn{5}{l}{\textit{Cancer genomics (TCGA)}} \\
Lung adenocarcinoma vs squamous \citep{TCGA2014LUAD, TCGA2012LUSC, Goldman2020} & tumour type (515 : 502) & binary, AUROC $\uparrow$ & 1,017 & tissue source site (68) \\
Breast cancer PAM50 subtype \citep{TCGA2012BRCA, Parker2009, Goldman2020} & subtype (421 : 192 : 141 : 67) & 4 classes, weighted F1 $\uparrow$ & 821 & tissue source site (19) \\
\bottomrule
\end{tabular}
\end{table}

% TODO: review, needs manual review
\paragraph{Pre-processing.} The steps shared across datasets are listed below. The per-dataset steps and the counts after each one are in Appendix~\ref{app:datasets}.
\begin{enumerate}[leftmargin=*]
    \item Every molecule is parsed with RDKit \citep{rdkit}, and molecules that fail are dropped. On hERG $K_i$, LogD and Caco-2 the structures are also standardised: largest fragment kept, charges neutralised, canonical tautomer, and molecules containing a metal removed.
    \item Repeat measurements of one molecule are merged into their median (hERG $K_i$, LogD, Caco-2). hERG $K_i$ molecules whose repeat measurements have a standard deviation above 1.0 p$K_i$ units are removed. $K_i$ values combined from different sources are a known source of large label noise \citep{landrum2024}.
    \item Duplicate molecules are removed by InChIKey on the six small-molecule datasets, keeping the first.
    \item LogD and Caco-2 labels are clipped at the 1st and 99th percentiles of the dataset.
    \item Two datasets are subsampled once with a fixed seed: 10,000 of the 129,428 parsable QM9 molecules, and 5,000 of the 524,346 trastuzumab variants, drawn in proportion to each combination of label and number of mutations.
\end{enumerate}

% TODO: review, needs manual review
\paragraph{Splitting rules.} Four rules hold for every dataset.
\begin{enumerate}[leftmargin=*]
    \item Folds are formed as in Chapter~\ref{chap:background}, ``Evaluation methodology'': no group is split, regression datasets use group $k$-fold, and classification datasets use stratified group $k$-fold.
    \item There are five outer folds, and the stratified folds use a fixed seed.
    \item No split is ever random. What happens when a dataset has fewer groups than folds is stated in Chapter~\ref{chap:background}, ``Evaluation methodology''. Every dataset here has at least five groups, so all outer splits have five folds.
    \item On a given dataset, every representation, model and hybrid is scored on the same five folds.
\end{enumerate}

% TODO: review, needs manual review
The hybrid construction step uses a nested internal train/validation split within each training fold to score representations for feature allocation. This internal split is group-aware, taking the first fold of a grouped 5-fold split (GroupKFold) with the same grouping variable as the outer cross-validation, to prevent the allocation evaluator from seeing related samples on both sides of its split.

% TODO: review, needs manual review
\paragraph{Groups by domain.} The small-molecule and QM9 datasets are grouped by Bemis--Murcko scaffold, defined in Chapter~\ref{chap:background}, ``Evaluation methodology''. In this chapter the scaffold keeps stereochemistry, and all acyclic molecules share one empty scaffold. On QM9 this acyclic group holds 1,041 of the 9,978 molecules and is the largest group. On the two Biogen datasets (rat liver microsome clearance and MDR1--MDCK efflux) there are 1.2 and 1.1 molecules per scaffold (the $n$ and group columns of Table~\ref{tab:datasets}). On the Biogen datasets, holding out a scaffold is therefore close to holding out single molecules. Chapter~\ref{chap:robustness} used a different scaffold key and different curation for its versions of LogD, Caco-2, hERG $K_i$ and QM9; Appendix~\ref{app:datasets} lists the differences.

% TODO: review, needs manual review
The reaction datasets are Buchwald--Hartwig C--N couplings run on high-throughput plates, in which aryl halide, ligand, base and additive are varied. They are grouped by additive, following the out-of-sample tests of \citet{Ahneman2018}, so each held-out fold contains four or five additives the model never saw in training. The two sets are not independent: they share 3,955 reactions (Appendix~\ref{app:datasets}).

% TODO: review, needs manual review
On the Koenig deep mutational scan of the anti-VEGF antibody G6, every variant carries one mutation. Variants are grouped by the position of that mutation modulo 5, as in ProteinGym \citep{Notin2023}, so each fold is one residue class and holds 836 to 874 variants. All mutations at one position then fall in the same fold, and a model cannot learn a per-position effect and look it up at test time \citep{Vieira2026}. The nanobody and humanness sets are grouped by single-linkage clusters at 95\% sequence identity, after \citet{Kenlay2024}. The identity measure, the cluster sizes and the humanness labels are given in Appendix~\ref{app:datasets}. The trastuzumab variants differ from the wild type only in the ten residues of CDR-H3, and are grouped by how many of those residues are mutated. A held-out fold therefore contains variants at mutation counts not seen in training.

% TODO: review, needs manual review
The two nucleotide tasks come from the revised Nucleotide Transformer benchmark \citep{DallaTorre2025}. Both are grouped by chromosome, so a held-out fold contains whole chromosomes the model never saw in training: between two and six of the 23 chromosome groups (Appendix~\ref{app:datasets} gives the task definitions and how the released files are regrouped).

% TODO: review, needs manual review
The TCGA datasets use RNA-seq expression and other measurements of primary tumours, one sample per patient, obtained through UCSC Xena \citep{Goldman2020}. Samples are grouped by tissue source site, the institution that supplied the sample, which is read from the TCGA barcode. A model therefore cannot score well by recognising the site a sample came from, and no patient spans training and test.

% TODO: review, needs manual review
\paragraph{Consequences of the grouping.} Group definitions differ between domains on purpose; the one shared rule is that no group spans training and test. Some groupings leave few or very unequal groups, and the folds inherit that. The trastuzumab held-out folds range from 13 to 3,035 variants, with binder shares from 15.5\% to 92.3\% against 32.8\% overall. The breast cancer folds range from 3 to 403 tumours. None of the 68 lung sites supplied both tumour types, so the lung folds are unbalanced: one holds 20 adenocarcinomas and 113 squamous tumours. The Koenig folds are fixed by its five groups. Per-fold scores on these datasets spread more, so results are read per dataset, in the dataset's own metric, and never pooled across datasets.
\subsection{Representations by domain}
% ---- NOTES: Representations by domain ----
% Takeaway: Each dataset gets up to ten base representations that encode the input in different ways, chosen before the search, plus a few augmentation feature sets.
% - Change: "ten" in the sentence on pre-selecting base representations should read "at most ten". H3K4me3 keeps 8 and TATA promoter keeps 9.
% - Result needed: per dataset, the most and least alike pair of base representations by correlation of out-of-fold predictions.
% - The five per-domain tables (about 100 rows) belong in an Additional file. The main text keeps the domain paragraphs and Table~\ref{tab:dataset_reps}.
% ---- END NOTES ----
% TODO: review, needs manual review

We combine a variety of component, or base representations to create these hybrids.

Any representation that gives one fixed-length vector per sample can serve as a base representation, as set out in the description of how the hybrid is built earlier in this chapter. The first version of the method used three molecular representations: PDV, ECFP4 and MHG-GNN. These three are now one small-molecule example among many. Every base representation is frozen: it is computed once from the raw input (a SMILES string, a sequence, a reaction or a patient profile) and is never trained on the labels.

For every data set in this study, we pre-selected ten base representations that spanned across different domains of information and showed some positive results when used on their own for the given prediction task.

To choose them, every candidate representation was first scored on its own with a gradient-boosted tree model \citep{ke2017lightgbm} under the dataset's grouped folds. Candidates were then compared by the Spearman correlation between their out-of-fold predictions, since two representations whose predictions correlate highly carry much the same signal. At most ten were then chosen by hand, looking at each candidate's own score next to these correlations and favouring representations that encode the input in different ways. Two nucleotide datasets have fewer. H3K4me3 had 12 candidates, and four of them predicted almost the same as a better-scoring candidate, with a Spearman correlation of 0.95 or above between their out-of-fold predictions; eight remain. TATA promoter had 10 candidates, one was dropped by the same rule, and nine remain. Augmentation features do not count towards the ten. They are added after the hybrid is built (see the Augmentation subsection above), and the same representation can be a base representation on one dataset and augmentation on another.

If an input cannot be parsed or embedded, it gets an all-zero row, so that every representation of a dataset keeps the same rows. In computed descriptors, a missing or infinite value is set to 0, and values are stored unscaled. A pretrained model gives one vector per input, the mean of its last-layer token vectors with padding excluded, unless the tables say otherwise. Scaling and filling of missing values happen inside each fold, on the training rows only, as described in the construction of the hybrid.

\paragraph{Small molecules.} The small-molecule representations (Table~\ref{tab:reps_smallmol}) are the fingerprint, descriptor, graph and language-model families described in Chapter~\ref{chap:background}. The only additions are four sets of 3D shape and pharmacophore features, used as augmentation only. All four are computed from a single conformer per molecule, embedded with a fixed seed and optimised with the MMFF force field. A single conformer may misdescribe flexible molecules.

\paragraph{Chemical reactions.} Reaction datasets (Table~\ref{tab:reps_reactions}) use three kinds of representation. The condensed graph of reaction merges reactants and products into one graph in which the bonds that change are marked \citep{Hoonakker2011}; here it is hashed into a circular fingerprint after atom mapping \citep{Schwaller2021rxnmapper}. Reaction language models read the reaction SMILES directly \citep{Schwaller2021rxnfp, Sagawa2023ReactionT5}. Any molecular representation can also be applied to the reactants, to the product, or to their difference. A one-hot encoding of which aryl halide, ligand, base and additive was used carries no chemistry. It is kept because on plate data such identity features can predict yield about as well as chemical descriptors \citep{Chuang2018}.

\paragraph{Antibodies.} Antibody representations (Table~\ref{tab:reps_antibodies}) start from the heavy- and light-chain amino-acid sequences. Besides one-hot and BLOSUM62 encodings \citep{Henikoff1992}, most are frozen language models, either trained on all proteins or on antibody repertoires only. Most models embed each chain separately and join the two vectors, so their width doubles on datasets with both chains. Several representations read one model in different ways, for instance averaging over earlier layers or over the CDR residues only, with CDRs found by IMGT numbering \citep{Lefranc2003} using ANARCI \citep{Dunbar2016ANARCI}. Datasets of point variants of one parent antibody also use features that describe each mutation.

\paragraph{Nucleotide sequences.} A $k$-mer frequency vector counts how often each DNA word of length $k$ occurs in a sequence, divided by the total number of words \citep{Vinga2003}; the canonical variant merges each word with its reverse complement \citep{Teeling2004}. DNA language models are networks pretrained on genomes to predict masked or next tokens \citep{Zhou2023DNABERT2, DallaTorre2025, Nguyen2023HyenaDNA}. Table~\ref{tab:reps_nucleotides} lists both kinds, along with the small sequence summaries used as augmentation.

\paragraph{Cancer patients.} For the two TCGA datasets (Table~\ref{tab:reps_tcga}), each representation is a different measurement of the same patient's tumour, not a different encoding of one input. The measurements are RNA-seq gene expression (read through foundation models or summarised per chromosome arm), DNA methylation on the HM450 array, gene-level copy number, microRNA abundance, protein levels from reverse-phase protein arrays, and embeddings of the diagnostic whole-slide image. Data were obtained from TCGA through UCSC Xena \citep{TCGA2013, Goldman2020}. The most-variable selections, the z-scoring and the principal components for methylation, copy number, microRNA and protein were fitted once on all patients in a cohort, held-out patients included. No labels are used in these steps. A patient with no measurement for one data type gets an all-zero row for that representation.

We list all of these representations alongside their associated data sets and prediction tasks below:

\begin{longtable}{p{3.2cm}p{5.6cm}p{1.6cm}p{1.6cm}p{1.4cm}}
\caption{Representations used for small molecules. One row is one representation. Role: B = base representation, A = augmentation, B/A = base on some datasets and augmentation on others.}\label{tab:reps_smallmol}\\
\toprule
Representation & What it encodes & Features & Source & Role \\
\midrule
\endfirsthead
\toprule
Representation & What it encodes & Features & Source & Role \\
\midrule
\endhead
\bottomrule
\endfoot
ECFP4 & hashed circular substructures, radius 2 & 2048 & \citep{rogers2010} & B \\
Atom pairs & hashed atom-type pairs with bond distance & 2048 & \citep{Carhart1985} & B \\
Topological torsions & hashed paths of four bonded atoms & 2048 & \citep{Nilakantan1987} & B/A \\
Avalon & hashed paths, rings and atom environments & 2048 & \citep{Gedeck2006} & B \\
MACCS keys & 166 predefined substructure keys & 167 & \citep{Durant2002} & B/A \\
WL graph kernel & Weisfeiler--Lehman label counts, 5 rounds, labels present in at least 5\% of molecules & 140--317 & \citep{Shervashidze2011} & B \\
PDV & RDKit 2D physicochemical descriptors & 200 & \citep{rdkit} & B \\
Mordred & 2D molecular descriptors & 1613 & \citep{Moriwaki2018} & B \\
GROVER & frozen pretrained graph transformer & 4800 & \citep{Rong2020GROVER} & B \\
MHG-GNN & frozen pretrained graph encoder & 1024 & \citep{kishimoto2023} & B \\
GraphMVP & frozen 2D graph encoder pretrained with 3D views & 300 & \citep{Liu2022GraphMVP} & B \\
MoLFormer & frozen SMILES language model; public checkpoint trained on 10\% of its corpus & 768 & \citep{ross2022} & B \\
SMI-TED & frozen SMILES encoder--decoder & 768 & \citep{Soares2025} & B \\
ChemBERT & frozen SMILES language model; community checkpoint trained on ChEMBL following the cited method & 256 & \citep{Zhang2022ChemBERT} & B \\
ChemBERTa-2 & frozen SMILES language model & 384 & \citep{Ahmad2022} & B \\
SELFormer & frozen SELFIES language model & 768 & \citep{Yuksel2023} & B \\
Mol2vec & summed substructure word vectors & 300 & \citep{Jaeger2018} & B \\
3D pharmacophore & counts of 7 pharmacophore types, and mean and spread of distances for the 21 type pairs, one conformer & 49 & \citep{rdkit} & A \\
USR & shape moments from 4 reference points, one conformer & 12 & \citep{Ballester2007} & A \\
USRCAT & USR per pharmacophore atom type, one conformer & 60 & \citep{Schreyer2012} & A \\
Coulomb matrix & sorted eigenvalues of the Coulomb matrix, one conformer & 50 & \citep{Rupp2012} & A \\
\end{longtable}

\begin{longtable}{p{3.2cm}p{5.6cm}p{1.6cm}p{1.6cm}p{1.4cm}}
\caption{Representations used for chemical reactions. One row is one representation. Role: B = base representation, A = augmentation, B/A = base on some datasets and augmentation on others.}\label{tab:reps_reactions}\\
\toprule
Representation & What it encodes & Features & Source & Role \\
\midrule
\endfirsthead
\toprule
Representation & What it encodes & Features & Source & Role \\
\midrule
\endhead
\bottomrule
\endfoot
CGR fingerprint & circular fingerprint, radius 2, hashed from the condensed graph of reaction & 2048 & \citep{Hoonakker2011} & B \\
rxnfp & frozen reaction language model & 256 & \citep{Schwaller2021rxnfp} & B \\
ReactionT5 & frozen reaction language model & 768 & \citep{Sagawa2023ReactionT5} & B \\
Component one-hot & which aryl halide, ligand, base and additive & 45 & \citep{Ahneman2018} & B \\
GROVER (reactants) & GROVER averaged over reactants & 4800 & \citep{Rong2020GROVER} & B \\
Mordred (reactants) & Mordred averaged over reactants & 1613 & \citep{Moriwaki2018} & B \\
Mol2vec (reactants) & Mol2vec averaged over reactants & 300 & \citep{Jaeger2018} & B \\
ECFP4 (reactants) & ECFP4, bitwise OR over reactants & 2048 & \citep{rogers2010} & B \\
Avalon (difference) & product minus reactant Avalon bits & 2048 & \citep{Gedeck2006} & B \\
MHG-GNN (difference) & product minus reactant embedding & 1024 & \citep{kishimoto2023} & B \\
MACCS (difference) & product minus reactant MACCS keys & 167 & \citep{Durant2002} & B \\
Mol2vec (difference) & product minus reactant Mol2vec & 300 & \citep{Jaeger2018} & B \\
MACCS (reactants) & MACCS keys, bitwise OR over reactants & 167 & \citep{Durant2002} & A \\
MACCS (product) & MACCS keys of the product & 167 & \citep{Durant2002} & A \\
3D pharmacophore (reactants) & 3D pharmacophore averaged over reactants & 49 & \citep{rdkit} & A \\
3D pharmacophore (product) & 3D pharmacophore of the product & 49 & \citep{rdkit} & A \\
\end{longtable}

\begin{longtable}{p{3.2cm}p{5.6cm}p{1.6cm}p{1.6cm}p{1.4cm}}
\caption{Representations used for antibodies. One row is one representation. LM = language model. Widths marked /chain double on datasets with heavy and light chains (Koenig expression, Trastuzumab--HER2) and are single on the nanobody and heavy-chain-only humanness datasets. -- in Source: a standard encoding or one built for this thesis. Role: B = base representation, A = augmentation, B/A = base on some datasets and augmentation on others.}\label{tab:reps_antibodies}\\
\toprule
Representation & What it encodes & Features & Source & Role \\
\midrule
\endfirsthead
\toprule
Representation & What it encodes & Features & Source & Role \\
\midrule
\endhead
\bottomrule
\endfoot
One-hot & one-hot residues, both chains padded to fixed length & 5600 & -- & B \\
BLOSUM62 & BLOSUM62 row per residue, both chains padded & 5600 & \citep{Henikoff1992} & B \\
Amino-acid composition & fraction of each amino acid, per chain & 40 & -- & A \\
ESM-2 & protein LM, last layer, mean over residues & 1280/chain & \citep{Lin2023} & B \\
ESM-2 (joined chains) & ESM-2 on heavy and light chain joined into one string & 1280 & \citep{Lin2023} & B \\
ESM-2 early layers & ESM-2, mean over layers 0--16 & 1280/chain & \citep{Lin2023} & B \\
ESM-2 middle layers & ESM-2, mean over layers 20--33 & 1280/chain & \citep{Lin2023} & B \\
ESM-2 CDR & ESM-2, mean over CDR residues only & 1280/chain & \citep{Lin2023} & B \\
ESM-2 per CDR & ESM-2, CDR1/2/3 pooled separately then joined & 3840/chain & \citep{Lin2023} & B \\
ESM-2 variable positions & ESM-2, mean over the 20 highest-entropy positions & 1280/chain & \citep{Lin2023} & B \\
ESM-2 shift from parent & ESM-2 variant minus parent embedding & 1280/chain & \citep{Lin2023} & B \\
ESM-2 log-likelihood ratio & masked mutant vs parent log-probability: sum, mean, min & 3 & \citep{Meier2021} & A \\
ESM3 & protein LM & 1536/chain & \citep{Hayes2025} & B \\
ESM3 CDR & ESM3, mean over CDR residues & 1536/chain & \citep{Hayes2025} & B \\
ESM C & protein LM & 1152/chain & \citep{Candido2026} & B \\
ProtT5 & protein LM & 1024/chain & \citep{Elnaggar2022} & B \\
ProtT5 CDR & ProtT5, mean over CDR residues & 1024/chain & \citep{Elnaggar2022} & B \\
SaProt & structure-aware protein LM & 1280/chain & \citep{Su2023SaProt} & B \\
AbLang-2 & paired antibody LM, both chains read together & 480 & \citep{Olsen2024} & B \\
IgT5 & paired antibody LM, both chains read together & 1024 & \citep{Kenlay2024IgT5} & B \\
AntiBERTy & antibody LM & 512/chain & \citep{Ruffolo2021AntiBERTy} & B \\
AntiBERTy CDR & AntiBERTy, mean over CDR residues & 512/chain & \citep{Ruffolo2021AntiBERTy} & B \\
Ab-RoBERTa & antibody LM & 768/chain & \citep{Huh2025} & B \\
Ab-RoBERTa CDR & Ab-RoBERTa, mean over CDR residues & 768/chain & \citep{Huh2025} & B \\
Mutation features & wild-type and mutant residue, substitution score, changes in hydrophobicity, weight and charge, position; averaged over mutations & 48 & -- & B/A \\
\end{longtable}

\begin{longtable}{p{3.2cm}p{5.6cm}p{1.6cm}p{1.6cm}p{1.4cm}}
\caption{Representations used for nucleotide sequences. One row is one representation. LM = language model. -- in Source: a simple summary built for this thesis. Role: B = base representation, A = augmentation, B/A = base on some datasets and augmentation on others.}\label{tab:reps_nucleotides}\\
\toprule
Representation & What it encodes & Features & Source & Role \\
\midrule
\endfirsthead
\toprule
Representation & What it encodes & Features & Source & Role \\
\midrule
\endhead
\bottomrule
\endfoot
4-mer counts & normalised counts of every 4-letter word & 256 & \citep{Vinga2003} & B \\
5-mer counts & normalised counts of every 5-letter word & 1024 & \citep{Vinga2003} & B \\
6-mer counts & normalised counts of every 6-letter word & 4096 & \citep{Vinga2003} & B \\
Canonical TNF & 4-mer counts merged with reverse complements & 136 & \citep{Teeling2004} & B \\
DNABERT-2 & DNA LM & 768 & \citep{Zhou2023DNABERT2} & B \\
Nucleotide Transformer v2 250M & DNA LM & 768 & \citep{DallaTorre2025} & B \\
Nucleotide Transformer v2 500M & DNA LM & 1024 & \citep{DallaTorre2025} & B \\
HyenaDNA & long-context DNA LM & 256 & \citep{Nguyen2023HyenaDNA} & B \\
GROVER (DNA) & DNA LM (not the molecular GROVER) & 768 & \citep{Sanabria2024} & B \\
Mistral-DNA & mixture-of-experts DNA LM & 768 & \citep{Mourad2024} & B \\
GC content & GC fraction, GC and AT skew, mean and spread of GC over 100-base windows & 5 & -- & A \\
Dinucleotide frequency & frequency of each 2-letter word & 16 & \citep{Vinga2003} & A \\
Codon usage & relative usage of each of the 61 sense codons, read in frame from the first base & 61 & \citep{Sharp1986} & A \\
Sequence complexity & entropy of 1-, 2- and 3-letter words, linguistic complexity, fraction of alphabet used & 5 & -- & A \\
DNA shape & minor groove width, propeller twist, helix twist and roll from a dinucleotide lookup table; mean and spread & 8 & -- & A \\
\end{longtable}

\begin{longtable}{p{3.2cm}p{5.6cm}p{1.6cm}p{1.6cm}p{1.4cm}}
\caption{Representations used for cancer patients (TCGA). One row is one representation. -- in Source: computed for this thesis from the patient's own data. Role: B = base representation, A = augmentation, B/A = base on some datasets and augmentation on others.}\label{tab:reps_tcga}\\
\toprule
Representation & What it encodes & Features & Source & Role \\
\midrule
\endfirsthead
\toprule
Representation & What it encodes & Features & Source & Role \\
\midrule
\endhead
\bottomrule
\endfoot
Geneformer (2,048 genes) & single-cell foundation model (V2-316M) on the 2,048 top-ranked genes, [CLS] vector & 1152 & \citep{Theodoris2023, Geneformer2024card} & B \\
Geneformer (4,096 genes) & the same model on the 4,096 top-ranked genes, [CLS] vector & 1152 & \citep{Theodoris2023, Geneformer2024card} & B \\
BulkRNABert & bulk RNA-seq foundation model & 256 & \citep{Gelard2024} & B \\
Top-variance methylation probes & beta values of the 5,000 most variable CpG probes & 5000 & \citep{Goldman2020} & B \\
Methylation PCA & 50 principal components of the standardised CpG probes above median variance & 50 & \citep{Goldman2020} & B \\
Copy number PCA & 50 principal components of gene-level copy number & 50 & \citep{Mermel2011} & B \\
miRNA & 200 most variable miRNAs, z-scored & 200 & \citep{Goldman2020} & B \\
Protein (RPPA) & reverse-phase protein array, z-scored & 258 & \citep{Li2013TCPA} & B \\
Chromosome-arm expression & mean expression of the genes on each chromosome arm & 42 & -- & B/A \\
TITAN & whole-slide image foundation model & 768 & \citep{Ding2024TITAN} & B \\
Prov-GigaPath & whole-slide image foundation model & 512 & \citep{Xu2024GigaPath} & B \\
CTransPath & histopathology image encoder & 512 & \citep{Wang2022CTransPath} & B \\
Phikon & histopathology image encoder & 512 & \citep{Filiot2023} & B \\
H-optimus-0 & histopathology image encoder & 512 & \citep{Saillard2024} & B \\
HRD score & homologous recombination deficiency: LOH, TAI, LST and their sum & 4 & \citep{Knijnenburg2018} & A \\
Tumour mutation burden & non-silent, missense and truncating mutation counts, log count, rate per megabase & 5 & \citep{Ellrott2018} & A \\
Tumour purity & purity, ploidy, genome doublings, subclonal genome fraction & 4 & \citep{Carter2012} & A \\
Expression statistics & summaries of the patient's expression distribution & 15 & -- & A \\
Immune signatures & 10 CIBERSORT cell fractions and Th1 and Th2 signature scores & 12 & \citep{Thorsson2018} & A \\
CIBERSORT & immune cell fractions deconvolved from bulk expression & 21 & \citep{Newman2015, Thorsson2018} & A \\
Stemness index & stemness indices from expression and from epigenetically regulated expression & 2 & \citep{Malta2018} & A \\
Hypermethylation summary & fraction of all methylation probes with beta above 0.3, and mean beta; a genome-wide stand-in for CIMP status \citep{Toyota1999} & 2 & -- & A \\
\end{longtable}

\begin{longtable}{p{2.8cm}p{2.8cm}p{5.4cm}p{3.4cm}}
\caption{Datasets and the representations each one uses. One row is one dataset. Base representations are the (up to) ten the hybrid search selects features from; augmentation features are added afterwards. Datasets marked \textsuperscript{\textdagger} are searched at the fast tier (Table~\ref{tab:search_tiers}). Sample counts, metrics and folds are given in the Datasets and splitting subsection.}\label{tab:dataset_reps}\\
\toprule
Dataset & Task & Base representations & Augmentation \\
\midrule
\endfirsthead
\toprule
Dataset & Task & Base representations & Augmentation \\
\midrule
\endhead
\bottomrule
\endfoot
\multicolumn{4}{l}{\textit{Small molecules}} \\
Caco-2 efflux \citep{openadmet} & efflux ratio (log), regression & WL graph kernel, MHG-GNN, ECFP4, PDV, GROVER, Avalon, Topological torsions, SMI-TED, Atom pairs, MoLFormer & MACCS keys, 3D pharmacophore, USRCAT, Coulomb matrix, USR \\
LogD\textsuperscript{\textdagger} \citep{openadmet} & LogD, regression & MHG-GNN, Mordred, GROVER, ECFP4, SMI-TED, MoLFormer, Avalon, Topological torsions, ChemBERT, GraphMVP & MACCS keys, 3D pharmacophore, USRCAT \\
hERG Ki \citep{Zdrazil2023} & hERG inhibition, pKi, regression & MHG-GNN, Mordred, WL graph kernel, Avalon, GROVER, PDV, ECFP4, Atom pairs, ChemBERT, SMI-TED & MACCS keys, 3D pharmacophore, Coulomb matrix, USR, USRCAT \\
MDR1-MDCK efflux \citep{fang2023} & efflux ratio (log), regression & Mordred, PDV, MHG-GNN, Avalon, GROVER, SMI-TED, Atom pairs, WL graph kernel, MACCS keys, ECFP4 & 3D pharmacophore, Coulomb matrix \\
Rat liver microsome clearance \citep{fang2023} & intrinsic clearance (log), regression & PDV, MHG-GNN, SMI-TED, Mordred, MoLFormer, GROVER, WL graph kernel, Atom pairs, Avalon, ECFP4 & 3D pharmacophore, MACCS keys, USRCAT, Coulomb matrix, USR \\
QM9\textsuperscript{\textdagger} \citep{Ramakrishnan2014} & HOMO--LUMO gap, regression & Mordred, Avalon, SMI-TED, ECFP4, SELFormer, ChemBERTa-2, Mol2vec, ChemBERT, GraphMVP, MoLFormer & Topological torsions, 3D pharmacophore, USRCAT, Coulomb matrix \\
\midrule
\multicolumn{4}{l}{\textit{Chemical reactions}} \\
Buchwald--Hartwig yield \citep{Ahneman2018} & yield (\%), regression & CGR fingerprint, GROVER (reactants), Avalon (difference), Mordred (reactants), MHG-GNN (difference), ReactionT5, Component one-hot, Mol2vec (reactants), rxnfp, MACCS (difference) & MACCS (product), MACCS (reactants), 3D pharmacophore (product), 3D pharmacophore (reactants) \\
Buchwald--Hartwig plates \citep{Ahneman2018} & yield (\%), regression & ECFP4 (reactants), GROVER (reactants), MHG-GNN (difference), CGR fingerprint, Component one-hot, Avalon (difference), ReactionT5, Mol2vec (difference), Mordred (reactants), rxnfp & MACCS (product), MACCS (reactants), 3D pharmacophore (product), 3D pharmacophore (reactants) \\
\midrule
\multicolumn{4}{l}{\textit{Antibodies}} \\
Koenig expression\textsuperscript{\textdagger} \citep{Koenig2017} & expression enrichment of variants, regression & ProtT5, ESM3 CDR, SaProt, IgT5, ESM-2 early layers, ESM C, AbLang-2, ESM-2 (joined chains), ESM-2, ESM-2 middle layers & Mutation features, ESM-2 log-likelihood ratio, Amino-acid composition \\
Nanobody thermostability \citep{ValdesTresanco2023} & melting temperature, regression & SaProt, ESM3, ESM C, ESM-2 variable positions, ESM-2 CDR, ESM-2 early layers, ESM-2 shift from parent, ProtT5 CDR, IgT5, One-hot & Mutation features \\
Humanness \citep{Prihoda2022} & human vs non-human, binary & IgT5, AntiBERTy CDR, Ab-RoBERTa, AbLang-2, BLOSUM62, ESM C, ESM-2 per CDR, ESM-2 early layers, ProtT5, ESM-2 middle layers & none \\
Trastuzumab--HER2 binding\textsuperscript{\textdagger} \citep{Chinery2024} & binder vs non-binder, binary & BLOSUM62, ESM-2 per CDR, ESM-2 middle layers, ProtT5, ESM3, ESM C, Ab-RoBERTa CDR, AntiBERTy, Mutation features, AbLang-2 & none \\
\midrule
\multicolumn{4}{l}{\textit{Nucleotide sequences}} \\
H3K4me3\textsuperscript{\textdagger} \citep{DallaTorre2025} & histone mark present vs absent, binary & GROVER (DNA), Mistral-DNA, Canonical TNF, Nucleotide Transformer v2 250M, Nucleotide Transformer v2 500M, HyenaDNA, 6-mer counts, DNABERT-2 & Codon usage, Sequence complexity, Dinucleotide frequency, DNA shape, GC content \\
TATA promoter\textsuperscript{\textdagger} \citep{DallaTorre2025} & promoter vs non-promoter, binary & Mistral-DNA, GROVER (DNA), HyenaDNA, 6-mer counts, 5-mer counts, Canonical TNF, 4-mer counts, Nucleotide Transformer v2 500M, Nucleotide Transformer v2 250M & Codon usage, Sequence complexity, Dinucleotide frequency, DNA shape, GC content \\
\midrule
\multicolumn{4}{l}{\textit{Cancer patients (TCGA)}} \\
BRCA PAM50 subtype \citep{TCGA2012BRCA} & breast cancer subtype, 4 classes & Geneformer (4,096 genes), BulkRNABert, Protein (RPPA), Copy number PCA, TITAN, H-optimus-0, Prov-GigaPath, CTransPath, Phikon, miRNA & HRD score, Tumour mutation burden, Tumour purity, Expression statistics, Chromosome-arm expression, Immune signatures \\
LUAD vs LUSC \citep{TCGA2014LUAD} & lung cancer type, binary & Geneformer (2,048 genes), Top-variance methylation probes, BulkRNABert, TITAN, Methylation PCA, Chromosome-arm expression, Copy number PCA, miRNA, CTransPath, Protein (RPPA) & Immune signatures, Expression statistics, Stemness index, CIBERSORT, Hypermethylation summary, Tumour mutation burden, HRD score, Tumour purity \\
\end{longtable}

\subsection{Models and metrics}
% ---- NOTES: Models and metrics ----
% Takeaway: Five downstream models from different families score every hybrid and every single representation, with fixed settings, so differences come from the input.
% Takeaway: Each dataset is scored in its field's own metric and read on its own; results are per-fold scores and their spread, with no significance tests.
% Paragraph 1: the downstream models (her paragraph; fix it)
% - Five: random forest, XGBoost, LightGBM, SVM and DNN. They cover bagged trees, boosted trees, a kernel method and a neural network.
% - A Gaussian process \citep{Rasmussen2006} is added on humanness, the one dataset under 1,000 molecules where it runs.
% - Change: "Eleven models are available", with CatBoost, NGBoost, the MC-Dropout BNN and the 1D CNN. None of these four is used; name only the five.
% - Change: kNN is not a downstream model. "model pool" should read "set of downstream models". "We use the full set for model-matching and noise experiments" no longer holds.
% - Why five: a hybrid that helped only one model family would be a property of that model, not a better input.
% - Tree ensembles are the baseline to beat: they still lead on medium-sized tabular data \citep{Grinsztajn2022}.
% Paragraph 2: fixed settings and size caps
% - Every model uses the same settings on every dataset, with no per-dataset tuning, so differences come from the input.
% - The Gaussian process is skipped above 1,000 molecules. The SVM is skipped above 10,000 (5,000 under uncapped augmentation).
% - Models skipped on named datasets for run time are listed in Additional files. Every model's settings are in Table~\ref{tab:app:model_settings}.
% Paragraph 3: k-nearest neighbours (one sentence)
% - Not a downstream model. It is used, at 5, 10 and 20 neighbours, in the comparison of importance-weighted against unweighted distances. State it here because Chapter 7, "Simple is often better", relies on that result.
% Paragraph 4: metrics
% - Each dataset uses its field's usual metric, so its numbers compare with published work.
% - RMSE for small molecules, QM9 and reaction yield, in the label's units. Spearman for antibody regression, whose assays rank candidates.
% - AUROC for binary tasks. Weighted F1 for PAM50's four unequal subtypes. MCC for the nucleotide sets \citep{Chicco2020}.
% - Definitions are in Chapter 2. Metrics are never averaged across datasets.
% Paragraph 5: how results are reported
% - Per-fold scores on outer test folds, each hybrid paired by fold with the best single representation, and their spread.
% - No significance tests and no confidence intervals. Fold scores share training molecules, and no unbiased estimator of the variance of $k$-fold cross-validation exists \citep{Bengio2004}.
% - Result needed: hours per dataset for the full search, and the hardware.
% ---- END NOTES ----

% TODO: review, needs manual review
For a fair comparison, the downstream model pool must cover a range of inductive biases, from methods that degrade in high dimensions (kNN, SVM, GP) to tree ensembles that tolerate irrelevant features. Eleven models are available: tree ensembles (random forest, XGBoost, LightGBM, CatBoost, NGBoost), kernel methods (SVM, Gaussian process), nearest-neighbor (kNN), and neural networks (three-layer DNN, MC-Dropout BNN for uncertainty quantification, and a multi-scale 1D CNN). Standard experiments use five of these: random forest, XGBoost, LightGBM, SVM and DNN, with a Gaussian process added on datasets under 1,000 molecules. We use the full set for model-matching and noise experiments, and skip models with cubic or higher scaling for large datasets: Gaussian processes above 1,000 samples, SVMs above 10,000. No per-dataset tuning is performed: all models use fixed hyperparameters throughout, so performance differences reflect the input, not the model configuration.




\subsection{Comparators, controls and ablations}
% ---- NOTES: Comparators, controls and ablations ----
% Takeaway: These experiments ask whether the hybrid's gain is real, or only the reward for searching many configurations.
% Takeaway: Where a search picks the best of many configurations, the comparator pays the same selection cost, by reselecting on the other folds.
% Paragraph 1: what these experiments are for
% - The hybrid picks the best of many configurations. A reader will ask whether the gain is the reward for searching hard. Each paragraph below answers one version of that question.
% Paragraph 2: paying the same selection cost on both sides
% - The best of many cross-validated scores is biased upwards \citep{Varma2006}.
% - Leave-one-fold-out reselection: for held-out fold $f$, pick the hybrid configuration and the single representation with the best mean score on the other folds. Report their difference on fold $f$, positive when the hybrid is better.
% - It is arithmetic on stored fold scores and needs no extra model fits.
% - Some experiments choose the configuration by a search over all molecules and redo only the feature selection per fold. Their scores are not fully held out; this is marked below.
% Paragraph 3: does the search overfit?
% - Nested cross-validation at every budget from 10 to 8,000 features, including budgets near the number of training molecules.
% - Selection bias = the inner search score minus the held-out score. Also report the budget the search picked against the budget with the best held-out score.
% - Figure: x = budget / training-set size (log scale), y = held-out score, one line per model. Second panel: selection bias.
% Paragraph 4: does a larger search give a better hybrid?
% - Does the hybrid's lead over each comparator grow with the number of configurations searched? Report it after reselection.
% - This needs the same dataset searched at two sizes. The full runs search each dataset at one tier only.
% Paragraph 5: fusing features against fusing predictions
% - Ensemble: one model per representation, predictions pooled (late fusion \citep{Stahlschmidt2022}). Its members, model and pooling rule are chosen on an inner split of the training molecules.
% - Controls: the best single representation; all representations joined and reduced by PCA to the hybrid's width.
% - Report cost beside score: representations computed at prediction time, and models deployed.
% - TODO(lit): read the Turkina2025 methods for whether high-level fusion pools one model per fingerprint, as the ensemble here does.
% Paragraph 6: the hybrid against standard feature reduction at the same width (not fully held out)
% - Every method gets $B$ features drawn from all base representations, fitted on the fold's training molecules, and one 100-tree random forest. Regression datasets only.
% - Methods: the hybrid, PCA, mRMR \citep{Peng2005}, recursive feature elimination \citep{Guyon2002}, and $B$ random features.
% - Question: does the gain come from drawing on several representations, or only from having fewer features?
% Paragraph 7: must the internal evaluator match the final model?
% - Nested cross-validation, one search per outer fold per evaluator setting: matched, a near family, a far family, a wider network for the DNN, and a fixed set of three tree evaluators.
% - Report each setting minus matched, paired per outer fold, at the same number of features. The answer says whether one cheap evaluator is enough.
% Paragraph 8: stability of the selected features across folds
% - Nogueira stability index \citep{Nogueira2018}: 1 means the same features every fold, about 0 is random selection. Its $d$ is the width of all base representations joined.
% - Unstable feature sets would weaken any chemical reading of the blueprints.
% Paragraph 9: gain against dataset size
% - Subsamples of 10 to 100% of the molecules, drawn as whole scaffold groups, 3 per size, each with a full search. Gain = hybrid minus best single representation, reselected.
% - Small subsamples search fewer budgets, so part of any rise in gain may come from searching more configurations.
% Paragraph 10: importance used during training (not fully held out)
% - The 100-feature hybrid, with random forest, LightGBM and XGBoost.
% - Five conditions: unweighted; features scaled by importance; molecules down-weighted when their label disagrees with their 10 nearest neighbours; both; only the features above median importance.
% - Tree ensembles ignore a positive rescaling of a feature \citep{Hastie2009}. So on these three models, scaling by importance only drops the zero-importance features.
% - Down-weighting by neighbour agreement is a noise mitigation method. It is a soft version of removing training molecules that disagree with their nearest neighbours \citep{Wilson1972}.
% Paragraph 11: a few features added to one strong representation, with decoys
% - Two decoys are planted among the candidates: pure Gaussian noise, and near-copies of 64 base features. Count the decoy features taken per fold; a working rule takes zero.
% - Also run on the multi-representation hybrid, to ask where extra features do more.
% Paragraph 12: are hybrids more robust to label noise than their base representations? (only if this stays in this chapter)
% - Chapter 3's normalised area under the retention curve (\aucnorm, Equation~\ref{eq:auc_norm}), in each dataset's own metric, at seven noise levels.
% - One downstream model, a random forest, so only the representation changes. The hybrid is searched again at every noise level on the noisy training labels; held-out labels stay clean.
% Paragraph 13: a concatenation of two representations against each alone
% - For each dataset's top three single representations, score each alone and each of the three pairs joined, with the same model and folds.
% - A concatenation that lands between its two parts supports combining and filtering at once.
% - Figure: one panel per dataset. For each pair: the score of A, of B and of the concatenation, with the band between the two single scores shaded.
% - Result needed: for each pair, the score of each representation alone and of their concatenation.
% ---- END NOTES ----
\todo[inline, color=yellow!40]{5.3.4: write the methods for each complexity experiment.}


\subsection{Transfer across datasets, antibody mutational scans and data held at separate sites}
% ---- NOTES: Transfer across datasets, antibody mutational scans and data held at separate sites ----
% Takeaway: Three setups test the search outside the single-dataset setting: reusing search settings on a second dataset, hybrids on antibody mutational scans, and feature selection with data held at separate sites.
% Takeaway: Each needs only several numeric representations of the same objects, one supervised target, and a grouped split.
% Paragraph 1: reusing search settings across datasets
% - Question: can search settings found on a large dataset be reused on a small one, instead of a new search?
% - The settings are decisions (importance method, allocation method, budget, internal evaluator), not a set of features.
% - Source: lipophilicity, 5,028 molecules. Target: hERG $\mathrm{p}K_i$, 1,412 molecules. Only representations listed for both, at the same width, are used. Result needed: the shared set the full run uses.
% - Why it matters: the search becomes a one-off cost paid on the largest dataset.
% - Closest prior work: Meta-QSAR chose the learning method and representation for a new QSAR dataset from results on over 2,700 others \citep{Olier2018}. It reuses decisions, not selected features.
% - Four comparisons inside every held-out fold of the target. The best single representation. A hybrid searched on the target (upper reference).
% - The source settings with features re-selected on the target. The source's selected features copied unchanged.
% - The last two separate whether the decisions transfer and whether the features transfer.
% Paragraph 2: hybrids on antibody mutational scans
% - Background, two sentences: a deep mutational scan measures a property for many single-point mutants of one antibody. It maps which positions tolerate change.
% - Four scans with a single wild type, including the expression scan used here \citep{Koenig2017}. Candidates are the antibody representations plus two representations of the change from the wild type.
% - Nested cross-validation; Spearman for continuous scans, AUROC for binarised binding scans. Report hybrid minus best single representation.
% Paragraph 3: data held at separate sites, background
% - Federated learning trains one model across data held by several parties without moving the data \citep{McMahan2017}. It has been run across pharmaceutical companies \citep{Heyndrickx2024}.
% - Structures and assay results are the asset, so pooling them is not possible \citep{Hanser2023,Bassani2023}.
% - The collaborator (Lhasa Limited) uses the teacher-student process of Chapter 2 \citep{Hanser2025}, unmodified here, so numbers stay comparable to theirs.
% - Each site's teacher labels shared molecules. Each vote is weighted by how close the molecule is to that site's training molecules in fingerprint distance. A student is trained on the most confident combined labels.
% - TODO(lit): check that \citet{Hanser2025} states this weighting before citing it for it; otherwise cite the collaborator's released code.
% Paragraph 4: feature selection at the sites (the addition in this thesis)
% - Each site runs the hybrid search on its own molecules and reports feature indices, importance values and one score. No molecule and no feature value leaves the site.
% - A coordinator splits the total budget $B$ across sites in proportion to their scores. A feature another site already took is skipped.
% - Only the proportional split is run, so the text cannot say what the site scores buy.
% - Why it matters: indices and importances are smaller, and harder to invert, than model weights. Splitting the budget between sites is the same question as splitting it between representations.
% - Related: clients run mutual-information feature selection locally and a server aggregates the results \citep{Cassara2022}.
% Paragraph 5: the hERG comparison
% - hERG blockade is a cardiotoxicity liability screened early in drug discovery \citep{Czodrowski2013}.
% - Three real sources act as sites: the public ChEMBL hERG $K_i$ set, and the collaborator's training set (7,772 molecules) and test set (4,383 molecules). They differ in curation rules, not only chemistry. Sites are usually simulated by clustering one dataset.
% - Blocker if $\mathrm{p}K_i \ge 6$. Five scaffold-group folds; MCC is reported.
% - Two by two: molecules pooled or kept at their sites, crossed with all features or a selected set.
% - Selected minus all, at fixed site structure, is what selection buys. Pooled minus site-held, at fixed features, is what pooling buys.
% - Caution: both all-features settings use one 2,048-bit fingerprint. Both selected settings use features from the hERG candidate representations. So selected minus all differs in two ways at once, width and source.
% - The same search on ECFP4 alone is the closest control for the source difference. The budget is swept from 50 to 1,000 features.
% - Figure: two rows, the same three sites. Top: sites, teachers, a shared set being labelled, one student. Bottom: each site's feature search, the coordinator splitting the budget, each site's model on the shared feature list. Mark what crosses a site boundary.
% Paragraph 6: other separate-site experiments, one line each
% - Sites with different amounts of label noise; representations as sites; families of representations as sites; antibody sources as sites.
% Paragraph 7: limits
% - One source is proprietary, so the comparison cannot be reproduced from public data alone (Appendix~\ref{app:availability}).
% - The published process labels a separate public set; here the training molecules play that role. State this as a compromise.
% ---- END NOTES ----
\todo[inline, color=yellow!40]{5.3.5: OPIG tasks and the federated setup.}


\section{Results: when the hybrid helps}
% ---- NOTES: Results: when the hybrid helps ----
% Takeaway: Dataset by dataset, where the hybrid scores better than the best single representation and than an ensemble, and where it does not.
% Takeaway: Where it does not help, one representation may already carry most of the signal; a tie there is a sign the search is not inflating scores.
% Paragraph 1: how to read the numbers
% - Outer held-out folds; leave-one-fold-out reselected numbers for both sides; each dataset in its own metric; per-fold scores and spread.
% Paragraph 2: hybrid against the best single representation
% - Grouped by domain. Report the count of datasets, per model, where the reselected hybrid is ahead on the fold mean, out of the datasets run.
% - Table (\label{tab:hybrid_vs_single}): one row per dataset and downstream model. Columns: domain; dataset; metric; best single representation per fold; its score, the hybrid's score and the difference per fold (positive = hybrid better); mean and spread of the difference.
% - Result needed: this table from the full runs.
% Paragraph 3: hybrid against an ensemble
% - Report hybrid minus ensemble and hybrid minus best single, with the cost beside the score.
% - If an ensemble matches the hybrid at lower cost, say so; it also belongs in Chapter 7, "Simple is often better".
% - Result needed: hybrid, ensemble and best single per dataset and model, with cost.
% Paragraph 4: gain against the number of molecules
% - A gain that rises with size can come from more data or from searching more budgets. State this next to the figure.
% - Figure: x = number of molecules (log scale), y = hybrid minus best single; one line per model, one panel per dataset; per-fold points.
% Paragraph 5: does a larger search give a larger gain? (only if paired runs at two sizes are made)
% - An early comparison is not usable. It differed in three ways at once, not one: configurations, folds and representations searched.
% Paragraph 6: where the hybrid does not help
% - Name the datasets and models where the reselected hybrid is behind or level.
% - For each, report what share of the hybrid's features comes from the best single representation. This tests one explanation, one representation carrying most of the signal. Do not state it as found.
% - Result needed: those datasets, with that share.
% ---- END NOTES ----
\todo[inline, color=yellow!40]{Results TODO(prelim): methods are final, numbers will be replaced by the server runs.}

\todo[inline, color=orange!40]{RESULTS CAVEAT: state the re-run explicitly, once, here, and cross-reference it from ch.6 and ch.7.}

\section{Results: why the hybrid helps}
% ---- NOTES: Results: why the hybrid helps ----
% Takeaway: The gain comes from which features are chosen, not only from having fewer: at the same number of features the hybrid is ahead of PCA and random features.
% Takeaway: Different downstream models want different numbers of features and different representations.
% Paragraph 1: several representations, or only fewer features?
% - If the hybrid leads PCA and random features at the same number of features, the gain is not only from having fewer.
% - Compare only where realised widths match. One sentence of caveat: the configuration was chosen using all molecules.
% - Preliminary, antibody thermostability at 50 features: hybrid 0.6131 against 0.5225 Spearman for PCA. Recursive feature elimination came within 0.003 Spearman of the hybrid at 250 features.
% - Preliminary, hERG Ki at 10 features: hybrid 0.6304 against 0.6776 RMSE (pKi) for PCA, better on all 5 folds. Other budgets go in the figure.
% - Figure: held-out score against $B$ (features, log scale), one line per method, one panel per dataset.
% Paragraph 2: must the internal evaluator be the downstream model?
% - Report each setting minus matched, paired per fold, at the same number of features. A difference at different widths is not a difference in evaluator.
% - If matched is not always best, read it as informative. The search already tries four evaluators and keeps the best.
% - Table: one row per downstream model and dataset. Columns: near family, far family, wider network (DNN only), three-evaluator set, each minus matched.
% Paragraph 3: does importance help when used during training?
% - On the three tree models, read the feature-scaling condition as "delete the zero-importance features". Report how many of the hybrid's features score zero.
% - Preliminary: LightGBM gain gave zero importance to 66 to 78 of the 250 features of a hybrid. On Caco-2 every weighting scored at or below unweighted.
% - The same scaling used as a distance is reported in Chapter 6.
% Paragraph 4: a few features added to one strong representation
% - Report the score with and without extras, against the same number of random extra features, and the decoys taken per fold.
% - If decoys make up a large share of what is taken, a score gain alone does not show the extras carry signal.
% - Preliminary, antibody thermostability, SaProt plus up to 20 mutation features, 3 folds, no admission rule. Random forest went from 0.5683 to 0.5762 Spearman.
% - In the same run, near-copy decoys outnumbered real features on every fold: 21, 11 and 21 decoys against 12, 7 and 18 real features.
% Paragraph 5: do tree ensembles and neural networks want different hybrids?
% - Number of features: the budget with the best held-out score per downstream model. The thermostability example is in "The search space and why each axis is swept".
% - Representations: rank the whole base representations by score for each model. Report the Spearman correlation of those rankings between model pairs, per dataset.
% - Features: Jaccard overlap $|A\cap B|/|A\cup B|$ between a tree-model hybrid and a DNN hybrid on the same dataset, by source representation.
% - Figure: heatmap of rank correlation between downstream models, one panel per dataset.
% - Result needed: best budget per model and dataset; rank correlation per model pair; Jaccard overlap.
% ---- END NOTES ----
\todo[inline, color=yellow!40]{Results TODO(prelim): methods are final, numbers will be replaced by the server runs.}


\section{Results: blueprints on specific datasets}
% ---- NOTES: Results: blueprints on specific datasets ----
% Takeaway: On datasets with named features, the best hybrid's selected features are read against the known chemistry of the endpoint.
% Takeaway: Which settings win differs by dataset, and the selection is only partly stable across folds.
% Paragraph 1: which datasets, and how features are read
% - Only named features can be read: ECFP bits, MACCS keys and the physicochemical descriptor vector (PDV). hERG Ki and Caco-2 carry two of them (ECFP4, PDV); lipophilicity carries one.
% - Per fold, random forest and XGBoost each get their own hybrid. SHAP values \citep{Lundberg2020} are computed for the better of the two, on held-out molecules.
% - Summed per representation, they give each representation's share of the hybrid.
% - An ECFP bit is drawn from the first molecule that sets it. One hashed bit can stand for several environments; say which was drawn.
% Paragraph 2: worked example, hERG Ki
% - The budget split across representations, and the top named features with their substructures or descriptor names.
% - Read them against hERG binding. Most blockers carry a basic nitrogen, protonated at physiological pH, and hydrophobic or aromatic groups. These meet Tyr652 and Phe656 in the channel pore.
% - Sources, neither yet in refs.bib: Aronov, Drug Discovery Today 10(2):149--155, 2005, doi 10.1016/S1359-6446(04)03278-7. Mitcheson et al., PNAS 97(22):12329--12333, 2000, doi 10.1073/pnas.210244497.
% - For each top feature: confirms known chemistry, suggests something unexpected, or cannot be read.
% - Result needed: per-representation feature counts and importance shares; the top 10 to 20 named features.
% Paragraph 3: a second dataset for contrast (Caco-2 or lipophilicity), read the same way, shorter.
% Paragraph 4: stability across folds
% - Preliminary, 3 folds, Nogueira index: 0.447 on Caco-2, 0.275 on hERG Ki, 0.074 on nanobody thermostability. All five representations searched on thermostability were learned embeddings.
% - A low index does not by itself mean overfitting. Correlated features can swap between folds and carry the same information.
% - Result needed: stability index per dataset and downstream model, 5 folds.
% Paragraph 5: does the budget split change across chemical space?
% - Importance and allocation are computed separately in 5 k-means clusters of the training molecules. Compare per-representation shares between clusters against random parts of the same sizes.
% - Result needed: per-cluster budget shares on hERG blockers (Lhasa Limited), against the random-part control.
% Paragraph 6: which settings won, per dataset
% - The one place the full-run winners are reported: allocation method, importance method, internal evaluator, representation count and budget, per dataset and downstream model.
% - Also: whether data complexity measures track the winning budget. Spearman coefficients only.
% - Figure: one small panel per setting; one bar group per dataset; bar = how many (downstream model, budget) combinations each value won.
% - Figure: for the hERG Ki hybrid, a stacked bar of features per representation, and the top named features drawn on an example molecule.
% ---- END NOTES ----
\todo[inline, color=yellow!40]{Results TODO(prelim): methods are final, numbers will be replaced by the server runs.}

\section{Results: transfer across datasets and antibody mutational scans}
% ---- NOTES: Results: transfer across datasets and antibody mutational scans ----
% Takeaway: Search settings found on lipophilicity, with features re-selected on hERG, beat hERG's best single representation on 3 of 5 downstream models; copying the features unchanged did so on 2 of 5 (preliminary).
% Takeaway: On an antibody expression scan the hybrid was ahead of the best single representation on all 5 downstream models (preliminary).
% Paragraph 1: reusing search settings (preliminary)
% - RMSE in pKi on hERG, 5 scaffold folds. Difference = the best single representation's RMSE minus the comparison's, so positive means lower error.
% - Re-selecting features with the source settings beat the best single representation on 3 of 5 downstream models. Copying the features did so on 2 of 5.
% - Re-selecting came within 0.01 RMSE of the hybrid searched on hERG itself on 4 of 5 models.
% - Give the spread with each number. On the DNN the fold standard deviation of the copied-features difference is 0.1823 RMSE, larger than any mean difference.
% - Table: one row per downstream model. Columns: best single; hybrid searched on hERG; source settings re-selected; source features copied; the last two minus best single, with fold standard deviation.
% Paragraph 2: the antibody mutational scan (preliminary)
% - Expression scan \citep{Koenig2017}, 4,275 variants, 3 unequal outer folds, reduced search (200 configurations, budgets of 100 and 500 features).
% - The best single representation was ProtT5 on every model and fold. The hybrid was above it on all 5 downstream models, and in 14 of 15 model-fold pairs. Example: LightGBM, +0.0431 Spearman.
% - Most of the budget went to ProtT5, IgT5 and AbLang-2; the shares go in the figure or table.
% - The two wild-type-difference representations took at most 6% of features. Report this as a hybrid-against-best-single result, not a result about them.
% - Result needed: the other three scans, and the full search.
% ---- END NOTES ----
\todo[inline, color=yellow!40]{Results TODO(prelim): methods are final, numbers will be replaced by the server runs.}

\section{Results: data held at separate sites}
% ---- NOTES: Results: data held at separate sites ----
% Takeaway: Does feature selection run separately at each site match feature selection on pooled molecules, when only feature indices, importances, one score per site and predictions cross between sites?
% Paragraph 1: how to read it
% - MCC on five scaffold folds, per-fold scores and spread, at every budget from 50 to 1,000 features.
% Paragraph 2: the hERG two by two
% - Report the selection difference separately for pooled and site-held molecules, and at which budgets it keeps its sign.
% - Selected minus all differs in width and source at once. With pooled molecules, pooled/selected minus the ECFP4-only search, at equal width, isolates what several representations buy.
% - Table (\label{tab:sites_two_by_two}): one row per budget (features). Columns: the four settings and the two selection differences; mean MCC and spread over folds.
% - Figure: the two by two, with the MCC of each setting at one budget in its box.
% Paragraph 3: context comparisons
% - The best single source, picked on its own training molecules. Its gap to the site-held selected setting is what combining sites is worth.
% - The collaborator's two teacher-weighting schemes, and one teacher per representation.
% - Each site's reported score and its share of the features, per budget.
% Paragraph 4: the other separate-site experiments, one line each, nulls included, in their own metric
% - The antibody sites' labels are in different units (kcal/mol, a normalised score, degrees Celsius), so the teacher-ensemble comparisons were not run there. One sentence.
% - Sites differ in curation as well as chemistry, so a gap between sites is not only a chemistry effect.
% - Result needed: MCC for the four settings and each context comparison, per budget and fold.
% ---- END NOTES ----
\todo[inline, color=yellow!40]{Results TODO(prelim): methods are final, numbers will be replaced by the server runs.}


\section{Discussion}
% ---- NOTES: Discussion ----
% Takeaway: Feature selection can overfit, and the search picks the best of many configurations. Nested cross-validation charges the hybrid for that, so the held-out folds, not the inner search, decide whether it helps.
% Takeaway: Where the hybrid does not help, a tie with the best single representation is a sign the search is not inflating scores.
% Paragraph 1: selection bias (her paragraph below; fix it)
% - Keep the argument. Selection outside the folds inflates the score \citep{Ambroise2002}. Selection is a form of model selection \citep{Cawley2010}. Searching adds degrees of freedom \citep{Tibshirani2015}.
% - Change: the Ambroise and McLachlan numbers. With random labels, selection outside the folds gave cross-validated error near zero at 128 genes; inside each fold, 0.40 to 0.45. "11 out of 40" does not match.
% - Change: Vabalas et al. say selection on pooled training and test data adds "considerably more" bias than parameter tuning \citep{Vabalas2019}.
% - Change: "KIRBy" should read "the hybrid search". It scores only the configurations in the grid, so "many possible feature subsets" should read "many configurations".
% - Change: the \citet{Ng2004} sentence misstates it. With $L_1$ regularisation, the training examples needed grow with the logarithm of the number of irrelevant features. For SVMs and neural networks they grow at least linearly in the worst case.
% - Change: \citet{Kumar2022} found fine-tuning on average 7\% less accurate out of distribution than a linear probe; check "up to 10\%".
% Paragraph 2: the case that the search overfits
% - Importance is computed on the same training labels the model is fitted to, so label noise can steer the pick.
% - The best of many noisy scores is biased upwards \citep{Cawley2010}. Small datasets are the most exposed.
% - Sort \& Slice, which uses no labels, beat two supervised selection methods \citep{sns}. Its authors call label-based overfitting "conceivable": a hedge, not a finding.
% Paragraph 3: the case that it does not, or that the cost is outweighed
% - Nested cross-validation scores the search's choice on an outer fold it never saw. It also charges the single representation for its own selection.
% - Two further checks: stability across folds, and settings reused on a second dataset keeping their score.
% - A hybrid does not depend on one representation's blind spots, such as hash collisions or molecules it cannot tell apart (Chapter 4).
% - Whether strong, varied features outweigh those picked on noise is answered by the held-out folds. Chapter 6 tests it with noisy labels.
% - Point to the selection-bias figure in "Comparators, controls and ablations"; no second figure here.
% Paragraph 4: the budget and the information bottleneck
% - The information bottleneck keeps a short code of the input that is informative about the target \citep{Tishby2000}.
% - The budget is an analogy only: it keeps original features rather than learning a mapping. The evidence is the matched-width comparison.
% Paragraph 5: whether the internal evaluator must match the final model
% - If the settings score alike, the model scoring representations during the search need not be the model finally trained. One cheap evaluator may then be enough. That is a finding, not a failed test.
% - Say how this fits with different downstream models preferring different settings.
% Paragraph 6: where the hybrid does not help
% - A tie with the best single representation is consistent with the search not inflating scores. A large gain there would be the warning sign.
% - Practical signal: score each representation alone first. If one stands well clear, expect a small gain.
% - Prior work reports ties without saying when to expect them. ECFP4 joined with a bioactivity fingerprint beat both on 18 of 24 assays and was "comparable" on the rest \citep{Laufktter2019}.
% - TODO(lit): a source that studies when combining gives no gain is still needed.
% ---- END NOTES ----

% TODO: review, needs manual review
Ambroise and McLachlan \citep{Ambroise2002} showed the extreme case of feature selection overfitting: an SVM achieved near-zero cross-validated error on random labels when gene selection was performed on the full dataset before cross-validation, but when selection was nested inside each fold the median error was 11 out of 40, close to the 50\% expected by chance. Vabalas et al.\ \citep{Vabalas2019} showed that the bias introduced by pre-selection can exceed the bias from hyperparameter tuning, and Cawley and Talbot \citep{Cawley2010} argued that feature selection is itself a form of model selection whose complexity cost must be accounted for. KIRBy follows the standard remedy: we compute all importance scores within each cross-validation fold, on training data only. But the problem of ``search degrees of freedom'' \citep{Tibshirani2015} remains. Searching over many possible feature subsets adds complexity even when each individual selection is unbiased. The budget acts as a hard capacity constraint: selecting 100 features from a pool of 5,000 limits the space of possible selections far more tightly than $L_1$ regularization over the same pool, because the feasible set grows polynomially in the budget while the continuous parameter space of a regularized model is unbounded \citep{Ng2004}. Important here is a distinction: the features being selected are frozen embeddings, not parameters being fitted. Choosing which pre-existing knowledge to trust is, we think, a less expressive operation than learning new features, and Kumar et al.\ \citep{Kumar2022} showed that frozen embeddings fed through a linear probe can outperform full fine-tuning by up to 10\% in out-of-distribution accuracy precisely because fine-tuning distorts pretrained structure. And stability can be measured directly: the Nogueira index \citep{Nogueira2018} quantifies the consistency of selected feature sets between folds. If selection were overfit, feature sets would differ wildly from fold to fold.

\section{Limitations}
% ---- NOTES: Limitations ----
% Takeaway: The search is costly, must be repeated when data, representations or models change, and its gain is unknown until it has run.
% Takeaway: On small datasets, frozen inputs and noisy endpoints the hybrid has a ceiling, and simpler selection may be enough.
% Paragraph 1: her paragraph below; fix it
% - Change: "KIRBy's limitations" should read "the hybrid method's limitations". "best single source" should read "best single representation".
% - Change: match the two numbers. A 200-molecule dataset leaves a held-out part of about 32 molecules per fold. About 20 corresponds to about 125 molecules.
% - Keep the small-dataset claim only if the full runs show it.
% - TODO(lit): source for SchNet outperforming frozen embeddings on 3D quantum properties. The SchNet paper does not compare with frozen embeddings.
% - Change: one-hot and $k$-mer counts are base representations here. The "disperses ... information" sentence needs a source or a result, or should be cut.
% - Change: "20 to 50 foundation models" describes the codebase, not one search. The dataset lists draw on 33 pretrained models; each search uses a handful.
% - Change: \citet{Kramer2012} gives a "standard deviation of 0.54 pK$_i$ units", and a maximum squared Pearson correlation of 0.81 "on large datasets". Use \citet, not "Kramer et al.\ \citep{}".
% - Change: "The budget constraint plausibly guards against overshoot" is measured by nothing. Keep it as a guess in words, or cut it.
% Paragraph 2: compute and run time
% - The search is a full grid, run again inside every outer fold. Importance and standalone scores are reused within a fold; the global pool adds one fit per configuration.
% - It costs orders of magnitude more than fitting one model on one representation, and needs a cluster.
% - Result needed: hours per dataset, and which settings take most of the time. The per-dataset hours belong in an Additional file table.
% Paragraph 3: the search must be repeated when inputs change
% - New molecules, new folds, a new representation or a new downstream model each mean a new search. A new representation reruns everything, because shares are relative.
% - Datasets grow, so the cost recurs. Reusing a blueprint on a related dataset is a partial workaround.
% Paragraph 4: high upfront cost, gain unknown until the search has run
% - No cheap test predicts whether the hybrid will win. A standout representation is a hint, not a guarantee.
% - Figure: per dataset, hybrid minus best single against search cost in hours; per-fold points.
% Paragraph 5: reproducibility
% - A hybrid depends on the representations and their checkpoint versions, the seed, every search setting and the library versions.
% - A blueprint stores the settings, not the features, so a rebuild on other molecules can pick others.
% Paragraph 6: simpler selection may be enough
% - The best single representation, concatenation with a filter or $L_1$ penalty, or a tree ensemble's own selection each cost about one fit. An honest estimate for any of them also needs nested cross-validation \citep{Varma2006}.
% - With limited compute or a routine task, start with the best single representation. The grid is finite; the search finds the best within it only.
% ---- END NOTES ----

% TODO: review, needs manual review
We want to be clear about KIRBy's limitations. For small datasets (below about 200 samples), the hybrid search cannot reliably outperform the best single source. We suspect this is because the internal validation set within each fold shrinks to roughly 20 samples, and at that size the evaluator cannot distinguish genuine improvements from noise. Tuning cannot fix this. And at the other extreme, frozen embeddings have a ceiling. No signal can be added beyond what the inputs already contain, and on tasks where the true signal requires end-to-end learning from raw inputs (certain 3D quantum-mechanical properties, for example, where SchNet outperforms any frozen embedding by a wide margin), this ceiling binds. Although the methodology naturally accommodates embeddings from sequence-trained models (ESM-2, ProtBERT, etc.), encoding raw sequences directly as one-hot or $k$-mer features disperses the positional and compositional information that makes sequences useful in the first place. Integrating 20 to 50 foundation models into a single pipeline, each with its own dependencies, GPU requirements, and version constraints, is genuinely difficult engineering. Embeddings need only be computed once per dataset and everything is cached. Yet the initial integration burden is real and scales with the number of models in the pool. Finally, on endpoints with high experimental noise, even a well-constructed hybrid may not improve on simpler baselines. Kramer et al.\ \citep{Kramer2012} documented that experimental uncertainty in heterogeneous public $K_i$ data is approximately 0.54 log units. This sets a natural accuracy ceiling of $R^2 \approx 0.81$ regardless of model sophistication; the budget constraint plausibly guards against overshoot, but it cannot push past the noise floor.


\section{Conclusions}
% ---- NOTES: Conclusions ----
% Takeaway: One paragraph per question this chapter answers, then a short cost line, then the bridge to Chapter 6.
% - Quote each question's wording from Section~\ref{sec:intro:questions} at the start of its paragraph; answer in one or two sentences.
% Paragraph 1: does the hybrid beat its inputs?
% - Against the best single representation at full width, and an ensemble.
% - Where it only matches the best single representation, the search still found that representation without it being named in advance.
% - Result needed: hybrid against each comparator per dataset; the datasets where it did not help, and which representation dominated.
% Paragraph 2: selection, not only reduction
% - At matched width, against PCA, mRMR, recursive feature elimination and random features. Ahead there means the gain comes from which features are chosen.
% - Whether tree ensembles and neural networks prefer different widths and representations.
% Paragraph 3: one search setting is not enough
% - Each setting picked different values on different datasets, so a fixed recipe would lose on some. This is why the search exists.
% - Result needed: per setting, the number of datasets on which each value won.
% Paragraph 4: is the search overfitting?
% - One sentence on nested cross-validation. Result needed: selection bias per budget, and the stability index per dataset.
% Paragraph 5: what blueprints tell a chemist
% - Result needed: one or two datasets where selected named features match known chemistry.
% Paragraph 6: beyond small molecules, transfer, and data held at separate sites
% - What a dataset needs: enough samples, and representations whose features carry importance. Label noise sets a separate ceiling.
% - One sentence each: antibodies, reactions, nucleotides, cancer genomics; transfer; mutational scans; separate sites.
% Paragraph 7: label-noise robustness and importance weighting (the parts this chapter answers)
% - Result needed: \aucnorm for the hybrid against single representations; each weighted training condition minus unweighted.
% - Distances and learning with noisy labels are answered in Chapter 6.
% Paragraph 8: cost, one or two sentences
% - The search is expensive and must be rerun for new folds, representations or models. Simpler selection may be enough.
% Paragraph 9: bridge to Chapter 6
% - This chapter used the hybrid only as model input.
% - Chapter 6 uses it for similarity, in two ways. It asks whether either improves noise detection and learning with noisy labels, which Chapter 4 showed failing under standard similarity.
% - Figure (optional): three boxes. Chapter 4 (standard representations and similarity fail), Chapter 5 (a feature set chosen for the task), Chapter 6 (that feature set used for similarity and label noise).
% ---- END NOTES ----
\chapter{Hybrid Representations for Similarity and Label Noise}\label{chap:similarity}
% ---- TAKEAWAYS ----
% - The hybrid began as a fix for similarity; this chapter tests whether features chosen for the task give better neighbours than the single representations they come from.
% - A similarity hybrid uses the same search but ranks candidates by predicting each molecule's activity from its five nearest neighbours.
% - Finding noisy labels, estimating how large the noise is, and improving the held-out score are three different tasks. A method can do the first and fail the third.
% - The type of noise decides what can be fixed. Censoring has a known direction and location and is largely correctable; Gaussian noise, the standard assumption, is the hardest case.
% - Under Gaussian noise the input moved the score more than the choice of LNL method did (preliminary).
% - Better similarity cannot remove the noise on the neighbours' own labels. No one method fits every noise type and dataset, so model the noise first.
% ---- END TAKEAWAYS ----

\section{Introduction}
% ---- NOTES: Introduction ----
% Takeaway: Chapter 4 showed where similarity breaks; Chapter 5 built a representation chosen for the task. This chapter asks whether that representation repairs similarity, noise detection and learning with noisy labels.
% Takeaway: Short, about one page: three paragraphs, the research questions as a list, and one status sentence.
% Paragraph 1: what the previous two chapters leave open
% - Chapter \ref{chap:flaws} showed where the similarity assumption fails on standard representations: ECFP4 bits compared by Tanimoto mix a few task-relevant features with many that are constant or irrelevant on the dataset.
% - Methods that trust a molecule's neighbours inherit that error, including finding noisy labels by their neighbours.
% - Chapter \ref{chap:method} built the adaptive hybrid representation: a budget of features picked from several base representations for one task and one model.
% Paragraph 2: the original motivation (one or two sentences)
% - The hybrid started as a fix for similarity: select the task-relevant features first, then compute distances only on those. It then turned out to work as an ordinary model input; this chapter returns to the original motivation.
% Paragraph 3: the methods under test (two or three sentences)
% - Nearest-neighbour editing \citep{Wilson1972}, filtering by classifiers trained on other folds \citep{Brodley1999}, and confident learning \citep{Northcutt2021}. Each depends on the feature space it works in.
% List: research questions, one per results section
% - Distances and nearest-neighbour prediction: is predicting each molecule's activity from its five nearest neighbours more accurate in a hybrid than in the single representations it is built from?
% - Noise detection and estimation: are noisy labels easier to find, or their size easier to estimate, in a hybrid?
% - Learning with noisy labels: does a hybrid change how well noise mitigation and noise-robust training work?
% Paragraph 4: status (one sentence)
% - Planned as the second hybrid paper. Methods are final; results are preliminary.
% - Result needed: one-sentence preview of the pattern found for each research question.
% - End with what is new in this chapter as a short numbered list, one line each.
% ---- END NOTES ----


\section{Methods}
% ---- NOTES: Methods ----
% Takeaway: Every hybrid comes from the Chapter 5 search. Two things are new: a second score to rank candidate hybrids on, and what is done to the training labels.
% Takeaway: The hybrid is always rebuilt inside each fold from the training labels the experiment gives it, so feature choice sees the noise.
% Paragraph 1: what the experiments share
% - Datasets, scaffold-group-aware folds and metrics are as in Chapter \ref{chap:method}. Noise touches training labels only.
% - Comparison: each downstream model's own top two single representations, ranked by that model's own score.
% Paragraph 2: which labels the hybrid saw
% - The hybrid's features are rebuilt within each fold from that fold's training labels: noisy, clipped or uncleaned.
% - The search settings (representations, budget, allocation method) come from different labels per experiment:
%   - Survey of twelve LNL methods: searched again in each fold on the noisy training labels, with the same neural network the methods use. Budgets 50, 100 and 250 features.
%   - Labels clipped at a detection limit: settings from the search on the unclipped labels, which a real user would not have. State it as a limit.
%   - Cleaning real labels: settings from the search on the same real labels.
% Paragraph 3: the three tasks kept apart (two sentences)
% - Finding noisy labels, estimating how large the noise is, and improving the held-out score are separate tasks. A method can do the first well and fail the third.
% Figure: one-panel flow. Left: dataset with injected noise, or none. Middle: the hybrid built per fold on the training labels, beside the top two single representations. Right: similarity score, and the held-out score after each LNL method.
% ---- END NOTES ----


\subsection{Similarity in a hybrid}
% ---- NOTES: Similarity in a hybrid ----
% Takeaway: Two ways to use a hybrid for similarity: compute distances in an ordinary hybrid, or search a similarity hybrid whose candidates are ranked by nearest-neighbour prediction instead of a downstream model.
% Takeaway: In a hybrid, width is weight: a representation's pull on the distance is roughly its share of the budget.
% Paragraph 1: two ways to use a hybrid for similarity
% - Ordinary hybrid: the Chapter \ref{chap:method} search picks features that help a downstream model. Distances are then computed in that hybrid.
% - Similarity hybrid: the same search, but each candidate is ranked by how well each molecule's activity is predicted from its five nearest neighbours in the hybrid. No downstream model is fitted.
% - Everything else is shared: importance methods, allocation methods, budgets and folds. The neighbour score also sets each representation's share of the budget.
% - Every representation in the dataset's set is a candidate; the number of representations is not searched.
% Paragraph 2: the distance inside a hybrid
% - Manhattan distance (Chapter \ref{chap:flaws}) on features z-scored with the training molecules' mean and standard deviation.
% - Width is weight: every z-scored feature adds one term of similar size, so a representation's pull on the distance is roughly its number of features divided by the budget. The allocation method therefore also weights the representations inside the distance.
% - Result needed: confirm Manhattan distance and one-over-distance neighbour weighting once the design experiment on distance, neighbour weighting and number of neighbours has run.
% Paragraph 3: the neighbour score
% - Computed on the held-out molecules of each fold. Each molecule's 5 nearest neighbours come from other scaffold groups, because a same-scaffold neighbour is an easy match and inflates the score.
% - Prediction, weighting each neighbour by one over its distance: $\hat{y}_i = \frac{\sum_{j\in N_5(i)} y_j / d_{ij}}{\sum_{j\in N_5(i)} 1/d_{ij}}$.
% - Regression score: Spearman $\rho(\hat{\mathbf{y}}, \mathbf{y})$. Binary classification: $2(\mathrm{AUROC} - 0.5)$ on the distance-weighted fraction of neighbours in class 1, so chance is 0 under both.
% - The score measures how far the similar property principle (Chapter~\ref{chap:background}) holds in the hybrid.
% Paragraph 4: where the score is undefined (one or two sentences)
% - The similarity hybrid is not searched on three datasets where the score is undefined. The datasets and reasons go in an Additional file.
% Figure: two pipelines sharing a first half (base representations -> importance -> allocation -> hybrid). Top: hybrid -> downstream model -> dataset metric. Bottom: hybrid -> 5 nearest neighbours from other scaffold groups -> distance-weighted mean activity -> Spearman or rescaled AUROC. Caption: same search, different score.
% ---- END NOTES ----
\todo[inline, color=yellow!40]{6.2.1: objective=similarity. [DIAGRAM D6.1]}

\paragraph{Feature importance weighting.}
% ---- NOTES: Feature importance weighting ----
% Takeaway: Scaling each feature by its importance before the distance lets features that matter for the task count more, and drops zero-importance features.
% Paragraph 1: what the weighted distance is (her paragraph below)
% - Change: her first sentence names noise detection as the only purpose. Add the second use: nearest-neighbour prediction, weighted against unweighted (results in "Results: similarity").
% - Importance is computed on each fold's training molecules only. Weights are applied after z-scoring; applied before, z-scoring divides them back out.
% - $w$ under Manhattan and $\sqrt{w}$ under Euclidean both multiply each feature's term in the distance by $w$ once.
% - Average 1 keeps weighted and unweighted distances on the same scale; neighbour rankings do not change under any rescaling of $w$.
% ---- END NOTES ----

Beyond simply using the feature importance scores for allocation, we investigate whether feature weighting by importance during distance computation improves noise detection. For a feature importance vector $w$, weighted distances are computed by scaling the standardised features: $x' = x \odot w$ under Manhattan distance ($x \odot \sqrt{w}$ under Euclidean), with negative importances set to 0 and $w$ rescaled to average 1. 


\subsection{Noise conditions}
% ---- NOTES: Noise conditions ----
% Takeaway: Three conditions drive the learning-with-noisy-labels experiments: censoring, Gaussian noise, and real labels with nothing injected. Outlier and grouped noise are used only to grade detection and estimation.
% Takeaway: Each condition stands for a real source of error, and noise is added to training labels only.
% Paragraph 1: which conditions this chapter uses
% - Defined in Section~\ref{sec:noise_injection} (Figure~\ref{fig:noise_conditions}). Noise is sized from each fold's clean training labels and added to training labels only.
% Paragraph 2: censoring at a detection limit
% - The level is the fraction of training labels clipped: 0 to 0.40. The clip has no random draw, so one run per level.
% - The error is one-sided, and its direction and location (the limit) are known.
% - Stands for assay range limits, e.g. "> 10 µM" recorded as 10 µM \citep{Svensson2025}. Regression only.
% Paragraph 3: Gaussian noise
% - Levels are multiples of the fold's clean training-label standard deviation, 0 to 1.5, three draws per level. On classification datasets, the fraction of labels flipped, 0 to 0.40.
% - Stands for the random error of repeating a measurement. It is also the standard assumption in noise-robustness studies.
% Paragraph 4: outlier, grouped-wider and grouped-shifted noise (detection and estimation only)
% - Outlier noise stands for annotation and transcription errors \citep{Kalliokoski2013}; here 1\%, 5\% and 10\% of training molecules.
% - Grouped-wider and grouped-shifted stand for differences between laboratories or sources \citep{Bentz2013}. Both need a scaffold per molecule, so run only on datasets with SMILES.
% Paragraph 5: real labels, nothing injected
% - No list of wrong labels exists, so success is only a better held-out score after cleaning. That shows cleaning helped the model, not that it found errors.
% - QM9 is the control: its labels are computed, with no experimental error \citep{Ramakrishnan2014}.
% Table 6.1: one row per condition. Columns: what it stands for; levels; which experiments use it; datasets.
% - Result needed: the datasets each condition covers in the full runs.
% ---- END NOTES ----
\todo[inline, color=yellow!40]{6.2.2: Table 6.1, one row per noise type.}


\subsection{Noise mitigation methods}
% ---- NOTES: Noise mitigation methods ----
% Takeaway: Three experiments, one per question: correcting labels clipped at a detection limit, a survey of twelve LNL methods under Gaussian noise, and cleaning real labels.
% Takeaway: Noise-robust training never has to say which direction a label is wrong in. Noise mitigation does, per molecule.
% Paragraph 1: the two kinds of LNL method
% - LNL, noise mitigation and noise-robust training are defined in Section~\ref{sec:lnl_background}. The first and third experiments are noise mitigation; the survey covers both.
% - Change: this heading says noise mitigation, but the survey also covers noise-robust training. A heading naming learning with noisy labels fits all three.
% Paragraph 2: correcting labels clipped at a detection limit
% - Training labels above a limit $c$ are recorded as $c$. The treatments read the limit off the training labels as the largest recorded value; they are never given it.
% - Reference prediction $p_i$: out-of-fold boosted-tree prediction, scaffold groups kept together. Residual spread $s$: SD of $y_i - p_i$ over molecules not at the limit.
% - Basis, the mean of a normal above a limit: $\mathbb{E}[Y \mid Y > c] = \mu + \sigma\frac{\phi(a)}{1-\Phi(a)}$, $a = (c-\mu)/\sigma$ \citep{Tobin1958}.
% - Treatments: per-molecule lift, $c + s\frac{\phi(a_i)}{1-\Phi(a_i)}$ with $a_i = (c-p_i)/s$; a shuffled control dealing those lifts to other clipped molecules; a constant lift, the mean of the per-molecule lifts; dropping the clipped molecules. Two further treatments belong in an Additional file.
% - Model comparators, not proposals (the thesis changes the input, not the model), each paired with the model it replaces: gradient boosting with a censored likelihood, a random forest with absolute-error splits, gradient boosting with the Huber loss.
% - Scored by recovery rate, defined below.
% Paragraph 3: a survey of twelve LNL methods
% - Question: does changing the input from a single representation to the hybrid change which methods work, and how well?
% - Base for every method: one plain fully connected network. Inputs: the top two single representations, the hybrid, and the importance-weighted hybrid.
% - Most of these methods were built for image classification with flipped labels. Here they are adapted to continuous labels.
% - Scored against the plain network on the same input and noise level.
% - Table: the twelve methods in three groups, one plain line each.
%   - Noise mitigation, remove then refit: small-loss selection \citep{Han2018}, confident learning \citep{Northcutt2021}, removal by distance to neighbours, staged removal \citep{zhao2025}.
%   - Noise-robust training that changes the training loop: per-epoch down-weighting (no source yet), curriculum by confidence \citep{Bengio2009}, DivideMix \citep{li2020dividemix}, Meta-Weight-Net \citep{shu2019metaweightnet}, MentorNet \citep{jiang2018mentornet}, Contrast to Divide \citep{zheltonozhskii2021contrast}.
%   - Noise-robust training that models the noise (continuous labels only): heteroscedastic GP \citep{liu2020}, deep evidential regression \citep{amini2020deep}.
% Paragraph 4: cleaning real labels, nothing injected
% - Per fold: fit on the uncleaned labels and score; clean the training labels; refit the same model and score again.
% - Two suspicion scores: neighbour residual (distance from the mean label of 10 nearest neighbours in other scaffold groups) and cross-validation residual (distance from an out-of-fold prediction).
% - Seven cleanups, one line each: remove flagged molecules; down-weight by suspicion; move every label toward its neighbours' mean in proportion to suspicion; move only flagged labels there; move labels toward an out-of-fold prediction; toward a Gaussian-process mean; clip labels to the range carried by similar training molecules.
% - Each is run at a few strengths (share flagged, or how far labels move).
% - Shuffled-score control: the same cleanups with suspicion scores dealt at random. It separates "which molecules" from "how many".
% - Scored as folds where the cleaned model beats the uncleaned one, out of folds run.
% Paragraph 5: what the four subsubsections below are for (one sentence)
% - The three experiments use only the neighbour residual, the cross-validation residual, proportional smoothing, inverse weighting and the recovery rate. The rest describes the exploratory study (Section~\ref{sec:lnl_motivation}).
% - The exploratory study's full lists of removal methods, noise scores and smoothing and weighting rules can move to an Additional file.
% Figure: true label against recorded label for the clipped-label experiment, with the pile-up at $c$ and arrows for the lift treatments.
% ---- END NOTES ----
\todo[inline, color=yellow!40]{6.2.3: methods for censoring correction, the LNL survey and cleaning real labels.}

In addition to evaluating hybrid representation quality for distance-based reasoning, our investigation evolved to explore alternative noise handling strategies. Initial experiments with sample removal methods revealed limitations when applied to Gaussian noise affecting all samples. This motivated a pivot toward sample weighting and label smoothing approaches, which do not require accurate binary classification of noisy versus clean samples but instead leverage continuous noise magnitude estimates.

\subsubsection{Sample removal methods}
% ---- NOTES: Sample removal methods ----
% Takeaway: Removal was the first approach tried in the exploratory study. Under Gaussian noise every label is wrong by some amount, so removal throws away labels that are only partly wrong. This led to weighting and smoothing.
% Paragraph 1: the exploratory study (two sentences)
% - QM9 HOMO-LUMO gap, 10,000 molecules, one scaffold split; Gaussian noise of SD 0.3 and 0.6 eV on training labels; ECFP4, PDV and MHG-GNN; random forest after removal.
% - Cite only the single-representation rows. Some hybrid rows used an allocation method no longer in the thesis.
% Paragraph 2: k-nearest-neighbour filtering
% - Change: the first sentence re-states the similar property principle with a different citation from Chapter~\ref{chap:background}. Cite \citet{Johnson1990} as that chapter does.
% - Change: the sentence beginning "Presumably" can state the idea directly. Editing a training set by its nearest neighbours goes back to \citep{Wilson1972}.
% Paragraph 3: Local Outlier Factor
% - Change: there is no "LOF greater than 2" rule. LOF uses 20 neighbours and flags the 10\% of training molecules with the highest score.
% - LOF sees only the features, never the labels, so it flags molecules in sparse regions, not molecules with unusual labels.
% Paragraph 4: the other removal methods
% - Change: rename "AC detection" to "activity-cliff detection".
% - Co-teaching \citep{Han2018} and DivideMix \citep{li2020} were cut down to their selection step, run with random forests, and their rejected molecules deleted. Say so, so a reader does not expect the original methods.
% - Cross-validation filtering follows \citep{Brodley1999}.
% Paragraph 5: why removal does not suit Gaussian noise (one or two sentences)
% - Under Gaussian noise every training label is wrong by some amount, so detection precision is 1 by construction. The useful question becomes how wrong each label is.
% ---- END NOTES ----

Distance-based approaches rely on the assumption that similar molecules have similar properties \citep{Cherkasov2014}. Presumably, noisy data would break this pattern, making it a method for detecting anomalous data. k-Nearest neighbor filtering examines the variance of a sample's neighborhood. For sample $x_i$ with $k$ nearest neighbors $\mathcal{N}_k(x_i)$, we compute the standard deviation of neighbor labels:
$$
\sigma_i = \text{std}(\{y_j : x_j \in \mathcal{N}_k(x_i)\})
$$
Samples with $\sigma_i$ exceeding the 90th percentile are marked as potentially noisy. We test $k \in \{5, 10, 20\}$.

Local Outlier Factor (LOF) identifies anomalous samples by detecting local density estimates between a sample and its neighbors \citep{Breunig2000}. These density estimates are based on the distance to its $k$ nearest neighbors. If a point's density differs from that of its neighbors such that the LOF is substantially greater than 2, it is flagged as an outlier.

We also tested the following noise removal methods:  distance-weighted prediction filtering, neighborhood consensus filtering, AC detection, Mahalanobis outlier detection, cross-validation disagreement, bootstrap ensemble filtering, co-teaching \citep{Han2018}, and DivideMix \citep{li2020}.

\subsubsection{Noise magnitude estimation}
% ---- NOTES: Noise magnitude estimation ----
% Takeaway: A score of how wrong each label is lets later steps down-weight or smooth in proportion. Removal is all or nothing.
% Paragraph 1: why score how noisy each label is
% - \citet{Martin2021} adapt classification noise filters to continuous labels and find they help regression models. Cite it for the general idea only.
% - Change: the first sentence calls estimation "the core of the noise mitigation research". Estimating noise is not itself noise mitigation; it is the step noise mitigation builds on.
% Paragraph 2: scores from each molecule's nearest neighbours
% - Neighbour residual is the continuous-label form of editing a training set by neighbour disagreement \citep{Wilson1972}.
% - Change: the prose writes the local variance ratio's denominator as std + epsilon. It is a floor: the neighbours' SD, never below 1\% of the dataset's label SD.
% - Change: spell out "activity cliff" instead of "AC" \citep{Stumpfe2012, Dablander2023}.
% Paragraph 3: scores from models that never saw the molecule
% - Cross-validation residual: cite \citep{Brodley1999}; scaffold groups kept together.
% Paragraph 4: scores from models that predict their own uncertainty
% - All are fitted out of fold, so no molecule is scored by a model that trained on it.
% - Change, Gaussian process: one learned noise level is shared by every molecule, so it cannot rank molecules. The score is the out-of-fold predictive SD plus the out-of-fold absolute residual.
% - Change, random forest uncertainty: the prose calls it "Bayesian random forest" and says out-of-bag. It is the spread across trees plus half the out-of-fold absolute residual \citep{taddy2015}.
% - Change: the last sentence says NGBoost was "selected". It is one score among the others. The correlation it cites from Section~\ref{sec:uq} is pooled across noise levels; within one noise level it is near zero.
% - Five more scores computed but not in the prose belong in an Additional file.
% Table (optional): one row per noise score; columns family, what is computed in one line, citation. It can replace the three itemised lists.
% ---- END NOTES ----

We now examine the core of the noise mitigation research: estimating the magnitude of noise. This enables proportional handling rather than complete removal. We used the following neighbor-based estimators:

\begin{itemize}
    \item Neighbor residual: Estimates noise as deviation from k-NN mean, normalized by global label standard deviation:
    $$\text{score}_i = \frac{|y_i - \bar{y}_{\mathcal{N}_k(x_i)}|}{\sigma_y}$$
    
    \item Local variance ratio: Measures deviation relative to local label variance, identifying samples that deviate from neighbors more than neighbors deviate from each other:
    $$\text{score}_i = \frac{|y_i - \bar{y}_{\mathcal{N}_k(x_i)}|}{\text{std}(y_{\mathcal{N}_k(x_i)}) + \epsilon}$$
    
    \item AC-aware: Distinguishes noise from activity cliffs by weighting deviation by neighborhood consistency. High deviation with high neighbor agreement suggests noise while high deviation with neighbor disagreement suggests an AC \citep{Dablander2023}:
    $$\text{score}_i = \frac{|y_i - \bar{y}_{\mathcal{N}_k(x_i)}|}{\sigma_y} \cdot \frac{1}{1 + \text{std}(y_{\mathcal{N}_k(x_i)})/\sigma_y}$$
\end{itemize}

And we used the following model-based estimators using prediction disagreement:

\begin{itemize}
    \item Cross-validation residual: Uses out-of-fold predictions; high residuals indicate labels inconsistent with the model learned from other samples.
    \item Ensemble disagreement: Trains multiple models and combines prediction variance with mean residual.
\end{itemize}

We used probabilistic methods that directly predict noise magnitude using uncertainty estimates:

\begin{itemize}
    \item Gaussian Process: Provides natural uncertainty estimates via posterior variance, with learned noise level $\sigma^2$ from the kernel \citep{gauche}.
    
    \item Heteroscedastic Neural Network: Predicts both $\mu(x)$ and $\sigma^2(x)$ using $\beta$-NLL loss \citep{Seitzer2022}:
    $$\mathcal{L} = \frac{1}{2}\left(\frac{(y - \mu)^2}{\sigma^2} + \log \sigma^2\right) \cdot \sigma^{2\beta}$$
    where $\beta = 0.5$ prevents the variance from hiding the mean errors.
    
    \item Conformal prediction: Uses split conformal prediction to obtain prediction intervals. Wider intervals should indicate higher uncertainty \citep{Shafer2008}.
    
    \item Deep ensemble: Trains $M$ independent networks and uses the variance of means as epistemic uncertainty plus the mean of variances as aleatoric uncertainty \citep{lakshminarayanan2017}.
    
    \item Bayesian Random Forest: Uses tree prediction variance and out-of-bag residuals as uncertainty estimates \citep{taddy2015}.
    
    \item Quantile regression: Fits models for multiple quantiles and uses interval width as noise estimate.
    
    \item NGBoost: Natural gradient boosting that predicts distributional parameters $(\mu, \sigma)$ directly \citep{Duan2020}.
\end{itemize}

NGBoost in particular was selected due to its promising uncertainty values with respect to the correlation between noise and prediction error, as discussed in Section~\ref{sec:uq}.  

\subsubsection{Label smoothing and sample weighting}
% ---- NOTES: Label smoothing and sample weighting ----
% Takeaway: Smoothing moves a label part of the way toward a target; weighting keeps the label and lowers how much the molecule counts. Only the proportional rule and the inverse weight are used in this chapter's experiments.
% Paragraph 1: before the lists, define the two inputs once
% - score_norm: the noise score rescaled to 0-1 over the fold's training molecules. Target: the mean label of the nearest training neighbours, or a model's prediction. alpha: strength, 0 (no change) to 1 (label replaced).
% Paragraph 2: changes to the prose
% - Threshold rule: the top alpha share by score is flagged, and each moves alpha of the way to its target, so alpha sets both how many labels move and how far.
% - Soft threshold: add the second line, $y + \text{strength}\cdot(\text{target} - y)$.
% - "to preserve effective sample size" is wrong. Say "so the total weight equals the number of training molecules".
% ---- END NOTES ----

Given the fact that all labels contain some magnitude of noise, we focused strongly on label smoothing over removal using the following methods:

\begin{itemize}
    \item Proportional: $y_{\text{smooth}} = y + \alpha \cdot \text{score}_{\text{norm}} \cdot (\text{target} - y)$
    
    \item Threshold: Only correct samples above the $(1-\alpha)$ percentile of scores.
    
    \item Soft threshold: Applies sigmoid-like correction strength:
    $$\text{strength} = \alpha \cdot \sigma\left(5 \cdot (\text{score}_{\text{norm}} - 0.5)\right)$$
    
    \item Confidence: For methods that output raw uncertainty $\sigma$ (e.g., NGBoost):
    $$\text{correction} = \frac{\alpha \cdot \sigma}{1 + \alpha \cdot \sigma}, \quad y_{\text{smooth}} = (1 - \text{correction}) \cdot y + \text{correction} \cdot \mu$$
\end{itemize}


Rather than directly modifying labels, we weighted more trustworthy samples during training:

\begin{itemize}
    \item Inverse: $w = \frac{1}{1 + \alpha \cdot \text{score}_{\text{norm}}}$
    
    \item Exponential: $w = \exp(-\alpha \cdot \text{score}_{\text{norm}})$
    
    \item Soft: $w = 1 - \alpha \cdot \text{score}_{\text{norm}}^2$, clipped to $[0.1, 1.0]$
\end{itemize}

All weights are normalized to mean 1 to preserve effective sample size.

\subsubsection{Recovery rate}
% ---- NOTES: Recovery rate ----
% Takeaway: Recovery is the share of the score lost to label noise that an LNL method gives back. 100\% restores clean performance; 0\% is no change; negative is worse than leaving the noise in.
% Takeaway: Detection is scored separately, with three numbers that need no fixed cut-off.
% Paragraph 1: how it is computed here
% - Change: the prose describes an older protocol, one random forest on ECFP4 scored by $R^2$. Replace the first-person opening ("We have a standard RF...") with a plain description: the same model fitted on clean, noisy and treated training labels, test labels never altered.
% - Every downstream model is scored in the dataset's own metric and divided by its own clean-to-noisy gap. The formula's direction cancels, so it works for RMSE too.
% - No recovery is reported when the gap is under 2\% of the clean score, or when the noisy score is no worse than the clean one.
% - Change: use \[ ... \] for the display equation instead of $$ ... $$.
% Paragraph 2: how detection is scored
% - Change: the prose says precision, recall and F1. The experiments report average precision (the headline), precision in the top k (k = number of changed molecules), and AUROC.
% - Detection is not reported under a noise type that changes every label.
% ---- END NOTES ----

For noise detection accuracy, we compute precision, recall, and F1 score, treating noise injection as ground truth.

We have a standard RF on ECFP4 train and test on the clean dataset, producing $R^2_{\text{clean}}$. Noise is then added to the clean data, RF on ECFP4 is then freshly trained and tested on this noisy data, producing $R^2_{\text{noisy}}$. Using the given estimation method and representation to test, noise estimation and label smoothing is applied to the noisy data, resulting in the denoised set. RF on ECFP4 is yet again run, this time on the denoised data, to produce $R^2_{\text{denoised}}$. From these performance values, we derive recovery rate, or the fraction of performance loss recovered after denoising, as follows: 

$$
\text{Recovery} = \frac{R^2_{\text{denoised}} - R^2_{\text{noisy}}}{R^2_{\text{clean}} - R^2_{\text{noisy}}}
$$

A recovery rate of 1.0 indicates complete restoration to clean performance; negative values indicate the denoising procedure made performance worse than the noisy baseline.


\section{Results: similarity}
% ---- NOTES: Results: similarity ----
% Takeaway: Answers the first research question: is predicting each molecule's activity from its nearest neighbours more accurate in a hybrid, or in a similarity hybrid, than in the single representations?
% Takeaway: Also asks whether scaling features by importance gives neighbours with closer activities.
% Paragraph 1: nearest-neighbour prediction in a hybrid against the single representations
% - Prediction: the mean activity of the 5 or 10 nearest training molecules.
% - Spaces compared: each single representation, the top three ordinary hybrids and the top three similarity hybrids, each hybrid with and without importance weighting.
% - Reference: for each held-out fold, the best single representation and best hybrid are chosen on the other folds. Report hybrid minus that representation, in the dataset's own metric, per fold.
% - Result needed: held-out score per dataset and fold for the best single representation, the ordinary hybrid and the similarity hybrid.
% Table: one row per dataset. Columns: best single representation (name and score), ordinary hybrid, similarity hybrid, each minus the single representation. Entry: mean over folds with the fold range.
% Paragraph 2: distance against activity difference
% - For sampled pairs of held-out molecules, absolute activity difference against distance, laid out like the hERG Ki figure in Chapter \ref{chap:flaws}.
% - Two panels: best single representation and similarity hybrid. One number per panel: Spearman correlation between distance and activity difference.
% - Result needed: the figure, and the correlation per space for each dataset.
% Figure: two scatter panels side by side, activity difference (dataset units) against distance. Caption gives the correlation for each panel.
% Paragraph 3: importance-weighted distance for k-nearest-neighbour prediction
% - Space: all of the dataset's representations joined side by side, not a hybrid. Each feature multiplied by the square root of its importance; Euclidean distance; k = 5, 10 and 20.
% - Scored in $R^2$ on regression datasets, not the dataset's own metric. One sentence.
% - No control shuffles the weights across features, so a gain cannot be split into "which features" and "unequal weights".
% - Result needed: unweighted and weighted $R^2$ per dataset, importance method and k.
% Paragraph 4: the similarity hybrid as an ordinary model input
% - In the cleaning of real labels, downstream models are also fitted on the similarity hybrid. With no cleaning, this compares the two hybrids as inputs to the same model on the same folds.
% - Result needed: held-out score with no cleaning, per dataset, model and fold, for the similarity hybrid, the ordinary hybrid and the top two single representations.
% ---- END NOTES ----
\todo[inline, color=yellow!40]{Results TODO(prelim): methods are final, numbers will be replaced by the server runs. [DIAGRAM D6.2]}


\section{Results: noise detection and estimation}
% ---- NOTES: Results: noise detection and estimation ----
% Takeaway: Two of the three tasks: finding which training labels are wrong (detection), and sizing how wrong they are (estimation).
% Takeaway: Finding noise is possible for outliers, and a systematic offset can be sized; correcting labels to improve the held-out score is a different problem.
% Paragraph 1: how detection and estimation are scored
% - Each noise score ranks a fold's training molecules from most to least suspect. Detection is graded against the molecules the injection singled out: widened under outlier and grouped-wider noise, clipped under censoring.
% - Chance level is the share singled out, e.g. about 0.1 at 10\% outlier noise. Report it beside every average precision.
% - Estimation is graded against the injected noise size in label units, or a size measured from repeat measurements.
% Paragraph 2: finding injected noisy labels, and whether the space matters
% - Question: does a neighbour-based score find the corrupted molecules more often when neighbours are found in a hybrid than in a single representation?
% - Ways to find neighbours: the single representation, the same weighted by importance, forest proximity, the hybrid, and the similarity hybrid. Model-based scores change only their input features.
% - Result needed: average precision per score, space and noise condition, with chance level and spread across folds.
% - Preliminary, will be replaced: antibody thermostability, outlier noise on 10\% of labels, one-hot, no hybrid. Average precision: plain Euclidean 0.1905, TreeSHAP-weighted Euclidean 0.2292, forest proximity 0.2251; chance about 0.09 (about 50 moved of 537 training molecules).
% - Preliminary, will be replaced, unverified: hERG Ki, outlier noise on 10\% of labels, one fold. Ten single representations reach average precision 0.4036 to 0.4326; every way of combining the ten falls inside that range.
% Paragraph 3: finding real suspect labels, nothing injected
% - Curator flags: ChEMBL curators mark suspect records \citep{Zdrazil2023}. On hERG IC50 in ChEMBL 37, 301 of 12,008 records carry a flag, on 215 of 9,802 molecules.
% - Most are flagged "outside typical range" because the value is extreme, so a score reading distance from neighbours is partly graded on the value itself. State this next to the number.
% - Inter-laboratory disagreement: BindingDB holds the same compound measured by different laboratories \citep{Gilson2016}. The Spearman correlation between a score and that spread grades it against real noise.
% - Result needed: average precision against curator flags; correlation with inter-laboratory spread.
% Paragraph 4: estimating how large the noise is
% - From repeat measurements: a molecule's label error variance is the per-measurement variance divided by its number of measurements $k$. Regressing squared out-of-fold residuals on $1/k$ gives the per-measurement noise variance as the slope.
% - From near-duplicate pairs: half the mean squared label difference of the closest 1\% of training pairs. An upper bound.
% - From the gamma test \citep{Stefansson1997}. Also an upper bound.
% - Put each estimate beside the model's own held-out mean squared error. Noise mitigation can only pay where label error is a large share of model error. Compare with the published size of Ki measurement error \citep{Kramer2012}.
% - Result needed: the three estimates per dataset and space, beside the model's mean squared error.
% Paragraph 5: the gap between finding noise and correcting it
% - Flag sets are built with precision set by hand, 0 to 1 in steps of 0.25, and each is acted on. This gives the precision a score would need before acting improves the held-out score. Compare it with the average precision the scores reach.
% - Result needed: the precision needed before acting pays, per dataset, against the precision reached.
% Table: one row per noise score. Columns: average precision under outlier noise at 10\%, grouped-wider noise and censoring, each in the best single representation and in the hybrid; chance level in the header row.
% ---- END NOTES ----
\todo[inline, color=yellow!40]{Results TODO(prelim): methods are final, numbers will be replaced by the server runs.}

\section{Results: noise mitigation by noise type}
% ---- NOTES: Results: noise mitigation by noise type ----
% Takeaway: Results go experiment by experiment (clipped labels, the twelve-method survey, cleaning real labels), then one closing subsection compares what each noise type allows.
% Paragraph 1: what every result is compared against
% - Each model's own top two single representations, the hybrid, and, where run, the importance-weighted or similarity hybrid. Scores in each dataset's own metric, per fold, no significance tests.
% - Change: the section heading says noise mitigation, but the survey also covers noise-robust training.
% - Caveat for the preliminary numbers: they compare against one fixed single representation per dataset, one-hot on antibody thermostability (its weakest representation for a random forest, 0.3471 Spearman) and ECFP4 on hERG Ki and Caco-2.
% ---- END NOTES ----


\subsection{Censoring: correcting labels clipped at a detection limit}
% ---- NOTES: Censoring: correcting labels clipped at a detection limit ----
% Takeaway: Raising every clipped label by one constant recovered most of the damage from clipping; per-molecule amounts did no better, and deleting the clipped molecules was worse than keeping them (preliminary).
% Paragraph 1: the question and why the clip is injected
% - Question: when training labels past an assay's limit are recorded at the limit, can they be corrected so the model predicts held-out molecules better?
% - Injected, not real: the hERG Ki data keep only ChEMBL records with an exact "=" relation, so clean training labels are clipped on purpose. Held-out labels keep their real values.
% Paragraph 2: hERG Ki at 10\% clipped (preliminary, will be replaced)
% - Setup: ECFP4, 3 folds, random forest, LightGBM, kNN and SVM. Limit 6.86 pKi, about 90 of 941 training molecules per fold.
% - Deleting the clipped molecules scores worse than keeping them clipped, on all four models, in every fold and at every clipping level.
% - Robust losses barely move from the clipped row: random forest with absolute-error splits 0.6574 pKi, LightGBM with the Huber loss 0.6480 pKi.
% Paragraph 3: other clipping levels and the per-molecule lift (preliminary, will be replaced)
% - Also run: hERG Ki at 25\% and 40\% clipped, and Caco-2 efflux ratio at 25\%.
% - The constant lift beat the clipped labels in 48 of 48 comparisons (4 models, 3 folds, 4 conditions). Recovery shrinks as more labels are clipped.
% - A separate lift per molecule scored higher RMSE than the constant in 15 of 16 model-condition pairs. Dealing the per-molecule lifts to the wrong molecules scored lower RMSE than dealing them correctly, in 16 of 16.
% - So per-molecule lift sizes did no better than one shared size, and no better than the same sizes given to the wrong molecules.
% - Result needed: recovery rate per treatment, model, feature space and clipping level; the three model comparators against their parent models.
% Table: hERG Ki, 10\% of training labels clipped at 6.86 pKi, ECFP4, RMSE (pKi), mean of 3 folds. One row per version of the training labels; columns random forest, LightGBM, kNN, SVM.
% - Never clipped: 0.6032, 0.6087, 0.6871, 0.6841.
% - Clipped: 0.6588, 0.6499, 0.7158, 0.7148.
% - Clipped, then constant lift: 0.6118, 0.6248, 0.6937, 0.6847.
% - Clipped molecules deleted: 0.7706, 0.7631, 0.8037, 0.8162.
% Table: recovery rate of the constant lift, ECFP4, mean of 3 folds. One row per dataset and clipping level; columns random forest, LightGBM, kNN, SVM.
% - hERG Ki 10\%: 85\%, 61\%, 77\%, 98\%.
% - hERG Ki 25\%: 56\%, 62\%, 54\%, 65\%.
% - hERG Ki 40\%: 45\%, 48\%, 51\%, 53\%.
% - Caco-2 25\%: 65\%, 71\%, 72\%, 68\%.
% - The raw RMSE pairs behind each share go in an Additional file.
% ---- END NOTES ----
\todo[inline, color=yellow!40]{Results TODO(prelim): methods are final, numbers will be replaced by the server runs.}


\subsection{Gaussian noise: learning-with-noisy-labels methods}
% ---- NOTES: Gaussian noise: learning-with-noisy-labels methods ----
% Takeaway: Under Gaussian noise the input moved the score more than the method: the hybrid with no LNL method beat the best method on one-hot (preliminary).
% Takeaway: The methods helped only at the noise level that actually lowered the plain network's score.
% Paragraph 1: what was run (preliminary, will be replaced)
% - Antibody thermostability, 672 antibodies, melting temperature. Spearman on test antibodies, mean of 3 folds.
% - Gaussian noise at 0.5 and 1.5 times the SD of the fold's clean training labels.
% - Inputs: one-hot (5600 features), the hybrid built on each fold's noisy labels, and the importance-weighted hybrid.
% Paragraph 2: the plain network before any method
% - No added noise: one-hot 0.4167, hybrid 0.5418, importance-weighted hybrid 0.5857 Spearman.
% - Noise at 0.5 times did not lower one-hot or the hybrid. At 1.5 times the hybrid fell from 0.5418 to 0.4859. So 1.5 times is the level that left damage to repair.
% Paragraph 3: how many of the twelve methods beat the plain network on the same input
% - At 0.5 times: 0 of 12 on one-hot, 0 of 12 on the hybrid, 5 of 12 on the importance-weighted hybrid.
% - At 1.5 times: 8 of 12 on one-hot, 8 of 12 on the hybrid, 6 of 12 on the importance-weighted hybrid.
% Paragraph 4: input against method
% - The hybrid scored above one-hot at both levels, for the plain network and every method.
% - At 1.5 times, the best method on one-hot (small-loss selection, 0.3568) is below the plain network on the hybrid (0.4859). The two ranges do not overlap.
% Paragraph 5: limits (two sentences)
% - 3 folds, one noise draw per level, and the hybrid was searched with a random forest while the methods train a network. The final runs fix both and replace one-hot with each model's top two single representations.
% - Result needed: per noise level and dataset, the count of methods beating the plain network on each input, and whether the input moves the score more than the method.
% Table: one row per method, the plain network first (0.4534, 0.5991, 0.5833 at 0.5 times; 0.2729, 0.4859, 0.5118 at 1.5 times), then the twelve in the order of the methods table. Columns: one-hot, hybrid and importance-weighted hybrid at 0.5 times, then the same three at 1.5 times. Cells: Spearman, mean of 3 folds. A line above the table says what one row is.
% ---- END NOTES ----
\todo[inline, color=yellow!40]{Results TODO(prelim): methods are final, numbers will be replaced by the server runs.}


\subsection{Real labels, nothing injected: cleaning}
% ---- NOTES: Real labels, nothing injected: cleaning ----
% Takeaway: On the two small-molecule datasets no cleanup improved on leaving the labels as they came (preliminary).
% Takeaway: On the antibody dataset whether cleaning helped depended on the feature space, not only on the cleanup (preliminary).
% Paragraph 1: what was run (preliminary, will be replaced)
% - hERG Ki and Caco-2 (RMSE, lower is better) and antibody thermostability (Spearman). Random forest, 3 folds; each number averaged over every setting of that cleanup.
% Paragraph 2: the two small-molecule datasets
% - hERG Ki, RMSE in pKi. No cleaning: ECFP4 0.6032, hybrid 0.6029, similarity hybrid 0.6235. Every cleanup, in every space, scored at or above no cleaning.
% - Caco-2, RMSE in log efflux ratio. No cleaning: ECFP4 0.3242, hybrid 0.3090, similarity hybrid 0.3226. Down-weighting matched no cleaning; every other cleanup was worse.
% Paragraph 3: the antibody dataset, where the direction depends on the feature space
% - Spearman. No cleaning: one-hot 0.4068, hybrid 0.6279, similarity hybrid 0.6087.
% - On one-hot every cleanup lowered the score (worst: removal, 0.3498). On the hybrid every cleanup raised it (largest: removal 0.6429). On the similarity hybrid the seven split three above, four below.
% - Do not call it a gain until the shuffled-score control is in. The range clip reads no suspicion score, so that control cannot test it.
% - Result needed: the shuffled-score control per cleanup and feature space, antibody dataset first.
% Paragraph 4: what the final runs report
% - Folds where the cleaned model beat the uncleaned one, out of folds run, with the mean difference beside it. QM9 reported apart as the control.
% - Result needed: folds won per cleanup, suspicion score, feature space and dataset, including QM9.
% Table: one block per dataset. Rows: no cleaning, then the seven cleanups. Columns: top two single representations, the hybrid, the similarity hybrid. Cell: folds won out of folds run, with the mean difference in the dataset's metric. The line above the table says what one row is and what one cell counts.
% ---- END NOTES ----
\todo[inline, color=yellow!40]{Results TODO(prelim): methods are final, numbers will be replaced by the server runs.}


\subsection{By noise type}
% ---- NOTES: By noise type ----
% Takeaway: Gaussian noise is the standard assumption and the hardest to fix: every label is wrong by an unknown amount in an unknown direction, and the neighbours used to judge it are noisy too.
% Takeaway: Noise can be found or sized and still not be removable. Better similarity does not solve this.
% Takeaway: Ask what kind of noise is likely before choosing a method.
% Paragraph 1: why Gaussian noise is the hardest to fix
% - Discussion, no new results.
% - Every label carries an error of unknown size and sign. No clean subset exists to learn from.
% - A detector that compares a molecule with its neighbours compares it with labels that are also noisy. A large disagreement can be noise or a real activity cliff.
% - Even a correctly flagged label has no known correction. Removing it loses a molecule. Replacing it with the neighbours' mean swaps its noise for theirs.
% Paragraph 2: finding the noise is not the same as removing it
% - Noise can be identified or estimated, yet a cleaned dataset can give a worse model than the noisy one it came from.
% - Evidence: the exploratory study ran 860 cleaning configurations on the QM9 HOMO-LUMO gap (Chapter~\ref{chap:flaws}). Only 8 beat training on the noisy labels at both noise levels (SD 0.3 and 0.6 eV).
% - The survey's preliminary numbers do not yet support this: at 0.5 times the noise did no damage to repair, and at 1.5 times 8 of 12 methods helped.
% - A better similarity can make neighbours closer in true activity. It cannot remove the noise on their labels.
% Paragraph 3: the other noise types, from easiest to fix to hardest
% - Censoring is the most fixable: the limit and the direction are known, only the size is not. That is why one constant lift recovered most of the damage.
% - Outlier noise is easier to find: the widened labels stick out against their neighbours. The fix is removal or down-weighting; neither recovers the true value.
% - Grouped-wider noise makes every neighbour in the family noisy too. Grouped-shifted noise moves a family together, so comparing with same-family neighbours cannot reveal it; knowing each record's source makes the offset estimable.
% - These three have detection and estimation results but no LNL result in this chapter. One sentence.
% Paragraph 4: practical takeaways
% - Data pooled from many laboratories: expect an offset per source. Keep the source of every record.
% - A single controlled experiment: expect censoring. Keep the "<" and ">" qualifiers.
% - Molecules from very different chemical domains: neighbour-based detection is least reliable.
% - Source and amount unknown: noise-robust training is safer than noise mitigation (Chapter~\ref{chap:discussion}).
% - Result needed: per noise type, whether any LNL method improved on the untreated model, on each feature space.
% Table: one row per noise type (Gaussian, outlier, grouped-wider, grouped-shifted, censoring). Columns: what is wrong in the labels; how easy the noisy labels are to find; whether the direction of the error is known; what makes a correction possible.
% Figure: five panels, one per noise type. x is the true label, y the observed label; the diagonal is no noise. Shade the part that can be detected; mark the part that can be corrected.
% ---- END NOTES ----
\todo[inline, color=yellow!40]{6.5.5: why Gaussian is hardest; practical takeaways. [DIAGRAM D6.3]}


\section{Conclusions}
% ---- NOTES: Conclusions ----
% Takeaway: One paragraph per research question, then what each noise type allows, the limits, and the bridge: model the noise before choosing a method.
% Takeaway: Short, 1-2 pages, no formulas.
% Paragraph 1: distances and nearest-neighbour prediction
% - Result needed: per dataset, nearest-neighbour prediction in the ordinary and similarity hybrid minus the best single representation; and the similarity hybrid as a downstream-model input.
% - If nearest-neighbour prediction improves when features are chosen for the property, then "similar" depends on the property, not on the two structures alone. Write this only if the results show it.
% Paragraph 2: noise detection and estimation
% - Finding a noisy label and correcting it are separate problems.
% - Result needed: per dataset, detection and estimation in the hybrid against the top two single representations.
% Paragraph 3: learning with noisy labels (one sentence per experiment)
% - Result needed: censoring, Gaussian-noise survey, and real-label cleaning, each against leaving the labels alone.
% Paragraph 4: what each noise type allows (two sentences)
% - The type of noise decides whether it can be corrected, not only found. Censoring has a known location and direction; Gaussian noise has neither.
% Paragraph 5: limit specific to this chapter
% - The clipped-label experiment takes its search settings from the unclipped labels, which a real user would not have.
% Paragraph 6: bridge to Chapter \ref{chap:discussion}
% - Model the noise before choosing a method. Which method helps depends on noise type, amount, dataset and feature space.
% ---- END NOTES ----


\chapter{Discussion and Future Directions}\label{chap:discussion}
% ---- TAKEAWAYS ----
% - Handle label noise realistically: ask where it comes from and what shape it has, not only how much there is.
% - Real assay noise is not only Gaussian, and Gaussian noise is the hardest case to correct.
% - No single method fits every dataset, for building the hybrid or for learning with noisy labels.
% - Simple is often better: tree ensembles under noise, fewer features, and neighbours on a small task-relevant feature set.
% - Correcting labels without ground truth is hard; prefer LNL methods that change training over noise mitigation, except for noise of known structure.
% - Next steps: classification and flipped labels, hybrids as inputs to generative models, and a faster search.
% ---- END TAKEAWAYS ----
% ---- NOTES: Discussion and Future Directions ----
% Paragraph 1: handle noise realistically (framing only, no new results).
% - Real assay labels carry shared offsets per laboratory \citep{Bentz2013} and labels clipped at detection limits \citep{Svensson2025}.
% - Gaussian noise is the hardest case: every label is off by an unknown amount in an unknown direction.
% - Cleaning labels can leave a model worse than the uncorrected labels would.
% - One sentence on what to ask of a dataset: pooled from many laboratories, one controlled experiment, or unknown source.
% Paragraph 2: the order of the chapter: a summary, two arguments (no single method fits; simple is often better), limitations, future directions.
% ---- END NOTES ----



\section{Summary}
% ---- NOTES: Summary ----
% Takeaway: each chapter's question and its answer, one paragraph per chapter, in the order of Section~\ref{sec:intro:questions}.
% - No numbers in this section; each chapter's conclusions hold the detail.
% Paragraph 1: the problem and the tools (Chapters~\ref{chap:intro} and~\ref{chap:background}).
% - QSAR models are trained on labels that carry experimental noise, and most work assumes that noise is Gaussian. One sentence on what Chapter~\ref{chap:background} defined.
% Paragraph 2: noise robustness of representations and models (Chapter~\ref{chap:robustness}).
% - Accuracy: model and representation explain comparable shares of the variance in R2.
% - Robustness (\aucnorm): the model explains far more of the variance than the representation.
% - Censored and grouped-shifted noise cost more robustness than Gaussian noise.
% - Mean predicted uncertainty rose with noise, but no model's per-molecule uncertainty found the corrupted labels. Under censoring the models were most confident on the clipped labels.
% Paragraph 3: why similarity fails on standard representations (Chapter~\ref{chap:flaws}).
% - Many features of a general-purpose fingerprint do not vary, or do not relate to the property, on a given dataset. A distance sums over all of them.
% - Some representations give different molecules identical feature vectors.
% - Judging a label by its neighbours' labels usually made the model worse.
% - What a representation would need: few features, chosen for the task, from several representations, traceable to their source.
% Paragraph 4: the adaptive hybrid representation (Chapter~\ref{chap:method}).
% - The method: choose features from several representations for one task, under a budget in features, with every choice made inside each fold.
% - Result needed: where the hybrid beat the best single representation and where it did not.
% - Result needed: whether importance weighting helped; transfer, antibody and separate-site results in one line each.
% - Limits: refer to Chapter~\ref{chap:method}'s limitations section; do not list them.
% Paragraph 5: hybrids for similarity and label noise (Chapter~\ref{chap:similarity}).
% - Result needed: whether the hybrid improved nearest-neighbour prediction, noise detection and LNL, and under which noise types any LNL method helped.
% Paragraph 6: how the pieces fit (short).
% - Chapter~\ref{chap:robustness} tells a reader which model to choose under noise. Chapters~\ref{chap:flaws} to~\ref{chap:similarity} tell them which features to give it.
% - Contributions in one line, then refer to "Summary of Research Contributions".
% ---- END NOTES ----

\todo[inline, color=yellow!40]{Write Sec 7.1. Draw on [NI 586], already used in the Chapter 3 conclusions -- do not simply repeat it.}


\section{No single method fits every dataset}
% ---- NOTES: No single method fits every dataset ----
% Paragraph 1: the claim.
% - No setting of the hybrid search was best on every dataset: downstream model, budget, number of base representations, importance method, allocation method, internal evaluator.
% - The same holds for LNL methods: which one helps, if any, depends on the data.
% Paragraph 2: what the best choice depends on.
% - The type and amount of label noise.
% - The number of molecules. A small dataset leaves few molecules inside a fold to score each setting on, so the chosen setting is less stable.
% - Whether any feature describes the property. The search cannot add signal the inputs do not carry.
% - The downstream model. The best hybrid for a tree ensemble need not be the best for a neural network.
% Paragraph 3: the evidence, pointing to each section rather than restating its numbers.
% - Chapter~\ref{chap:robustness}: the most robust representation for a model changed with the noise type, and the most robust model on QM9 was not the most robust on hERG Ki.
% - Chapter~\ref{chap:method} (preliminary): a fixed share of features per representation won some combinations of dataset, model and budget; one pooled ranking won others.
% - Chapter~\ref{chap:method}: which search settings won varied by dataset. Point to its tables; do not restate them.
% - Chapter~\ref{chap:similarity}: which LNL methods helped under each noise type.
% Paragraph 4: what the thesis offers instead of one recipe.
% - Not a fixed recipe or an architecture tuned to one problem. A procedure that gives one task its best chance by choosing its features.
% - For a practitioner: there is no default to trust. The useful question is which settings to try and how to compare them fairly.
% - The grid of winning search settings per dataset belongs in Chapter~\ref{chap:method} or an Additional file; refer to it here.
% ---- END NOTES ----


\subsection{Simple is often better}
% ---- NOTES: Simple is often better ----
% Takeaway: in four places the simpler choice did as well as the complex one, or better: the model under noise, the number of features, nearest-neighbour prediction, and LNL methods. No new results here.
% Paragraph 1: simple models under label noise (Chapter~\ref{chap:robustness}; robustness is \aucnorm).
% - On QM9 under Gaussian noise, forests, boosted trees, SVM and the GP were accurate on clean labels and robust.
% - The plain and Bayesian networks were among the most accurate on clean labels and the least robust. The variational networks sat in the middle, so write "tree ensembles and SVM, not plain neural networks".
% - Robust is not accurate: the most robust model on QM9, NGBoost, was among the least accurate on clean labels.
% Paragraph 2: fewer features (Chapters~\ref{chap:method} and~\ref{chap:similarity}).
% - The search sweeps the budget in features and does not assume more features are better.
% - Evidence: point to Chapter~\ref{chap:method}'s budget results and Chapter~\ref{chap:similarity}'s results for predicting each molecule's activity from its five nearest neighbours on a small hybrid. No per-dataset numbers here.
% Paragraph 3: LNL methods that did not help (no method-level detail).
% - Methods that used a distance between molecules, and neighbour-based label cleaning, made models worse.
% - Evidence: point to Chapter~\ref{chap:similarity} for which LNL methods beat the uncorrected model, and whether cleaning real labels helped. No per-dataset counts here.
% - Clipped labels are the exception (see Limitations). Do not write that every method failed.
% Paragraph 4: why LNL methods from other fields do not carry over to assay data.
% - Most were built and tested on large image datasets with flipped class labels \citep{Song2022}.
% - Assay datasets are small. Their labels are continuous, and every label carries some error.
% - Many methods rely on neighbours, or on how easily a sample is learned. Both need a space where close molecules have similar activity, which Chapter~\ref{chap:flaws} shows standard representations do not give.
% Paragraph 5: the balance.
% - A method tuned closely to one dataset loses generality; the tasks here span binding affinity, ADME, quantum properties and antibodies.
% - Fix the simple, robust parts (a tree ensemble, or neighbours on a small feature set). Search only the choices that vary by dataset. The search can over-fit too (Chapter~\ref{chap:method}).
% ---- END NOTES ----


\section{Limitations}
% ---- NOTES: Limitations ----
% Paragraph 1: correcting label noise without ground truth.
% - One sentence: limits of the hybrid search itself (run time, repeating the search, reproducibility) are in Chapter~\ref{chap:method}'s "Limitations".
% - The thesis set out to mitigate label noise. It is a hard problem with no clear general solution.
% - Every experimental label carries unknown noise, so there is no clean copy to check a correction against. In the assay experiments "clean" means "no added noise"; only QM9 is close to noise-free.
% - Finding a noisy label is easier than correcting it. Possible: outlier detection, and systematic errors of known structure (a shift per source, clipping at a detection limit).
% - The injected noise conditions are models. Real data mix them in unknown amounts.
% Paragraph 2: recommendation.
% - Noise mitigation writes a correction into the data before training. A wrong guess stays in the data for every later model.
% - LNL methods that change training adapt throughout training and use all the data \citep{Song2022}. Noise mitigation is cheaper: done once, reused by any model.
% - Correct labels directly only when the noise has a known structure: censoring (limit and direction known), gross outliers, or an offset tied to one source.
% - Evidence: the clipped-label result in Chapter~\ref{chap:similarity}, one clause.
% - Use uncertainty to judge how noisy a dataset is, not which labels are wrong (Chapter~\ref{chap:robustness}).
% Paragraph 3: other limits of scope.
% - The label-noise experiments are regression only.
% - The controlled label-noise experiments use QM9 and three small assay datasets. No dataset has repeated measurements with a known true value for every molecule.
% - Compute limited how many noise conditions, noise levels, folds and datasets could be crossed.
% - Results are per-fold scores and their spread over 3 to 5 folds, with no significance tests. Some patterns are a direction only.
% ---- END NOTES ----
\todo[inline, color=yellow!40]{7.5: noise mitigation is hard; recommend LNL over direct mitigation. Method limitations are in 5.10.}


\section{Future directions}
% ---- NOTES: Future directions ----
% Paragraph 1: classification and flipped labels (the first next step).
% - The noise work here is almost all regression. In classification a noisy label is a flipped class, not a shifted value.
% - Many QSAR class labels come from cutting a continuous value at a threshold, so flips concentrate near the cut-off.
% - NoiseInject already has six flip conditions, and the hybrid search already handles classification. Missing: a study of the molecule-dependent flips at the scale of the regression study.
% - Open question: do the regression findings carry over? Most LNL methods were designed for flipped labels.
% Paragraph 2 (her paragraph "I also plan to extend..."): edits.
% - Cut GNN embeddings, sequence encoders, TreeSHAP and gradient-based attribution: the search already includes them.
% - Cut small-loss selection: it scores training molecules, not features, and is one of the twelve LNL methods surveyed.
% - Keep as open: transformed inputs (normalised or binarised features), features from graph distances, and importance methods that treat the model as a black box (KernelSHAP, dropping one feature and refitting).
% - Keep "and datasets" only if a domain outside "Datasets and splitting" is named.
% Paragraph 3 (her paragraph "As mentioned earlier..."): edits.
% - Cut the plan to adapt the hybrid to datasets and models: the search already keeps one best configuration per dataset, downstream model and budget.
% - Cut the planned trade-off study: Chapter~\ref{chap:method} reports held-out score against budget per model.
% - Cut "a very small default hybrid with RF-based importance" and "I do not currently have a full set of results on hybrid".
% - Keep one sentence: a search objective that also rewards robustness to label noise. It is not built.
% Paragraph 4 (her paragraph "Various estimation methods..."): other noise types.
% - Chapter~\ref{chap:similarity} compares LNL methods under censoring, Gaussian noise and real labels. Still open: outliers and grouped offsets.
% - Keep her two questions: does neighbour-based cleaning do better on outliers, and do label smoothing and sample weighting still lead there?
% Paragraph 5: hybrids as inputs to generative models.
% - The model that scores molecules for a molecule generator could take the hybrid as its input.
% - Open question: what should the search optimise? Held-out accuracy on known molecules may not reward the right thing for new ones.
% Paragraph 6: a faster search.
% - The full search is too heavy for routine use. The fast tier (Table~\ref{tab:search_tiers}) and adaptive pruning already cut cost.
% - Next step: a warm start, searching near the blueprint of a similar dataset.
% Paragraph 7: one sentence each.
% - New embeddings: does the hybrid improve a weak new embedding, not only absorb a strong one?
% - Blueprints across domains, at the level of representation families.
% - Data at separate sites: a fuller federated-learning deployment with Lhasa, with real sites running the selection on their own data \citep{McMahan2017}.
% Paragraph 8 (her paragraph "What does this mean in practice?"): edits.
% - Say "the hybrid search" instead of "KIRBy", and "the representations searched" instead of "the pool".
% - Tone down "a substantive change" and "a principled combination", or tie them to Chapter~\ref{chap:method}'s results.
% Figure: hub-and-spoke with the hybrid search in the centre. Spokes: classification, generative scoring, separate sites, a faster search, new embeddings or domains.
% ---- END NOTES ----

I also plan to extend the hybrid research to additional base representations and datasets. Some alternative representations include additional GNN embeddings, sequential encoders, augmented features such as embedded distances from molecular graphs, and transformed representations (i.e., normalization, binarization of continuous features). Additionally, I want to explore alternative feature importance scoring methods, especially ones that involve different families of models. Specifically, I plan to look at TreeSHAP, KernelSHAP, drop-column importance for noise robustness, gradient-based attribution for NN embeddings, and small-loss selection, as discussed in Section~\ref{sec:lnl_motivation}.

As mentioned earlier, although I do not currently have a full set of results on hybrid, they have been showing promise throughout a rank of QSAR and noise mitigation tasks. I would like to explore the complexity tradeoff between model and representation complexity, as I've seen a trend of complicated models like NNs having stronger predictive performance with smaller hybrids and more simplistic tree-based methods benefiting from larger hybrids. This noise mitigation experiment relied on a very small default hybrid and used RF-based feature importance. I would like to leverage my allocation strategies to adapt hybrid models to datasets \textit{and} models. These strategies can also be modified to optimize for noise robustness in addition to predictive performance.

Various estimation methods showed promise, meriting further investigation. Additionally, neighbor-based noise mitigation strategies which are not fit for Gaussian noise may perform differently on outlier-based strategies. The dominance of estimators with label smoothing and sample weighting may be challenged for different types of noise.

What does this mean in practice? A new foundation model appears every few months in molecular, protein, or antibody modeling. A researcher using KIRBy does not need to decide whether to replace their current encoding or retrain their pipeline; they compute the new model's embeddings, add the matrix to the pool, and re-run the search. The search determines whether the new embedding contributes useful features and, if so, how many and which ones. In a field where encoding choice is often determined by familiarity or lab convention rather than by systematic comparison, this is a substantive change. It replaces a single bet with a principled combination, tailored to each task, that updates as the pool grows.


\pagebreak
\chapter*{Summary of Research Contributions}
\addcontentsline{toc}{chapter}{Summary of Research Contributions}

% ---- NOTES: Summary of Research Contributions ----
% Order: papers in chapter order, the conference presentation, the blog posts, then a new \section{Software}. One or two sentences per entry.
% Software: three pieces of open-source software. Each entry: name, one sentence, link. Licence and version: refer once to Appendix~\ref{app:availability}.
% - NoiseInject: a Python package that adds label noise to training labels and scores how a model's performance and uncertainty change as the noise grows. github.com/adpunt/noise_inject; https://doi.org/10.5281/zenodo.20532710; \citep{noiseinject}.
% - The QM9 pipeline: github.com/adpunt/qsar_qm_models. It produced the QM9 results of Chapter~\ref{chap:robustness}.
% - KIRBy (to be released): the Python repository holding the hybrid search, the representation generators, the dataset loaders and the experiments of Chapters~\ref{chap:flaws} to~\ref{chap:similarity}. It runs without the QM9 pipeline.
% ---- END NOTES ----

\section{Research Papers}
% ---- NOTES: Research Papers ----
% Paragraph 1: the NoiseInject paper (Chapter~\ref{chap:robustness}).
% - Change: title to "NoiseInject: A Benchmarking Framework to Study Noise Robustness in QSAR Models"; author order to Punt, Hanser, Werner, Morris. Journal of Cheminformatics, in preparation.
% - Change: "prepration" to "preparation", in both entries.
% Paragraph 2: split the hybrid entry into two papers, each with a title and one sentence on its status.
% - Building the adaptive hybrid representation (Chapter~\ref{chap:method}): when and why it helps, whether it is explainable, and whether it extends to new embeddings and separate sites.
% - Hybrids for molecular similarity and label noise (Chapter~\ref{chap:similarity}). The current title's "noise prediction robustness and mitigation" is too narrow: the chapter covers noise detection and LNL.
% - Working titles and venue for both: not yet chosen.
% ---- END NOTES ----
\textbf{In prepration:} Punt, A., Morris, G., Werner, S., \& Hanser, T. (2026). Noise-Robust QSAR: evaluating probabilistic methods across molecular representations \\
My supervisors and I are currently iterating on an advanced draft for this research paper based on the material in Chapter~\ref{chap:robustness}, which we intend to submit to the \textit{Journal of Cheminformatics} within the next few months.

\textbf{In prepration:} Punt, A., Morris, G., Werner, S., \& Hanser, T. (2026). Improving noise prediction robustness and mitigation with feature-importance hybrid representations in multi-modal contexts \\ I am in the early stages of writing a draft and extending the experiments in Chapters~\ref{chap:method} and~\ref{chap:similarity} within the next few months.

\section{Conference Presentations}
% ---- NOTES: Conference Presentations ----
% - Change: "Posted presented" to "Poster presented".
% ---- END NOTES ----
Punt, A., Morris, G., Werner, S., \& Hanser, T. (2024). Investigating the impact and mitigation of noise in QSAR models. Posted presented at the Autumn UK-QSAR 2024 Meeting.


\section{Blogposts}

I have contributed to the \textit{Blopig} blog of the Oxford Protein Informatics Group (OPIG):

\begin{itemize}
    \item \href{https://www.blopig.com/blog/2025/10/are-you-addicted-to-dopamine/}{Are you addicted to dopamine?}
    \item \href{https://www.blopig.com/blog/2025/05/demystifying-git-and-merge-conflicts/}{Demystifying Git and merge conflicts}
    \item \href{https://www.blopig.com/blog/2025/03/ai-in-academic-writing-ethical-considerations-and-best-practices/}{AI in Academic Writing: Ethical Considerations and Best Practices}
    \item \href{https://www.blopig.com/blog/2024/08/memory-mapped-files-for-efficient-data-processing/}{Memory-mapped files for efficient data processing}
    \item \href{https://www.blopig.com/blog/2024/04/exploring-multilingual-programming/}{Exploring multilingual programming}
\end{itemize}

%%========================= APPENDICES ======================================

\appendix

\chapter{Datasets and curation}\label{app:datasets}

% ---- TAKEAWAYS ----
% - Every dataset has its source, what its label measures, and each curation step in the order it runs, so the exact sets can be rebuilt.
% - Small molecules are standardised and deduplicated by InChIKey, so one molecule cannot sit in both a training fold and a test fold.
% - Some curation steps already change the labels before any noise is added: repeats merged into a median, extremes clipped, repeat measurements kept.
% - Chapter 3 used the same sources with a lighter curation, which explains the different molecule counts between Chapters 3 and 5.
% - What one sample is, and how folds are grouped, differs by domain: molecule, whole reaction, antibody, DNA sequence, tumour.
% ---- END TAKEAWAYS ----

% TODO: review, needs manual review
This appendix records where each dataset came from, what its label measures, and every curation step applied to it, in the order the steps run. It covers the four datasets of Chapter~\ref{chap:robustness}, the 16 datasets of Chapter~\ref{chap:method}, and the further sets used only for the separate-site experiments of Chapter~\ref{chap:method} and in Chapter~\ref{chap:similarity}. The label, metric and fold grouping of each of the 16 datasets are summarised in Table~\ref{tab:datasets}, and the splitting rules are given in the Datasets and splitting subsection of Chapter~\ref{chap:method}; neither is repeated here. Table~\ref{tab:curation_counts} gives the number of samples in each release and after curation. All counts were obtained by running the curation code on the downloaded files.

% TODO: review, needs manual review
\begin{table}[htbp]
\centering
\small
\setlength{\tabcolsep}{3pt}
\caption{Sample counts before and after curation. One row is one dataset. ``Released'' is the number of rows in the downloaded file; ``Used'' is the number after every step in this appendix. The last column names the steps that change the count.}
\label{tab:curation_counts}
\begin{tabular}{p{3.4cm}p{1.6cm}p{1.4cm}p{7.0cm}}
\toprule
Dataset & Released & Used & Steps that remove samples \\
\midrule
\multicolumn{4}{l}{\textit{Small molecules}} \\
hERG $K_i$ & 1,815 & 1,412 & no pChEMBL value (3); not a binding assay (75); repeats merged into one median per compound (322); InChIKey duplicates (3) \\
LogD & 5,326 & 5,028 & no value (287); repeats merged after standardisation (11) \\
Caco-2 efflux & 5,326 & 2,158 & no value (3,165); repeats merged after standardisation (3) \\
Rat liver microsome clearance & 3,521 & 3,053 & no value (467); InChIKey duplicates (1) \\
MDR1--MDCK efflux & 3,521 & 2,641 & no value (879); InChIKey duplicates (1) \\
QM9 & 130,831 & 9,978 & not parsable by RDKit (1,403); random draw of 10,000; InChIKey duplicates (22) \\
hERG blockers, training set & 7,772 & 7,606 & also in the hERG $K_i$ set (166) \\
hERG blockers, test set & 4,383 & 4,213 & also in the hERG $K_i$ set (170) \\
hERG IC$_{50}$ & 12,010 & 12,008 & missing or non-positive value (2) \\
\midrule
\multicolumn{4}{l}{\textit{Chemical reactions}} \\
Buchwald--Hartwig yield & 3,955 & 3,955 & none \\
Buchwald--Hartwig plates & 4,599 & 4,312 & no aryl halide (287) \\
\midrule
\multicolumn{4}{l}{\textit{Antibodies}} \\
Koenig expression & 4,275 & 4,275 & none \\
Nanobody thermostability & 672 & 672 & none \\
Humanness & 564 & 564 & none \\
Trastuzumab--HER2 binding & 524,346 & 5,000 & stratified random draw of 5,000 \\
\midrule
\multicolumn{4}{l}{\textit{Nucleotide sequences}} \\
H3K4me3 & 18,244 & 18,244 & none (training and test files merged) \\
TATA promoter & 5,274 & 5,274 & none (training and test files merged) \\
\midrule
\multicolumn{4}{l}{\textit{Cancer genomics (TCGA)}} \\
Lung adenocarcinoma vs squamous & 1,129 & 1,017 & not a primary tumour (112) \\
Breast cancer PAM50 subtype & 1,097 & 821 & no PAM50 call, or called normal-like (276); count is of primary tumours \\
\bottomrule
\end{tabular}
\end{table}
% Result needed: re-count Table~\ref{tab:curation_counts} on the copies of the files the full runs read on the cluster (these counts come from the local copies).

\section{Small molecules}\label{app:datasets:smallmol}

% TODO: review, needs manual review
Three steps are shared by several small-molecule datasets. The first is structure standardisation with RDKit \citep{rdkit}. The steps run in this order: cleanup, keep the largest fragment (removing salts and solvents), neutralise charges, and choose the canonical tautomer. Any molecule that then contains a metal atom is rejected, and the result is written as canonical SMILES. Standardisation is applied to hERG $K_i$, LogD and Caco-2. The Biogen sets, QM9 and the hERG blocker sets are only checked to parse in RDKit.

% TODO: review, needs manual review
The second shared step removes duplicate molecules by InChIKey \citep{Heller2015}, keeping the first occurrence. Two different SMILES strings can describe the same molecule, and without this step one molecule could sit in both a training fold and a test fold. The step is applied to hERG $K_i$, LogD, Caco-2, the two Biogen sets and QM9. The hERG blocker sets are deduplicated by InChIKey only against the other sites, as described below. The step is skipped on purpose for hERG IC$_{50}$, which keeps one row per measurement. The third shared step is the fold grouping by Bemis--Murcko scaffold \citep{Bemis1996}, whose form in the hybrid experiments is given in the Datasets and splitting subsection of Chapter~\ref{chap:method}.

% TODO: review, needs manual review
\paragraph{\texorpdfstring{hERG $K_i$.}{hERG Ki.}} Records for the hERG potassium channel (ChEMBL target CHEMBL240) were retrieved from ChEMBL \citep{Zdrazil2023} through its web interface. The query keeps records with standard type $K_i$, a pChEMBL value present, an exact relation (``$=$''), and no curator data-validity comment. Only binding assays (ChEMBL assay type B) are kept. The label is the pChEMBL value, $\mathrm{p}K_i = -\log_{10} K_i$ with $K_i$ in molar units. Repeat measurements of one compound are merged into their median, and compounds whose repeats have a standard deviation above 1.0 p$K_i$ unit are removed, following \citet{landrum2024}. The structures are then standardised and deduplicated by InChIKey. Because repeats are merged on ChEMBL's own SMILES before standardisation, two ChEMBL entries that standardise to the same molecule are not merged by median; the InChIKey step keeps the first of them. On a pull of September 2026, 1,815 $K_i$ records reduce to 1,812 with a pChEMBL value and 1,737 from binding assays. These give 1,415 compounds, 226 of them measured more than once and none removed by the standard-deviation rule. The InChIKey step removes 3 more, leaving 1,412 compounds with p$K_i$ from 4.07 to 9.41. The file read by the experiments was pulled on 1 February 2026; the ChEMBL release was not recorded at the time. Re-running the same query in September 2026 against ChEMBL release 37 gives a byte-identical file of 1,415 compounds.
% Result needed: confirm the hERG $K_i$ file the full runs read on the cluster is the same file (1,415 compounds, pulled 1 February 2026).

% TODO: review, needs manual review
\paragraph{LogD and Caco-2 efflux.} Both endpoints come from the training split of the OpenADMET--ExpansionRx challenge data, released on Hugging Face under a CC-BY-4.0 licence \citep{openadmet}. The file holds 5,326 molecules; 5,039 have a LogD value and 2,161 a Caco-2 efflux ratio. For each endpoint, rows without a value are dropped and the structures are standardised. Rows that standardise to the same SMILES are merged into their median label. This merges 11 LogD pairs and 3 Caco-2 pairs. The efflux ratio, which has no unit, is then transformed to $\log_{10}$. Labels are clipped at the dataset's own 1st and 99th percentiles: $-1.0$ and $4.5$ for LogD, and $-0.319$ and $1.679$ for $\log_{10}$ efflux ratio. Clipping sets 116 LogD labels (54 low, 62 high) and 45 Caco-2 labels (23 low, 22 high) to a limit. The extreme labels of both sets are therefore flattened before any noise is added in Chapter~\ref{chap:similarity}. The InChIKey step removes nothing on either set, leaving 5,028 and 2,158 molecules.
% TODO(lit): pH at which the OpenADMET LogD values were measured (not stated on the Hugging Face page).

% TODO: review, needs manual review
\paragraph{Rat liver microsome clearance and MDR1--MDCK efflux.} These are two endpoints of the Biogen ADME set of 3,521 compounds \citep{fang2023}, downloaded from Polaris as \texttt{biogen/adme-fang-v1} \citep{polaris}. Rat liver microsome clearance is the intrinsic clearance $\mathrm{CL_{int}}$ in mL/min/kg, and the MDR1--MDCK endpoint is an efflux ratio. Both are released already as $\log_{10}$ values and are used as released. Rows without a value, or whose SMILES RDKit cannot parse, are dropped, and duplicates are removed by InChIKey. There is no standardisation, no median merge and no clipping. Clearance has 3,054 values and efflux 2,642; the InChIKey step removes one molecule from each, leaving 3,053 and 2,641.

% TODO: review, needs manual review
\paragraph{QM9.} QM9 \citep{Ramakrishnan2014, qm9_dataset} was read through PyTorch Geometric \citep{pytorchGeometric}, which distributes 130,831 molecules. Of these, 129,428 parse in RDKit; the other 1,403 carry an over-valent carbon atom in the released structure file. This list of parsable molecules is identical to the one used in Chapter~\ref{chap:robustness}. From it, 10,000 molecules were drawn once at random with seed 42, so the subsample is a random draw of molecules, not a random split. The InChIKey step removes 22 duplicates, leaving 9,978. The label is the HOMO--LUMO gap in eV, from 1.88 to 11.37 eV in this sample.

% TODO: review, needs manual review
\paragraph{hERG blockers held at separate sites.} The separate-site experiments of Chapter~\ref{chap:method} use two further hERG sets from Lhasa Limited: a training set of 7,772 molecules (3,630 blockers, 46.7\%) and a test set of 4,383 molecules (1,084 blockers). Both were provided through the Effiris and FLuID collaboration \citep{Bassani2023, Hanser2023}. Each label is Lhasa's own binary blocker call, and the two sets were labelled under different rules. The molecules are stored as RDKit molecule objects and converted to SMILES. Entries that give no valid SMILES are dropped, which removes none. Neither set is standardised. The sets are proprietary and cannot be redistributed (Appendix~\ref{app:availability}). In the same experiments the hERG $K_i$ set acts as a third site. Its labels are made binary, with a compound counted as a blocker when $\mathrm{p}K_i \ge 6.0$ ($K_i \le 1\,\mu$M), which gives 467 blockers and 945 non-blockers. Molecules are then matched by InChIKey across the three sites, taken in the order hERG $K_i$, training set, test set, and only the first copy is kept. This removes 166 molecules from the training set and 170 from the test set, all of them also in the hERG $K_i$ set, leaving 7,606 and 4,213. The two Lhasa sets share no molecule. For 56 of the 166 and 57 of the 170 shared molecules, the Lhasa blocker call disagrees with the call from p$K_i$.

% TODO: review, needs manual review
\paragraph{\texorpdfstring{hERG IC$_{50}$.}{hERG IC50.}} This set is used only in Chapter~\ref{chap:similarity}, where the disagreement between repeat measurements and the curator flags serve as real evidence of label error. It keeps every measurement as its own row. The query takes ChEMBL records for CHEMBL240 with standard type IC$_{50}$ in nM and an exact relation \citep{Zdrazil2023}. Records with a missing or non-positive value are dropped, and the label is $\mathrm{pIC_{50}} = 9 - \log_{10}(\mathrm{IC_{50}}/\mathrm{nM})$. No assay-type filter is applied. Of the 12,010 records, 10,219 come from binding assays. Records with a curator data-validity comment are kept and marked; there are 301 of them, on 215 molecules. There is no standardisation, no median merge and no InChIKey step. The result is 12,008 measurements of 9,802 molecules. Folds are grouped by scaffold, and every measurement of a molecule shares that molecule's scaffold, so no molecule is split between training and test.

\section{\texorpdfstring{The Chapter~\ref{chap:robustness} versions of LogD, Caco-2, hERG $K_i$ and QM9}{The Chapter 3 versions of LogD, Caco-2, hERG Ki and QM9}}\label{app:datasets:ch3}

% ---- NOTES: The Chapter 3 versions of LogD, Caco-2, hERG K_i and QM9 ----
% - Change: the clause "and is also redistributed through the OpenADMET benchmark suite" is wrong. OpenADMET does not carry QM9; the QM9 runs read the PyTorch Geometric copy \citep{pytorchGeometric}.
% - Change: "Section~\ref{sec2}" points to the Background chapter. The hERG $K_i$ curation is in Section~\ref{sec:robustness:methods}.
% ---- END NOTES ----
The datasets used in this study are publicly available. The QM9 dataset is available on Figshare \citep{qm9_dataset} and is also redistributed through the OpenADMET benchmark suite \citep{Ramakrishnan2014, openadmet}. The LogD and Caco-2 datasets are available in the OpenADMET repository \citep{openadmet}. The hERG-Ki dataset was derived from ChEMBL \citep{Zdrazil2023} following the curation protocol described in Section~\ref{sec2}, and can be regenerated from ChEMBL using that protocol.

% TODO: review, needs manual review
Chapter~\ref{chap:robustness} used the same sources as Chapter~\ref{chap:method} for LogD, Caco-2 and hERG $K_i$, but a lighter curation. The differences are listed in Table~\ref{tab:curation_ch3_ch5}. They explain why the sample counts differ between the two chapters. The QM9 experiments of Chapter~\ref{chap:robustness} draw ten replicates of 10,000 molecules from the same 129,428 parsable molecules. Each replicate has its own seed and its own 80/10/10 scaffold split into 8,000, 1,000 and 1,000 molecules. Chapter~\ref{chap:robustness} also merges hERG $K_i$ repeats a second time, on its own standardised SMILES; on this file that merges nothing, so its set keeps 1,415 compounds.

% TODO: review, needs manual review
\begin{table}[htbp]
\centering
\small
\caption{Curation of the three assay datasets in Chapters~\ref{chap:robustness} and~\ref{chap:method}. One row is one curation step or one count.}
\label{tab:curation_ch3_ch5}
\begin{tabular}{p{3.4cm}p{4.8cm}p{4.8cm}}
\toprule
 & Chapter~\ref{chap:robustness} & Chapter~\ref{chap:method} \\
\midrule
Structure standardisation & largest fragment and canonical SMILES only & cleanup, largest fragment, neutralised charges, canonical tautomer, metals rejected \\
Label clipping (LogD, Caco-2) & none & 1st and 99th percentile \\
InChIKey deduplication & none & first occurrence kept \\
Scaffold key & stereochemistry removed; each acyclic molecule its own group & stereochemistry kept; all acyclic molecules one group \\
Molecules: LogD, Caco-2, hERG $K_i$ & 5,039; 2,161; 1,415 & 5,028; 2,158; 1,412 \\
Scaffold groups: LogD, Caco-2, hERG $K_i$ & 2,245; 1,038; 749 & 2,413; 1,086; 769 \\
\bottomrule
\end{tabular}
\end{table}

\section{Chemical reactions}\label{app:datasets:reactions}

% TODO: review, needs manual review
Both reaction datasets come from the high-throughput Buchwald--Hartwig C--N coupling experiments of \citet{Ahneman2018}, in which aryl halide, palladium ligand, base and additive were varied on reaction plates. The label is the reaction yield in \%, from 0 to 100. Here one sample is a whole reaction, written as its reactants and product, so the representations and the grouping differ from those of the single-molecule sets.

% TODO: review, needs manual review
\paragraph{Buchwald--Hartwig yield.} This is the version released with \citet{Schwaller2021a}: the first cross-validation sheet of the input workbook in that paper's public repository. All 16 sheets of the workbook hold the same 3,955 reactions. Each reaction is written as aryl halide, 4-methylaniline, ligand, base and additive going to a product. The product is generated by applying a C--N coupling template to the aryl halide. A reaction would be dropped if the template failed, which it never does, so all 3,955 reactions are kept. They cover 15 aryl halides, 4 ligands, 3 bases and 22 additives, and fall into 22 additive groups of 178 to 180 reactions.

% TODO: review, needs manual review
\paragraph{Buchwald--Hartwig plates.} This is the data table of the public code repository of \citet{Ahneman2018}, which holds 4,599 rows over three plates. Its product SMILES are taken from the file. The 287 rows without an aryl halide are dropped. The 180 remaining rows without an additive are kept and form their own group, so the set has 24 additive groups of 177 to 180 reactions. The 4,312 reactions kept are not independent of the first set. Every one of the 3,955 reactions in the first set appears here, with the same aryl halide, ligand, base, additive and yield. The plates set adds 357 reactions: the 180 without an additive and 177 with 5-phenyl-1,2,4-oxadiazole as the additive.

\section{Antibodies}\label{app:datasets:antibodies}

% TODO: review, needs manual review
The Koenig and nanobody sets were downloaded as CSV files from the FLAb benchmark repository \citep{Chungyoun2024}. In both, rows with a missing heavy chain, a heavy chain shorter than 11 residues, or a missing or non-numeric label are dropped; no row is removed by this rule. The humanness set was downloaded from the BioPhi publication repository and the trastuzumab set from its own repository. The FLAb and BioPhi files are taken from each repository's main branch, with no commit pinned. Each file has been changed only once: the Koenig file in FLAb commit \texttt{c41e56e} and the nanobody file in commit \texttt{005e6c1}, both on 6 August 2025, and the two humanness files in BioPhi commit \texttt{c0797f2} on 24 August 2021. The local copies of the Koenig and humanness files are identical to those commits.
% Result needed: confirm on the cluster that the nanobody file matches FLAb commit 005e6c1 (there is no local copy to compare).

% TODO: review, needs manual review
The nanobody and humanness sets are grouped by sequence identity. The identity of two sequences $a$ and $b$ is
\begin{equation}
\mathrm{id}(a,b) = 1 - \frac{\mathrm{lev}(a,b)}{\max(|a|,|b|)},
\end{equation}
where $\mathrm{lev}$ is the Levenshtein (edit) distance and $|a|$ is the length in residues. Sequences are clustered by single linkage at $\mathrm{id} \ge 0.95$: two sequences share a cluster if a chain of pairs, each at identity 0.95 or more, joins them. For a set with paired chains, heavy and light chain would be joined into one string before clustering, as in \citet{Kenlay2024}; both clustered sets here carry one chain per antibody. The clusters are not one antibody each. The 672 nanobodies fall into 309 clusters, of which 185 hold one nanobody and the largest holds 32. The 564 humanness antibodies fall into 461 clusters, of which 404 hold one antibody and the largest holds 13.

% TODO: review, needs manual review
\paragraph{Koenig expression.} This is a deep mutational scan of the anti-VEGF antibody G6 \citep{Koenig2017}, taken from FLAb. The label is the expression enrichment ratio, which has no unit and runs from 0.00098 to 3.178. All 4,275 variants carry both chains, with heavy chains of 120 residues and light chains of 108. Each variant carries exactly one mutation: 2,261 in the heavy chain and 2,014 in the light chain, at 225 distinct positions. No two variants have the same pair of sequences.

% TODO: review, needs manual review
\paragraph{Nanobody thermostability.} This is the NbThermo set \citep{ValdesTresanco2023}, taken from FLAb. The label is the melting temperature $T_m$ in $^\circ$C, from 39.0 to 98.2 $^\circ$C, for 672 single-domain antibodies with no light chain. The melting temperatures were measured with several methods and pooled into one label. The 672 rows hold 556 distinct sequences. Each of 116 sequences appears twice, and 106 of those pairs carry two different melting temperatures, differing by up to 8.8 $^\circ$C. These repeats are kept, so this set carries disagreement between measurements in its own labels before any noise is added.

% TODO: review, needs manual review
\paragraph{Humanness.} This is the list of therapeutic antibodies from IMGT/mAb-DB released with BioPhi \citep{Prihoda2022}, 564 antibodies with one sequence each. The sequence used is the first one listed for each antibody in the released sequence file. An antibody is labelled human (1) only when its species field is exactly ``Human'', which gives 198 human antibodies. The other 366, most of them humanised or chimeric, are labelled non-human (0). The 564 antibodies hold 506 distinct sequences, mostly because an antibody and its drug conjugate share a sequence. One sequence is shared by three antibodies with different labels, two humanised and one human. All three sit in one identity cluster, so they never fall on both sides of a split.

% TODO: review, needs manual review
\paragraph{Trastuzumab--HER2 binding.} This is the library of variants of the ten-residue CDR-H3 (the third complementarity-determining region of the heavy chain) of trastuzumab, screened for binding to HER2 by \citet{Chinery2024}. The released training, validation and test files are joined and re-split. Variants whose CDR-H3 is not exactly 10 residues would be dropped; there are none, so all 524,346 are kept: 172,149 binders (32.8\%) and 352,197 non-binders. From these, 5,000 variants were drawn once with seed 42. The draw is stratified on each combination of label and number of mutated CDR-H3 residues, so each combination keeps its share of the full set. The sample holds 1,639 binders (32.8\%). Each full heavy chain of 120 residues is rebuilt by inserting the variant CDR-H3 into the trastuzumab framework. The light chain is the fixed trastuzumab light chain, the same on every row.

\section{Nucleotide sequences}\label{app:datasets:nucleotides}

% TODO: review, needs manual review
Both nucleotide tasks come from the revised downstream tasks of the Nucleotide Transformer \citep{DallaTorre2025}, downloaded from Hugging Face (\texttt{InstaDeepAI/nucleotide\_transformer\_downstream\_tasks\_revised}). The download is not pinned to a revision; the files used are those of revision \texttt{851f994}, dated 30 June 2025, the latest at the time of writing. The released training and test files are merged. Each sequence's chromosome is read from its name, and the folds are grouped by chromosome. Nothing else is filtered, and no sequence contains a letter other than A, C, G or T. Both tasks cover chromosomes 1 to 21, X and Y; chromosome 22 is absent.

% TODO: review, needs manual review
\paragraph{H3K4me3.} Each sample is a 1,000-base sequence, labelled by whether it carries the H3K4me3 histone mark. The released files hold 17,468 training and 776 test sequences, 18,244 in all, with 9,122 in each class.

% TODO: review, needs manual review
\paragraph{TATA promoter.} Each sample is a 300-base sequence, labelled by whether it is a promoter with a TATA box. The released files hold 5,062 training and 212 test sequences, 5,274 in all, with 2,637 in each class.

\section{Cancer genomics (TCGA)}\label{app:datasets:tcga}

% TODO: review, needs manual review
The two cancer datasets use data from The Cancer Genome Atlas \citep{Weinstein2013}, downloaded through UCSC Xena \citep{Goldman2020}. Expression is read from each cohort's Illumina HiSeq RNA-seq matrix on the Xena TCGA hub (version 2017-10-13). It holds 20,530 genes as $\log_2(x+1)$ of the RSEM-normalised count \citep{Li2011RSEM} from the UNC RNA-seq pipeline. The same 20,530 genes appear in every cohort, and no gene filter is applied. Only primary tumours are kept (sample-type code 01 in the TCGA barcode). Barcodes are cut to the sample level, which leaves one sample per patient. Missing expression values are set to 0. There is no batch correction. Every fold holds out whole tissue source sites. A correction fitted on the training sites has no parameter for a site it never saw. A correction fitted on all sites has already seen the test samples.
% Result needed: confirm the full runs on the cluster read the per-cohort expression matrix described above, not the fallback matrix the loader uses when that download fails.

% TODO: review, needs manual review
\paragraph{Lung adenocarcinoma against lung squamous cell carcinoma.} The adenocarcinoma \citep{TCGA2014LUAD} and squamous cell carcinoma \citep{TCGA2012LUSC} cohorts are stacked, and the cohort is the label. The two matrices hold 576 and 553 samples. Dropping 110 normal-tissue and 2 recurrent-tumour samples leaves 515 adenocarcinomas and 502 squamous tumours, 1,017 in all, from 68 tissue source sites.

% TODO: review, needs manual review
\paragraph{Breast cancer PAM50 subtype.} The label is the PAM50 subtype call \citep{Parker2009} from the breast cancer clinical matrix on Xena \citep{TCGA2012BRCA}. Four classes are kept: luminal A, luminal B, basal-like and HER2-enriched. Tumours called normal-like, and tumours without a call, are dropped. Of 1,097 primary tumours with expression data, 821 remain: 421 luminal A, 192 luminal B, 141 basal-like and 67 HER2-enriched, from 19 tissue source sites.
\todo[inline, color=orange!40]{Appendix A: add per-dataset provenance, accession numbers, DOIs.}
\chapter{Additional files}\label{app:additional}\label{app:noiseinject}\label{app:results}
% ---- TAKEAWAYS ----
% - Every supporting table and figure too long for a main chapter is here, as one numbered Additional file, grouped by the chapter it supports.
% - The Chapter 3 files are the NoiseInject paper's Additional files; ten of the eleven come from the earlier version of the study and must be replaced.
% - The Chapter 5 files let a reader rebuild any hybrid: every model's settings, the models left out per dataset, the size of the search, and what a blueprint saves.
% - The design-experiment file shows where each search default came from. Those numbers are preliminary design runs, not results of the thesis.
% - The Chapter 4 and 6 files hold the full per-dataset tables that those chapters summarise. Most need the full runs.
% ---- END TAKEAWAYS ----
\newcounter{addfile}
\providecommand{\theHaddfile}{addfile.\arabic{addfile}}
\newcommand{\addfile}[1]{\refstepcounter{addfile}\par\medskip\noindent\textbf{Additional file~\theaddfile: #1}\par\nopagebreak\smallskip}

% TODO: review, needs manual review
This appendix holds the tables and figures that support the main chapters but are too long to include in them. They are grouped by the chapter they support and numbered in order as Additional files. The files for Chapter~\ref{chap:robustness} follow the Additional files of the paper that chapter is adapted from.

\section{Noise robustness across molecular representations and model architectures}\label{app:add:robustness}
% ---- NOTES: Noise robustness across molecular representations and model architectures ----
% Takeaway: the supporting tables and figures for Chapter 3, in the NoiseInject paper's grouping.
% - Out of date: every file but the first. They use the noise degradation slope, six noise strategies and Mol2Vec. Chapter 3 now uses the normalised area under the retention curve.
% Result needed: replace the ten out-of-date tables and figures once the Chapter 3 results are updated.
% - Exclusion threshold: clean $R^2 < 0.3$, per replicate on QM9 and per fold on the assay datasets.
% - Change: the excluded-configurations caption and the Chapter 3 Results paragraph (66 configurations) still use the earlier threshold of 0.6.
% Result needed: excluded configurations at clean $R^2 < 0.3$, one row per model, representation and noise condition.
% - The redundancy tables also list five representations Chapter 3 does not use (Morgan, randomised SMILES, SMILES, binary PDV, Mol2Vec).
% - The redundancy and ICC files belong to the correlation filter before the ANOVA; the paper replaced it with model families. Which one to keep is still open.
% - One new file is planned, cited by Chapter 3: base models against their probabilistic counterparts on LogD, Caco-2 efflux and hERG $K_i$, laid out like the QM9 table in Chapter 3.
% - Each entry: median change in AUC_norm, then clean $R^2$ in brackets, counterpart minus base, over 5 folds. One column says whether the pair shares one setting.
% Result needed: that table, on the three assay datasets.
% ---- END NOTES ----

\addfile{Hyperparameters for every model configuration, and the descriptors that make up PDV}
\label{app:noise:sec:hyperparameters_default}\label{app:noise:sec:hyperparameters_tuned}\label{app:noise:sec:pdv_descriptors}
% TODO: review, needs manual review
Table~A lists the default settings of every model configuration, Table~B the tuned settings that replace a default on each dataset, and Table~C the 200 descriptors that make up PDV.
% Thesis appendix item: hyperparameters_default
% Seeded on 2026-09-27 from qsar_qm_models/supplementary/items/hyperparameters_default.tex by thesis_appendix/seed_from_paper.py.
% From here on this copy is edited independently of the paper. Rules: thesis_appendix/README.md.
% The lines below, down to the first blank line, are the PAPER copy's header, kept for
% provenance. They describe the paper file, not this one: in particular a GENERATED
% notice there does not apply here, and this copy may be edited by hand.
% Supplementary item: hyperparameters_default
% GENERATED by scripts/generate_supp_table1.py --write. Do not edit;
% scripts/test_supp_table1.py fails on a hand edit. Rules: supplementary/README.md.

{\small
% TODO: REVIEW MODIFICATION (30 Sep) -- was p{5.0cm}p{5.0cm}p{4.6cm}, 12pt wider than the text; the first column is wider so
% the configuration names fit on one line.
\begin{longtable}{@{}p{5.9cm}p{4.6cm}p{3.6cm}@{}}
\caption{\textbf{Table A: default hyperparameters for every model configuration in the study.}
Every value is read from \texttt{models/model\_defaults.py} at spec version 1.7.0, spec hash \texttt{9249926c38dc}, which is the file both the QM9 pipeline and the assay pipeline load their settings from. Every QM9 results row carries the spec version and the spec hash as two of its columns, so a QM9 table can be traced back to the settings that produced it; the three assay datasets are written by a second pipeline whose results files do not carry those two columns. The configuration list is the model roster the job generator submits. One setting is used for every representation and for every dataset, except where Table B gives a tuned replacement. The support vector machine uses an RBF kernel on every representation, so no result in this study carries a kernel that changes with the representation. Values not held in the shared spec name their source file in the third column.}
\label{af1:hyperparameters}\\
\toprule
\textbf{Configuration} & \textbf{Hyperparameter} & \textbf{Default value} \\
\midrule
\endfirsthead
\toprule
\textbf{Configuration} & \textbf{Hyperparameter} & \textbf{Default value} \\
\midrule
\endhead
\multicolumn{3}{@{}l}{\textit{Shared by every neural network below}} \\
\midrule
Training, all networks & optimizer & adam \\
 & lr & 0.001 \\
 & weight\_decay & 0.0 \\
 & batch\_size & 32 \\
 & epochs & 100 \newline {\footnotesize\raggedright cap; early stopping usually reaches it first} \\
 & patience & 10 \newline {\footnotesize\raggedright epochs without a strictly better mean validation loss} \\
 & restore\_best\_weights & True \newline {\footnotesize\raggedright the reported fit is the best epoch, not the last} \\
 & val\_loss\_reduction & mean \\
 & improvement\_tolerance & 0.0 \\
 & mc\_passes & 100 \newline {\footnotesize\raggedright stochastic forward passes at inference} \\
 & standardise\_targets & True \newline {\footnotesize\raggedright fitted on training labels alone} \\
\midrule
\multicolumn{3}{@{}l}{\textit{Tree and boosting models}} \\
\midrule
RF \newline \footnotesize\texttt{rf} & n\_estimators & 100 \\
 & max\_depth & None \newline {\footnotesize\raggedright no limit} \\
 & min\_samples\_leaf & 5 \\
 & min\_samples\_split & 2 \\
 & max\_features & 0.3 \\
 & bootstrap & True \\
\addlinespace
RF300 \newline \footnotesize\texttt{rf300} & n\_estimators & 300 \\
 & max\_depth & None \newline {\footnotesize\raggedright no limit} \\
 & min\_samples\_leaf & 5 \\
 & min\_samples\_split & 2 \\
 & max\_features & 0.3 \\
 & bootstrap & True \\
\addlinespace
XGBoost \newline \footnotesize\texttt{xgboost} & n\_estimators & 100 \\
 & max\_depth & 6 \\
 & learning\_rate & 0.1 \\
 & subsample & 1.0 \\
 & colsample\_bytree & 1.0 \\
 & colsample\_bylevel & 1.0 \\
 & min\_child\_weight & 1 \\
 & gamma & 0.0 \\
 & reg\_alpha & 0.0 \\
 & reg\_lambda & 1.0 \\
\addlinespace
LightGBM \newline \footnotesize\texttt{lgb} & n\_estimators & 100 \\
 & learning\_rate & 0.1 \\
 & num\_leaves & 31 \\
 & max\_depth & $-1$ \newline {\footnotesize\raggedright no limit} \\
 & subsample & 1.0 \\
 & colsample\_bytree & 1.0 \\
 & min\_child\_samples & 20 \\
 & reg\_alpha & 0.0 \\
 & reg\_lambda & 0.0 \\
\addlinespace
NGBoost \newline \footnotesize\texttt{ngboost} & n\_estimators & 500 \\
 & learning\_rate & 0.01 \\
 & natural\_gradient & True \\
 & dist & Normal \newline {\footnotesize\raggedright ngboost.distns.Normal} \\
 & score & MLE \newline {\footnotesize\raggedright ngboost.scores.MLE, the log scoring rule} \\
 & minibatch\_frac & 1.0 \\
 & col\_sample & 1.0 \\
 & base\_max\_depth & 3 \newline {\footnotesize\raggedright depth of one boosting stage} \\
 & base\_criterion & friedman\_mse \\
 & early\_stopping\_rounds & 50 \newline {\footnotesize\raggedright stages without a better validation log likelihood} \\
 & use\_best\_iteration & True \newline {\footnotesize\raggedright predictions use the best stage, not every stage fitted} \\
 & ngboost\_ensemble\_seeds & 1 \newline {\footnotesize\raggedright 1 = one fit, no seed ensemble} \\
\addlinespace
QRF \newline \footnotesize\texttt{qrf} & n\_estimators & 300 \\
 & max\_depth & None \newline {\footnotesize\raggedright no limit} \\
 & min\_samples\_leaf & 5 \\
 & min\_samples\_split & 2 \\
 & max\_features & 0.3 \\
 & bootstrap & True \\
\midrule
\multicolumn{3}{@{}l}{\textit{Kernel models}} \\
\midrule
SVM \newline \footnotesize\texttt{svm} & C & 1.0 \\
 & gamma & scale \\
 & kernel & rbf \\
\addlinespace
GP \newline \footnotesize\texttt{gauche\_rbf} & kernel & rbf \newline {\footnotesize\raggedright set by this job, not by the shared spec} \\
 & likelihood\_noise & 0.001 \newline {\footnotesize\raggedright starting value; it is learned} \\
 & outputscale & 1.0 \\
 & apply\_outputscale & True \\
 & max\_train\_n & 5000 \newline {\footnotesize\raggedright molecules the process is fitted on; an exact Gaussian process is cubic in this. The test set is untouched and the number actually fitted is written into every results row} \\
 & fallback\_adam\_lr & 0.1 \newline {\footnotesize\raggedright used only when the standard fitter refuses the model} \\
 & fallback\_adam\_iters & 100 \newline {\footnotesize\raggedright used only when the standard fitter refuses the model} \\
 & single\_thread\_fit & False \newline {\footnotesize\raggedright a threading workaround, off since the environment check found one runtime} \\
 & init\_lengthscale\_from\_data & True \newline {\footnotesize\raggedright the RBF length scale starts at the median distance between training molecules, not at the library default} \\
 & lengthscale\_probe\_n & 500 \newline {\footnotesize\raggedright molecules sampled to estimate that median} \\
 & collapse\_fraction & 0.05 \newline {\footnotesize\raggedright a fit whose predictions vary by less than this fraction of the training label spread is recorded as collapsed rather than scored} \\
\addlinespace
GP-Tanimoto \newline \footnotesize\texttt{gauche} \newline \footnotesize run on ecfp4 only & kernel & tanimoto \newline {\footnotesize\raggedright set by this job, not by the shared spec} \\
 & likelihood\_noise & 0.001 \newline {\footnotesize\raggedright starting value; it is learned} \\
 & outputscale & 1.0 \\
 & apply\_outputscale & True \\
 & max\_train\_n & 5000 \newline {\footnotesize\raggedright molecules the process is fitted on; an exact Gaussian process is cubic in this. The test set is untouched and the number actually fitted is written into every results row} \\
 & fallback\_adam\_lr & 0.1 \newline {\footnotesize\raggedright used only when the standard fitter refuses the model} \\
 & fallback\_adam\_iters & 100 \newline {\footnotesize\raggedright used only when the standard fitter refuses the model} \\
 & single\_thread\_fit & False \newline {\footnotesize\raggedright a threading workaround, off since the environment check found one runtime} \\
 & init\_lengthscale\_from\_data & True \newline {\footnotesize\raggedright the RBF length scale starts at the median distance between training molecules, not at the library default} \\
 & lengthscale\_probe\_n & 500 \newline {\footnotesize\raggedright molecules sampled to estimate that median} \\
 & collapse\_fraction & 0.05 \newline {\footnotesize\raggedright a fit whose predictions vary by less than this fraction of the training label spread is recorded as collapsed rather than scored} \\
\addlinespace
GP-Hetero \newline \footnotesize\texttt{heteroscedastic\_gp} & kernel & rbf \newline {\footnotesize\raggedright set by this job, not by the shared spec} \\
 & likelihood\_noise & 0.001 \newline {\footnotesize\raggedright starting value; it is learned} \\
 & outputscale & 1.0 \\
 & apply\_outputscale & True \\
 & max\_train\_n & 5000 \newline {\footnotesize\raggedright molecules the process is fitted on; an exact Gaussian process is cubic in this. The test set is untouched and the number actually fitted is written into every results row} \\
 & fallback\_adam\_lr & 0.1 \newline {\footnotesize\raggedright used only when the standard fitter refuses the model} \\
 & fallback\_adam\_iters & 100 \newline {\footnotesize\raggedright used only when the standard fitter refuses the model} \\
 & single\_thread\_fit & False \newline {\footnotesize\raggedright a threading workaround, off since the environment check found one runtime} \\
 & init\_lengthscale\_from\_data & True \newline {\footnotesize\raggedright the RBF length scale starts at the median distance between training molecules, not at the library default} \\
 & lengthscale\_probe\_n & 500 \newline {\footnotesize\raggedright molecules sampled to estimate that median} \\
 & collapse\_fraction & 0.05 \newline {\footnotesize\raggedright a fit whose predictions vary by less than this fraction of the training label spread is recorded as collapsed rather than scored} \\
 & noise network hidden sizes & [64, 64] \newline {\footnotesize\raggedright models/models.py, not in the shared spec} \\
 & noise network output & Softplus, plus $10^{-4}$ \newline {\footnotesize\raggedright models/models.py, not in the shared spec} \\
 & Adam learning rate, process & 0.1 \newline {\footnotesize\raggedright models/models.py, not in the shared spec} \\
 & Adam learning rate, noise network & 0.001 \newline {\footnotesize\raggedright models/models.py, not in the shared spec} \\
 & noise network epochs & 100 \newline {\footnotesize\raggedright the shared epoch cap} \\
\midrule
\multicolumn{3}{@{}l}{\textit{Neural networks}} \\
\midrule
NN-$\alpha$ \newline \footnotesize\texttt{dnn} & hidden\_sizes & [128, 64] \\
 & dropout\_rate & 0.2 \\
 & activation & relu \\
\addlinespace
NN-$\beta$ \newline \footnotesize\texttt{mlp} & hidden\_size & 128 \\
 & num\_hidden\_layers & 2 \\
 & dropout\_rate & 0.2 \\
 & activation & relu \\
\addlinespace
BNN-$\alpha$ \newline \footnotesize\texttt{dnn\_bnn\_full} & hidden\_sizes & [128, 64] \\
 & dropout\_rate & 0.2 \\
 & activation & relu \\
 & bnn\_prior\_mu & 0.0 \\
 & bnn\_prior\_sigma & 0.1 \\
 & bnn\_kl\_weight & elbo \newline {\footnotesize\raggedright 1 / number of training molecules} \\
\addlinespace
BNN-$\beta$ \newline \footnotesize\texttt{mlp\_bnn\_full} & hidden\_size & 128 \\
 & num\_hidden\_layers & 2 \\
 & dropout\_rate & 0.2 \\
 & activation & relu \\
 & bnn\_prior\_mu & 0.0 \\
 & bnn\_prior\_sigma & 0.1 \\
 & bnn\_kl\_weight & elbo \newline {\footnotesize\raggedright 1 / number of training molecules} \\
\addlinespace
VBLL-$\alpha$ \newline \footnotesize\texttt{dnn\_bnn\_full\_variational} & hidden\_sizes & [128, 64] \\
 & dropout\_rate & 0.2 \\
 & activation & relu \\
 & vbll\_prior\_mu & 0.0 \\
 & vbll\_prior\_sigma & 1.0 \\
 & vbll\_init\_log\_sigma & $-3.0$ \\
 & vbll\_init\_log\_noise\_var & $-2.0$ \\
\addlinespace
VBLL-$\beta$ \newline \footnotesize\texttt{mlp\_bnn\_full\_variational} & hidden\_size & 128 \\
 & num\_hidden\_layers & 2 \\
 & dropout\_rate & 0.2 \\
 & activation & relu \\
 & vbll\_prior\_mu & 0.0 \\
 & vbll\_prior\_sigma & 1.0 \\
 & vbll\_init\_log\_sigma & $-3.0$ \\
 & vbll\_init\_log\_noise\_var & $-2.0$ \\
\addlinespace
VBLL-$\alpha$ (noise head) \newline \footnotesize\texttt{dnn\_bnn\_full\_variational\_hetero} & hidden\_sizes & [128, 64] \\
 & dropout\_rate & 0.2 \\
 & activation & relu \\
 & vbll\_prior\_mu & 0.0 \\
 & vbll\_prior\_sigma & 1.0 \\
 & vbll\_init\_log\_sigma & $-3.0$ \\
 & vbll\_init\_log\_noise\_var & $-2.0$ \\
 & observation noise & predicted per molecule \newline {\footnotesize\raggedright a variational noise head on the same features, started at the single-number value above} \\
\addlinespace
VBLL-$\beta$ (noise head) \newline \footnotesize\texttt{mlp\_bnn\_full\_variational\_hetero} & hidden\_size & 128 \\
 & num\_hidden\_layers & 2 \\
 & dropout\_rate & 0.2 \\
 & activation & relu \\
 & vbll\_prior\_mu & 0.0 \\
 & vbll\_prior\_sigma & 1.0 \\
 & vbll\_init\_log\_sigma & $-3.0$ \\
 & vbll\_init\_log\_noise\_var & $-2.0$ \\
 & observation noise & predicted per molecule \newline {\footnotesize\raggedright a variational noise head on the same features, started at the single-number value above} \\
\addlinespace
BNN-$\alpha$ (variance head) \newline \footnotesize\texttt{dnn\_bnn\_full\_mve} & hidden\_sizes & [128, 64] \\
 & dropout\_rate & 0.2 \\
 & activation & relu \\
 & bnn\_prior\_mu & 0.0 \\
 & bnn\_prior\_sigma & 0.1 \\
 & bnn\_kl\_weight & elbo \newline {\footnotesize\raggedright 1 / number of training molecules} \\
 & output head & value and log variance \newline {\footnotesize\raggedright two outputs} \\
 & loss & Gaussian negative log likelihood \\
 & log variance clamp & $[-10, 10]$ \newline {\footnotesize\raggedright scripts/loss\_functions.py, not in the shared spec} \\
\addlinespace
BNN-$\beta$ (variance head) \newline \footnotesize\texttt{mlp\_bnn\_full\_mve} & hidden\_size & 128 \\
 & num\_hidden\_layers & 2 \\
 & dropout\_rate & 0.2 \\
 & activation & relu \\
 & bnn\_prior\_mu & 0.0 \\
 & bnn\_prior\_sigma & 0.1 \\
 & bnn\_kl\_weight & elbo \newline {\footnotesize\raggedright 1 / number of training molecules} \\
 & output head & value and log variance \newline {\footnotesize\raggedright two outputs} \\
 & loss & Gaussian negative log likelihood \\
 & log variance clamp & $[-10, 10]$ \newline {\footnotesize\raggedright scripts/loss\_functions.py, not in the shared spec} \\
\bottomrule
\end{longtable}
}

% Thesis appendix item: hyperparameters_tuned
% Seeded on 2026-09-27 from qsar_qm_models/supplementary/items/hyperparameters_tuned.tex by thesis_appendix/seed_from_paper.py.
% From here on this copy is edited independently of the paper. Rules: thesis_appendix/README.md.
% The lines below, down to the first blank line, are the PAPER copy's header, kept for
% provenance. They describe the paper file, not this one: in particular a GENERATED
% notice there does not apply here, and this copy may be edited by hand.
% Supplementary item: hyperparameters_tuned
% GENERATED by scripts/generate_supp_table1.py --write. Do not edit;
% scripts/test_supp_table1.py fails on a hand edit. Rules: supplementary/README.md.

{\small
\begin{longtable}{@{}llll@{}}
\caption{\textbf{The tuned settings that replace a default, and the four datasets they apply to.}
A setting is chosen once per model per dataset, not per representation, and is then used for all six representations: a two-way comparison across model and representation cannot be read if each cell carries its own hyperparameters. Only the Bayesian and variational networks are tuned; every other configuration in Table~\ref{af1:hyperparameters} trains at the default. A model with no row for a dataset kept the default on that dataset. Learning rates are printed to three significant figures; the source files hold the full value. Source files: \texttt{results/master\_tuned\_hyperparameters.json} for QM9 and \texttt{results/master\_tuned\_hyperparameters\_lab.json} for the three assay datasets.}
\label{af1:tuned}\\
\toprule
\textbf{Dataset} & \textbf{Configuration} & \textbf{Hyperparameter} & \textbf{Tuned value} \\
\midrule
\endfirsthead
\toprule
\textbf{Dataset} & \textbf{Configuration} & \textbf{Hyperparameter} & \textbf{Tuned value} \\
\midrule
\endhead
QM9 & NN-$\alpha$ & activation & tanh \\
 &  & hidden\_size1 & 64 \\
 &  & hidden\_size2 & 64 \\
\addlinespace
QM9 & BNN-$\alpha$ & activation & tanh \\
 &  & hidden\_size1 & 64 \\
 &  & hidden\_size2 & 64 \\
\addlinespace
QM9 & VBLL-$\alpha$ & activation & tanh \\
 &  & hidden\_size1 & 64 \\
 &  & hidden\_size2 & 64 \\
\addlinespace
QM9 & NN-$\beta$ & dropout\_rate & 0.401 \\
 &  & hidden\_size & 64 \\
 &  & lr & 0.00587 \\
 &  & num\_hidden\_layers & 3 \\
\addlinespace
QM9 & BNN-$\beta$ & dropout\_rate & 0.401 \\
 &  & hidden\_size & 64 \\
 &  & lr & 0.00587 \\
 &  & num\_hidden\_layers & 3 \\
\addlinespace
QM9 & VBLL-$\beta$ & dropout\_rate & 0.357 \\
 &  & hidden\_size & 64 \\
 &  & lr & 0.00119 \\
 &  & num\_hidden\_layers & 1 \\
\addlinespace
Caco-2 & NN-$\alpha$ & activation & tanh \\
 &  & hidden\_size1 & 64 \\
 &  & hidden\_size2 & 64 \\
\addlinespace
Caco-2 & BNN-$\alpha$ & activation & tanh \\
 &  & hidden\_size1 & 64 \\
 &  & hidden\_size2 & 64 \\
\addlinespace
Caco-2 & VBLL-$\alpha$ & activation & relu \\
 &  & hidden\_size1 & 64 \\
 &  & hidden\_size2 & 32 \\
\addlinespace
Caco-2 & NN-$\beta$ & dropout\_rate & 0.401 \\
 &  & hidden\_size & 64 \\
 &  & lr & 0.00587 \\
 &  & num\_hidden\_layers & 3 \\
\addlinespace
Caco-2 & BNN-$\beta$ & dropout\_rate & 0.401 \\
 &  & hidden\_size & 64 \\
 &  & lr & 0.00587 \\
 &  & num\_hidden\_layers & 3 \\
\addlinespace
hERG K$_i$ & NN-$\alpha$ & activation & tanh \\
 &  & hidden\_size1 & 64 \\
 &  & hidden\_size2 & 64 \\
\addlinespace
hERG K$_i$ & BNN-$\alpha$ & activation & tanh \\
 &  & hidden\_size1 & 64 \\
 &  & hidden\_size2 & 64 \\
\addlinespace
hERG K$_i$ & VBLL-$\alpha$ & activation & tanh \\
 &  & hidden\_size1 & 64 \\
 &  & hidden\_size2 & 64 \\
\addlinespace
hERG K$_i$ & NN-$\beta$ & dropout\_rate & 0.401 \\
 &  & hidden\_size & 64 \\
 &  & lr & 0.00587 \\
 &  & num\_hidden\_layers & 3 \\
\addlinespace
hERG K$_i$ & BNN-$\beta$ & dropout\_rate & 0.401 \\
 &  & hidden\_size & 64 \\
 &  & lr & 0.00587 \\
 &  & num\_hidden\_layers & 3 \\
\addlinespace
LogD & NN-$\alpha$ & activation & tanh \\
 &  & hidden\_size1 & 64 \\
 &  & hidden\_size2 & 64 \\
\addlinespace
LogD & BNN-$\alpha$ & activation & tanh \\
 &  & hidden\_size1 & 64 \\
 &  & hidden\_size2 & 64 \\
\addlinespace
LogD & VBLL-$\alpha$ & activation & tanh \\
 &  & hidden\_size1 & 256 \\
 &  & hidden\_size2 & 32 \\
\addlinespace
LogD & NN-$\beta$ & dropout\_rate & 0.401 \\
 &  & hidden\_size & 64 \\
 &  & lr & 0.00587 \\
 &  & num\_hidden\_layers & 3 \\
\addlinespace
LogD & BNN-$\beta$ & dropout\_rate & 0.401 \\
 &  & hidden\_size & 64 \\
 &  & lr & 0.00587 \\
 &  & num\_hidden\_layers & 3 \\
\addlinespace
LogD & VBLL-$\beta$ & dropout\_rate & 0.079 \\
 &  & hidden\_size & 128 \\
 &  & lr & 0.00011 \\
 &  & num\_hidden\_layers & 2 \\
\bottomrule
\end{longtable}
}

% Thesis appendix item: pdv_descriptors
% Seeded on 2026-09-27 from qsar_qm_models/supplementary/items/pdv_descriptors.tex by thesis_appendix/seed_from_paper.py.
% From here on this copy is edited independently of the paper. Rules: thesis_appendix/README.md.
% The lines below, down to the first blank line, are the PAPER copy's header, kept for
% provenance. They describe the paper file, not this one: in particular a GENERATED
% notice there does not apply here, and this copy may be edited by hand.
% Supplementary item: pdv_descriptors
% GENERATED by scripts/generate_supp_table1.py --write. Do not edit;
% scripts/test_supp_table1.py fails on a hand edit. Rules: supplementary/README.md.

% [2026-10-10 fix] kept \scriptsize (30 Sep change; \small was 70pt wider than the text); additional_files_body.tex also sets \small to \scriptsize around this \input
{\scriptsize
\begin{longtable}{@{}llll@{}}
\caption{\textbf{The 200 descriptors that make up the PDV representation.}
One entry is one feature of the vector, in the order the pipeline computes them. The list is fixed in the code rather than taken from RDKit's current default set, which grows between releases, so the same 200 features are built on every dataset and in every run. They were computed with RDKit 2022.09.5. Descriptors that return a non-finite value on a molecule have that value replaced by zero before the standardization constants are computed. Source: \texttt{DEFAULT\_DESCRIPTOR\_LIST} in \texttt{scripts/process\_and\_train.py}.}
\label{af1:pdv}\\
\toprule
\endfirsthead
\toprule
\endhead
BalabanJ & NOCount & SlogP\_VSA7 & fr\_bicyclic \\
BertzCT & NumAliphaticCarbocycles & SlogP\_VSA8 & fr\_diazo \\
Chi0 & NumAliphaticHeterocycles & SlogP\_VSA9 & fr\_dihydropyridine \\
Chi0n & NumAliphaticRings & TPSA & fr\_epoxide \\
Chi0v & NumAromaticCarbocycles & VSA\_EState1 & fr\_ester \\
Chi1 & NumAromaticHeterocycles & VSA\_EState10 & fr\_ether \\
Chi1n & NumAromaticRings & VSA\_EState2 & fr\_furan \\
Chi1v & NumHAcceptors & VSA\_EState3 & fr\_guanido \\
Chi2n & NumHDonors & VSA\_EState4 & fr\_halogen \\
Chi2v & NumHeteroatoms & VSA\_EState5 & fr\_hdrzine \\
Chi3n & NumRadicalElectrons & VSA\_EState6 & fr\_hdrzone \\
Chi3v & NumRotatableBonds & VSA\_EState7 & fr\_imidazole \\
Chi4n & NumSaturatedCarbocycles & VSA\_EState8 & fr\_imide \\
Chi4v & NumSaturatedHeterocycles & VSA\_EState9 & fr\_isocyan \\
EState\_VSA1 & NumSaturatedRings & fr\_Al\_COO & fr\_isothiocyan \\
EState\_VSA10 & NumValenceElectrons & fr\_Al\_OH & fr\_ketone \\
EState\_VSA11 & PEOE\_VSA1 & fr\_Al\_OH\_noTert & fr\_ketone\_Topliss \\
EState\_VSA2 & PEOE\_VSA10 & fr\_ArN & fr\_lactam \\
EState\_VSA3 & PEOE\_VSA11 & fr\_Ar\_COO & fr\_lactone \\
EState\_VSA4 & PEOE\_VSA12 & fr\_Ar\_N & fr\_methoxy \\
EState\_VSA5 & PEOE\_VSA13 & fr\_Ar\_NH & fr\_morpholine \\
EState\_VSA6 & PEOE\_VSA14 & fr\_Ar\_OH & fr\_nitrile \\
EState\_VSA7 & PEOE\_VSA2 & fr\_COO & fr\_nitro \\
EState\_VSA8 & PEOE\_VSA3 & fr\_COO2 & fr\_nitro\_arom \\
EState\_VSA9 & PEOE\_VSA4 & fr\_C\_O & fr\_nitro\_arom\_nonortho \\
ExactMolWt & PEOE\_VSA5 & fr\_C\_O\_noCOO & fr\_nitroso \\
FpDensityMorgan1 & PEOE\_VSA6 & fr\_C\_S & fr\_oxazole \\
FpDensityMorgan2 & PEOE\_VSA7 & fr\_HOCCN & fr\_oxime \\
FpDensityMorgan3 & PEOE\_VSA8 & fr\_Imine & fr\_para\_hydroxylation \\
FractionCSP3 & PEOE\_VSA9 & fr\_NH0 & fr\_phenol \\
HallKierAlpha & RingCount & fr\_NH1 & fr\_phenol\_noOrthoHbond \\
HeavyAtomCount & SMR\_VSA1 & fr\_NH2 & fr\_phos\_acid \\
HeavyAtomMolWt & SMR\_VSA10 & fr\_N\_O & fr\_phos\_ester \\
Ipc & SMR\_VSA2 & fr\_Ndealkylation1 & fr\_piperdine \\
Kappa1 & SMR\_VSA3 & fr\_Ndealkylation2 & fr\_piperzine \\
Kappa2 & SMR\_VSA4 & fr\_Nhpyrrole & fr\_priamide \\
Kappa3 & SMR\_VSA5 & fr\_SH & fr\_prisulfonamd \\
LabuteASA & SMR\_VSA6 & fr\_aldehyde & fr\_pyridine \\
MaxAbsEStateIndex & SMR\_VSA7 & fr\_alkyl\_carbamate & fr\_quatN \\
MaxAbsPartialCharge & SMR\_VSA8 & fr\_alkyl\_halide & fr\_sulfide \\
MaxEStateIndex & SMR\_VSA9 & fr\_allylic\_oxid & fr\_sulfonamd \\
MaxPartialCharge & SlogP\_VSA1 & fr\_amide & fr\_sulfone \\
MinAbsEStateIndex & SlogP\_VSA10 & fr\_amidine & fr\_term\_acetylene \\
MinAbsPartialCharge & SlogP\_VSA11 & fr\_aniline & fr\_tetrazole \\
MinEStateIndex & SlogP\_VSA12 & fr\_aryl\_methyl & fr\_thiazole \\
MinPartialCharge & SlogP\_VSA2 & fr\_azide & fr\_thiocyan \\
MolLogP & SlogP\_VSA3 & fr\_azo & fr\_thiophene \\
MolMR & SlogP\_VSA4 & fr\_barbitur & fr\_unbrch\_alkane \\
MolWt & SlogP\_VSA5 & fr\_benzene & fr\_urea \\
NHOHCount & SlogP\_VSA6 & fr\_benzodiazepine & qed \\
\bottomrule
\end{longtable}
}


\addfile{Configurations excluded from the robustness analysis}
\label{app:noise:sec:excluded_configurations}
% TODO: review, needs manual review
For each model and representation on the QM9 HOMO--LUMO gap, the number of noise strategies for which the configuration was excluded because its clean-label $R^2$ was too low.
% Thesis appendix item: excluded_configurations
% Seeded on 2026-09-27 from qsar_qm_models/supplementary/items/excluded_configurations.tex by thesis_appendix/seed_from_paper.py.
% From here on this copy is edited independently of the paper. Rules: thesis_appendix/README.md.
% The lines below, down to the first blank line, are the PAPER copy's header, kept for
% provenance. They describe the paper file, not this one: in particular a GENERATED
% notice there does not apply here, and this copy may be edited by hand.
% Supplementary item: excluded_configurations
% Split out of additional_files.tex on 2026-09-27, where it was Additional file 5 (Excluded configurations (was Supp Table S2)).
% Content unchanged apart from: the heading line removed (the build writes it), the old
% 'Table S*/Figure S*' caption prefix replaced by \SuppFile{}, and the label renamed to supp:excluded_configurations.
% Rules: supplementary/README.md.

\begin{table}[H]
\centering
\caption{\textbf{Configurations excluded from robustness analysis due to baseline R$^2 \leq 0.6$ on the QM9 HOMO--LUMO gap, summarized by model and representation.}
Each cell shows the number of noise strategies (out of 6) for which the configuration was excluded.}
\label{app:noise:excluded_configurations}
\small
\begin{tabular}{lcccc}
\toprule
\textbf{Model} & \textbf{Mol2Vec} & \textbf{MHGGNN} & \textbf{SMILES} & \textbf{Rand.~SMILES} \\
\midrule
BNN-$\alpha$ & 6 & 6 & --- & --- \\
BNN-$\beta$ & 6 & 6 & --- & --- \\
VBLL-$\alpha$ & 6 & 6 & --- & --- \\
VBLL-$\beta$ & 6 & 6 & --- & --- \\
RF & --- & --- & --- & 4 \\
QRF & --- & --- & --- & 4 \\
NGBoost & --- & --- & 6 & 4 \\
\bottomrule
\end{tabular}
\end{table}


\addfile{Spearman correlations between the robustness profiles of the models}
\label{app:noise:sec:model_redundancy}
% TODO: review, needs manual review
Pairs of models whose noise degradation slope profiles on the QM9 HOMO--LUMO gap correlate at $\rho \geq 0.95$, and which model of each pair was left out of the ANOVA.
% Thesis appendix item: model_redundancy
% Seeded on 2026-09-27 from qsar_qm_models/supplementary/items/model_redundancy.tex by thesis_appendix/seed_from_paper.py.
% From here on this copy is edited independently of the paper. Rules: thesis_appendix/README.md.
% The lines below, down to the first blank line, are the PAPER copy's header, kept for
% provenance. They describe the paper file, not this one: in particular a GENERATED
% notice there does not apply here, and this copy may be edited by hand.
% Supplementary item: model_redundancy
% Split out of additional_files.tex on 2026-09-27, where it was Additional file 2 (Pairwise model NDS correlations (was Supp Table S3)).
% Content unchanged apart from: the heading line removed (the build writes it), the old
% 'Table S*/Figure S*' caption prefix replaced by \SuppFile{}, and the label renamed to supp:model_redundancy.
% Rules: supplementary/README.md.

{\scriptsize
\begin{longtable}{llccc}
\caption{\textbf{Pairwise Spearman rank correlations between model noise degradation slope (NDS) profiles on the QM9 HOMO--LUMO gap ($\rho \geq 0.95$).}
``Model A'' and ``Model B'' label each pair, both treated symmetrically. Checkmarks indicate pairs where one model was excluded from the ANOVA due to near-redundancy. Gauche (RBF) pairs omitted (all $\rho = 1.0$, $n = 6$, PDV-only artifact).}
\label{app:noise:model_redundancy}\\
\toprule
\textbf{Model A} & \textbf{Model B} & \textbf{$n$} & \textbf{$\rho$} & \textbf{Excluded} \\
\midrule
\endfirsthead
\toprule
\textbf{Model A} & \textbf{Model B} & \textbf{$n$} & \textbf{$\rho$} & \textbf{Excluded} \\
\midrule
\endhead
VBLL-$\beta$ & NGBoost & 18 & 0.998 &  \\
LightGBM & XGBoost & 46 & 0.998 &  \\
Gauche & VBLL-$\beta$ & 18 & 0.996 & \checkmark \\
QRF & RF & 40 & 0.995 & \checkmark \\
VBLL-$\alpha$ & LightGBM & 24 & 0.992 &  \\
RF & XGBoost & 46 & 0.989 &  \\
NGBoost & XGBoost & 40 & 0.988 &  \\
NGBoost & RF & 40 & 0.987 &  \\
LightGBM & RF & 46 & 0.986 &  \\
VBLL-$\beta$ & SVM & 24 & 0.986 &  \\
QRF & XGBoost & 40 & 0.986 & \checkmark \\
NGBoost & QRF & 34 & 0.986 & \checkmark \\
VBLL-$\alpha$ & NGBoost & 18 & 0.986 &  \\
VBLL-$\alpha$ & XGBoost & 24 & 0.984 &  \\
LightGBM & NGBoost & 40 & 0.984 &  \\
LightGBM & QRF & 40 & 0.984 & \checkmark \\
VBLL-$\alpha$ & VBLL-$\beta$ & 24 & 0.983 &  \\
NN-$\alpha$ & NN-$\beta$ & 52 & 0.983 &  \\
LightGBM & SVM & 46 & 0.982 &  \\
BNN-$\beta$ & SVM & 34 & 0.981 &  \\
VBLL-$\alpha$ & SVM & 24 & 0.979 &  \\
BNN-$\alpha$ & VBLL-$\beta$ & 24 & 0.979 &  \\
VBLL-$\alpha$ & QRF & 24 & 0.978 & \checkmark \\
SVM & XGBoost & 46 & 0.978 &  \\
BNN-$\alpha$ & BNN-$\beta$ & 36 & 0.977 &  \\
BNN-$\alpha$ & SVM & 34 & 0.976 &  \\
VBLL-$\alpha$ & Gauche & 18 & 0.975 & \checkmark \\
VBLL-$\alpha$ & RF & 24 & 0.975 &  \\
QRF & SVM & 40 & 0.974 & \checkmark \\
BNN-$\alpha$ & VBLL-$\alpha$ & 24 & 0.974 &  \\
LightGBM & VBLL-$\beta$ & 24 & 0.973 &  \\
VBLL-$\alpha$ & BNN-$\beta$ & 24 & 0.972 &  \\
Gauche & SVM & 40 & 0.972 & \checkmark \\
Gauche & XGBoost & 44 & 0.970 & \checkmark \\
NGBoost & SVM & 40 & 0.970 &  \\
NN-$\alpha$ & BNN-$\beta$ & 36 & 0.970 &  \\
Gauche & LightGBM & 42 & 0.969 & \checkmark \\
RF & SVM & 46 & 0.969 &  \\
BNN-$\beta$ & VBLL-$\beta$ & 24 & 0.968 &  \\
VBLL-$\beta$ & QRF & 24 & 0.966 & \checkmark \\
NN-$\alpha$ & BNN-$\alpha$ & 36 & 0.965 &  \\
VBLL-$\beta$ & XGBoost & 24 & 0.965 &  \\
NN-$\beta$ & SVM & 46 & 0.962 &  \\
NN-$\alpha$ & SVM & 46 & 0.960 &  \\
Gauche & NGBoost & 34 & 0.959 & \checkmark \\
LightGBM & BNN-$\beta$ & 36 & 0.959 &  \\
BNN-$\alpha$ & NN-$\beta$ & 36 & 0.958 &  \\
BNN-$\alpha$ & QRF & 28 & 0.956 & \checkmark \\
Gauche & QRF & 34 & 0.956 & \checkmark \\
VBLL-$\beta$ & RF & 24 & 0.955 &  \\
Gauche & RF & 40 & 0.954 & \checkmark \\
BNN-$\beta$ & XGBoost & 34 & 0.953 &  \\
BNN-$\beta$ & QRF & 28 & 0.952 & \checkmark \\
NN-$\alpha$ & VBLL-$\beta$ & 24 & 0.952 &  \\
NN-$\beta$ & BNN-$\beta$ & 36 & 0.952 &  \\
BNN-$\alpha$ & LightGBM & 36 & 0.951 &  \\
NN-$\beta$ & VBLL-$\beta$ & 24 & 0.950 &  \\
BNN-$\alpha$ & Gauche & 30 & 0.950 & \checkmark \\
\bottomrule
\end{longtable}}


\addfile{Spearman correlations between the robustness profiles of the representations}
\label{app:noise:sec:representation_redundancy}
% TODO: review, needs manual review
Pairwise correlations between the noise degradation slope profiles of the representations on the QM9 HOMO--LUMO gap, and which representation of each pair was left out of the ANOVA.
% Thesis appendix item: representation_redundancy
% Seeded on 2026-09-27 from qsar_qm_models/supplementary/items/representation_redundancy.tex by thesis_appendix/seed_from_paper.py.
% From here on this copy is edited independently of the paper. Rules: thesis_appendix/README.md.
% The lines below, down to the first blank line, are the PAPER copy's header, kept for
% provenance. They describe the paper file, not this one: in particular a GENERATED
% notice there does not apply here, and this copy may be edited by hand.
% Supplementary item: representation_redundancy
% Split out of additional_files.tex on 2026-09-27, where it was Additional file 3 (Pairwise representation NDS correlations (was Supp Table S4)).
% Content unchanged apart from: the heading line removed (the build writes it), the old
% 'Table S*/Figure S*' caption prefix replaced by \SuppFile{}, and the label renamed to supp:representation_redundancy.
% Rules: supplementary/README.md.

{\scriptsize
\begin{longtable}{llccc}
\caption{\textbf{Pairwise Spearman rank correlations between representation noise degradation slope (NDS) profiles on the QM9 HOMO--LUMO gap.}
``Rep A'' and ``Rep B'' label each pair of molecular representations, both treated symmetrically. Checkmarks indicate pairs where one representation was excluded from the ANOVA.}
\label{app:noise:representation_redundancy}\\
\toprule
\textbf{Rep.\ A} & \textbf{Rep.\ B} & \textbf{$n$} & \textbf{$\rho$} & \textbf{Excluded} \\
\midrule
\endfirsthead
\toprule
\textbf{Rep.\ A} & \textbf{Rep.\ B} & \textbf{$n$} & \textbf{$\rho$} & \textbf{Excluded} \\
\midrule
\endhead
ECFP4 & Morgan & 60 & 0.995 & \checkmark \\
MHGGNN & Morgan & 48 & 0.994 & \checkmark \\
Rand.~SMILES & SMILES & 16 & 0.994 & \checkmark \\
Rand.~SMILES & SNS & 16 & 0.994 & \checkmark \\
ECFP4 & MHGGNN & 54 & 0.993 &  \\
ECFP4 & SNS & 56 & 0.992 & \checkmark \\
Morgan & SNS & 52 & 0.992 & \checkmark \\
MHGGNN & SNS & 44 & 0.991 & \checkmark \\
PDV & ECFP4 & 72 & 0.990 &  \\
PDV & MHGGNN & 48 & 0.989 &  \\
SMILES & SNS & 52 & 0.988 & \checkmark \\
Morgan & Rand.~SMILES & 16 & 0.988 & \checkmark \\
Morgan & SMILES & 54 & 0.988 & \checkmark \\
PDV & Morgan & 54 & 0.986 & \checkmark \\
PDV & Rand.~SMILES & 12 & 0.986 & \checkmark \\
PDV & SNS & 50 & 0.984 & \checkmark \\
ECFP4 & SMILES & 72 & 0.984 &  \\
Mol2Vec & Morgan & 48 & 0.982 & \checkmark \\
MHGGNN & Mol2Vec & 54 & 0.982 &  \\
ECFP4 & Mol2Vec & 54 & 0.982 &  \\
PDV & PDV (binary) & 72 & 0.980 & \checkmark \\
MHGGNN & SMILES & 48 & 0.980 &  \\
PDV & Mol2Vec & 48 & 0.979 &  \\
ECFP4 & Rand.~SMILES & 16 & 0.979 & \checkmark \\
PDV & SMILES & 66 & 0.978 &  \\
MHGGNN & PDV (binary) & 54 & 0.972 & \checkmark \\
Mol2Vec & SNS & 44 & 0.969 & \checkmark \\
Mol2Vec & SMILES & 48 & 0.966 &  \\
ECFP4 & PDV (binary) & 78 & 0.965 & \checkmark \\
Morgan & PDV (binary) & 60 & 0.962 & \checkmark \\
PDV (binary) & SNS & 56 & 0.961 & \checkmark \\
Mol2Vec & PDV (binary) & 54 & 0.956 & \checkmark \\
MHGGNN & Rand.~SMILES & 16 & 0.956 & \checkmark \\
PDV (binary) & SMILES & 72 & 0.941 & \checkmark \\
PDV (binary) & Rand.~SMILES & 16 & 0.929 & \checkmark \\
Mol2Vec & Rand.~SMILES & 16 & 0.891 & \checkmark \\
\bottomrule
\end{longtable}}


\addfile{Intraclass correlation coefficients for model pairs}
\label{app:noise:sec:model_icc}
% TODO: review, needs manual review
ICC(1,1) of the noise degradation slope profiles for the model pairs with ICC $\geq 0.5$, on the QM9 HOMO--LUMO gap.
% Thesis appendix item: model_icc
% Seeded on 2026-09-27 from qsar_qm_models/supplementary/items/model_icc.tex by thesis_appendix/seed_from_paper.py.
% From here on this copy is edited independently of the paper. Rules: thesis_appendix/README.md.
% The lines below, down to the first blank line, are the PAPER copy's header, kept for
% provenance. They describe the paper file, not this one: in particular a GENERATED
% notice there does not apply here, and this copy may be edited by hand.
% Supplementary item: model_icc
% Split out of additional_files.tex on 2026-09-27, where it was Additional file 4 (ICC(1,1) (was Supp Table S5)).
% Content unchanged apart from: the heading line removed (the build writes it), the old
% 'Table S*/Figure S*' caption prefix replaced by \SuppFile{}, and the label renamed to supp:model_icc.
% Rules: supplementary/README.md.

\begin{table}[H]
\centering
\caption{\textbf{Intraclass correlation coefficients (ICC(1,1)) for model pairs with ICC $\geq 0.5$, assessing within-family consistency of noise degradation slope (NDS) profiles on the QM9 HOMO--LUMO gap.}
ICC(1,1) measures the proportion of total variance attributable to between-configuration differences; values near 1.0 indicate that two models rank configurations almost identically. Lower (more negative) NDS indicates faster degradation under noise. The Family column groups each pair: ``Tree'' (both tree-based), ``$\alpha$/$\beta$'' (architecture-pair partners such as BNN-$\alpha$/BNN-$\beta$), ``$\alpha$'' (both share the $\alpha$ neural-network backbone), or ``Cross'' (cross-family pair).}
\label{app:noise:model_icc}
\small
\begin{tabular}{llcccc}
\toprule
\textbf{Model A} & \textbf{Model B} & \textbf{$n$} & \textbf{ICC(1,1)} & \textbf{Mean $|\Delta$NDS$|$} & \textbf{Family} \\
\midrule
LightGBM & XGBoost & 8 & 0.944 & 0.0060 & Tree \\
BNN-$\alpha$ & BNN-$\beta$ & 6 & 0.846 & 0.0101 & $\alpha$/$\beta$ \\
QRF & RF & 7 & 0.810 & 0.0216 & Tree \\
VBLL-$\alpha$ & VBLL-$\beta$ & 4 & 0.782 & 0.0108 & $\alpha$/$\beta$ \\
QRF & XGBoost & 7 & 0.754 & 0.0157 & Tree \\
LightGBM & QRF & 7 & 0.743 & 0.0133 & Tree \\
VBLL-$\alpha$ & LightGBM & 4 & 0.698 & 0.0115 & Cross \\
VBLL-$\beta$ & SVM & 4 & 0.689 & 0.0133 & Cross \\
RF & XGBoost & 8 & 0.678 & 0.0212 & Tree \\
VBLL-$\alpha$ & XGBoost & 4 & 0.656 & 0.0133 & Cross \\
BNN-$\alpha$ & VBLL-$\alpha$ & 4 & 0.604 & 0.0134 & $\alpha$ \\
LightGBM & RF & 8 & 0.586 & 0.0244 & Tree \\
LightGBM & VBLL-$\beta$ & 4 & 0.577 & 0.0174 & Cross \\
LightGBM & SVM & 8 & 0.561 & 0.0117 & Cross \\
BNN-$\alpha$ & VBLL-$\beta$ & 4 & 0.552 & 0.0123 & $\alpha$/$\beta$ \\
VBLL-$\beta$ & XGBoost & 4 & 0.544 & 0.0195 & Cross \\
VBLL-$\alpha$ & QRF & 4 & 0.540 & 0.0229 & Cross \\
VBLL-$\alpha$ & SVM & 4 & 0.522 & 0.0124 & Cross \\
Gauche & VBLL-$\beta$ & 3 & 0.519 & 0.0113 & Cross \\
\bottomrule
\end{tabular}
\end{table}


\addfile{Global overview of noise robustness on ECFP4}
\label{app:noise:sec:ecfp4_overview}
% TODO: review, needs manual review
Noise robustness of every model on the ECFP4 representation, on the QM9 HOMO--LUMO gap.
% Thesis appendix item: ecfp4_overview
% Seeded on 2026-09-27 from qsar_qm_models/supplementary/items/ecfp4_overview.tex by thesis_appendix/seed_from_paper.py.
% From here on this copy is edited independently of the paper. Rules: thesis_appendix/README.md.
% The lines below, down to the first blank line, are the PAPER copy's header, kept for
% provenance. They describe the paper file, not this one: in particular a GENERATED
% notice there does not apply here, and this copy may be edited by hand.
% Supplementary item: ecfp4_overview
% Split out of additional_files.tex on 2026-09-27, where it was Additional file 6 (ECFP4 global overview (was Supp Figure S1)).
% Content unchanged apart from: the heading line removed (the build writes it), the old
% 'Table S*/Figure S*' caption prefix replaced by \SuppFile{}, and the label renamed to supp:ecfp4_overview.
% Rules: supplementary/README.md.

\begin{figure}[H]
\centering
\includegraphics[width=\textwidth,height=0.78\textheight,keepaspectratio]{fig1_supp_ecfp4_overview.png}
\caption{\textbf{Global overview of noise robustness on the ECFP4 representation, evaluated on QM9 HOMO--LUMO gap.}
a) Performance degradation curves (R$^2$ versus noise level $\sigma$) under Gaussian noise. b) Noise degradation slope (NDS) heatmap across model architectures and noise strategies; lower (more negative) NDS indicates faster degradation under noise.}
\label{app:noise:ecfp4_overview}
\end{figure}


\addfile{Clean $R^2$ against robustness on PDV under Gaussian noise}
\label{app:noise:sec:clean_r2_vs_robustness}
% TODO: review, needs manual review
Clean-label $R^2$ against the noise degradation slope for every model on PDV under Gaussian noise, on the QM9 HOMO--LUMO gap.
% Thesis appendix item: clean_r2_vs_robustness
% Seeded on 2026-09-27 from qsar_qm_models/supplementary/items/clean_r2_vs_robustness.tex by thesis_appendix/seed_from_paper.py.
% From here on this copy is edited independently of the paper. Rules: thesis_appendix/README.md.
% The lines below, down to the first blank line, are the PAPER copy's header, kept for
% provenance. They describe the paper file, not this one: in particular a GENERATED
% notice there does not apply here, and this copy may be edited by hand.
% Supplementary item: clean_r2_vs_robustness
% Split out of additional_files.tex on 2026-09-27, where it was Additional file 7 (Baseline R^2 vs NDS scatter (was Supp Figure S2)).
% Content unchanged apart from: the heading line removed (the build writes it), the old
% 'Table S*/Figure S*' caption prefix replaced by \SuppFile{}, and the label renamed to supp:clean_r2_vs_robustness.
% Rules: supplementary/README.md.

\begin{figure}[H]
\centering
\includegraphics[width=\textwidth]{fig3_ranking_consistency.png}
\caption{\textbf{Baseline R$^2$ versus noise degradation slope (NDS) scatter for all models on the physicochemical descriptor vector (PDV) representation under Gaussian noise, on the QM9 HOMO--LUMO gap.}
Lower (more negative) NDS indicates faster performance loss under noise.}
\label{app:noise:clean_r2_vs_robustness}
\end{figure}


\addfile{Bayesian transformation effects by representation}
\label{app:noise:sec:bayesian_transformation_tests}
% TODO: review, needs manual review
The effect of each Bayesian transformation of the neural networks on the noise degradation slope, by representation, on the QM9 HOMO--LUMO gap.
% Thesis appendix item: bayesian_transformation_tests
% Seeded on 2026-09-27 from qsar_qm_models/supplementary/items/bayesian_transformation_tests.tex by thesis_appendix/seed_from_paper.py.
% From here on this copy is edited independently of the paper. Rules: thesis_appendix/README.md.
% The lines below, down to the first blank line, are the PAPER copy's header, kept for
% provenance. They describe the paper file, not this one: in particular a GENERATED
% notice there does not apply here, and this copy may be edited by hand.
% Supplementary item: bayesian_transformation_tests
% Split out of additional_files.tex on 2026-09-27, where it was Additional file 8 (NN Bayesian transformation effects (was Supp Table S6)).
% Content unchanged apart from: the heading line removed (the build writes it), the old
% 'Table S*/Figure S*' caption prefix replaced by \SuppFile{}, and the label renamed to supp:bayesian_transformation_tests.
% Rules: supplementary/README.md.

\begin{table}[H]
\centering
\caption{\textbf{Neural network Bayesian transformation effects on noise degradation slope (NDS) by representation, on the QM9 HOMO--LUMO gap.}
Positive $\Delta$ NDS indicates the Bayesian variant is more robust (less negative slope) than its deterministic counterpart.}
\label{app:noise:bayesian_transformation_tests}
\small
\begin{tabular}{llrrrr}
\toprule
\textbf{Model} & \textbf{Representation} & \textbf{Baseline R$^2$} & \textbf{NN-$\beta$ NDS} & \textbf{BNN-$\beta$ NDS} & \textbf{$\Delta$ NDS} \\
\midrule
NN-$\beta$ & ECFP4 & 0.791 & $-0.464$ & $-0.388$ & $+0.076$ \\
NN-$\beta$ & PDV & 0.858 & $-0.487$ & $-0.382$ & $+0.106$ \\
NN-$\beta$ & SMILES & 0.721 & $-0.558$ & $-0.365$ & $+0.193$ \\
\bottomrule
\end{tabular}
\end{table}


\addfile{Uncertainty metrics for probabilistic models on ECFP4}
\label{app:noise:sec:uncertainty_metrics_ecfp4}
% TODO: review, needs manual review
Uncertainty metrics of the probabilistic models on ECFP4 under Gaussian noise, on the QM9 HOMO--LUMO gap.
% Thesis appendix item: uncertainty_metrics_ecfp4
% Seeded on 2026-09-27 from qsar_qm_models/supplementary/items/uncertainty_metrics_ecfp4.tex by thesis_appendix/seed_from_paper.py.
% From here on this copy is edited independently of the paper. Rules: thesis_appendix/README.md.
% The lines below, down to the first blank line, are the PAPER copy's header, kept for
% provenance. They describe the paper file, not this one: in particular a GENERATED
% notice there does not apply here, and this copy may be edited by hand.
% Supplementary item: uncertainty_metrics_ecfp4
% Split out of additional_files.tex on 2026-09-27, where it was Additional file 9 (Uncertainty metrics ECFP4 (was Supp Table S8)).
% Content unchanged apart from: the heading line removed (the build writes it), the old
% 'Table S*/Figure S*' caption prefix replaced by \SuppFile{}, and the label renamed to supp:uncertainty_metrics_ecfp4.
% Rules: supplementary/README.md.

\begin{table}[H]
\centering
\caption{\textbf{Uncertainty quantification metrics for probabilistic models on the ECFP4 representation (Gaussian noise strategy) on the QM9 HOMO--LUMO gap.}
Unc-Error $\rho$: Spearman correlation between predicted uncertainty and absolute prediction error. Unc-Noise $\rho$: Spearman correlation between predicted uncertainty and injected noise level $\sigma$. ECE: expected calibration error (lower = better-calibrated). Coverage targets: 68\% ($1\sigma$), 95\% ($2\sigma$).}
\label{app:noise:uncertainty_metrics_ecfp4}
\small
\begin{tabular}{lrrrrr}
\toprule
\textbf{Model} & \textbf{Unc-Error $\rho$} & \textbf{Unc-Noise $\rho$} & \textbf{ECE} & \textbf{Cov.\ $1\sigma$} & \textbf{Cov.\ $2\sigma$} \\
\midrule
\textbf{QRF}     & \textbf{0.25} & 0.27 & 0.19 & 74.7\% & 94.6\% \\
\textbf{NGBoost} & 0.22 & \textbf{0.37} & \textbf{0.11} & 69.2\% & 94.2\% \\
GP               & 0.18 & 0.34 & 0.23 & 78.6\% & 96.0\% \\
BNN-$\beta$      & 0.20 & 0.37 & 0.29 & 78.6\% & 95.2\% \\
BNN-$\alpha$     & 0.15 & 0.32 & 0.24 & 76.6\% & 95.1\% \\
VBLL-$\beta$     & 0.19 & 0.34 & 0.19 & 73.9\% & 95.5\% \\
VBLL-$\alpha$    & 0.11 & 0.31 & 0.22 & 74.5\% & 95.1\% \\
\bottomrule
\end{tabular}
\end{table}


\addfile{Variance decomposition on the three assay datasets}
\label{app:noise:sec:assay_variance_decomposition}
% TODO: review, needs manual review
ANOVA variance decomposition ($\eta^2$, in per cent) on LogD, Caco-2 efflux and hERG $K_i$.
% Thesis appendix item: assay_variance_decomposition
% Seeded on 2026-09-27 from qsar_qm_models/supplementary/items/assay_variance_decomposition.tex by thesis_appendix/seed_from_paper.py.
% From here on this copy is edited independently of the paper. Rules: thesis_appendix/README.md.
% The lines below, down to the first blank line, are the PAPER copy's header, kept for
% provenance. They describe the paper file, not this one: in particular a GENERATED
% notice there does not apply here, and this copy may be edited by hand.
% Supplementary item: assay_variance_decomposition
% Split out of additional_files.tex on 2026-09-27, where it was Additional file 10 (Validation ANOVA (was Supp Figure S3)).
% Content unchanged apart from: the heading line removed (the build writes it), the old
% 'Table S*/Figure S*' caption prefix replaced by \SuppFile{}, and the label renamed to supp:assay_variance_decomposition.
% Rules: supplementary/README.md.

\begin{figure}[H]
\centering
\includegraphics[width=\textwidth]{fig_validation_anova.png}
\caption{\textbf{ANOVA variance decomposition ($\eta^2$, \%) on external validation datasets: a) LogD, b) Caco-2 Efflux, and c) hERG K$_i$.}
Model, representation, and interaction contributions to noise degradation slope (NDS) variance for each dataset; lower (more negative) NDS indicates faster performance loss under noise.}
\label{app:noise:assay_variance_decomposition}
\end{figure}


\addfile{Random forest against quantile regression forest on the three assay datasets}
\label{app:noise:sec:rf_vs_qrf_assay}
% TODO: review, needs manual review
The noise degradation slope of random forest and of quantile regression forest on LogD, Caco-2 efflux and hERG $K_i$.
% Thesis appendix item: rf_vs_qrf_assay
% Seeded on 2026-09-27 from qsar_qm_models/supplementary/items/rf_vs_qrf_assay.tex by thesis_appendix/seed_from_paper.py.
% From here on this copy is edited independently of the paper. Rules: thesis_appendix/README.md.
% The lines below, down to the first blank line, are the PAPER copy's header, kept for
% provenance. They describe the paper file, not this one: in particular a GENERATED
% notice there does not apply here, and this copy may be edited by hand.
% Supplementary item: rf_vs_qrf_assay
% Split out of additional_files.tex on 2026-09-27, where it was Additional file 11 (RF vs QRF on validation datasets (was Supp Table S7)).
% Content unchanged apart from: the heading line removed (the build writes it), the old
% 'Table S*/Figure S*' caption prefix replaced by \SuppFile{}, and the label renamed to supp:rf_vs_qrf_assay.
% Rules: supplementary/README.md.

\begin{table}[H]
\centering
\caption{\textbf{Random forest (RF) vs quantile regression forest (QRF) noise degradation slope (NDS) on the three external validation datasets (LogD, Caco-2 Efflux, hERG K$_i$).}
Lower (more negative) NDS indicates faster degradation under noise.}
\label{app:noise:rf_vs_qrf_assay}
\small
\begin{tabular}{lrrr}
\toprule
\textbf{Dataset} & \textbf{RF NDS} & \textbf{QRF NDS} & \textbf{$\Delta$ NDS} \\
\midrule
ChEMBL-hERG-Ki     & $-0.231$ & $-0.359$ & $-0.128$ \\
OpenADMET-Caco2    & $-0.534$ & $-0.627$ & $-0.094$ \\
OpenADMET-LogD     & $-0.101$ & $-0.185$ & $-0.083$ \\
\bottomrule
\end{tabular}
\end{table}


\section{The flaws of the similarity assumption}\label{app:add:flaws}
% ---- NOTES: The flaws of the similarity assumption ----
% Takeaway: the per-dataset, per-representation counts behind Chapter 4's claims, in four tables.
% - Table, constant features. One row is one dataset and one representation. Columns: total features; constant features; near-constant features.
% - Table, molecules a representation cannot tell apart. One row is one representation and one dataset. Columns: molecules in groups with an identical feature vector; how many of those differ only in stereochemistry; lowest RMSE reachable when every group gets one prediction; largest label spread within a group.
% - Table, distance against property difference. One row is one dataset and one representation. Columns: rank correlation between pairwise distance and absolute property difference; score of predicting each molecule's activity from its five nearest neighbours.
% - Table, activity cliffs. One row is one dataset and one representation. Column: molecules in at least one activity cliff.
% Result needed: all four tables, from the full runs.
% ---- END NOTES ----
\todo[inline, color=orange!40]{Additional files for Chapter~\ref{chap:flaws}: all four need the full runs.}

\section{Building adaptive hybrid representations}\label{app:add:method}\label{app:searchspace}
% ---- NOTES: Building adaptive hybrid representations ----
% Takeaway: these files are what let a reader rebuild any hybrid in the thesis. Chapter 5 gives the reasons; this section holds only the reference tables.
% - Change: the design-experiment file's allocation row gives only counts. Add the reason the global pool with squared weights was kept over the other pooled candidates, once Chapter 5 states it.
% Result needed: the best representation count per (downstream model, budget) from the design run, for the design-experiment file's representation-count row.
% Result needed: per dataset and downstream model, per-fold score of the best single representation against the hybrid; the winning configuration; how often each setting value won.
% Result needed: augmentation per dataset: how often no extra features were chosen; per-fold base against base plus extras, uncapped and capped; which extra feature sets were kept; decoy features admitted.
% Result needed: per dataset and budget, held-out score against the score the inner search reported, and how stable the selected features are across folds.
% Result needed: hybrid minus each comparator (concatenation, one model per representation, PCA, mRMR, recursive feature elimination, random features) at each budget.
% Result needed: the smaller comparator, control and ablation experiments, one table each, if they leave Chapter 5.
% Result needed: per task, outcome of the transfer, antibody-scan and separate-site experiments, with which representations the selected features came from.
% Result needed: hours and model fits per dataset at each tier, with the hardware.
% ---- END NOTES ----

\addfile{Settings of every model, as downstream model and as internal evaluator}
\label{app:add:model_settings}
% TODO: review, needs manual review
One row is one model. The middle two columns give its settings when it scores a hybrid (downstream model) and when it scores one representation inside the search (internal evaluator). The models fitted only to compute feature importance are described in Chapter~\ref{chap:method}.

% TODO: review, needs manual review (whole table and caption, written from the current code)
{\small
\begin{longtable}{@{}p{1.9cm}p{4.6cm}p{3.6cm}p{1.6cm}p{2.3cm}@{}}
\caption{Settings of every model. $p$ is the number of features. ``Molecules'' is the number of training molecules the internal evaluator is fitted on; above it, a random subset of that size is drawn with a fixed seed. Settings not listed are the library's defaults. Every setting is the same on every dataset. The Gaussian process runs only on the datasets named in Chapter~\ref{chap:method}, ``Models and metrics''.}
\label{tab:add:model_settings}\\
\toprule
Model & As downstream model & As internal evaluator & Molecules & Library; source \\
\midrule
\endfirsthead
\toprule
Model & As downstream model & As internal evaluator & Molecules & Library; source \\
\midrule
\endhead
Random forest & 100 trees, no depth limit; $\sqrt{p}$ features tried at each split & 50 trees, depth at most 10; all $p$ features tried at each split on regression, $\sqrt{p}$ on classification & all & scikit-learn \citep{scikit-learn}; \citet{Breiman2001} \\
\addlinespace
LightGBM & 100 trees & 50 trees, depth at most 4 & all & lightgbm; \citet{ke2017lightgbm} \\
\addlinespace
XGBoost & 100 trees; half of the features sampled for each tree & 50 trees, depth at most 4 & all & xgboost; \citet{Chen2016} \\
\addlinespace
SVM & RBF kernel, $C = 1$. Not fitted above 10{,}000 molecules, or above 5{,}000 under uncapped augmentation & RBF kernel, $C = 1$ & at most 500 & scikit-learn; \citet{Vapnik1995} \\
\addlinespace
Gaussian process & Regression: RBF plus white-noise kernel, labels normalised, 3 optimiser restarts. Classification: RBF kernel. Not fitted above 1{,}000 molecules & Regression: scaled RBF kernel, fixed noise variance 0.01, no restarts. Classification: scaled RBF kernel, no restarts & at most 500 & scikit-learn; \citet{Rasmussen2006} \\
\addlinespace
DNN & Hidden layers of 256, 128 and 64 units, ReLU, dropout 0.2 \citep{Srivastava2014}; Adam \citep{Kingma2015} at learning rate 0.001, halved after 5 epochs without improvement; batch size 256; up to 200 epochs; stops after 20 epochs without improvement on 20\% of the training molecules that share no group with the rest, using the outer split's grouping; labels z-normalised & The same, except that the early-stopping 20\% is drawn at random & all & PyTorch \citep{Paszke2019} \\
\bottomrule
\end{longtable}
}

\addfile{Downstream models left out, by dataset and augmentation setting}
\label{app:add:models_left_out}
% TODO: review, needs manual review
One row is one rule that leaves a downstream model out. Only the 16 datasets of Table~\ref{tab:datasets} are listed; a dataset and setting not named in any row runs every downstream model in Table~\ref{tab:search_tiers}.

% TODO: review, needs manual review (whole table and caption, written from the current code)
{\small
\begin{longtable}{@{}p{2.4cm}p{3.0cm}p{8.6cm}@{}}
\caption{Downstream models left out, by dataset and augmentation setting. Every rule was set so that a search finishes within the time limit of one job; none was set from scores. hERG $K_i$ drops no model under any setting.}
\label{tab:add:models_left_out}\\
\toprule
Model left out & Augmentation setting & Datasets \\
\midrule
\endfirsthead
\toprule
Model left out & Augmentation setting & Datasets \\
\midrule
\endhead
XGBoost & all three & H3K4me3, TATA promoter \\
SVM & all three & datasets above 10{,}000 molecules: H3K4me3 \\
SVM & uncapped & datasets above 5{,}000 molecules: QM9, LogD$^{a}$, TATA promoter$^{a}$ \\
SVM & uncapped & Koenig expression \\
SVM & capped & Caco-2 efflux, LogD, rat liver microsome clearance, MDR1--MDCK efflux, Buchwald--Hartwig plates \\
Random forest and SVM & capped & QM9 \\
Random forest & uncapped & QM9, LogD \\
Gaussian process & all three & every dataset except humanness (Chapter~\ref{chap:method}, ``Models and metrics'') \\
\bottomrule
\end{longtable}
{\footnotesize $^{a}$Counted on the whole dataset: LogD has 5{,}028 molecules and TATA promoter 5{,}274. A search run inside one outer fold is given about 80\% of them and keeps the SVM.}
}

\addfile{Size of the search}
\label{app:add:search_size}
% TODO: review, needs manual review
One row is one downstream model. The counts are per dataset, objective and augmentation setting, with adaptive pruning off.

% TODO: review, needs manual review (whole table and caption)
\begin{table}[htbp]
\centering
\small
\caption{Number of configurations searched per dataset, objective and augmentation setting. A configuration is one choice of downstream model, internal evaluator, allocation method, importance method, representation count and budget. At the full tier each (downstream model, internal evaluator) pair runs 5 allocation methods $\times$ 3 importance methods $\times$ 4 representation counts $\times$ 8 budgets $= 480$ configurations; at the fast tier it runs $3 \times 1 \times 2 \times 8 = 48$ (Table~\ref{tab:search_tiers}). ``Matched'' is counted once where it repeats an evaluator already in the list. Each configuration is scored on 5 folds, so the full tier fits a downstream model 40{,}800 times (50{,}400 on humanness) and the fast tier 4{,}080 times. Where a dataset leaves a downstream model out (Additional file~\ref{app:add:models_left_out}), its pairs drop out: a fast search on H3K4me3 has 9 pairs, so 432 configurations.}
\label{tab:add:search_size}
\begin{tabular}{@{}lp{4.2cm}rr@{}}
\toprule
Downstream model & Internal evaluators & Full tier & Fast tier \\
\midrule
Random forest & 3 (matched is random forest) & 1{,}440 & 144 \\
LightGBM & 3 (matched is LightGBM) & 1{,}440 & 144 \\
DNN & 3 (matched is DNN) & 1{,}440 & 144 \\
XGBoost & 4 & 1{,}920 & 192 \\
SVM & 4 & 1{,}920 & 192 \\
Gaussian process (humanness only) & 4 & 1{,}920 & not run \\
\midrule
Total & 17 pairs (21 on humanness) & 8{,}160 (10{,}080 on humanness) & 816 \\
\bottomrule
\end{tabular}
\end{table}

\addfile{Design experiments behind the search settings}
\label{app:add:design_experiments}
% TODO: review, needs manual review
One row is one setting of the search, or one rule that was tested and then fixed instead of searched. Every number comes from preliminary design runs made before the main experiments; they are not results of the thesis, and several used other budgets than the final search.

% TODO: review, needs manual review (whole table and caption; numbers from preliminary design runs)
{\footnotesize
\begin{longtable}{@{}p{1.9cm}p{2.5cm}p{2.7cm}p{4.4cm}p{2.5cm}@{}}
\caption{Design experiments behind the search settings. Preliminary design runs, not the thesis's full runs.}
\label{tab:add:design_experiments}\\
\toprule
Setting & Values compared & Datasets, folds, budgets & What was measured, and the result & Values kept (full / fast) \\
\midrule
\endfirsthead
\toprule
Setting & Values compared & Datasets, folds, budgets & What was measured, and the result & Values kept (full / fast) \\
\midrule
\endhead
Downstream model & Nine models, among them random forest, LightGBM, XGBoost, CatBoost, DNN, $k$-nearest neighbours and a Gaussian process, each on each of 10 base representations at full width & Buchwald--Hartwig yield, Caco-2 efflux, Koenig expression, LogD, breast cancer PAM50 subtype; 5 folds & Best held-out score over the 10 representations, and seconds per fit. The DNN scored best on 4 of the 5 datasets and third on Caco-2 (0.2936 against 0.2851 RMSE, log$_{10}$ efflux ratio) & Full: random forest, LightGBM, XGBoost, SVM, DNN, and the Gaussian process on humanness / fast: the same five \\
\addlinespace
Budget & None compared & -- & -- & 10, 25, 50, 100, 250, 500, 1{,}000 and 2{,}000 features at both tiers \\
\addlinespace
Representation count & 2, 3, 5 and all 10 & 9 datasets, 3 tree-ensemble downstream models, budgets 500 and 1{,}000 features; 5 folds & The best count for each (downstream model, budget) combination & Full: 2, 3, 5, all / fast: 3, all \\
\addlinespace
Allocation method & 8 candidates: the five searched and 3 further pooled candidates, each at every representation count & The same 9 datasets, models and budgets: 54 (dataset, model, budget) combinations & Each candidate was scored at its best representation count, and the number of combinations in which it was in the top 3 of the 8 was counted. Fixed shares: linear 30, squared 27, rank decay 26, softmax 25; global pool with squared weights 14 & Full: the five described in Chapter~\ref{chap:method} / fast: linear, squared, rank decay \\
\addlinespace
Importance method, agreement & 8 methods & LogD, hERG $K_i$, Buchwald--Hartwig yield, nanobody thermostability, antibody FcRn retention, breast cancer PAM50 subtype, lung adenocarcinoma vs squamous; each method fitted once on the whole dataset & Spearman correlation between two methods' scores over every feature. Kept pairs: LightGBM gain and TreeSHAP 0.20--0.64, LightGBM gain and DeepLIFT 0.21--0.67, TreeSHAP and DeepLIFT 0.56--0.99 across the 7 datasets & One method kept from each group whose rankings agree above 0.95 \\
\addlinespace
Importance method, score & 28 methods, among them three controls: feature variance, Laplacian score and random ranking & Caco-2 efflux; 3 folds; budget 100 features; random forest and LightGBM downstream & RMSE of the best single representation minus RMSE of the hybrid (log$_{10}$ efflux ratio; positive means the hybrid has the lower error): DeepLIFT $+0.0144$, LightGBM gain $+0.0100$, TreeSHAP $+0.0067$, random ranking $-0.0118$ & Full: LightGBM gain, TreeSHAP, DeepLIFT / fast: LightGBM gain \\
\addlinespace
Internal evaluator settings & 20 hyperparameter settings from 5 model families & Buchwald--Hartwig yield, Caco-2 efflux, Koenig expression; 3 folds; budget 100 features & Spread of mean scores across all 20 settings: 0.0253 $R^2$, 0.0179 $R^2$ and 0.0085 Spearman. Median standard deviation across the 3 folds for one setting, allocation method and downstream model: 0.0462 $R^2$, 0.0693 $R^2$ and 0.0186 Spearman & Matched, random forest, LightGBM and DNN, at fixed settings, both tiers \\
\addlinespace
Rescaling of each representation's score (a rule) & 5 rescalings & 9 datasets; 5 folds; budgets 50 and 300 features & Whether any rescaling beat the one in use by more than that rule's own spread across folds: in 0 of 90 (dataset, allocation method, budget) combinations & One fixed rule per metric (Chapter~\ref{chap:method}) \\
\bottomrule
\end{longtable}
{\footnotesize The representation-count and allocation rows come from one design run on 10 datasets, 3 downstream models and 2 budgets (60 combinations); QM9 is left out, which gives the 54.}
}

\addfile{Fields a blueprint saves}
\label{app:add:blueprint_fields}
% TODO: review, needs manual review
One row is one field. Together the fields are the recipe that rebuilds a hybrid.

% TODO: review, needs manual review (whole table and caption, written from the current code)
{\small
\begin{longtable}{@{}p{4.6cm}p{9.4cm}@{}}
\caption{Fields a blueprint saves. A rebuild computes the internal-evaluator scores that split the budget again, on the molecules it is given.}
\label{tab:add:blueprint_fields}\\
\toprule
Field & What it holds \\
\midrule
\endfirsthead
\toprule
Field & What it holds \\
\midrule
\endhead
\multicolumn{2}{@{}l}{\textit{The recipe}} \\
\midrule
Representations & The representations offered to the search \\
Representation count & How many of them the configuration used; empty when it used all \\
Allocation method & One of the five described in Chapter~\ref{chap:method} \\
Importance method & LightGBM gain, TreeSHAP or DeepLIFT \\
Budget & Number of features in the hybrid \\
Internal evaluator & Matched, random forest, LightGBM or DNN \\
Objective & Predictive, or predicting each molecule's activity from its five nearest neighbours \\
Metric & The dataset's metric, and whether lower is better \\
Task type & Regression or classification \\
Folds and seed & The number of folds, and the random seed \\
\midrule
\multicolumn{2}{@{}l}{\textit{The recipe: augmentation}} \\
\midrule
Augmentation setting & None, uncapped or capped \\
Candidate ranking & How candidate extra features are ranked \\
Candidate list length & 50 features \\
Cap per source & Optional; the most candidates taken from one source \\
Margin & How much better than adding none a block of extra features must score \\
Fraction of the budget & Optional; used instead of the fixed block sizes \\
\bottomrule
\end{longtable}
}

\todo[inline, color=orange!40]{Additional files for Chapter~\ref{chap:method}: the per-dataset results, augmentation, held-out against inner-search scores, comparators, smaller experiments, transfer experiments and compute need the full runs.}

\section{Hybrid representations for similarity and label noise}\label{app:add:similarity}
% ---- NOTES: Hybrid representations for similarity and label noise ----
% Takeaway: the full tables behind Chapter 6's summaries, and the settings of the twelve LNL methods.
% - Table, finding injected noisy labels. One row is one noise score. Columns: average precision under outlier noise at 1\%, 5\% and 10\%, grouped-wider noise and censoring, in each feature space; chance level in the header row.
% - The feature spaces: each model's top two single representations, the same features weighted by importance, forest proximity, the hybrid, and the hybrid ranked on the neighbour score.
% - Table, finding real suspect labels. One row is one noise score and feature space. Columns: average precision against ChEMBL curator flags on hERG IC50, per record and per molecule, with chance level; Spearman correlation with the spread between laboratories on the BindingDB molecules with repeat measurements.
% - Table, estimated noise size. One row is one dataset and feature space. Columns: the estimate from repeat measurements (with the measured variance between repeats where the dataset keeps them); the estimate from near-duplicate pairs; the gamma-test estimate; the model's held-out mean squared error. All in label units squared.
% Result needed: the three tables above, from the full runs.
% Result needed: per dataset and fold, importance-weighted against unweighted features for predicting each molecule's activity from its five nearest neighbours; the hybrid searched for that objective against the ordinary hybrid and the two best single representations.
% Result needed: censoring: share of the lost score recovered, per treatment, model, feature space and clipping level; deleting clipped molecules against keeping them.
% Result needed: the survey of twelve LNL methods: per noise level, the methods that beat the plain network on each input.
% Result needed: cleaning real labels: folds won per cleaning method, noise score and feature space.
% Result needed: per noise type, whether any LNL method, noise mitigation or noise-robust training, improved on the untreated model, on each feature space.
% ---- END NOTES ----

\addfile{Settings of the twelve LNL methods and the plain network they share}
\label{app:add:lnl_settings}
% TODO: review, needs manual review
One row is one method; the first row is the plain network every method starts from. The kind of each method follows Chapter~\ref{chap:similarity}.

% TODO: review, needs manual review (whole table and caption, written from the current code)
{\small
\begin{longtable}{@{}p{2.9cm}p{2.0cm}p{7.3cm}p{1.6cm}@{}}
\caption{Settings of the twelve LNL methods. On classification datasets the squared-error loss is replaced by cross-entropy.}
\label{tab:add:lnl_settings}\\
\toprule
Method & Kind & Settings & Labels \\
\midrule
\endfirsthead
\toprule
Method & Kind & Settings & Labels \\
\midrule
\endhead
Plain network & -- & The DNN of Table~\ref{tab:add:model_settings}, with the gradient norm clipped at 1. After a method has removed molecules, the 20\% held out for early stopping is drawn at random & both \\
\midrule
Small-loss selection & noise mitigation & 10 warm-up epochs on all molecules; keeps the 70\% with the smallest loss and fits a new plain network on them & both \\
Confident learning & noise mitigation & 3 plain networks; removes molecules whose error is above the 85th percentile and whose spread across the 3 networks is below the 25th percentile, then refits & both \\
Removal by distance to neighbours & noise mitigation & 10 nearest neighbours in the standardised features; score is the distance of a label from its neighbours' mean label, divided by their standard deviation (on classification, the share of neighbours with another label); removes molecules above the 85th percentile, then refits & both \\
Staged removal & noise mitigation & 3 stages; each removes the 10\% of the remaining molecules with the largest training loss & both \\
\midrule
Per-epoch down-weighting & noise-robust training & After each epoch, the weight of every molecule whose loss is above the median is multiplied by 0.95 & both \\
Curriculum by confidence & noise-robust training & 20 warm-up epochs, then trains on the 25\%, 50\%, 75\% and 100\% of molecules with the smallest loss, in stages & both \\
DivideMix & noise-robust training & Two networks, 10 warm-up epochs; a two-component Gaussian mixture on each network's squared errors splits the molecules & both \\
Meta-Weight-Net & noise-robust training & The weighting network is updated every 5 epochs at learning rate 0.001, on 20\% of the training molecules held out for it & both \\
MentorNet & noise-robust training & The mentor network reads each molecule's current loss, its last 5 losses and its mean distance to its 10 nearest neighbours & both \\
Contrast to Divide & noise-robust training & 5 perturbed copies of each molecule, Gaussian noise of standard deviation 0.01 on the features & both \\
\midrule
Heteroscedastic GP & noise-robust training & RBF kernel, 2 optimiser restarts, fitted on at most 1{,}000 training molecules; a small network predicts each molecule's noise variance & regression \\
Deep evidential regression & noise-robust training & Weight of the evidence penalty 0.01 & regression \\
\bottomrule
\end{longtable}
}

\todo[inline, color=orange!40]{Additional files for Chapter~\ref{chap:similarity}: finding injected and real noisy labels, estimated noise size and the other results need the full runs.}
\chapter{Data and code availability}\label{app:availability}
% TODO: needs manual review (whole chapter written by Claude)
% ---- TAKEAWAYS ----
% - Three separate pieces of software lie behind the thesis: NoiseInject, the QM9 pipeline and the KIRBy repository.
% - Each can be reused on its own; the hybrid search runs on any precomputed feature matrices.
% - Every dataset is public except Lhasa's curated hERG sets, which are proprietary.
% - Every result can be reproduced from public data and code, except the hERG work with data held at separate sites.
% - Chapter 3 is joint work; everything else was done by A.P. alone.
% ---- END TAKEAWAYS ----
% ---- NOTES: Data and code availability ----
% - Change: the NoiseInject block says version 1.0.0, but the archived DOI holds v0.3.0 (3 June 2026). Replace the DOI once a 1.0.0 release is archived.
% - Change: the QM9 pipeline block says Linux and macOS; neither has been checked.
% ---- END NOTES ----

% TODO: review, needs manual review
Three pieces of software lie behind this thesis. NoiseInject is the label-noise package released with Chapter~\ref{chap:robustness}. A separate pipeline ran the QM9 experiments of that chapter. The KIRBy repository implements the adaptive hybrid representation and holds the experiments of Chapters~\ref{chap:flaws}, \ref{chap:method} and~\ref{chap:similarity}, together with the assay-dataset runs of Chapter~\ref{chap:robustness}. The three are separate, so each can be reused on its own. For example, the hybrid search runs on feature matrices a reader has computed elsewhere, without the Rust part of the QM9 pipeline. The source and curation of each dataset are given in Appendix~\ref{app:datasets}.

NoiseInject implements the six regression and six classification noise injection strategies and the robustness and uncertainty-calibration metrics used in this work. It is available as an open-source Python package under an MIT license \citep{noiseinject}. It is model-agnostic and operates on label arrays.

% TODO: review, needs manual review
Its noise conditions and metrics are described in Section~\ref{sec:noise_injection}.

\begin{itemize}
    \item Project name: NoiseInject
    \item Project home page: \url{https://github.com/adpunt/noise_inject}
    \item Archived version: \url{https://doi.org/10.5281/zenodo.20532710}
    \item Version: 1.0.0
    \item Operating system(s): Platform independent
    \item Programming language: Python 3.7 or later
    \item Other requirements: NumPy, pandas, SciPy and scikit-learn; see the repository's \texttt{requirements.txt}
    \item License: MIT
    \item Any restrictions to use by non-academics: None
\end{itemize}

% TODO: review, needs manual review
The QM9 results of Chapter~\ref{chap:robustness} were produced by a separate repository, written in Python and Rust, which also holds the hyperparameter tuning, the variance decomposition and the figure scripts for that chapter. It is described in Chapter~\ref{chap:robustness}, under ``The QM9 precompute pipeline''.

\begin{itemize}
    \item Project name: qsar\_qm\_models
    \item Project home page: \url{https://github.com/adpunt/qsar_qm_models}
    \item Archived version: not yet archived
    \item Operating system(s): Linux and macOS
    \item Programming language: Python and Rust (2021 edition)
    \item Other requirements: RDKit installed through conda, to which the Rust program links; PyTorch and PyTorch Geometric
    \item License: GNU GPL v3.0
    \item Any restrictions to use by non-academics: those of the GNU GPL v3.0
\end{itemize}

% TODO: review, needs manual review
The adaptive hybrid representation is implemented in the KIRBy repository. It is written in Python and works on feature matrices computed beforehand, so the hybrid search can be applied to representations from any domain. The repository contains the hybrid search, loaders for the datasets in Appendix~\ref{app:datasets} and generators for the representations of each application area. It also holds the neighbour-based similarity analyses, the learning-with-noisy-labels methods of Chapter~\ref{chap:similarity} and every experiment script. Label noise is added through NoiseInject, so the noise definitions are the same as in Chapter~\ref{chap:robustness}. A GPU is not needed for the hybrid search, but it speeds up the pretrained encoders.

\begin{itemize}
    \item Project name: KIRBy
    \item Project home page: \url{https://github.com/adpunt/KIRBy}
    \item Archived version: [Zenodo DOI]
    \item Operating system(s): Linux and macOS
    \item Programming language: Python 3.8 or later; several pretrained encoders need Python 3.10 or later
    \item Other requirements: NumPy, pandas, SciPy, scikit-learn, and RDKit installed through conda; each application area has an optional install
    \item License: [licence]
    \item Any restrictions to use by non-academics: [depends on the licence]
\end{itemize}


% TODO: review, needs manual review
All datasets but one are publicly available. QM9 \citep{Ramakrishnan2014} is deposited on Figshare \citep{qm9_dataset}. The LogD and Caco-2 efflux data come from the OpenADMET--ExpansionRx challenge release, under a CC-BY-4.0 licence \citep{openadmet}. The hERG $K_i$ data were extracted from ChEMBL \citep{Zdrazil2023}, whose data are released under CC BY-SA 3.0. They can be regenerated with the curation protocol in Appendix~\ref{app:datasets}. The rat liver microsome clearance and MDR1--MDCK efflux data are the Biogen ADME set \citep{fang2023}, released on Polaris \citep{polaris} under CC-BY-4.0. The reaction-yield, antibody, nucleotide and tumour gene-expression datasets were downloaded from their public releases, which are cited in Appendix~\ref{app:datasets}. Pretrained model weights were downloaded from their original sources and are not redistributed.

% TODO: review, needs manual review
The exception is the hERG blocker data used for feature selection with data held at separate sites (Chapter~\ref{chap:method}). These are Lhasa Limited's curated hERG training set of 7,772 molecules and test set of 4,383 molecules \citep{Hanser2025}. They are proprietary and are not redistributed. The comparison with FLuID also calls Lhasa's FLuID code and cannot be run without it. Every other result in this thesis can be reproduced from public data and the code above.

% TODO: review, needs manual review
% TODO: update
Chapter~\ref{chap:robustness} is adapted from a jointly authored manuscript by Punt, Hanser, Werner and Morris, in preparation for the \textit{Journal of Cheminformatics}. For that work, A.P. implemented the NoiseInject framework, conducted all experiments, and drafted the manuscript. T.H., S.W., and G.M. supervised the project, contributed to the experimental design, and revised the manuscript. All authors read and approved the final manuscript. The adaptive hybrid representation, the KIRBy repository, and the experiments of Chapters~\ref{chap:flaws}, \ref{chap:method} and~\ref{chap:similarity} were designed, implemented, run and written up by A.P. alone. Lhasa Limited supplied the curated hERG data and the FLuID code \citep{Hanser2025}. The use of generative AI is described in the declaration in the front matter.
